% The TandemLLM manual. Build: latexmk -lualatex tandemllm-manual.tex (or `make` in manual/).
\documentclass[paper=a4,fontsize=11pt,DIV=11,twoside=false,open=any,parskip=half-,
  toc=listof,toc=index,headings=standardclasses,numbers=noenddot]{scrbook}

% The manual's preamble: fonts, colours, page furniture, tables, code, figures, index, links.
% Built with LuaLaTeX (latexmk -lualatex); no shell escape is needed.

\usepackage{fontspec}
\setmainfont{LibertinusSerif}[
  Extension=.otf, UprightFont=*-Regular, BoldFont=*-Semibold,
  ItalicFont=*-Italic, BoldItalicFont=*-SemiboldItalic, Ligatures=TeX]
\setsansfont{LibertinusSans}[
  Extension=.otf, UprightFont=*-Regular, BoldFont=*-Bold, ItalicFont=*-Italic,
  Ligatures=TeX]
\setmonofont{Inconsolatazi4}[
  Extension=.otf, UprightFont=*-Regular, BoldFont=*-Bold, Scale=MatchLowercase]

\usepackage[english]{babel}
\usepackage{microtype}
\usepackage[dvipsnames,table]{xcolor}
\usepackage{graphicx}
\usepackage{amsmath}
\usepackage{array,booktabs,longtable,tabularx,multirow}
\usepackage{enumitem}
\usepackage{caption}
\usepackage{float}
\graphicspath{{screenshots/}}
% \screenshot[includegraphics options]{file}{caption}{label}: a screenshot of the running application
\newcommand{\screenshot}[4][]{%
  \begin{figure}[htbp]\centering
  \setlength{\fboxsep}{0pt}\setlength{\fboxrule}{.4pt}%
  \fcolorbox{rule}{white}{\includegraphics[width=\dimexpr\linewidth-1pt\relax,#1]{#2}}%
  \caption{#3}\label{#4}\end{figure}}
\usepackage{csquotes}
\usepackage[obeyspaces,spaces]{url}
\usepackage{listings}
\usepackage{tikz}
\usepackage{pgfplots}
\usepackage{pgfplotstable}
\usepackage[noautomatic]{imakeidx}
\usepackage{scrlayer-scrpage}

% ---------------------------------------------------------------------------- colours
% One palette for the whole book: cobalt for structure and links, signal orange for the one
% thing a figure is about, ink and slate for text, rule and paper for furniture.
\definecolor{cobalt}{HTML}{1F4FCC}
\definecolor{cobaltlight}{HTML}{E8EEFB}
\definecolor{sky}{HTML}{8FB3FF}
\definecolor{signal}{HTML}{E8641B}
\definecolor{signallight}{HTML}{FCEBDD}
\definecolor{ink}{HTML}{1C2230}
\definecolor{slate}{HTML}{5A6475}
\definecolor{rule}{HTML}{C9D0DC}
\definecolor{paper}{HTML}{F5F7FA}
\definecolor{moss}{HTML}{2E8B57}

% ---------------------------------------------------------------------------- page and headings
\setlength{\emergencystretch}{2em}
\clubpenalty=10000
\widowpenalty=10000
\displaywidowpenalty=10000
\addtokomafont{disposition}{\color{ink}}
\setkomafont{chapterprefix}{\normalfont\sffamily\Large\color{slate}}
\addtokomafont{chapter}{\color{cobalt}}
\setkomafont{part}{\normalfont\bfseries\fontsize{30}{36}\selectfont\color{cobalt}}
\setkomafont{partnumber}{\normalfont\sffamily\LARGE\color{slate}}
\setkomafont{descriptionlabel}{\normalfont\bfseries}
\addtokomafont{caption}{\small}
\addtokomafont{captionlabel}{\sffamily\bfseries\color{cobalt}}
\setkomafont{pageheadfoot}{\small\sffamily\color{slate}}
\setkomafont{pagenumber}{\small\sffamily\color{ink}}
\renewcommand*{\sectionformat}{\textcolor{slate}{\thesection\autodot}\enskip}
\renewcommand*{\subsectionformat}{\textcolor{slate}{\thesubsection\autodot}\enskip}
\RedeclareSectionCommand[beforeskip=-2.4ex plus -1ex minus -.2ex,afterskip=1.1ex plus .2ex]{section}
\RedeclareSectionCommand[beforeskip=-1.8ex plus -1ex minus -.2ex,afterskip=.7ex plus .2ex]{subsection}
% a signal rule closes every chapter heading
\renewcommand*{\chapterheadendvskip}{\par\nobreak\vspace{1.6ex}{\color{signal}\hrule height .9pt}\nobreak\vspace{4.5ex}}
\renewcommand*{\partpagestyle}{empty}
% \partintro{text}: the short preamble under a part title, after a signal rule
\newcommand{\partintro}[1]{\setpartpreamble[c][.78\textwidth]{\vspace*{3ex}{\centering\color{signal}\rule{3cm}{1.2pt}\par}\vspace{3.5ex}\small\noindent #1}}
\setkomafont{partentry}{\sffamily\bfseries\color{cobalt}}
\setkomafont{chapterentry}{\bfseries}
\setkomafont{chapterentrypagenumber}{\normalfont}

% running heads: the chapter on the inside, the section on the outside, a hairline under both
\pagestyle{scrheadings}
\clearpairofpagestyles
\ihead{\leftmark}
\ohead{\rightmark}
\ofoot*{\pagemark}
\automark[section]{chapter}
\KOMAoptions{autooneside=false}
\renewcommand*{\chaptermarkformat}{\thechapter\autodot\enskip}
\renewcommand*{\sectionmarkformat}{\thesection\autodot\enskip}
\KOMAoptions{headsepline=.4pt}
\setkomafont{headsepline}{\color{rule}}
\makeatletter\renewcommand\@pnumwidth{2.2em}\makeatother
\setcounter{tocdepth}{1}
\setcounter{secnumdepth}{2}
\setlist{itemsep=.2ex,topsep=.6ex}
\setlist[description]{style=nextline, font=\normalfont\bfseries}
\captionsetup{font=small,labelfont={sf,bf,color=cobalt},format=plain,justification=raggedright,singlelinecheck=false,skip=6pt}

\renewcommand{\topfraction}{.9}
\renewcommand{\bottomfraction}{.7}
\renewcommand{\textfraction}{.08}
\renewcommand{\floatpagefraction}{.75}
\setcounter{topnumber}{3}
\setcounter{bottomnumber}{2}
\setcounter{totalnumber}{5}
% floats on a float page sit at the top, not in the middle
\makeatletter\setlength{\@fptop}{0pt}\setlength{\@fpbot}{0pt plus 1fil}\makeatother

% ---------------------------------------------------------------------------- inline code and terms
\urlstyle{tt}
\makeatletter
\g@addto@macro{\UrlBreaks}{\do\_\do\-\do\.\do\=\do\,\do\:}
\makeatother
\Urlmuskip=0mu plus 1mu
% \code{...}: identifiers, flags, variables, routes and paths, set in monospace and broken
% after punctuation; for text with braces, %, # or \ use \texttt with escapes instead.
\DeclareUrlCommand\code{\urlstyle{tt}}
% \term{word}: a defined term, set in italics and put in the index
\newcommand{\term}[1]{\emph{#1}\index{#1}}
% the three kinds of number (About this manual, "How numbers are labelled")
\newcommand{\measured}{\textsf{\small measured}}
\newcommand{\arith}{\textsf{\small arithmetic}}
% \hd{...}: a table header cell
\newcommand{\hd}[1]{{\sffamily\color{slate}#1}}

% ---------------------------------------------------------------------------- tables
\renewcommand{\arraystretch}{1.15}
\arrayrulecolor{rule}
\newcolumntype{L}[1]{>{\raggedright\arraybackslash}p{#1}}
\newcolumntype{Y}{>{\raggedright\arraybackslash}X}
% reftable: a long reference table with three columns, a caption, a label and a header
\newenvironment{reftable}[3]{%
  \footnotesize
  \setlength{\LTpre}{.8ex}\setlength{\LTpost}{.8ex}%
  \longtable{@{}L{.29\linewidth}L{.19\linewidth}L{.445\linewidth}@{}}
  \caption{#1}\label{#2}\\
  \toprule #3\\ \midrule \endfirsthead
  \multicolumn{3}{@{}l}{\textit{\small continued}}\\ \toprule #3\\ \midrule \endhead
  \bottomrule \endlastfoot
}{\endlongtable}
% httptable: the HTTP routes, three columns: method and path, who may call, what it does
\newenvironment{httptable}[2]{%
  \footnotesize
  \setlength{\LTpre}{.8ex}\setlength{\LTpost}{.8ex}%
  \longtable{@{}L{.31\linewidth}L{.09\linewidth}L{.55\linewidth}@{}}
  \caption{#1}\label{#2}\\
  \toprule \hd{route} & \hd{who} & \hd{what, body, answer}\\ \midrule \endfirsthead
  \multicolumn{3}{@{}l}{\textit{\small continued}}\\ \toprule \hd{route} & \hd{who} & \hd{what, body, answer}\\ \midrule \endhead
  \bottomrule \endlastfoot
}{\endlongtable}
% reftabletwo: the same with two columns, a name and what it is
\newenvironment{reftabletwo}[3]{%
  \footnotesize
  \setlength{\LTpre}{.8ex}\setlength{\LTpost}{.8ex}%
  \longtable{@{}L{.33\linewidth}L{.62\linewidth}@{}}
  \caption{#1}\label{#2}\\
  \toprule #3\\ \midrule \endfirsthead
  \multicolumn{2}{@{}l}{\textit{\small continued}}\\ \toprule #3\\ \midrule \endhead
  \bottomrule \endlastfoot
}{\endlongtable}

% ---------------------------------------------------------------------------- code listings
% One style for every block of code, commands, output and status lines.
\lstdefinestyle{tandem}{
  basicstyle=\ttfamily\small,
  columns=fullflexible, keepspaces=true,
  breaklines=true, breakatwhitespace=false,
  postbreak=\mbox{\textcolor{slate}{$\hookrightarrow$}\space},
  frame=leftline, framerule=1.2pt, rulecolor=\color{cobalt},
  backgroundcolor=\color{paper},
  xleftmargin=1.2em, framexleftmargin=.6em, aboveskip=1.2ex, belowskip=1.2ex,
  showstringspaces=false, upquote=true, extendedchars=true,
  literate={›}{{\guilsinglright}}1 {µ}{{\textmu}}1 {~}{{\textasciitilde}}1,
  commentstyle=\color{slate},
  keywordstyle=\color{cobalt},
  stringstyle=\color{ink},
}
\lstset{style=tandem}
\lstdefinelanguage{shellx}{morekeywords={export,set,bash,cd,pip,curl,hf,python,python3,docker,sudo,tandem-app,git,vllm,flock,printf,echo,grep,security,rm,cp,cat,sqlite3,npm,npx},
  sensitive=true, morecomment=[l]{\#}, morestring=[b]"}
\lstnewenvironment{shell}{\lstset{language=shellx}}{}
\lstnewenvironment{plain}{\lstset{language={}}}{}
\lstnewenvironment{lines}{\lstset{language={}}}{}
\lstnewenvironment{jsonc}{\lstset{language={},morestring=[b]"}}{}

% ---------------------------------------------------------------------------- figures
\usetikzlibrary{positioning,arrows.meta,shapes.geometric,shapes.misc,fit,backgrounds,calc,
  decorations.pathreplacing,patterns,matrix,chains}
\pgfplotsset{compat=1.18}
\tikzset{
  >={Stealth[length=2.2mm]},
  box/.style={draw=cobalt, fill=cobaltlight, rounded corners=2pt, align=center,
    font=\sffamily\small, inner sep=4pt, minimum height=8mm},
  hot/.style={draw=signal, fill=signallight, rounded corners=2pt, align=center,
    font=\sffamily\small, inner sep=4pt, minimum height=8mm},
  quiet/.style={draw=rule, fill=paper, rounded corners=2pt, align=center,
    font=\sffamily\small, inner sep=4pt, minimum height=8mm},
  lbl/.style={font=\sffamily\scriptsize, text=slate, align=center},
  flow/.style={->, thick, draw=ink},
  dflow/.style={->, thick, draw=slate, dashed},
  group/.style={draw=rule, rounded corners=4pt, inner sep=6pt},
  % the instrument-readout pieces of the dashboard figures
  panel/.style={draw=rule, fill=white, rounded corners=3pt, inner sep=5pt},
  readout/.style={font=\sffamily\bfseries\Large, text=ink},
  unit/.style={font=\sffamily\scriptsize, text=slate},
}
% ---------------------------------------------------------------------------- the mark
% The TandemLLM mark (Scope: a dial, the StairCut step trace inside it, the orange now-rate
% breaking out on the right), drawn in TikZ from the geometry of docs/media/logo/mark.svg so the
% build needs no SVG converter. Its own colours, not the book's: brandcobalt and brandsignal.
\definecolor{brandcobalt}{HTML}{3D63FF}
\definecolor{brandsignal}{HTML}{FF6B1F}
% \tandemmark[<unit>]{<dial and trace>}{<now-rate>}: the 64-unit grid at one unit = <unit>;
% the box is the mark's ink, units 2 to 62.
\newcommand{\tandemmark}[3][.25mm]{%
  \begin{tikzpicture}[x=#1, y=-#1]
    \useasboundingbox (2,2) rectangle (62,62);
    \draw[draw=#2, line width=\dimexpr#1*4\relax] (32,32) circle[radius=26];
    \draw[draw=#2, line width=\dimexpr#1*6\relax, line cap=butt, line join=round,
      rounded corners=\dimexpr#1*3/2\relax]
      (8,42) -- (22,42) -- (22,30) -- (38,30) -- (38,18) -- (41,18);
    \draw[draw=#3, line width=\dimexpr#1*6\relax, line cap=round] (40,18) -- (60,18);
  \end{tikzpicture}}
% built once, outside any other picture (a tikzpicture inside the cover's overlay would inherit it)
\newsavebox{\tandemcoverbox}
\newsavebox{\tandempartbox}
\AtBeginDocument{%
  \sbox{\tandemcoverbox}{\tandemmark[.32mm]{brandcobalt}{brandsignal}}%
  \sbox{\tandempartbox}{\tandemmark[.14mm]{slate}{slate}}}
% a small mono mark above "Part I" on each part page
\renewcommand*{\partformat}{\usebox{\tandempartbox}\\[2.2ex]\partname~\thepart\autodot}

\pgfplotsset{
  tandem/.style={
    width=\linewidth, height=6.2cm,
    axis lines*=left, grid=major, grid style={draw=rule!60},
    tick label style={font=\sffamily\footnotesize}, label style={font=\sffamily\small},
    legend style={font=\sffamily\footnotesize, draw=none, fill=none},
    every axis plot/.append style={thick},
  },
  % a single stacked bar: no y axis, no empty-range warning
  onebar/.style={ytick=\empty, axis y line=none, ymin=-1, ymax=1, enlarge y limits=false},
}

% ---------------------------------------------------------------------------- index and links
\makeindex[intoc, columns=2, title=Index, options=-s tandem.ist]
\usepackage[colorlinks=true,linkcolor=cobalt,urlcolor=cobalt,citecolor=cobalt,
  bookmarksnumbered=true,bookmarksopen=true,bookmarksopenlevel=0,unicode=true,
  pdftitle={TandemLLM Manual},pdfauthor={The TandemLLM authors},
  pdfsubject={An inference engine for Qwen3.8-27B on one DGX Spark, and its conductor},
  pdfkeywords={speculative decoding, inference, NVFP4, DGX Spark, StairCut}]{hyperref}
\usepackage{bookmark}
\newcommand{\version}{0.16.0}
\newcommand{\manualdate}{October 2026}

% "see" entries of the index
\index{lookup store|see{suffix store}}
\index{n-gram drafter|see{lookup drafter}}
\index{KV cache|see{resident prefix, state store}}
\index{queue|see{admission}}


\begin{document}

\frontmatter
% the cover's logo hook (chapters/titlepage.tex): the mark, top right on the ink page
\newcommand{\coverlogo}{\usebox{\tandemcoverbox}}
% The cover, drawn with TikZ on a full-bleed ink page.
% The hero is the verify-cost staircase from data/verify-cost.csv, the measured table StairCut
% prices its cuts on: the one chart the whole engine is built around.
%
% Logo hook: define \coverlogo before this file is read and it is placed in the top-right
% corner; the title block does not move. tandemllm-manual.tex sets it to the mark (\tandemmark,
% preamble.tex); without a definition the corner stays empty.
\providecommand{\coverlogo}{}
\begin{titlepage}
\thispagestyle{empty}
\begin{tikzpicture}[remember picture, overlay]
  \fill[ink] (current page.south west) rectangle (current page.north east);
  % a faint instrument grid behind the chart
  \begin{scope}
    \clip ([yshift=-11.2cm]current page.north west) rectangle ([yshift=7.2cm]current page.south east);
    \foreach \y in {0,1,...,11} \draw[ink!94!white, line width=.3pt]
      ([yshift=-11.2cm-\y cm]current page.north west) -- ([yshift=-11.2cm-\y cm]current page.north east);
  \end{scope}
  % the staircase: tree verify plus commit, ms, by rows, at five context lengths (measured)
  \pgfplotstableread[col sep=comma, comment chars={\#}]{data/verify-cost.csv}\covercost
  \begin{axis}[at={($(current page.west)+(2.4cm,-0.4cm)$)}, anchor=west,
      width=16.8cm, height=10.4cm, axis lines=none, clip=false,
      xmin=0, xmax=34, ymin=68, ymax=116,
      every axis plot/.append style={line cap=round, line join=round}]
    % glow, then line, for each context
    \addplot[color=sky!45!ink, line width=4pt, opacity=.22] table[x=rows, y=tree1024] {\covercost};
    \addplot[color=sky!45!ink, line width=1.1pt] table[x=rows, y=tree1024] {\covercost};
    \addplot[color=sky!60!ink, line width=4pt, opacity=.22] table[x=rows, y=tree4096] {\covercost};
    \addplot[color=sky!60!ink, line width=1.1pt] table[x=rows, y=tree4096] {\covercost};
    \addplot[color=sky!78!ink, line width=4pt, opacity=.22] table[x=rows, y=tree8192] {\covercost};
    \addplot[color=sky!78!ink, line width=1.1pt] table[x=rows, y=tree8192] {\covercost};
    \addplot[color=sky, line width=4.5pt, opacity=.25] table[x=rows, y=tree16384] {\covercost};
    \addplot[color=sky, line width=1.2pt] table[x=rows, y=tree16384] {\covercost};
    \addplot[color=signal, line width=6pt, opacity=.28] table[x=rows, y=tree32768] {\covercost};
    \addplot[color=signal, line width=1.5pt, mark=*, mark size=1.5pt] table[x=rows, y=tree32768] {\covercost};
    % the step
    \draw[sky!50!ink, dashed, line width=.5pt] (axis cs:16.5,70) -- (axis cs:16.5,114);
    \node[font=\sffamily\scriptsize, text=sky!70!ink, anchor=south west, align=left] at (axis cs:16.9,108)
      {17 rows need a second\\tensor-core tile};
    % context labels at the right end of each line
    \node[font=\sffamily\scriptsize, text=sky!45!ink, anchor=west] at (axis cs:32.6,87.76) {1k context};
    \node[font=\sffamily\scriptsize, text=sky!78!ink, anchor=west] at (axis cs:32.6,92.05) {8k};
    \node[font=\sffamily\scriptsize, text=sky, anchor=west] at (axis cs:32.6,99.60) {16k};
    \node[font=\sffamily\scriptsize, text=signal, anchor=west] at (axis cs:32.6,111.56) {32k};
    % a bare scale
    \node[font=\sffamily\scriptsize, text=slate, anchor=north] at (axis cs:8,69.5) {8};
    \node[font=\sffamily\scriptsize, text=slate, anchor=north] at (axis cs:16,69.5) {16};
    \node[font=\sffamily\scriptsize, text=slate, anchor=north] at (axis cs:24,69.5) {24};
    \node[font=\sffamily\scriptsize, text=slate, anchor=north] at (axis cs:32,69.5) {32};
    \node[font=\sffamily\scriptsize, text=slate, anchor=east] at (axis cs:-0.6,75) {75};
    \node[font=\sffamily\scriptsize, text=slate, anchor=east] at (axis cs:-0.6,90) {90};
    \node[font=\sffamily\scriptsize, text=slate, anchor=east] at (axis cs:-0.6,105) {105};
    \node[font=\sffamily\scriptsize, text=slate, anchor=north west] at (axis cs:0,66.2)
      {rows in the verified tree \quad $\cdot$ \quad verify + commit, ms, published NVFP4 weights (measured)};
  \end{axis}
  % the title block
  \node[anchor=north west, text width=15cm, align=left, text=white]
    at ([xshift=2.4cm,yshift=-2.5cm]current page.north west)
    {{\sffamily\fontsize{46}{50}\selectfont\bfseries TandemLLM}\\[5mm]
     {\sffamily\fontsize{17}{20}\selectfont\color{sky} The manual}\\[9mm]
     {\fontsize{12.5}{16}\selectfont\color{white!85!ink} An environment for the agents you run on your own machine:\\
      several models working on one request, caches that survive a restart,\\
      memory notes, parallel requests, and one dashboard for this engine,\\
      other engines and cloud APIs}};
  % the logo hook, top right: the mark
  \node[anchor=north east] at ([xshift=-2.4cm,yshift=-2.5cm]current page.north east) {\coverlogo};
  % the foot: version, date, where the code is
  \draw[signal, line width=1.2pt] ([xshift=2.4cm,yshift=4.2cm]current page.south west) -- ++(3cm,0);
  \node[anchor=north west, text width=15cm, align=left, text=white]
    at ([xshift=2.4cm,yshift=3.7cm]current page.south west)
    {{\sffamily\large Version \version}\quad{\sffamily\color{sky}\manualdate}\\[3mm]
     {\small\color{white!75!ink} The text of this manual is CC BY 4.0; the code it describes is AGPL-3.0-only or a commercial license.}\\[1.5mm]
     {\ttfamily\small\color{sky} github.com/TandemLLM/TandemEngine}};
\end{tikzpicture}
\end{titlepage}

\input{chapters/about}

{\hypersetup{linkcolor=ink}\tableofcontents\listoffigures\listoftables}

\mainmatter

\partintro{%
Two chapters for every reader. Chapter~\ref{ch:intro} says what TandemLLM is, names the machine and the model every number in this book comes from, says what the repository contains and what runs today, and states the two rules every change has to pass. Chapter~\ref{ch:tour} follows one chat request through all of it, from the client to the stream going back, and names the chapter that explains each step.}
\part{Introduction}
\chapter{What TandemLLM is}\label{ch:intro}

TandemLLM\index{TandemLLM} runs the models your agents need on your own hardware, behind one OpenAI-compatible address. A coding agent such as opencode or a chat client such as Open WebUI sends its request to that address. TandemLLM picks the model that answers, plans the turn, and watches the turn while it runs.

Several models work on one request. Qwen3.8-27B writes the answer. Two small drafters guess its next tokens and a lookup drafter offers text it has seen before, so one step of the big model commits 4.58 tokens on average on the benchmark row (\measured, the paper). More models can run beside them in the same process. Laya's decision models\index{Laya} answer typed questions such as \enquote{is this turn code or prose}. A Qwen3.5-2B narrator\index{narrator} writes one short line about what the big model is doing, and a Qwen3.5-0.8B shadow reader\index{shadow reader} keeps its own copy of each coding session. A residency planner refuses any model that would not fit in memory (Section~\ref{sec:residency}).

You get on your own machine what a cloud provider gives you. Clients use the OpenAI API with streaming and tool calls. Every request leaves one row in a usage ledger with its tokens and its timings. TandemApp, a separate project with its own manual, puts this engine and other backends (vLLM, SGLang, cloud APIs) behind one address.

An agent's next turn does not read its whole prompt again. The resident prefix keeps the last prompt's rows with a snapshot of the recurrent state every 1,024 tokens, and the end anchor keeps the end of the last answer. A parked session holds a second conversation while another one runs, and the state store keeps whole snapshots. Two caches live on the NVMe and survive a shutdown. The disk tier keeps prefixes that come back. The session spill writes live conversations at a stop and reads them back at the next start: a 41,587-token conversation came back in 0.94~s (\measured, the 0.4.0 series). Last, the persistent suffix store keeps the n-gram tables the lookup drafter copies from. Chapter~\ref{ch:caches} describes every cache.

With memory notes on, the model keeps notes as files, one fact per file and at most 400 words. Notes about the user and about each project last across sessions. Session notes hold what is true for one conversation only (TandemApp's manual, \enquote{Memory notes}).

The engine says what each request is doing right now under \code{/v1/dashboard/*} and counts it in \code{/metrics} (Chapter~\ref{ch:ref-metrics}). TandemApp's dashboard shows it, beside the other engines it watches.

Requests run in parallel. The shipped profile's 8 lanes decode up to eight requests in the same rounds (Section~\ref{sec:parallel-requests}), and the proxy relays any number of requests to other backends at once.

TandemLLM is not tied to one model. Its main model is Qwen3.8-27B, and a second, Qwen3.8-Flash-Next, runs on the same board in its place (Chapter~\ref{ch:flashnext}). The small models are Qwen3.5 and Laya checkpoints, the proxy serves any model behind an OpenAI-compatible API, and Chapter~\ref{ch:adding-model} lists the seams where a main model plugs in.

\section{The test environment}\label{sec:test-env}\index{test environment}\index{DGX Spark}\index{GB10}

Every number in this book comes from one machine with one main model: an NVIDIA DGX Spark running Qwen3.8-27B. The Spark has a GB10 GPU and 128~GB of memory that the CPU and the GPU share. That is the project's test hardware and its test model, and the project has tested nothing else so far. Its kernels' tile tables and StairCut's price tables come from that board. Much of the engine follows that model's mix of 48 Gated DeltaNet layers and 16 attention layers (Chapter~\ref{ch:architecture}).

The project is growing past both. Other CUDA GPUs run the engine with the one Spark-only kernel turned off; their speed has no measurement yet. The proxy and the dashboard run on any machine with Python 3.11. The roadmap (Chapter~\ref{ch:roadmap}) names the next models and boards. When this book says \enquote{the board} or \enquote{the model} without a name, it means the DGX Spark and Qwen3.8-27B.

\section{Two levers}

A decode step of Qwen3.8-27B reads about 27~GB of weights in the vendor's FP8 format. The DGX Spark reads its memory at 273~GB/s on paper and about 240~GB/s in practice, so every step costs at least 100~ms before any code runs (\arith). That is a ceiling of 10 tokens a second for one request. The engine is built to get under that floor while writing the model's own greedy text.

With the published NVFP4 weights and StairCut on, it writes 49.89 tokens a second for one request on the benchmark row, against 13.69 for plain greedy decoding on the same box (\measured, the StairCut paper, 50 requests each). The rest of this section says how.

The engine pulls two levers on the same cost.

The first lever cuts the bytes a step reads. The served weights store every projection in NVFP4\index{NVFP4}, four bits a weight with a scale per 16 weights, and the output head in FP8\index{FP8}. A step then reads 15.01~GB instead of 26.93~GB (\arith, from the checkpoint's tensor sizes). Chapter~\ref{ch:quantisation} shows how the 4-bit copies were made and what they cost in quality.

The second lever makes each step produce several tokens. Small \emph{drafters}\index{drafter} guess the next tokens, and the big model checks all the guesses in one forward pass. A guess the model agrees with is a free token. At the first wrong guess the model puts in its own token, so nothing a drafter proposes reaches the output unless the model would have written it. On the benchmark row a round takes 93.4~ms and commits 4.58 tokens (\measured, the paper). That is where the name comes from: two models working in tandem, one guessing and one checking.

The verify is cheap because of the memory floor. A forward pass over 16 tree rows reads the same weights as a pass over one row, so on the published weights at 1k tokens of context it takes 74.4~ms against 75.5~ms for 8 rows (\measured, \code{ops/stair-tables-nvfp4.json}). Extra guesses in the same pass are nearly free until the row count crosses a kernel's tile boundary. StairCut\index{StairCut}, the router that sizes each round's draft tree, is built on that staircase (Chapter~\ref{ch:speculation}).

\section{What is in the repository}

TandemLLM is several programs that share one repository and one release number.

\begin{description}[style=nextline]
\item[The engine] \code{engine/}, \code{tools/} (kernels) and \code{server/app.py}: the model, its kernels, the drafters and routers, the caches, and an OpenAI-compatible HTTP server built on Python's standard library. It needs a CUDA GPU (measured on the DGX Spark only), PyTorch, Triton and the weights. The small models run on its own code: the Qwen3.5 narrator and shadow reader on its Qwen forward, Laya's decision models in the same process.
\item[The conductor] \code{server/conductor/}: a layer in front of the model. It classifies each turn, sets the thinking effort, adds hidden notes and skills, watches the generation for loops and broken tool calls, keeps memory notes and shows the client what it does. It is off unless a configuration directory is given.
\item[The dashboard] \code{dashboard/}: a web page the server serves at \code{/dashboard/}. It shows what each request is doing now, a year of usage, the log, a playground, the conductor's settings, the backends and the machine.
\item[The proxy] \code{proxy/}: the same dashboard in front of other OpenAI-compatible engines and cloud APIs, routed by model id. It installs with \code{pip} and needs no GPU.
\item[The plugins] \code{plugins/}: two small packages that run inside vLLM and SGLang and report each request's queue, prefill and decode times to the proxy.
\item[The installer] \code{install.sh}: one command that sets up the served profile on the test board, and best effort on another CUDA GPU. With \code{--setup} it starts setup mode instead: the Setup Guide in the browser, or \code{ops/use-profile.sh} from a shell, then downloads and starts any weight profile (Section~\ref{sec:new-box}).
\end{description}

\section{What runs today}

The roadmap (Chapter~\ref{ch:roadmap}) widens each of these, one step at a time.

\begin{itemize}
\item Models: Qwen3.8-27B is the main model the engine serves, a hybrid of 48 Gated DeltaNet layers and 16 attention layers; Chapter~\ref{ch:architecture} explains why much of the engine follows that mix. Qwen3.8-Flash-Next, a mixture of experts on the same hybrid, is the second, served from turboderp's EXL3 packs one model at a time (Chapter~\ref{ch:flashnext}). Beside it run the two fine-tuned DFlash2 drafters, the lookup drafter and, when you turn them on, Laya, the narrator and the shadow reader. Behind the proxy, any model a backend serves.
\item Hardware: the DGX Spark is the only board with measured kernels and prices (Section~\ref{sec:test-env}). Other CUDA GPUs are best effort.
\item Requests: \code{--parallel N} gives the engine N sequences, and requests that decode at the same time share one draft pass and one verify pass a round (Section~\ref{sec:parallel-requests}). The shipped profile runs 8 lanes since 0.15.0 (SRV-707, \code{QSE_PARALLEL=8} in \code{ops/serve.env}), and every weight card has a speed row with 1, 2, 4 and 8 requests at once. With one sequence a request that finds the engine busy waits in a bounded queue, and a short request is served at a safe point inside the running one (the side lane, Section~\ref{sec:side-lane}). The proxy relays any number of requests to other backends at once (TandemApp's manual, \enquote{The conductor in front of any backend}).
\end{itemize}

\section{Two rules}

Every change to the engine keeps two rules, and the release gate checks both before anything ships.

The first rule is agreement with greedy decoding\index{greedy agreement}. For every greedy request the engine's output tokens equal what plain decoding on the same weights writes, one token at a time, up to the first near tie where the two best tokens are within one unit in the last place. There the batched verify can round the other way. The paper measured how often: of 30 prompts, 17 outputs equal plain greedy decoding to the end of the answer, and 13 diverge at a near tie with a gap of 0.04 to 0.76 ulp (\measured). Chapter~\ref{ch:exactness} has the full statement and the checks.

The second rule is that no measured number gets worse\index{never worse}. A candidate runs against the current release in the same hour on the same box, alternated, and ships only if no figure resolves worse: not the mean, the median, the 90th percentile or the maximum, not the time to first token, and not tokens or milliseconds per round. Chapter~\ref{ch:measurement} describes the method and the traps it avoids.

\section{The headline numbers}

Table~\ref{tab:headline} collects the numbers that describe the served engine on the test board. Each comes from the StairCut paper's measurements on the published NVFP4 weights, with the lookup drafter's persistent store off.

\begin{table}[htbp]
\centering\small
\caption{The served engine's headline numbers. Source: the StairCut paper, published NVFP4 weights, \measured.}\label{tab:headline}
\begin{tabularx}{\linewidth}{@{}Yr@{}}
\toprule
\hd{what} & \hd{value} \\
\midrule
single request, StairCut, benchmark row (mean of 50) & 49.89 tok/s \\
the same drafter at a fixed block of 8 & 47.13 tok/s \\
the same drafter at a fixed block of 16 & 44.80 tok/s \\
plain greedy decoding, speculation off & 13.69 tok/s \\
teacher-forced bench, 25 workloads, StairCut & 90.45 tok/s \\
one round on the benchmark row & 93.4 ms, 4.58 tokens \\
time to first token, cold, at 32k / 64k / 128k tokens & 44.4 / 93.7 / 208.3 s \\
the same prompt sent again, resumed from the cache & 0.3 / 0.6 / 1.1 s \\
held-out loss against Qwen's BF16 release, prose / code & +0.0221 / +0.0426 nats \\
\bottomrule
\end{tabularx}
\end{table}

\begin{figure}[htbp]
\centering
\begin{tikzpicture}
\begin{axis}[tandem, xbar, height=4.6cm, width=.86\linewidth, bar width=8pt,
    xmin=0, xmax=62, xlabel={tokens a second, one request (mean of 50)},
    symbolic y coords={greedy,block16,block8,staircut}, ytick={greedy,block16,block8,staircut}, bar shift=0pt,
    yticklabels={plain greedy decoding, fixed block 16, fixed block 8, StairCut},
    nodes near coords, nodes near coords style={font=\sffamily\scriptsize, /pgf/number format/fixed, /pgf/number format/precision=2, /pgf/number format/fixed zerofill},
    enlarge y limits=0.2, grid=none, axis lines*=left]
  \addplot[fill=cobaltlight, draw=cobalt] coordinates {(13.69,greedy) (44.80,block16) (47.13,block8)};
  \addplot[fill=signal, draw=signal] coordinates {(49.89,staircut)};
\end{axis}
\end{tikzpicture}
\caption{Single-request speed on the benchmark row \texttt{serve-single-i256-o256-v1}: StairCut against the same drafter at fixed block widths and against plain greedy decoding. \measured: the StairCut paper, published NVFP4 weights, persistent store off.}\label{fig:headline}
\end{figure}

These are single-request numbers. Parallel requests on the engine are measured in Section~\ref{sec:parallel-requests}. Long prompts are measured separately in Chapter~\ref{ch:caches}. No task benchmark suite has been run, so this manual makes no claim about task accuracy.

\section{The paper}\label{sec:paper}\index{StairCut paper}\index{Zenodo}

Khaled Bakeer's paper describes the engine's decoding method:

\begin{quote}
Khaled Bakeer. \emph{StairCut: sizing speculative draft trees on a measured verification-cost staircase.} Zenodo, 29 September 2026. \href{https://doi.org/10.5281/zenodo.23045275}{doi:10.5281/zenodo.23045275}.
\end{quote}

The paper measured the release tagged \code{v0.2.0-staircut} (commit \code{c111fa7}) with the published NVFP4 weights on the test board. It backs these parts of this book:

\begin{itemize}
\item the draft tree merged from the block drafter and the lookup drafter, and StairCut, which cuts it each round on the measured verify-cost staircase (Chapter~\ref{ch:speculation});
\item the commit of an accepted path through the Gated DeltaNet layers in closed form, and the NVFP4 kernels with a fixed reduction order (Chapter~\ref{ch:kernels});
\item the NVFP4 quantisation from Qwen's BF16 release and its held-out loss (Chapter~\ref{ch:quantisation});
\item the speeds of Table~\ref{tab:headline}: 49.89~tok/s for StairCut, 47.13 and 44.80 at fixed blocks of 8 and 16, 13.69 for plain greedy decoding, 50 requests each on the benchmark row, and the teacher-forced bench of 25 workloads;
\item the agreement with plain greedy decoding: of 30 outputs, 17 identical to the end of the answer and 13 that part at a near tie (Chapter~\ref{ch:exactness});
\item the time to first token of a 32k to 128k-token prompt from cold, and again when it resumes from the anchors (Chapter~\ref{ch:caches}).
\end{itemize}

What came after that release is not in the paper: image input, the conductor, the caches on the NVMe, memory notes, parallel requests, the proxy and the plugins. Their numbers come from the release series, and the text labels them that way. The release gate holds every later release to the paper's single-request numbers and to the agreement rule. \code{CITATION.cff} in the repository carries the citation; in BibTeX:

\begin{plain}
@misc{bakeer2026staircut,
  author    = {Bakeer, Khaled},
  title     = {StairCut: sizing speculative draft trees on a measured
               verification-cost staircase},
  year      = {2026},
  publisher = {Zenodo},
  doi       = {10.5281/zenodo.23045275},
  url       = {https://doi.org/10.5281/zenodo.23045275}
}
\end{plain}

\section{Credits}

Qwen made the model weights and their reference implementation. The base block drafter is DFlash2 for Qwen3.8-27B, released by its vendor under Apache-2.0. The engine runs on PyTorch and Triton over CUDA, and its own kernels are written in Triton and CUDA. The lookup corpus that the installer builds comes from public, permissively licensed text. The decision models the conductor can ask are Laya's published checkpoints (Apache-2.0).

\chapter{One request, end to end}\label{ch:tour}

This chapter follows one chat completion from a coding agent through every part of TandemLLM. Each step names the chapter that explains it. Figure~\ref{fig:request-path} is the map.

\begin{figure}[htbp]
\centering
\input{figures/request-path}
\caption{The request path. A client sends an OpenAI request; the proxy routes it by model id; the conductor's router plans the turn; admission decides when it runs; the caches decide where the prefill starts; then drafters and the target model take turns, round after round, while the conductor's monitor watches the output at fixed safe points. Requests for other backends leave at the proxy.}\label{fig:request-path}
\end{figure}

\section{The client sends a request}

The client is any program that speaks the OpenAI chat completions API: the OpenAI SDK\index{OpenAI SDK}, opencode\index{opencode}, the Hermes agent\index{Hermes agent}, Open WebUI\index{Open WebUI} or \code{curl}. It sends \code{POST /v1/chat/completions} with its messages, its tools and its options, usually streamed. The engine's server reads every OpenAI field it can honour, accepts the ones that have no effect here, and refuses the rest with a 400 that names the field (Section~\ref{sec:api-fields}).

\section{The proxy routes it}

When the server runs with \code{--backends}\index{backends@\texttt{-\kern0pt-backends}}, or when the standalone proxy (\code{tandem-app}) is in front, the request is routed by its \code{model} field. A model of another backend (vLLM, SGLang, a cloud API) is relayed there: the answer's bytes go back to the client unchanged, and the proxy writes one ledger row with the tokens, the time to first token, the decode rate and the cost. A request for \code{TandemLLM} stays in the engine's server. Chapter~\ref{ch:install-app} sets this up and Chapter~\ref{ch:harness-proxy} covers the conductor in front of other backends.

\section{The conductor plans the turn}

With a conductor configuration the server serves one model name, \code{TandemLLM}\index{TandemLLM (model name)}: the weights with a conductor in front. Before the request queues, the conductor resolves a \term{profile} (by header, by client, or the default), and its router decides what kind of task the turn is. A tool-result turn keeps the session's class; only a new user message asks a classifier, and the default classifier is Qwen itself under a choice grammar. The class sets the thinking budget and adds a hidden note at the end of the last message: the effort sentence, rules and \emph{packages}\index{package} (skill texts). Later turns replay that note at the same token position, so the cached prefix keeps matching. Chapter~\ref{ch:conductor} explains the router and Chapter~\ref{ch:packages} the packages.

A client's side request, such as opencode's title call, is recognised by its shape and served small: no classifier, no note, at most 64 output tokens for a title.

\section{Admission decides when it runs}

The engine holds one sequence unless \code{--parallel N} gives it more (Section~\ref{sec:parallel-requests}). A request that finds it busy waits in a bounded queue (8 by default). A streamed request that waits more than 10~s gets its stream opened and an SSE comment every 10~s, with an estimate of the wait. A short request (at most 2,048 prompt tokens and 64 output tokens, thinking off) does not wait at all: the running generation serves it at its next safe point, between two decode blocks or two prefill chunks, and gets its state back byte for byte. Before every prefill a memory gate checks that the prompt's workspace fits under the declared peak, giving caches back cheapest first. Chapter~\ref{ch:serving} covers the queue, the side lane and the memory gate.

\section{The caches decide where the prefill starts}

A prefill is the most expensive thing the engine does: 208.3~s for a 128k-token prompt from cold (\measured, the paper). Before it prefills, the engine looks for state it already has. In an agent session each turn's prompt is the previous prompt plus the answer plus a tool result, and the \term{resident prefix} still holds the previous turn's rows in the engine's own buffers, with recurrent-state \emph{anchors}\index{anchor} every 1,024 tokens and an \term{end anchor} at the end of the last generation. The next turn resumes there and forwards only what is new. Further down are a state store of snapshots, a disk tier on the NVMe for prefixes that come back, and a session spill that keeps whole conversations through a restart. Chapter~\ref{ch:caches} describes them all.

\section{The decode loop}

Then rounds run until the answer ends. In each round the block drafter proposes up to 15 tokens and the lookup drafter proposes continuations it has seen before; StairCut merges them into one tree and cuts it to the size that maximises expected tokens per millisecond on the measured verify cost. One forward pass of the target model verifies the whole tree through all 64 layers. The engine walks the tree with the model's own tokens, keeps the longest agreeing path plus one token of the model's own, and commits that path's KV rows and recurrent state. Chapter~\ref{ch:speculation} explains the round and Chapter~\ref{ch:kernels} the kernels that make the verify cheap.

\section{The monitor watches the output}

While the model writes, the conductor's \term{session monitor} reads the output at four safe points: before the prefill, after every committed block, at the end token, and when the answer's tool calls are complete. It can close a reasoning block that loops, continue a turn that ended empty, repair or regenerate a malformed tool call, and stop a tool call whose arguments repeat. Everything it adds is hidden from the client and never counted in \code{usage}. Chapter~\ref{ch:monitor} lists its tasks.

\section{The stream goes back}

Committed tokens stream to the client as they are verified, several at a time. A tool call streams as argument deltas while the model writes it, with values typed by the request's JSON schema. With a profile's status lines on, the client also sees one short line for each step the conductor takes (\enquote{› reading 20,745 tokens, about 18 s left}), which the engine strips again from the next prompt (Chapter~\ref{ch:lines}).

\section{After the request}

The request leaves one log line (counts, finish reason, rate, never text) and one row in the usage ledger. The conductor writes its decisions to a journal, a transcript block for a request with tools, and, when memory notes are on, any \code{<tandem_memory>} block the model wrote. When the session goes quiet, a background writer may write the session's facts as notes (Chapter~\ref{ch:memory}). The dashboard shows all of it live (Chapter~\ref{ch:dashboard}).


\partintro{%
How to install and run it. Chapter~\ref{ch:install-engine} starts with the setups that work today, then installs the engine with one command, says who can reach it, and adds the small models and other backends. TandemApp's manual installs the dashboard and proxy on any machine and puts other engines and cloud APIs behind it. Chapter~\ref{ch:operations} is the operator's chapter: the memory rules, the service scripts, the watchdog, how to hold the board for a measurement, where the logs are, and how to upgrade and roll back. Chapter~\ref{ch:community} runs weights other people published, and says how to list your own.}
\part{Installing and running}
\chapter{Installing TandemLLM}\label{ch:install-engine}

This chapter installs the engine. It runs the models and needs a CUDA GPU; its test board is the DGX Spark (Section~\ref{sec:test-env}). The dashboard, the Chat tab and the proxy are TandemApp, a separate project with its own manual: it needs only Python, runs on any machine, and talks to the engine over HTTP. Find your setup in Table~\ref{tab:setups}, then follow the sections it names.

\section{Setups, and which ones work}\label{sec:setups}\index{setups}

The engine runs alone, or behind TandemApp together with other engines and cloud APIs. Table~\ref{tab:setups} says which combinations work today, and how far each one is tested.

\begin{table}[htbp]
\centering\small
\caption{Setups and their state in this release. \emph{Live}: run on the test board in the 0.4.0 series. \emph{Built}: the code supports it and its tests pass, but no live run is recorded.}\label{tab:setups}
\begin{tabularx}{\linewidth}{@{}L{.27\linewidth}YL{.13\linewidth}@{}}
\toprule
\hd{setup} & \hd{how} & \hd{state} \\
\midrule
the engine alone & \code{install.sh} on the board (this chapter) & live \\
the engine and cloud APIs, one address & \code{tandem-app} on the board or another machine, with the engine and the cloud API as backends (TandemApp's manual) & live, with DeepSeek \\
the engine and other engines on other machines & \code{tandem-app} as above, the other engines as \code{vllm}, \code{sglang} or \code{generic} backends (TandemApp's manual) & built; vLLM and SGLang live behind the proxy on their own \\
the engine and a second GPU engine on one larger machine & each engine on its own GPU, \code{tandem-app launch} for the second & built \\
the engine and a second GPU engine on one DGX Spark & the engine plans its memory as if it owns the board (up to 77~GiB during a prefill, plus its small models); a board out of memory stops answering until it is power cycled & not supported \\
\bottomrule
\end{tabularx}
\end{table}

\begin{figure}[htbp]
\centering
% Which install path: a decision tree over the three paths of Table tab:install-paths.
\begin{tikzpicture}[node distance=6mm and 8mm,
    q/.style={draw=slate, fill=paper, shape=rounded rectangle, font=\sffamily\small,
      align=center, inner sep=4pt, minimum height=9mm},
    path/.style={draw=cobalt, fill=cobaltlight, rounded corners=2pt, text width=50mm, align=left,
      font=\sffamily\small, inner sep=5pt}]
  \node[q] (spark) {a machine with a CUDA GPU?\\{\scriptsize tested: DGX Spark}};
  \node[q, below right=8mm and 6mm of spark] (others) {other engines or\\cloud APIs too?};
  \node[q, below left=8mm and 6mm of spark] (py) {Python 3.11 or newer,\\or docker?};
  \node[path, below=9mm of others] (a) {\textbf{Path A: the engine}\\[2pt]
    {\normalfont\footnotesize the engine with its models at \texttt{:8000}, and the conductor in chat mode}\\[3pt]
    {\ttfamily\scriptsize curl -fsSL .../install.sh | bash}\\[2pt]
    {\footnotesize\color{slate} Chapter~\ref{ch:install-engine}}};
  \node[path, below=9mm of py] (b) {\textbf{Path B: TandemApp, the dashboard and proxy}\\[2pt]
    {\normalfont\footnotesize one address, one ledger and one dashboard for backends you already run; no GPU, no torch. In docker: TandemApp's \texttt{packaging/Dockerfile.app}}\\[3pt]
    {\ttfamily\scriptsize pip install '.[host]' \&\& tandem-app}\\[2pt]
    {\footnotesize\color{slate} TandemApp's manual}};
  \node[path, draw=signal, fill=signallight, text width=70mm, below=10mm of $(a.south)!0.5!(b.south)$] (c) {\textbf{Path C: several engines and cloud APIs behind one address}\\[2pt]
    {\normalfont\footnotesize A or B, plus the backends. The engine's own server relays (\texttt{QSE\_BACKENDS}), or \texttt{tandem-app} lists the engine as one backend among others. TandemApp's manual}};
  \draw[flow] (spark) -- node[right, lbl, pos=.55] {yes} (others);
  \draw[flow] (spark) -- node[left, lbl, pos=.55] {no} (py);
  \draw[flow] (others) -- node[right, lbl] {no} (a);
  \draw[flow] (py) -- node[left, lbl] {yes} (b);
  \draw[flow] (others.east) -- ++(8mm,0) |- node[pos=.25, right, lbl] {yes} (c.east);
  \draw[dflow] (a.south) -- node[right, lbl] {add backends} (a.south |- c.north);
  \draw[dflow] (b.south) -- node[left, lbl] {add backends} (b.south |- c.north);
\end{tikzpicture}

\caption[Which install path]{Which install path. A machine with the GPU gets the engine (path A); any other machine gets TandemApp, the dashboard and proxy (path B, TandemApp's manual); either grows into path C by adding backends behind TandemApp, which lists the engine as one backend among others (Section~\ref{sec:engine-backends}).}\label{fig:install-paths}
\end{figure}

\section{What the engine needs}\index{requirements}

The engine is tested on one NVIDIA DGX Spark, a GB10 GPU (\code{sm_121}) with 128~GB of memory that the CPU and the GPU share, running Linux on aarch64. On another CUDA GPU the installer turns off the one kernel built for the GB10, and the engine runs best effort. The installer checks four things before it changes anything: an NVIDIA GPU and driver, the CUDA compiler (one kernel builds on its first use), about 70~GB of free disk, and 64~GiB of free memory at the start. With \code{--setup} it downloads only Wren's model, so it asks for about 5~GiB of disk beside the Python environment and 8~GiB of free memory; \code{use-profile.sh} checks the disk and the memory a profile needs before it downloads or starts it.

The software is pinned. The engine imports PyTorch 2.13 built for CUDA 13.0, torchvision 0.28 (the CUDA 13.0 build that matches that torch), Triton 3.7, transformers 5.12, Pillow 11.3, safetensors and numpy, and nothing else. Python 3.11 is the tested version; under 3.12 one grammar test fails. The installer's lock lists every package with its version, for aarch64 and x86\_64.

Two of those packages are there for images. Pillow decodes a request's images, and torchvision is the backend of the image processor that the checkpoint ships. The engine pins that backend (\code{--image-backend torchvision}) and refuses to start without it, because transformers would otherwise fall back to its Pillow backend, whose patches differ by up to two 8-bit levels after a resize (\measured, 0.4~\% of the values of one resized image). Every cache keys on those patches, so the backend has to be fixed. \code{--image-backend pil} is the explicit choice without torchvision, and \code{--vision off} serves text only.

Run one engine per board\index{one engine per board}. Two engines loading side by side run the board out of memory, and a DGX Spark in that state stops answering until it is power cycled. Section~\ref{sec:one-engine} explains how the scripts enforce it.

\section{A new box: setup mode, then the Setup Guide}\label{sec:new-box}\index{setup mode}\index{use-profile.sh@\texttt{use-profile.sh}}

Since 0.8.1 a new DGX Spark needs two commands and no Hugging Face token. Every file a profile reads is public on Hugging Face: our own packs, the vendors' releases, and turboderp's Flash-Next packs (Chapter~\ref{ch:flashnext}).

\begin{shell}
curl -fsSLO https://raw.githubusercontent.com/0xBakeer/TandemLLM/v0.15.4/install.sh
bash install.sh --setup --yes
\end{shell}

The first line fetches the installer from the release tag. With \code{--setup} the installer puts the code and the Python environment in place and downloads only Wren's model (Qwen3.5-2B, 4.3~GB), the setup guide's assistant, and the two Wikipedia shards of the lookup corpus (1.2~GB). The corpus itself needs the model's tokenizer, so \code{ops/use-profile.sh} builds it once the model is on the disk. It writes the configuration with no weights and starts setup mode on \code{127.0.0.1:8000}: the Setup Guide's API (\code{/v1/setup/*}), Wren, and no model loaded. The engine serves no page.

Then pick a model with the Setup Guide. The guide shows each model with its weight packs, checks that the choice fits in memory, downloads it and starts the engine; a new browser opens its simple form, three pages: the model, its weights, and one page of folded settings. The switch at the top brings back all eight steps, through the model, the weights, the tools and the parallel slots. \code{docs/setup.md} describes the page.

\begin{itemize}
\item In a terminal on the box, run \code{~/TandemLLM/bin/tandem setup}: the same guide, step by step.
\item In a browser, use TandemApp's Setup page and point it at the engine, \texttt{http://127.0.0.1:8000} on the box itself.
\item From another computer, open an SSH tunnel to the box first, \code{ssh -L 8000:127.0.0.1:8000 <user>@<box>}, and use the same address on that computer.
\item On your network, install with \code{--host 0.0.0.0} and use \texttt{http://<box>:8000} from any computer that reaches the box. Until an admin password is set, anyone who reaches the port can run the setup.
\end{itemize}

The installer prints the addresses when it finishes, and while the engine is in setup mode TandemApp's page opens on the Setup Guide by itself.

\subsection{Advanced: from the command line}\label{sec:new-box-cli}

The Setup Guide runs \code{ops/use-profile.sh} for you. To pick the weights from a shell instead, run the copy in \code{~/TandemLLM/run/ops/}, next to the \code{serve.env} the installer wrote:

\begin{shell}
bash ~/TandemLLM/run/ops/use-profile.sh --list
bash ~/TandemLLM/run/ops/use-profile.sh <profile> --dry-run
bash ~/TandemLLM/run/ops/use-profile.sh <profile>
\end{shell}

\code{--list} prints every profile and marks the active one. With \code{--dry-run} the script prints every file's URL and size and changes nothing. Without it, the script downloads what the profile needs and checks every weight file against the sha256 the hub lists. It then leaves setup mode, starts the engine on port 8000 and sends one smoke request. Since 0.8.2, \code{community/<id>} runs a pack someone else published (Chapter~\ref{ch:community}).

\begin{table}[htbp]
\centering\small
\caption{What two of the profiles download on a new box. Free disk: the download plus the 60~GB the script keeps free after it.}\label{tab:new-box}
\begin{tabularx}{\linewidth}{@{}L{.17\linewidth}YL{.14\linewidth}L{.1\linewidth}@{}}
\toprule
\hd{profile} & \hd{downloads} & \hd{size} & \hd{free disk} \\
\midrule
\code{flashnext-3.05} & the pack's branch of \code{turboderp/Qwen3.8-Flash-Next-exl3}; the drafts' corrector (\code{0xBakeer/TandemLLM-Qwen3.8-Flash-Next-EXL3-3.05-corrector}, into \code{~/TandemLLM/src/ops/correctors/qwen3.8-flash-next}) & 85.1 + 0.035~GB & 145~GB \\
\code{exl3-4.5} & \code{0xBakeer/TandemLLM-Qwen3.8-27B-EXL3-4.5}; the K3 drafters and the K4 head (\code{...-EXL3-drafters}); the BF16 drafters (\code{...-DFlash2-b8}, \code{-b16}, into \code{~/TandemLLM/src/train}); the base checkpoint \code{Qwen/Qwen3.8-27B-FP8}, without \code{mtp.safetensors} & 18.4 + 1.9 + 7.7 + 30.4~GB & 118~GB \\
\bottomrule
\end{tabularx}
\end{table}

A folder appears only when all its files are there. If a download stops, run the same command again: it keeps the files that finished and continues. The corrector file of \code{flashnext-3.05} is checked against the sha256 the profile pins. If it does not download, the script prints a warning and goes on, and Flash-Next drafts without it: the same text, fewer drafts accepted. The engine then needs about 80~GiB of free memory for a 27B profile and 70~GiB for \code{flashnext-3.05}. If the profile does not come up, the script brings setup mode back.

The installer checks out the tag \code{v0.15.4} from the repository \code{--repo} names, the public one by default. Where the tag is missing, \code{curl} fails with a 404, and the installer stops at its clone step with \enquote{no tag, branch or commit 'v0.15.4'}. \code{--ref} picks another release, and \code{--repo} another repository.

\section{Updating a box in setup mode}\label{sec:update-setup}\index{upgrade}

A box with a running profile counts as configured, and \code{--setup} alone refuses to touch it. To update, add \code{--force-setup}: the installer stops the engine, checks out the new tag and comes back in setup mode. Then switch to the profile again; its files are on disk already, so nothing downloads twice.

\begin{shell}
bash ~/TandemLLM/src/install.sh --setup --force-setup --yes --ref <new tag>
bash ~/TandemLLM/run/ops/use-profile.sh flashnext-3.05
\end{shell}

\section{The installer}\index{install.sh}

Without \code{--setup}, one command sets up the NVFP4 profile, StairCut included:

\begin{shell}
curl -fsSL https://raw.githubusercontent.com/0xBakeer/TandemLLM/v0.15.4/install.sh | bash
\end{shell}

The installer works in seven steps, and each step is skipped when it is already done, so running it again updates an install and resumes a download that stopped.

\begin{enumerate}
\item Checks. The GPU needs its driver and the CUDA compiler, the disk and the memory need room, and Hugging Face must answer.
\item The engine: a clone of the repository \code{--repo} names into \code{~/TandemLLM/src} at the pinned release, \code{v0.15.4}.
\item A Python venv in \code{~/TandemLLM/venv} with the pinned packages; torch and torchvision come from the CUDA 13.0 index. If neither \code{uv} nor Python 3.11 is on the machine, the installer fetches \code{uv} first.
\item The weights, into the normal Hugging Face cache: the vendor's FP8 checkpoint (it supplies the tensors the NVFP4 overlays do not replace, and the tokenizer), the four NVFP4 files and both fine-tuned drafters. With \code{--setup}, only Wren's model; \code{use-profile.sh} fetches the rest later.
\item A lookup corpus built from public text (one English and one German Wikipedia shard, plus Python sources from the venv's own packages, about 31 million tokens). \code{--no-corpus} skips it. With \code{--setup}, only the two shards: \code{use-profile.sh} builds the corpus for the model it installs, under \code{corpus/tokenizers/}, with no network.
\item The configuration: \code{~/TandemLLM/run/ops/serve.env}, the release's \code{ops/serve.env}\index{ops/serve.env@\texttt{ops/serve.env}} with this machine's paths, and a control script, \code{~/TandemLLM/bin/tandem}.
\item The start: the engine starts, the script waits for \code{/health}, and a smoke test prints tokens a second. \code{--no-start} stops after step 6.
\end{enumerate}

The weights are public. \code{--hf-token} (or \code{HF_TOKEN}) is needed only for a repository that is private or gated, and a stored \code{hf auth login} works too. Table~\ref{tab:ref-install-flags} lists every option; the most used ones are \code{--setup} (Section~\ref{sec:new-box}), \code{--repo} (another repository, such as a fork), \code{--dir} (default \code{~/TandemLLM}), \code{--port} and \code{--host} (default \code{127.0.0.1:8000}), \code{--ref} (any git ref; \code{v0.2.0-staircut} is the paper's release), \code{--hf-token}, \code{--systemd} for a systemd user unit, \code{--dry-run} to print every step and change nothing, and \code{--uninstall}.

% Generated by manual/tools/gen_reference.py from the code. Do not edit by hand.

\begin{reftabletwo}{install.sh options}{tab:ref-install-flags}{\hd{Option} & \hd{What it does}}
\texttt{-{}-dir DIR} & install directory (default \textasciitilde{}/TandemLLM) \\
\texttt{-{}-port PORT} & the API port (default 8000) \\
\texttt{-{}-host ADDR} & the address to bind (default 127.0.0.1; 0.0.0.0 serves the LAN, with no auth on the OpenAI routes) \\
\texttt{-{}-ref REF} & the engine release: tag, branch or commit (default v0.15.4) \\
\texttt{-{}-repo URL} & the git repository to clone (default the public repository) \\
\texttt{-{}-hf-\allowbreak{}token TOKEN} & a Hugging Face token (or set HF\_TOKEN; a stored 'hf auth login' works too) \\
\texttt{-{}-no-\allowbreak{}corpus} & skip the lookup corpus (the lookup drafter then copies from the conversation only) \\
\texttt{-{}-no-\allowbreak{}start} & install and configure, but do not start the engine \\
\texttt{-{}-with-\allowbreak{}chat} & the chat (beta) too: the conductor (router, user memory, packages), its small models, and TandemApp's TandemNode with chatd, the Chat tab's service; needs -\kern0pt-app-wheel (a re-run keeps the TandemApp it installed). Without it a new install serves the engine alone (no page, no chat) and a re-run keeps the mode \\
\texttt{-{}-app-\allowbreak{}wheel FILE} & with -\kern0pt-with-chat: TandemApp's wheel, a local file (tandemapp-*.whl), into this install's venv (the engine's pins held) \\
\texttt{-{}-app-\allowbreak{}wheel-\allowbreak{}sha256 HEX} & the wheel's sha256, checked before it is installed \\
\texttt{-{}-engine-\allowbreak{}only} & the engine alone, a new install's default; on a re-run it turns the chat off \\
\texttt{-{}-setup} & set up instead of a model now: install the code and Wren's model only, then start setup mode (the Setup API, /v1/setup/*, for TandemApp's Setup page; docs/setup.md); 127.0.0.1 unless -\kern0pt-host says otherwise \\
\texttt{-{}-force-\allowbreak{}setup} & with -\kern0pt-setup, on a directory that already has a configured engine \\
\texttt{-{}-systemd} & also install a systemd user unit, tandemllm.service (off by default) \\
\texttt{-{}-runtime-\allowbreak{}only} & the engine's source and Python environment only: no weights, no corpus, no config, no start (what TandemApp's Add Spark wizard installs on a new node) \\
\texttt{-{}-yes,\allowbreak{} -\allowbreak{}y} & non-interactive: answer yes to every question \\
\texttt{-{}-dry-\allowbreak{}run} & print every step and command, change nothing \\
\texttt{-{}-uninstall} & stop the engine, remove the unit and DIR (the weights stay in the HF cache) \\
\texttt{-h,\allowbreak{} -{}-help} & this text \\
\end{reftabletwo}
\noindent Environment variables the installer reads: \texttt{TANDEM\_\allowbreak{}REF}, \texttt{TANDEM\_\allowbreak{}REPO}, \texttt{TANDEM\_\allowbreak{}DIR}, \texttt{TANDEM\_\allowbreak{}PORT}, \texttt{TANDEM\_\allowbreak{}HOST}, \texttt{HF\_\allowbreak{}TOKEN}, \texttt{HF\_\allowbreak{}HOME}.


The installer's corpus is not the corpus the paper's numbers come from; that one is not published. With the installer's corpus, or with \code{--no-corpus}, the lookup drafter proposes different continuations, so speeds can differ from the paper's. How much is not measured. On a CUDA GPU other than the GB10 the installer turns off the skinny NVFP4 kernel, which is built for \code{sm_121a} only, and StairCut's cost tables still come from the Spark, so such GPUs are best effort.

\section{After the install}

The install directory holds everything the installer wrote:

\begin{table}[htbp]
\centering\small
\caption{The layout of an installed engine.}\label{tab:install-layout}
\begin{tabularx}{\linewidth}{@{}L{.32\linewidth}Y@{}}
\toprule
\hd{path} & \hd{what} \\
\midrule
\code{~/TandemLLM/src} & the repository at the installed release \\
\code{~/TandemLLM/venv} & the Python environment, about 8~GiB with the CUDA wheels \\
\code{~/TandemLLM/corpus} & the lookup corpus \\
\code{~/TandemLLM/run/ops/serve.env} & the served profile with this machine's paths, rewritten on every run of the installer \\
\code{~/TandemLLM/run/local.env} & your own settings; \code{serve.env} reads it last and the installer never touches it \\
\code{~/TandemLLM/state} & the tokens file (the metrics token, and the admin password if you set one), once you create it \\
\code{~/TandemLLM/bin/tandem} & the control script \\
\code{~/.qwen38-spark-engine/} & the engine's state: the usage ledger, the conductor's documents and journal, memory notes, the disk tier, small models \\
\bottomrule
\end{tabularx}
\end{table}

The control script takes one command:

\begin{shell}
~/TandemLLM/bin/tandem start      # ops/start.sh with run/ops/serve.env; waits for /health
~/TandemLLM/bin/tandem stop 120   # SIGTERM, wait for the drain, kill only after it
~/TandemLLM/bin/tandem restart
~/TandemLLM/bin/tandem status
~/TandemLLM/bin/tandem smoke      # /v1/models and one short completion, with tok/s
~/TandemLLM/bin/tandem logs       # follow the engine's log
\end{shell}

\section{Engine only, or the chat (beta)}\label{sec:mode}\index{qsemode@\texttt{QSE\_MODE}}

A new install serves the engine alone: the OpenAI API under the served model's own name, the dashboard reads under \code{/v1/dashboard/*} and the caches. Nothing runs beside it: no extension (no \code{TandemLLM} model, router, user memory, packages, library or session monitor), no decision models, narrator, shadow reader or roles. One setting decides it, \code{QSE_MODE} in \code{serve.env}: \code{engine} or \code{chat}.

In chat mode (beta) the engine loads TandemApp's conductor as its extension (\code{QSE_EXT}), with the small models the weight profile names; TandemApp's TandemNode runs chatd, the Chat tab's service, and the engine runs none. Turn it on with \code{install.sh --with-chat --app-wheel FILE}, TandemApp's wheel (a re-run works too), with the Setup Guide's \emph{Chat (beta)} switch, which writes \code{QSE_MODE=chat} to \code{run/setup.env} and turns on Laya's decision models with it, or with \code{QSE_MODE=chat} in \code{run/local.env} and \code{tandem restart}. \code{--engine-only} turns it off again, and TandemNode stops a running chatd once the engine comes up in engine mode. On a box set up in the Setup Guide its switch decides: \code{run/setup.env} is read after the installer's \code{QSE_MODE}, so the two flags there only warn. The last block of \code{ops/serve.env} applies the mode after the weight profile, \code{run/setup.env} and \code{run/local.env}; \code{ops/use-profile.sh} keeps it.

An upgrade keeps each box's mode. Left empty, \code{QSE_MODE} is \code{chat} on a box whose conductor directory (\code{QSE_CONDUCTOR_CONFIG}, by default \code{~/.qwen38-spark-engine/conductor}) exists, so a box that ran the conductor stays a chat box with its documents, memory and sessions; \code{QSE_CONDUCTOR_CONFIG=off} set by hand stays engine only. \code{/health} and \code{/v1/models} say \code{"mode": "engine"} or \code{"chat"} and what is \code{loaded}: the conductor and its monitor's mode, the decision models, the narrator, the loaded roles, and whether chatd answers; \code{extensions} names the extension that runs.

\section{The first request}

The server speaks the OpenAI API on port 8000. In engine-only mode it serves the engine under its own model name (\code{/v1/models} lists it). With the chat on, it serves one model called \code{TandemLLM}, and any model name a client sends is served as \code{TandemLLM}:

\begin{shell}
curl -s localhost:8000/v1/chat/completions -H 'Content-Type: application/json' \
    -d '{"model": "TandemLLM", "messages": [{"role": "user", "content": "Say hello."}]}'
\end{shell}

Point other clients at the same address:

\begin{table}[htbp]
\centering\small
\caption{Pointing a client at the engine.}\label{tab:clients}
\begin{tabularx}{\linewidth}{@{}L{.2\linewidth}Y@{}}
\toprule
\hd{client} & \hd{setting} \\
\midrule
an OpenAI SDK & \code{OPENAI_BASE_URL=http://<host>:8000/v1} and any key \\
opencode & an OpenAI-compatible provider with that base URL and the model \code{TandemLLM}; with the chat (beta) on it gets the \code{coding} profile by itself \\
Open WebUI & an OpenAI connection to that base URL, with the connection header \code{X-OpenWebUI-Chat-Id: {{CHAT_ID}}} so each chat is its own conductor session with the chat (beta) on (TandemApp's manual, \enquote{Profiles and the router}); it gets the \code{chat} profile \\
\code{curl}, a script & as above; with the chat (beta) on it gets the \code{general} profile \\
\bottomrule
\end{tabularx}
\end{table}

An engine-only box (Section~\ref{sec:mode}) has no profiles: every client gets the raw path under the served model's own name, so opencode and Open WebUI name that model instead of \code{TandemLLM}. With the chat (beta) on, the header \code{X-Tandem: raw}\index{X-Tandem header} gives the engine without the conductor for one request. That is the release's raw path, the one the benchmarks and the paper measure. TandemApp's manual explains the profiles (\enquote{Profiles and the router}).

\screenshot[trim=0 640 0 0, clip]{client-open-webui}{Open WebUI with \texttt{TandemLLM} as its model: a demo request on the test board in October 2026, served under the \texttt{chat} profile, with the reasoning folded into \enquote{Thought for 1 seconds}. Only the message pane is shown.}{fig:shot-open-webui}

\screenshot{client-opencode}{opencode with \texttt{TandemLLM} as its provider: the end of a demo task on the test board in October 2026, under the \texttt{coding} profile. The model added a function, tested it, asked Laya for a verdict on the test log, removed its scratch files and reported what it did. The screenshot is from a build that sent the reasoning to opencode as text between \texttt{<think>} tags; every client but the Hermes agent now gets \texttt{reasoning\_content} by default (\texttt{QWEN38\_CLIENT\_REASONING}), which opencode shows as folded reasoning parts.}{fig:shot-opencode}

The engine asks for no password; read Section~\ref{sec:access} before anyone else can reach the port.

\section{Access and sign-in}\label{sec:access}\label{sec:dash-login}\index{dashboard login}\index{admin password}\index{sign-in}

TandemLLM has no sign-in turned on by default. The engine's OpenAI routes are always open, and so is everything else until you set an admin password: the dashboard, the conductor's settings, the memory notes, the transcripts, the decision models, the configuration and the cache routes. Anyone who reaches the port can read and change them. The installer binds the engine to \code{127.0.0.1}, so only the machine itself reaches it. With \code{--host 0.0.0.0}, or when you start \code{ops/serve.env} from a checkout (its \code{HOST} is \code{0.0.0.0}), the whole local network does.

Table~\ref{tab:access} has the switches. Only \code{/metrics} has a token of its own, for a Prometheus scraper, and \code{ops/make-secrets.sh} writes that one alone unless \code{--admin} asks for the password too.

\begin{table}[htbp]
\centering\small
\caption{Who can reach what, by default and with the switches set.}\label{tab:access}
\begin{tabularx}{\linewidth}{@{}L{.27\linewidth}L{.27\linewidth}Y@{}}
\toprule
\hd{what} & \hd{by default} & \hd{with the switch} \\
\midrule
the engine's OpenAI routes & anyone who reaches the port & no switch on the engine's own port; behind TandemApp's proxy, a user key from the first key on \\
the dashboard's reads & anyone & \code{--dashboard-login on} and a password: the token or a session \\
admin routes: conductor writes, memory, decision models, configuration, caches, the full \code{/health} & anyone & \code{QSE_ADMIN_TOKEN} (\code{ops/make-secrets.sh --admin}): the token or a session, or a tool on the machine itself \\
\code{/metrics} & anyone & \code{QSE_METRICS_TOKEN}: that token \\
\bottomrule
\end{tabularx}
\end{table}

With a password set, TandemApp's page has no sign-in screen in this release (TandemApp's manual). A request from the machine itself (loopback, without proxy headers) needs no token, because the watchdog and the measurement tools rely on that.

TandemLLM serves one user. Every client that reaches the engine is that user, so a \enquote{remember that} from any client lands in the same memory notes. A request that must not touch them sends \code{X-Tandem-Memory: off}. Until sign-in and several users arrive, keep the port on loopback or a trusted network. To reach it from elsewhere, put a reverse proxy with its own sign-in, or a VPN, in front of it.

\section{The small models}\label{sec:small-models}

Three small models can run beside the main model. They are off in the shipped profile, because each needs its weights on the disk. Each is one line in \code{local.env} and a restart.

\begin{table}[htbp]
\centering\small
\caption{The small models and the line that turns each on.}\label{tab:optional-parts}
\begin{tabularx}{\linewidth}{@{}L{.22\linewidth}L{.34\linewidth}Y@{}}
\toprule
\hd{part} & \hd{set in \code{local.env}} & \hd{the model} \\
\midrule
Laya, the decision models & \code{QWEN38_DECIDER_MODELS=laya-english,laya-multilingual,laya-typed-decisions} & \code{convaiinnovations/laya} under \code{~/.qwen38-spark-engine/deciders/laya} \\
the narrator's voice & \code{QWEN38_NARRATOR_MODEL=qwen3.5-2b} & \code{Qwen/Qwen3.5-2B} under \code{~/.qwen38-spark-engine/models/qwen3.5-2b} \\
the shadow reader & \code{QWEN38_SHADOW_MODEL=qwen3.5-0.8b} & \code{Qwen/Qwen3.5-0.8B} under \code{~/.qwen38-spark-engine/models/qwen3.5-0.8b} \\
\bottomrule
\end{tabularx}
\end{table}

Download a model with \code{hf download <repo> --local-dir <dir>}. The small models load after the engine, one at a time, through the residency planner\index{residency planner}, which refuses a load that would leave the board short of memory (Section~\ref{sec:residency}). The dashboard's System view shows what is loaded.

\section{Other engines and cloud APIs beside the engine}\label{sec:engine-backends}

The engine serves its own model only. Other engines and cloud APIs sit behind TandemApp's controller, \code{tandem-app}, on the board or another machine, with the engine as a backend of kind \code{tandemllm} (TandemApp's manual). Since engine API 1.8 the engine takes \code{--backends}\index{backends@\texttt{-\kern0pt-backends}} (\code{QSE_BACKENDS}\index{qsebackends@\texttt{QSE\_BACKENDS}} in \code{serve.env}) only as \code{off}: any other value stops the start with one line that names \code{tandem-app}, before anything loads. On a DGX Spark the other backends are cloud APIs or engines on other machines; a second GPU engine on the same Spark is not supported (Table~\ref{tab:setups}).

\section{Running from a checkout}

The installer does nothing you cannot do by hand. In an environment with the packages above, download the published weights and drafters with \code{hf}\index{Hugging Face}, the Hugging Face CLI (\code{pip install -U huggingface_hub}):

\begin{shell}
hf download Qwen/Qwen3.8-27B-FP8
hf download 0xBakeer/TandemLLM-Qwen3.8-27B-NVFP4 --local-dir ~/tandem/nvfp4
hf download 0xBakeer/TandemLLM-Qwen3.8-27B-DFlash2-b8 --local-dir ~/tandem/ft-b8
hf download 0xBakeer/TandemLLM-Qwen3.8-27B-DFlash2-b16 --local-dir ~/tandem/ft-b16
\end{shell}

Then edit the path lines at the top of \code{ops/serve.env}: \code{REPO} (the checkout), \code{PY} (a Python with the packages), \code{NV} (the three NVFP4 files, \code{mlp}, \code{gdn} and \code{attn}, comma-separated), \code{HEAD} (\code{head-fp8.safetensors}), \code{CKPT8}, \code{CKPT16} and \code{CORPUS}. Start it with \code{bash ops/start.sh}. That is the served profile.

In a PyTorch container that already has torch 2.13.0+cu130, add the two image packages:

\begin{shell}
pip install pillow==11.3.0 torchvision==0.28.0+cu130 \
    --extra-index-url https://download.pytorch.org/whl/cu130
\end{shell}

A server started directly with \code{python server/app.py} and no further flags is the raw engine. The conductor, the disk tier and the usage ledger come only with \code{--conductor-config}\index{conductor-config@\texttt{-\kern0pt-conductor-config}}, \code{--disk-cache}\index{disk-cache@\texttt{-\kern0pt-disk-cache}} and \code{--usage-ledger}\index{usage-ledger@\texttt{-\kern0pt-usage-ledger}}, which \code{ops/start.sh}\index{ops/start.sh@\texttt{ops/start.sh}} passes from \code{QSE_CONDUCTOR_CONFIG}\index{qseconductorconfig@\texttt{QSE\_CONDUCTOR\_CONFIG}}, \code{QSE_DISK_CACHE}\index{qsediskcache@\texttt{QSE\_DISK\_CACHE}} and \code{QSE_USAGE_LEDGER}\index{qseusageledger@\texttt{QSE\_USAGE\_LEDGER}}. To try the engine with the vendor's base drafter instead of the fine-tuned ones:

\begin{shell}
hf download z-lab/Qwen3.8-27B-DFlash2
python server/app.py --port 8000 --max-len 32768 --drafter dflash2 --dflash2-path greedy \
    --nvfp4 ~/tandem/nvfp4/mlp.safetensors,~/tandem/nvfp4/gdn.safetensors,~/tandem/nvfp4/attn.safetensors \
    --fp8-head ~/tandem/nvfp4/head-fp8.safetensors
\end{shell}

\section{Building the weights yourself}

The published NVFP4 set was quantised once from Qwen's BF16 release with GPTQ on the NVFP4 grid. Rebuilding it needs the BF16 checkpoint and a GPU with about 80~GB of memory at the default \code{--h-budget-gb 22}; it took 25 minutes on an RTX PRO 6000 (\measured):

\begin{shell}
hf download Qwen/Qwen3.8-27B
python tools/calib_corpus.py --out ~/calib
python tools/quant_nvfp4.py build --model <the BF16 snapshot> --methods gptq --out-dir ~/nvfp4-bf16 \
    --corpus ~/calib/calib-code.txt:80000,~/calib/calib-en.txt:60000,~/calib/calib-de.txt:16000
python tools/quant_head.py build --out ~/nvfp4-bf16/head-fp8.safetensors --ratios 1.0,0.95,0.90
\end{shell}

The build writes \code{gptq/mlp.safetensors}, \code{gdn.safetensors} and \code{attn.safetensors}; \code{tools/quality_gate.py --plan} scores them against the BF16 model. The engine loads them over the FP8 checkpoint, as it does the published files. Chapter~\ref{ch:quantisation} explains the method and the gate.

\section{Profiles}\index{profile (serve.env)}

A serving profile is a \code{serve.env} file: which weights, which drafters, which settings. Three ship with the NVFP4 set. They differ only in their weights; drafters, StairCut, caches and the exactness gate are the same in all three. The weight profiles under \code{ops/profiles/} are the other kind: \code{nvfp4}, six \code{exl3-<bpw>} sizes of the 27B and five \code{flashnext-<bpw>} packs, which \code{ops/use-profile.sh} switches between one engine at a time (Section~\ref{sec:two-models}).

\begin{table}[htbp]
\centering\small
\caption{The three serving profiles. Only the NVFP4 profile's speed is measured on this release; the bytes a step are \arith{} from the tensor sizes.}\label{tab:profiles-env}
\begin{tabularx}{\linewidth}{@{}L{.17\linewidth}L{.22\linewidth}Yr@{}}
\toprule
\hd{profile} & \hd{file} & \hd{weights} & \hd{bytes a step} \\
\midrule
NVFP4 (served) & \code{ops/serve.env} & the published set: every projection NVFP4, FP8 head & 15.0 GB \\
balanced & \code{ops/serve-balanced.env} & the published NVFP4 MLPs and FP8 head; the checkpoint's FP8 recurrent and attention layers & about 18.2 GB \\
FP8 & \code{ops/serve-fp8.env} & the checkpoint's own FP8 projections and BF16 head & 26.9 GB \\
\bottomrule
\end{tabularx}
\end{table}

The balanced and FP8 profiles source \code{ops/serve.env} and change two lines, \code{NV} and \code{HEAD}. They inherit StairCut's verify prices, which were measured on the NVFP4 profile. The router corrects their level online, but the staircase's shape for those two profiles is not measured, and neither is their speed. Do not confuse these files with the conductor's profiles (\code{coding}, \code{chat} and the others), which are documents the dashboard edits.

\section{Upgrading and uninstalling}\index{upgrade}\index{uninstall}

To upgrade, run the installer again with the same \code{--dir}. It checks out the new release, installs the new lock and rewrites \code{serve.env} from the release's template; \code{local.env} stays, and so does the box's mode (Section~\ref{sec:mode}). The usage ledger migrates in place at the first start: columns are added and every row is kept. A box installed with \code{--setup} and switched to a profile since then upgrades with \code{--setup --force-setup} and the profile switch after it, as Section~\ref{sec:update-setup} shows.

Coming from 0.3.0, a client sees these changes:

\begin{itemize}
\item The conductor is on in the shipped profile. \code{/v1/models} lists \code{TandemLLM} alone, every response carries the headers \code{X-Tandem}\index{x-tandem@\texttt{X-Tandem}} and \code{X-Tandem-Profile}\index{x-tandem-profile@\texttt{X-Tandem-Profile}}, and every request runs through the router, with a real cap on the thinking per effort level. \code{X-Tandem: raw} gives the release path of 0.3.0. A benchmark client that cannot send a header should run against a server with \code{CONDUCTOR=raw} and an empty \code{QSE_CONDUCTOR_CONFIG} in \code{local.env}.
\item The disk tier and the session spill are on: up to 50~GB under \code{~/.qwen38-spark-engine/cache/disk}, never below 50~GB free on the disk. \code{QSE_DISK_CACHE=off} turns both off.
\item A streamed request that finds the engine busy waits as long as a request may run, with an SSE comment every 10~s, instead of a 429 after 240~s. A short request is served between two blocks of the running turn.
\item The image processor's backend is pinned to torchvision.
\item A malformed \code{tools} field is a 400 that names the item.
\item The memory limits are enforced: 70~GiB between requests, 77~GiB during a prefill.
\end{itemize}

To run 0.4.0 as 0.3.0 was served, put \code{CONDUCTOR=raw}, an empty \code{QSE_CONDUCTOR_CONFIG=} and \code{QSE_DISK_CACHE=off} in \code{local.env}.

To roll back, run the installer with \code{--ref} set to the earlier tag; the ledger keeps its added columns and the older release ignores them. To uninstall:

\begin{shell}
bash ~/TandemLLM/src/install.sh --uninstall --yes
\end{shell}

This stops the engine, removes the systemd unit if one was installed, and deletes the install directory. It keeps the weights in the Hugging Face cache, which other tools may share, and \code{~/.qwen38-spark-engine/} with the conductor's documents, the journal, memory notes, transcripts, the disk tier and the small models. Delete that directory to remove them too.

\input{chapters/operations}
\input{chapters/community}

\partintro{%
Why the engine is fast, and why its output is still the model's own text. The chapters go from the model's layers and the byte budget (\ref{ch:architecture}) through the weights (\ref{ch:quantisation}), the kernels (\ref{ch:kernels}), the drafters and StairCut (\ref{ch:speculation}), the exactness contract (\ref{ch:exactness}), the caches (\ref{ch:caches}), the server (\ref{ch:serving}) and the measurement method (\ref{ch:measurement}) to the seams a new model plugs into (\ref{ch:adding-model}), and ends with the second and third models served through them, Qwen3.8-Flash-Next (\ref{ch:flashnext}) and LiquidAI's LFM2.5-24B-A2B (\ref{ch:lfm25}). Every number is labelled measured or arithmetic.}
\part{How the engine works}
\chapter{Architecture}\label{ch:architecture}

This part explains how the engine works, from the model's layers to the measurement method. This first chapter is the tour: the model, the byte budget that sets the speed limit, the decode loop, and a map of the code.

\section{The model}\index{Qwen3.8-27B}

\code{engine/config.py} reads the checkpoint's \code{config.json} and nothing else. The text model has 64 decoder layers, a hidden size of 5,120, a dense SiLU MLP of width 17,408 in every layer, and a vocabulary of 248,320 tokens. It has no mixture of experts. A 0.92~GB vision tower sits in the same checkpoint; the language-model loader skips it by name, and \code{engine/vision.py} loads it on its own for image input.

Every fourth layer (3, 7, 11 and so on up to 63) is full attention, 16 layers in all. Gated DeltaNet\index{Gated DeltaNet} fills the other 48. It is a recurrent layer: it keeps a fixed-size state per head instead of a KV cache that grows with every token.

\begin{figure}[htbp]
\centering
\begin{tikzpicture}[x=2.3mm, y=1mm]
  \foreach \i in {0,...,63} {
    \pgfmathtruncatemacro{\att}{mod(\i,4)==3 ? 1 : 0}
    \ifnum\att=1
      \fill[signal] (\i,0) rectangle +(0.82,9);
    \else
      \fill[cobalt!70] (\i,0) rectangle +(0.82,9);
    \fi
  }
  \foreach \i in {0,15,31,47,63} \node[lbl, below] at (\i+0.4,0) {\i};
  \node[lbl, anchor=west] at (0,13) {\textcolor{cobalt}{\rule{2.5mm}{2.5mm}} Gated DeltaNet, 48 layers: a 128$\times$128 fp32 state per head};
  \node[lbl, anchor=west] at (36,13) {\textcolor{signal}{\rule{2.5mm}{2.5mm}} full attention, 16 layers: KV rows};
\end{tikzpicture}
\caption[The layer pattern of Qwen3.8-27B]{The layer pattern of Qwen3.8-27B: three Gated DeltaNet layers, then one full-attention layer, sixteen times.}\label{fig:layers}
\end{figure}

\begin{table}[htbp]
\centering\small
\caption{The two layer kinds and the state each keeps per sequence (\arith, from the configuration's sizes).}\label{tab:layer-kinds}
\begin{tabularx}{\linewidth}{@{}L{.22\linewidth}rY@{}}
\toprule
\hd{layer kind} & \hd{count} & \hd{state per sequence} \\
\midrule
Gated DeltaNet & 48 & an fp32 state of 48 heads, each 128 by 128 (3.15 MB), and a 4-wide convolution tail (80 KB), per layer \\
full attention & 16 & BF16 keys and values, 4 KV heads of 256: 4 KB a token per layer \\
\bottomrule
\end{tabularx}
\end{table}

That split shapes most of the code. Published work on speculation and caching assumes attention, where undoing a token means moving a pointer back by one row. A recurrent state has absorbed every token and has no inverse. The engine needs exact ways to undo it, to branch it for a tree of guesses, and to restore it from a cache. Chapter~\ref{ch:speculation} shows the undo and the branch, and Chapter~\ref{ch:caches} the restore.

\section{The byte budget}\index{byte budget}

A decode step reads every weight once. Table~\ref{tab:bytes} sums the bytes from the checkpoint's safetensors headers.

\begin{table}[htbp]
\centering\small
\caption[The bytes one decode step reads, as shipped]{The bytes one decode step reads, as the vendor ships the model (\arith, from the safetensors headers).}\label{tab:bytes}
\begin{tabular}{@{}lr@{}}
\toprule
\hd{group} & \hd{bytes a step} \\
\midrule
48 Gated DeltaNet layers & 18.43 GB \\
16 attention layers & 5.96 GB \\
the output head, BF16, 248,320 rows by 5,120 & 2.54 GB \\
total, FP8 as shipped & 26.93 GB \\
\bottomrule
\end{tabular}
\end{table}

The three MLP projections are 63.6~\% of that. Quantising them first, then the rest, moved the budget down in steps, and each step had to pass the quality gate (Chapter~\ref{ch:quantisation}). Figure~\ref{fig:byte-budget} shows the steps and the floor each puts under a step at the board's paper bandwidth of 273~GB/s.

\begin{figure}[htbp]
\centering
\begin{tikzpicture}
\begin{axis}[tandem, ybar, height=6cm, width=.92\linewidth, bar width=22pt,
    ymin=0, ymax=31, ylabel={GB read per decode step},
    symbolic x coords={fp8,mlp,head,all}, xtick={fp8,mlp,head,all}, bar shift=0pt,
    xticklabels={{FP8 as shipped},{NVFP4 MLPs},{+ FP8 head},{+ NVFP4 GDN, attention}},
    x tick label style={align=center, text width=2.6cm},
    nodes near coords, nodes near coords style={font=\sffamily\scriptsize, align=center},
    point meta=explicit symbolic, enlarge x limits=0.18, grid=none]
  \addplot[fill=cobaltlight, draw=cobalt] coordinates
    {(fp8,26.93) [26.93 GB\\98.6 ms] (mlp,19.44) [19.44 GB\\71.2 ms] (head,18.17) [18.17 GB\\66.5 ms]};
  \addplot[fill=signal, draw=signal, bar shift=0pt] coordinates {(all,15.01) [15.01 GB\\55.0 ms]};
\end{axis}
\end{tikzpicture}
\caption[Bytes a step and the floor they set]{Bytes a step for each weight set, and the floor each sets under one step at 273~GB/s (\arith). The served set is the last bar. The board reads about 240~GB/s in practice, so the real floors are about 14~\% higher.}\label{fig:byte-budget}
\end{figure}

State is small beside the weights. The 48 recurrent states\index{recurrent state} are 151~MB together. The KV rows\index{KV rows} cost 64~KB a token across the 16 attention layers: 2.1~GB at 32k tokens and about 17~GB at the full 262,144-token context.

\section{The decode loop}\index{decode loop}\index{tokens per round}

Speed is one fraction:
\[
\text{tok/s} = \frac{\text{tokens committed a round}}{\text{draft ms} + \text{verify ms} + \text{commit ms}}.
\]
A round starts from the last committed token. The drafters propose a tree of guesses after it. One forward pass runs over the whole tree through all 64 layers. The engine reads the model's own choice at every node and keeps the longest path where the model agrees with the guesses, plus one more token, the model's own choice after the last accepted guess, so a round yields at least one token. Then it commits the state of that path and starts again.

On the served profile a round on the benchmark row takes 93.4~ms and commits 4.58 tokens, and the row's rate is 49.89 tok/s (\measured, the paper). The rate is the mean of the 50 requests' own rates, not the quotient of the two means (4.58 tokens in 93.4~ms is 49.0 tok/s). Without speculation the same box writes 13.69 tok/s at 72.8~ms a token. Chapter~\ref{ch:speculation} follows a round step by step.

\section{The parts}

Table~\ref{tab:code-map} maps each part of the engine to the files that implement it and the chapter that explains it.

\begin{table}[htbp]
\centering\small
\caption{A map of the code.}\label{tab:code-map}
\begin{tabularx}{\linewidth}{@{}L{.3\linewidth}YL{.14\linewidth}@{}}
\toprule
\hd{part} & \hd{where} & \hd{chapter} \\
\midrule
checkpoint layout and loading & \code{engine/loader.py}, \code{engine/config.py} & \ref{ch:adding-model} \\
the quantiser and its gate & \code{tools/quant_nvfp4.py}, \code{tools/quant_head.py}, \code{tools/quality_gate.py} & \ref{ch:quantisation} \\
matrix-product kernels & \code{tools/nvfp4_*.py}, \code{tools/fp8_linear.py}, \code{tools/head_gemv.py} & \ref{ch:kernels} \\
recurrent-layer kernels & \code{tools/gdn_*_kernels.py}, \code{engine/gdn.py} & \ref{ch:kernels} \\
target model, verify and commit & \code{engine/model.py}, \code{engine/tree.py} & \ref{ch:speculation} \\
drafters & \code{engine/drafters/} & \ref{ch:speculation} \\
routers: StairCut, length router, tree & \code{engine/lenrouter.py}, \code{engine/router.py}, \code{engine/prices.py} & \ref{ch:speculation} \\
sampling and penalties & \code{engine/sample.py}, \code{engine/penalty.py} & \ref{ch:exactness} \\
structured outputs & \code{engine/grammar.py} & \ref{ch:serving} \\
caches & \code{engine/cache.py}, \code{engine/disk_cache.py}, \code{engine/spill.py} & \ref{ch:caches} \\
settings registry & \code{engine/settings.py} & \ref{ch:operations} \\
the server & \code{server/app.py}, \code{server/stream.py}, \code{server/toolcall.py} & \ref{ch:serving} \\
the conductor & \code{server/conductor/} & \ref{ch:conductor} \\
live activity, metrics, ledger & \code{server/live.py}, \code{server/activity.py}, \code{server/metrics.py}, \code{server/ledger.py} & \ref{ch:dashboard} \\
the dashboard & \code{dashboard/}, \code{docs/contract/dashboard-v1/} & \ref{ch:dashboard} \\
the proxy and the collector & \code{proxy/} & \ref{ch:other-engines} \\
service scripts and the gate & \code{ops/} & \ref{ch:operations} \\
measurement tools & \code{tools/row3.py}, \code{tools/block_ab.py}, \code{tools/verify_spec.py} & \ref{ch:measurement} \\
\bottomrule
\end{tabularx}
\end{table}

\section{One sequence, or several}

The numbers in this part describe the loop of one sequence. The shipped profile runs 8 lanes since 0.15.0 (SRV-707). With \code{--parallel N} the engine holds N sequences over the one copy of the weights, and requests that decode at the same time verify their trees in one forward pass, so the weights are read once for all of them (Section~\ref{sec:parallel-requests}). With one sequence, requests wait for the engine's lock in a bounded queue, the side lane serves a short request inside the running one (Section~\ref{sec:side-lane}), and the proxy relays other backends' requests without touching the engine's lock (Chapter~\ref{ch:install-app}).

\input{chapters/quantisation}
\input{chapters/kernels}
\chapter{Drafters, StairCut and tree verify}\label{ch:speculation}

A verify of 16 tree rows costs about the same as a verify of 8. On the published NVFP4 set at 1k tokens of context, a tree verify with its commit takes 75.5~ms at 8 rows, 74.4~ms at 16, 83.1~ms at 24 and 87.8~ms at 32 (\measured, \code{ops/stair-tables-nvfp4.json}, median of 8). That flat stretch is why this engine speculates: the model's forward pass is bound by reading weights, and extra rows in the same pass are nearly free. TandemLLM therefore lets cheap drafters\index{drafter} guess many tokens and checks them all at once. This chapter covers the drafters, how StairCut decides what to spend each round, how the guesses become one tree, and how the tree is verified through recurrent layers.

\section{The round}\index{round}

Each round has four steps (Figure~\ref{fig:round}).

\begin{enumerate}
\item Draft. The block drafter proposes a block of tokens, and the lookup drafter proposes continuations it has seen before. The router merges them into one draft tree of at most 32 nodes.
\item Verify\index{verify}. One forward pass of the target model runs over the whole tree and reads the model's own next token at every node.
\item Accept. Walk down from the root. At each node take the model's token and move to the child that carries it. Stop where no child carries it. The model's token at the stopping point is kept too, so a round yields at least one token.
\item Commit\index{commit}. Keep the KV rows and the recurrent state of the accepted path, and drop the rest.
\end{enumerate}

\begin{figure}[htbp]
\centering
\input{figures/round}
\caption[One speculative round]{One speculative round. Both drafters propose; StairCut merges their proposals and cuts the tree to the node count with the most expected tokens per millisecond; one verify pass checks every node; the longest agreeing path and one token of the model's own are committed.}\label{fig:round}
\end{figure}

On the benchmark row a StairCut round takes 93.4~ms and commits 4.58 tokens (\measured, the paper).

\section{The block drafter}\index{DFlash2}\index{block drafter}

\code{engine/drafters/dflash2.py} runs DFlash2, a small drafter that proposes a whole block of tokens at once. It has five decoder layers of its own (hidden size 5,120) and reads five hidden states of the target, the residual stream entering layers 5, 19, 33, 47 and 61, fused into one input. It has no embedding and no head of its own: it borrows the target's. Inside a block every row attends to every other row, which is what lets it guess 15 tokens in one pass instead of 15 chained steps.

The engine holds two fine-tunes of the vendor's drafter, trained on the target's own output and published on Hugging Face: one trained at block 8 (the \emph{narrow} drafter: it proposes 7 tokens a block) and one at block 16 (the \emph{wide} drafter, 15 tokens). StairCut, the served behaviour, drafts with the block 16 one only. The served profile loads the block 8 one too and releases it at the first round; it drafts only with StairCut off. A draft moves about 6~GB (3.6~GB of drafter layers plus one read of the head for all rows), about 25~ms.

The drafter's head gives the top 16 candidates per slot, and a small bigram selector scores how each candidate follows the one before. That is a lattice of 15 slots by 16 candidates. The greedy path through it becomes the spine of the tree, and the next-best candidates become branches.

\section{The lookup drafter}\index{lookup drafter}

\code{engine/drafters/ngram.py} copies text instead of predicting it. It looks for the longest match of the last few tokens in three places:

\begin{itemize}
\item the request itself, prompt and output so far (orders 2 to 8);
\item a corpus of public text, memory-mapped with a suffix array;
\item the persistent suffix store\index{suffix store}, a log of what this engine has read and written before (Section~\ref{sec:suffix-store}).
\end{itemize}

Every continuation found is a vote, and the votes form a tree whose node scores are path probabilities. A lookup costs about 0.2~ms against about 25~ms for the block drafter. On fresh prose it rarely fires. On an edit of a file already in the context, or on a quotation, it can commit 15 tokens in a round. It fires only when the tree it would propose is worth its verify cost, because a drafter that proposes on every step and is usually wrong costs the round. After two wide rounds in a row that commit their whole width, the lookup drafter may propose one deep chain of up to 32 rows (\code{QWEN38_DEEP}\index{qwen38deep@\texttt{QWEN38\_DEEP}}).

\section{StairCut}\index{StairCut}

StairCut\index{StairCut} is the served behaviour: \code{QWEN38_LEN_SWITCH=1} and \code{QWEN38_LEN_MODE=wide}, off in the code and on in \code{ops/serve.env}\index{ops/serve.env@\texttt{ops/serve.env}}. Every round drafts with the wide drafter and builds four candidates of up to 31 nodes: its lattice tree, its chain, the lookup tree, and their merge. Each candidate is cut to the node count that maximises its expected committed tokens per millisecond,
\[
\frac{1 + \min\bigl(\sum_{\text{kept}} q,\; n\bigr)}{\text{verify}(\text{rows}) + \text{draft} + \text{commit}},
\]
and the round runs the best of the four.

A node's $q$ is its path probability from the drafter's lattice (read at temperature 1.4), or the lookup's vote share, scaled by a calibration the router learns online and capped at 1. Nodes are admitted in the order of the drafter's own score, each with the ancestors it needs. The price of a tree is verify plus draft plus commit; a chain pays the rollback instead of the commit, weighted by the chance it is rejected. The verify price is a staircase\index{staircase}, not a line, so the router tries every count instead of a search that assumes a smooth cost. A cut ends at 7, 15, 23 or 31 nodes (8, 16, 24 or 32 rows), or at the whole candidate when it is smaller.

\begin{figure}[htbp]
\centering
\begin{tikzpicture}
\pgfplotstableread[col sep=comma, comment chars={\#}]{data/verify-cost.csv}\verifycost
\begin{axis}[tandem, height=7cm, width=.92\linewidth,
    xlabel={rows in the verified tree}, ylabel={verify + commit, ms},
    xmin=0, xmax=34, ymin=70, ymax=115, xtick={2,4,8,12,16,20,24,28,32},
    legend pos=north west, legend cell align=left]
  \addplot[color=cobalt!35, mark=*, mark size=1.5pt] table[x=rows, y=tree1024] {\verifycost};
  \addplot[color=cobalt!55, mark=square*, mark size=1.4pt] table[x=rows, y=tree4096] {\verifycost};
  \addplot[color=cobalt!75, mark=triangle*, mark size=1.8pt] table[x=rows, y=tree8192] {\verifycost};
  \addplot[color=cobalt, mark=diamond*, mark size=1.8pt] table[x=rows, y=tree16384] {\verifycost};
  \addplot[color=signal, mark=*, mark size=1.6pt] table[x=rows, y=tree32768] {\verifycost};
  \legend{1k context, 4k, 8k, 16k, 32k}
  \draw[slate, dashed] (axis cs:16.5,70) -- (axis cs:16.5,115);
  \node[lbl, anchor=north west] at (axis cs:16.7,114) {17 rows need a second\\tensor-core tile};
\end{axis}
\end{tikzpicture}
\caption[Verify cost against tree rows]{The verify-cost staircase: a tree verify with its commit, by rows, at five context lengths. Flat to 16 rows, then steps. \measured: \texttt{ops/stair-tables-nvfp4.json}, published NVFP4 set, the medians StairCut prices its cuts on; this chart is drawn from that file.}\label{fig:verify-cost}
\end{figure}

\code{QWEN38_STAIR_TABLES}\index{qwen38stairtables@\texttt{QWEN38\_STAIR\_TABLES}} names the JSON file of verify prices by context length, chain and tree apart (\code{tools/verify_curve.py} measures them). Figure~\ref{fig:verify-cost} draws the served file. At 32k a 24-row tree costs about 16~ms more than a 16-row one, against 9~ms at 1k, and a tree costs 2 to 4~ms more than a chain of the same rows. \code{QWEN38_STAIR_RHO}\index{qwen38stairrho@\texttt{QWEN38\_STAIR\_RHO}} names a prior for the lookup's continuation rate on 8-token matches (\code{ops/lookup-rho.json}), which lets a long verbatim copy keep its whole line; the router updates it online against every committed block. \code{QWEN38_STAIR_SKIP=1} decides on the same prices whether to draft the head at all once a copy run has started: on a copy round the lookup line alone often wins, and at 32k a draft costs more than at 1k.

The draft cost is the drafter's own measured time at the current context length, learned online per context class. What the router learns (the draft cost per context class, the calibration, the copy counts) belongs to the router object, and the server keeps one router for its whole life, so what one request teaches reaches the next. \code{--no-learn-cost}\index{no-learn-cost@\texttt{-\kern0pt-no-learn-cost}} keeps the served price table instead, so what earlier requests took no longer moves a block's choice and a speed measurement repeats much more closely (Chapter~\ref{ch:measurement}). The timings inside the request still move it, so two runs can still part at a near tie (Section~\ref{sec:two-runs}).

With the node count calculated, the wide drafter's lattice cut to 16 nodes is a bush near the root on fresh text and a long line on copies, and the narrow drafter has nothing left to add. The router releases it at the first round.

\begin{figure}[htbp]
\centering
\begin{tikzpicture}
\begin{axis}[tandem, ybar, height=5.4cm, width=.8\linewidth, bar width=26pt,
    ymin=0, ymax=105, ylabel={tok/s, mean of 25 workloads},
    symbolic x coords={b8,b16,sc}, xtick={b8,b16,sc}, bar shift=0pt,
    xticklabels={fixed block 8, fixed block 16, StairCut},
    nodes near coords, nodes near coords style={font=\sffamily\scriptsize,
      /pgf/number format/fixed, /pgf/number format/precision=2, /pgf/number format/fixed zerofill},
    enlarge x limits=0.3, grid=none]
  \addplot[fill=cobaltlight, draw=cobalt] coordinates {(b8,59.31) (b16,81.95)};
  \addplot[fill=signal, draw=signal, bar shift=0pt] coordinates {(sc,90.45)};
\end{axis}
\end{tikzpicture}
\caption[The teacher-forced bench]{The teacher-forced bench: 25 workloads where every configuration commits the same reference text, so only rounds and their time differ. StairCut is above block 16 on 23 of the 25 workloads and above block 8 on 22. \measured: the StairCut paper.}\label{fig:forced}
\end{figure}

\section{Longer contexts}\index{context length}

Every token of context adds 64~KB of KV rows that each round reads, and the drafter reads the whole context too, so a round costs more as the conversation grows. The verify itself shows it: a 16-row tree with its commit takes 74.4~ms at 1k tokens of context and 81.5~ms at 16k, and 88.1~ms at 32k (\measured, the price table). Figure~\ref{fig:decode-context} puts the two measured decode rates the public record holds beside that curve and beside the bandwidth bound of plain decoding.

\begin{figure}[htbp]
\centering
\begin{tikzpicture}
\pgfplotstableread[col sep=comma, comment chars={\#}]{data/verify-cost.csv}\verifycost
\begin{axis}[tandem, height=6.2cm, width=.47\linewidth, name=left,
    xlabel={prompt length, tokens}, ylabel={tok/s, one request},
    xmode=log, log basis x=2, xmin=180, xmax=50000,
    xtick={256,1024,4096,16384}, xticklabels={256, 1k, 4k, 16k},
    ymin=0, ymax=60, legend style={at={(0.03,0.03)}, anchor=south west, font=\sffamily\scriptsize},
    legend cell align=left]
  \addplot[only marks, mark=*, mark size=2.4pt, color=signal] coordinates {(256,49.89) (16350,36.02)};
  \addplot[domain=256:40000, samples=40, color=slate, dashed] {240/(15.01 + 0.0000655*x)};
  \legend{StairCut (measured), plain decoding bound}
\end{axis}
\begin{axis}[tandem, height=6.2cm, width=.47\linewidth, at={(left.east)}, anchor=west, xshift=1.6cm,
    xlabel={context, tokens}, ylabel={16-row tree verify, ms},
    xmode=log, log basis x=2, xmin=800, xmax=40000,
    xtick={1024,4096,16384,32768}, xticklabels={1k,4k,16k,32k}, ymin=70, ymax=95]
  \pgfplotstableread[col sep=comma, comment chars={\#}]{data/verify-by-context.csv}\byctx
  \addplot[color=cobalt, mark=*] table[x=context, y=tree16] {\byctx};
\end{axis}
\end{tikzpicture}
\caption[Decode speed and verify cost against context length]{Left: single-request decode rate with StairCut at two prompt lengths (\measured: 256 tokens, the paper's benchmark row; 16,350 tokens, the shadow reader's A/B with the shadow off, the 0.4.0 series), and the bandwidth bound of plain one-token decoding when every step reads 15.01~GB of weights plus the context's KV at 240~GB/s (\arith). Right: the 16-row tree verify by context length (\measured, \texttt{ops/stair-tables-nvfp4.json}).}\label{fig:decode-context}
\end{figure}

Two points are all the public record holds for decode speed against prompt length on the published weights, and they come from two builds of the same series. A sweep over prompt lengths, under the rules of Chapter~\ref{ch:measurement}, is not published yet.

\section{The length router}\index{length router}

Before StairCut, the served profile picked the block width once per request, and with the StairCut flags off \code{engine/lenrouter.py} still does. It runs four wide rounds first. A wide round prices both widths at once, because a narrow block is a prefix of a wide one and the target's token at row $i$ does not depend on rows after $i$. Where that says the narrow width still has room, it runs up to four narrow probes. Then it decides, keeps the decision for the rest of the request, and releases the losing drafter (\code{--drop-idle}\index{drop-idle@\texttt{-\kern0pt-drop-idle}}), which then stops syncing its cache. It decides once because switching costs acceptance: a drafter keeps its own cache of the context, and a switched-in drafter starts behind. \code{--len-fixed 8} or \code{16} pins one width; those are the fixed baselines the paper compares StairCut with.

\section{Building the tree}\index{draft tree}

\code{engine/router.py} merges the block drafter's lattice and the lookup tree into one \code{DraftTree} (\code{engine/tree.py}), where a shared prefix becomes a shared node. With StairCut on, the merge is built at up to 31 nodes and the staircase cut sets its size. With StairCut off, the router prunes it to a node budget: a node stays only while its expected gain in accepted tokens pays for its share of the verify cost, priced on a verify curve (\code{QWEN38_TREE_MS}, or a price table). The budget is 16 nodes for either drafter, and the wide drafter gets 24 once the request has committed 32 tokens, so a short answer never pays for branches it cannot use.

The router's cost estimates are learned in the loop. It scores the lookup tree against the tokens the target actually wrote on every round, chosen or not, so its calibration settles before the lookup drafter has had a turn.

\section{Tree verify through recurrent layers}\index{tree verify}

Verifying a tree through attention layers is known work: give each node an attention mask of its ancestors and a position equal to its depth. The 48 Gated DeltaNet layers are harder. Their state after node $t$ is the state after $t$'s parent, updated by $t$, and a plain left-to-right walk overwrites the parent's state before $t$'s sibling needs it.

\code{engine/tree.py} stores a tree in depth-first pre-order: every node's parent has a lower index, and every subtree is a contiguous range (Figure~\ref{fig:verify-tree}). In that order the ancestor relation is part of the strict lower triangle, and every prefix operation of a chain verify has an exact tree counterpart.

\begin{figure}[htbp]
\centering
\input{figures/verify-tree}
\caption[A draft tree, its accepted path and its mask]{A draft tree of 8 nodes in depth-first pre-order. The model's own token at node 0 is \enquote{cat}, which node 1 carries; at node 1 it is \enquote{sat}, which node 2 carries; at node 2 it is \enquote{on}, which no child carries. The round commits \enquote{cat sat} and the model's \enquote{on}. Right: the ancestor mask the verify applies in the attention layers, and the relation the recurrent layers walk.}\label{fig:verify-tree}
\end{figure}

\begin{table}[htbp]
\centering\small
\caption{Every chain operation has an exact tree counterpart.}\label{tab:chain-tree}
\begin{tabular}{@{}ll@{}}
\toprule
\hd{chain} & \hd{tree} \\
\midrule
causal mask & the committed prefix plus the node's ancestors \\
position = index & position = depth, so two siblings share a position \\
the convolution slides over neighbours & it gathers each node's 3 nearest ancestors \\
cumulative gate along the sequence & the sum of gates along the node's path \\
the delta rule's lower-triangular solve & the same solve on the ancestor relation \\
accept a prefix, roll back to a length & accept a path, commit it by a gather \\
\bottomrule
\end{tabular}
\end{table}

The state after any accepted path can be written from quantities the verify already computed. Let $S_0$ be the state before the round, and for each node $t$ on the path let $k_t$ be its key, $u_t$ its pseudo-value and $g_t$ its cumulative gate. Then
\[
S_\text{path} = e^{g_\text{last}}\, S_0 + \sum_{t \in \text{path}} e^{g_\text{last} - g_t}\; k_t \otimes u_t .
\]
$u_t$ depends only on $t$'s own path, because the inverse of an ancestor-masked matrix is ancestor-masked too. The formula holds for every root-to-node path, not only for prefixes. The commit is therefore one read of the state and one small matrix product per layer, and it reads no weights. For a chain the same identity is the rollback: keeping the first $n$ tokens of a block is the path of length $n$. The KV side is simpler: the cache holds every node at its depth-first slot, and a commit gathers the accepted rows down into place.

\section{Sampling}\index{sampling}

Greedy requests (temperature 0) take the path above. A sampled request stays exact in distribution: the target draws its own token at each node, and a drafted token survives exactly while the draw lands on it. For a drafter that proposes one fixed token this is the standard rejection sampler, with no second draw. With \code{--sampled-tree det} (served) a sampled request keeps the tree. A request with a \code{seed} draws with noise keyed by position (the Gumbel-max trick), so the same seed gives the same text whatever the drafters and the router did, up to the same near-tie rule as greedy decoding. Chapter~\ref{ch:exactness} covers both.

\section{Reasoning}\index{reasoning budget}

The model thinks inside \code{<think>} tags. A reasoning budget closes the block when it is spent: the engine appends the vendor's closing sentence and \code{</think>}, runs those tokens through the model at their real positions, and the model answers from there. This changes the answer, by design. The budget comes from the server's \code{--think-budget}, from the client in any of the shapes clients send (Section~\ref{sec:client-budgets}), or from the conductor's effort level. With the conductor on, the closing phrase is hidden from the client, which receives \code{</think>} and the answer; the conductor's loop guards close a reasoning block that goes in circles (TandemApp's manual, \enquote{The session monitor}).

\section{The same method on Flash-Next}\label{sec:spec-flashnext}\index{MTP drafter}\index{StairCut!Flash-Next}

Flash-Next drafts with the multi-token-prediction layer its checkpoint ships, and its StairCut is written once for every model (\code{engine/methods/staircut.py}). Each round the MTP head drafts a chain of 4 tokens with its top two a depth, from one captured graph. The cut keeps the spine's prefix and, with siblings, the second choices that maximise
\[
\frac{1 + \sum_{\text{kept}} q}{V(\text{rows}) + \text{draft}},
\]
where $q$ comes from \code{ProbDeciles}, the head's probability mapped to the acceptance measured per tenth of probability and learned from every verified round. With \code{QWEN38_FN_LOOKUP=1} the lookup's best continuation joins the same tree: where its tokens match the head's chain the two share nodes, where they differ it hangs under the last shared node, and the cut admits nodes best-first up to 32 rows. A round still runs one verify, and on a copy the lookup lifted the rate from 102.3 to 132.9~tok/s (Table~\ref{tab:fn-lookup}).

$V(\text{rows})$ comes from \code{PriceModel}: a prior from a price file in the format of the 27B's stair tables, \code{ops/prices/<model>/<profile>.<device>.json}, followed online and never cheaper for more rows. The device is \code{dgx-spark} for the GB10 (\code{engine/methods/device.py}); a profile without a file for its device learns every price from the first rounds. On this model the verify is a ramp, not a staircase: every row can route to other experts, and at about 1k of context 2 rows cost 27.8~ms and 16 rows 68.8~ms (\measured, learned online).

The cut also decides whether to draft at all. Prices, the plain step's cost and a drafted round's cost a token are kept per context class, by wall time; a class whose drafts cost more a token than its plain steps takes the plain step, with probe rounds at a doubling interval, and a turn drafts only when the saving pays for the prompt rows the head has not read yet. Section~\ref{sec:fn-plain} has the rule and the runs that set it. A verified row on Flash-Next is the decode step's row bit for bit, so none of this moves the text.

\section{Knobs}

\begin{table}[H]
\centering\small
\caption{The speculative decoding knobs.}\label{tab:spec-knobs}
\begin{tabularx}{\linewidth}{@{}L{.5\linewidth}Y@{}}
\toprule
\hd{knob} & \hd{what} \\
\midrule
\code{--drafter lenrouter}, \code{--dflash2-ckpt}, \code{--dflash2-ckpt16} & the block drafters (both loaded; StairCut drafts with block 16) \\
\code{--narrow-arm on|off} & build the block 8 drafter at all \\
\code{QWEN38_LEN_SWITCH=1}, \code{QWEN38_LEN_MODE=wide} & StairCut \\
\code{QWEN38_STAIR_TABLES}, \code{_RHO}, \code{_SKIP} & its prices, the lookup prior, the head-skip rule \\
\code{--len-fixed 8|16}, \code{--len-latch}, \code{--drop-idle} & with StairCut off: pin a width, decide once a request, release the loser \\
\code{--tree}, \code{--budget}, \code{QWEN38_TREE_NODES*}, \code{QWEN38_TREE_MS} & the tree and its node budgets (StairCut off) \\
\code{--corpus}, \code{--suffix-store} & the lookup drafter's corpus and persistent store \\
\code{QWEN38_DEEP}, \code{QWEN38_DEEP_AFTER} & the deep chain after long accepted paths \\
\code{--sampled-tree det|mixed}, \code{QWEN38_DRAFT_TEMP} & sampled requests on the tree, and the drafter's temperature \\
\code{--learn-cost}, \code{--no-learn-cost} & learn verify and draft costs online, or keep the price table \\
\code{--think-budget}, \code{--reasoning-effort} & the server's reasoning defaults \\
\code{--drafter mtp}, \code{QWEN38_FN_DRAFT_HOT} & Flash-Next: its MTP layer as the drafter, through a head of the 32,768 most frequent ids \\
\code{QWEN38_FN_LOOKUP} & Flash-Next: the lookup in the MTP round's tree (on in its profiles) \\
\code{ops/prices/<model>/<profile>.<device>.json} & the shared StairCut's price file \\
\bottomrule
\end{tabularx}
\end{table}

\input{chapters/exactness}
\input{chapters/caches}
\chapter{Serving}\label{ch:serving}

\code{server/app.py}\index{server/app.py@\texttt{server/app.py}} speaks the OpenAI HTTP API, so any OpenAI client uses the engine by changing its base URL. It is built on the Python standard library (\code{ThreadingHTTPServer}), with no web framework. This chapter covers the request fields, reasoning and thinking budgets, tool calls, structured outputs, streaming, images, and how the server behaves under load: the queue, the side lane, the memory gate, the wait for memory and the loop guard. Chapter~\ref{ch:ref-http} lists every route.

\section{What the API supports}\label{sec:api-fields}\index{OpenAI API}

\code{server/compat.py} gives every OpenAI request field one of three answers, and \code{tests/test_compat.py} pins the table.

\begin{itemize}
\item Honoured: \code{messages}, \code{model}, \code{stream}, \code{stream_options}, \code{max_tokens}, \code{max_completion_tokens}, \code{stop}, \code{temperature}, \code{top_p}, \code{seed}, \code{presence_penalty}, \code{frequency_penalty}, \code{logit_bias}, \code{logprobs}, \code{top_logprobs} (up to 20), \code{n} (up to 16, not streamed), \code{tools}, \code{tool_choice}, \code{parallel_tool_calls}, \code{reasoning_effort}, \code{response_format}, \code{modalities} (text output only). Messages may carry \code{image_url} parts.
\item Accepted, without effect by their own definition: \code{user}, \code{metadata}, \code{store}, \code{service_tier}, \code{prompt_cache_key}, \code{safety_identifier}, \code{prediction}.
\item Refused with a 400 that names the field: \code{audio}, \code{functions} and \code{function_call} (use \code{tools}), \code{web_search_options}, \code{verbosity}, and for completions \code{suffix} and \code{echo}. A malformed \code{tools} (not a list, an item without a function name, \code{parameters} that is not an object) is a 400 that names the item.
\end{itemize}

The request's own shape is checked first, and a wrong one is a 400 that names the field before anything reads it: no \code{messages} or an empty list, an unknown role (\code{messages[0].role}), a system prompt alone or only tool results (no user message to answer), a system message after the conversation began, \code{max_tokens} or \code{max_completion_tokens} below 1, an empty completion \code{prompt}, a body that is not a JSON object, and a prompt that renders to no token. OpenAI's \code{developer} role is the system prompt; a system and a developer message at the start become one system prompt, because the chat template reads one.

Fields that are not OpenAI's pass through without a 400, because clients send extensions meant for other servers. The engine reads its own: \code{top_k}, \code{min_p}, \code{repetition_penalty}, \code{no_repeat_ngram_size}, \code{max_reasoning_tokens}, \code{reasoning_format}, \code{draft_temperature}, \code{chat_template_kwargs}, \code{cache_prompt}, \code{wait_for_memory} and \code{tandem.priority} (Section~\ref{sec:ref-tandem-fields}), and the other servers' thinking-budget shapes.

The prompt caches are on for every request. \code{cache_prompt: false} reads the whole prompt, so \code{timings.cache_n} is 0, and leaves every prompt cache as it was: the conversation in place is set aside as for another client's and is there for the next request without the flag. Any other value, or none, keeps the caches on. \code{docs/server.md} has the details.

\section{Reasoning}\index{reasoning format}

The model thinks inside \code{<think>} tags. The chat template opens the block in the prompt, so the model only ever writes \code{</think>}; the server re-emits \code{<think>} as the first delta so that clients which fold on a matched pair see one. \code{--reasoning-format}\index{reasoning-format@\texttt{-\kern0pt-reasoning-format}}, or \code{reasoning_format} in the request, picks how the thinking reaches the client:

\begin{table}[htbp]
\centering\small
\caption{Reasoning formats.}\label{tab:reasoning-format}
\begin{tabular}{@{}lll@{}}
\toprule
\hd{value} & \hd{\texttt{content}} & \hd{\texttt{reasoning\_content}} \\
\midrule
\code{tags} (default) & the thinking inside \code{<think>...</think>}, then the answer & empty \\
\code{reasoning_content} & the answer only & the thinking, as deltas \\
\code{both} & as \code{tags} & as \code{reasoning_content} \\
\bottomrule
\end{tabular}
\end{table}

\section{Thinking budgets from the client}\label{sec:client-budgets}\index{thinking budget}

A client may cap the thinking in any of the shapes clients send, and \code{server/thinking.py} reads them all into one budget.

\begin{table}[htbp]
\centering\small
\caption{The thinking-budget shapes the server reads.}\label{tab:budget-shapes}
\begin{tabularx}{\linewidth}{@{}L{.44\linewidth}Y@{}}
\toprule
\hd{shape} & \hd{read as} \\
\midrule
\code{max_reasoning_tokens}, \code{thinking_budget} & a cap in tokens \\
\code{chat_template_kwargs.thinking_budget}, \code{.max_thinking_tokens}, \code{.reasoning_budget}, \code{.thinking_token_budget} & a cap in tokens \\
\code{reasoning: {max_tokens, effort, enabled}} & \code{max_tokens} a cap, \code{effort} a word, \code{enabled: true} alone medium, \code{enabled: false} off \\
\code{thinking: {type, budget_tokens}} & \code{enabled} with \code{budget_tokens} a cap, \code{enabled} or \code{adaptive} alone on without a cap, \code{disabled} off \\
\code{reasoning_budget}, \code{max_thinking_tokens}, \code{thinking_token_budget}, \code{thinking_budget_tokens} & a cap in tokens \\
\code{reasoning_effort}, \code{reasoning.effort}, \code{chat_template_kwargs.reasoning_effort} & the level's cap: \code{minimal} and \code{low} 1,024, \code{medium} 8,192, \code{high} 16,384, \code{xhigh} and \code{max} none; \code{none} is thinking off \\
\code{enable_thinking: false} (top level or in \code{chat_template_kwargs}) & thinking off \\
\bottomrule
\end{tabularx}
\end{table}

Whether the model thinks is decided first: by \code{chat_template_kwargs.enable_thinking} when it is set, else by the top-level \code{enable_thinking}, else by any off form. The cap is then the first number in the table's order, else the first effort word's cap. A number beats a word, and 0 is no cap. An unknown word or a negative number is a 400 that names the field.

With the conductor deciding the effort, a client's number always wins, and its effort word is a ceiling on the router's level (Chapter~\ref{ch:conductor}). When a client asked for a budget in any shape and the reasoning would otherwise run into the output limit, the reasoning closes early enough to leave the answer \code{--answer-reserve}\index{answer-reserve@\texttt{-\kern0pt-answer-reserve}} tokens: \code{auto}, the default, is the smaller of 2,048 and a quarter of the limit. A request with none of these fields is served as before, token for token.

\section{Tool calls}\index{tool calls}

The request's \code{tools} go into the chat template\index{chat template}, and the model answers a call in its own XML form:

\begin{plain}
<tool_call>
<function=read>
<parameter=offset>
150
</parameter>
</function>
</tool_call>
\end{plain}

\code{server/toolcall.py} turns that into OpenAI \code{tool_calls} with \code{finish_reason: "tool_calls"}, on the streamed and the JSON path. It reads the strict form, a lenient form and a JSON call object, and only in the answer, never in the thinking. On a stream, a call goes out as argument deltas while the model writes it: the id and the name first, then each string value in pieces, so a whole-file write shows progress instead of minutes of silence.

Argument values carry their schema's types. The model writes every value as text, so \code{offset} arrives as the characters \code{150}. The parser looks up the request's own schema, and a parameter whose schema does not allow a string is returned as JSON: \code{150}, \code{true}, a list. A typed value is held back on the stream until it is complete, then sent whole. On an 8-task agent run, a strict client that validates arguments against the schema refused 58 of 80 calls before this, and 0 of 26 after (\measured).

\code{tool_choice} works as OpenAI defines it. \code{none} renders the prompt without tools. \code{required} and a named function become a constraint on the answer, the same token masks structured outputs use. \code{parallel_tool_calls: false} keeps one call.

\section{Structured outputs}\index{structured outputs}

\code{response_format} with \code{json_object} or \code{json_schema}, and \code{structured_outputs} with a \code{regex}, a \code{choice} list or a JSON schema, constrain the answer to a regular language. \code{engine/grammar.py} compiles each form to a finite automaton over bytes, because the tokenizer is byte-level and a token can end halfway through a character. A token is allowed in a state when walking its bytes never reaches the dead state; the engine computes that for the whole vocabulary at once per state and caches it, so a state costs a few milliseconds once. Each mask applies to every verified row of a draft tree, not only to one row per step, so constrained output under speculation is the text plain constrained decoding writes. JSON schemas with recursion are refused, and \code{json_object} allows a fixed nesting depth.

With tools offered and \code{tool_choice} auto, a document constrains the answer to the calls or the document (SRV-619). The automaton is the product of the tool-call language and the document's, and where both are still open the mask follows the row's argmax before the mask: a call's first token commits to the calls, whitespace leaves the choice open, and any other token (the document's opening brace, a fence, prose, the end token) commits to the document. The mask stays a function of the prefix and the row, so speculation changes no token. The log line \code{[structured] <id> branch=calls|document} and \code{extra.structured_branch} say which side won. \code{tool_choice: none} is the document alone, \code{required} and a named function the calls alone. Through 0.11.4 such a request was a 400; \code{--tool-format pythonic} keeps it.

A \code{</think>} the model writes again in its answer, alone on a line, is not sent under a conductor profile: Open WebUI closes a reasoning block at the first one and showed the second as text between two tool rows. The tag in a sentence, in code, in a tool call or closing a block the answer opened itself stays, and the raw path sends what the model wrote.

Open WebUI's terminal tools end \code{run_command}'s description with the terminal's live working directory, inside the system prompt's tool block, so an agent's first \code{cd} changed the prompt about 4,800 tokens in and the next user message read the whole conversation again (measured: 41,749 tokens, 56~s). Under a conductor plan that is on, the sentence moves into a note at the turn where the directory changed and replays there, so the tool block stays the same for the session (\code{QWEN38_CLIENT_CWD}, \code{note}; \code{off}: the tools as sent).

A client kind can have a format of its own when it names none and the server's is \code{tags}, under a conductor plan (\code{QWEN38_CLIENT_REASONING}, \code{kind=format,...}, \code{*} for every kind not named). The default, \code{*=reasoning_content,hermes=tags}, gives every client but the Hermes agent \code{reasoning_content}: an OpenAI SDK or curl client printed \code{message.content} and showed the model's thinking as its answer, and opencode with the installer's provider configuration showed it as literal \code{<think>} text. A client that wants the tags inline sends \code{reasoning_format: tags}. For Open WebUI the reason is its live view: Open WebUI 0.11.3 folds a \code{<think>} block out of the content after the fact and shows the stream piece that carries \code{</think>} as the text of a new message until the round's stream ends, so a long agent chat showed a literal \code{</think>} in 15 of 70 page scans; with \code{reasoning_content} three agent chats showed it in none of 683 scans (\measured). For that client the template's blank lines go out as one content piece where the reasoning ends, so Open WebUI closes and times each tool round's block even when the round goes straight to a tool call; a blank message is not shown. For the other clients the answer starts with its first character. A request's own \code{reasoning_format} wins. On the raw path an OpenAI SDK, curl or an unnamed client gets \code{reasoning_content} too, with stop strings matched on the answer, and every other client kind keeps the server's format (SRV-617): tool-eval-bench read the inline \code{<think>} block as the answer and failed its exact-answer scenarios.

Under a conductor plan a request's \code{stop} strings apply to the answer, never inside the reasoning, and the blank lines that start the answer are not matched: \code{stop: ["\textbackslash n\textbackslash n"]} used to end a response at the first blank line of the thinking, with no answer. A plain API client (curl, an OpenAI SDK, an unnamed client) gets no status line in its content. A request that offered no \code{tools} never gets \code{tool_calls}: a call the model writes anyway is read out of the answer and dropped, and the finish stays \code{stop}.

A JSON answer (\code{json_object}, \code{json_schema}) is parsed whole by its client, so a request that names no \code{reasoning_format} gets \code{reasoning_content} for it: the thinking goes to its own field and \code{message.content} is the JSON alone. A client that asks for \code{tags} keeps them. Such a request also gets the answer reserve a client's own thinking budget gets (Section~\ref{sec:client-budgets}), so the reasoning cannot spend \code{max_tokens} and cut the JSON, and it carries no status lines.

\section{Streaming}\index{streaming}

A streamed response sends one chunk per token, carries \code{usage} on the last chunk when the client asks for it (and by default, \code{--usage-default on}), then a chunk with \code{finish_reason} and \code{data: [DONE]}. A verify round accepts several tokens at once, so they reach the socket together: half the gaps between tokens are microseconds and the other half a whole round.

\code{usage} counts what the client sent and what it received. \code{prompt_tokens} is the client's messages as the chat template renders them; the tokens the conductor adds (its new and replayed notes, and the closed think block of a turn it turned thinking off) are \code{usage.prompt_tokens_details.hidden_tokens}, absent when there are none, and \code{cached_tokens} is the client's tokens served from a cache: the cached span less the hidden tokens inside it, which the server tracks by position (an opencode turn that re-read 3,827 of its 24,706 client tokens used to report all 24,706 cached; it now reports 20,879). The engine's own counts of the whole prompt are in \code{timings}, and the ledger keeps them as \code{engine_cached} and \code{forwarded_tokens}.

\code{timings} and \code{metrics} also carry the drafting of the answer. \code{timings.blocks} counts the decode forwards, one a verify round or plain step, and \code{tokens_per_block} is \code{(completion_tokens - 1) / blocks}: the drafted tokens a round kept plus the one the verify decides itself, so 1.0 means no draft was kept. \code{metrics.speculative_decoding.mean_acceptance_length} is the same number. \code{draft_n} and \code{draft_n_accepted} count the drafted tokens the verify checked and kept, and \code{draft_acceptance_rate} is their ratio. A stream carries both objects in the one chunk that carries \code{usage}. \code{docs/server.md} lists every field.

\code{finish_reason} is \code{stop}, \code{length}, \code{tool_calls}, \code{timeout} (the request's wall-clock cap) or \code{error}. On \code{error} the partial text has already been sent and an \code{error} object sits beside the finish reason. A client that hung up is logged as \code{abandoned}, not as an error. \code{server/stream.py} turns tokens into characters and holds back a trailing partial character until the next token completes it, so every character it sends is final and no replacement character reaches the client.

\section{A long prefill, and a client that leaves}

A long prompt prefills for minutes: 208.3~s at 128k tokens (\measured, the paper). After every prefill chunk the handler checks its socket. If the client has gone, the prefill stops there, the request ends \code{abandoned}, and the engine is free at once; the rows already prefilled stay resident, so a retry continues from there. A streamed prefill that has run for 5~s also sends an SSE comment after each chunk, at most once a second (\code{: prefill 12288/95131}). Clients ignore comment lines; proxies see a live stream and do not time it out.

\section{Images}\index{images}

A user, assistant or tool message may carry \code{image_url} parts with an \code{https://} or \code{data:image/...;base64,} URL. The server fetches and decodes them before the request queues and runs the checkpoint's own image processor: 16-pixel patches, 2 by 2 merge, 65,536 to 16,777,216 pixels. Each merged patch is one prompt row, so a 640 by 480 image is 300 rows, and those rows count as prompt tokens in \code{usage}. The vision tower\index{vision tower} (27 blocks, 0.92~GB, BF16) runs inside the prefill, when the first chunk that holds the image runs; the live view reads \enquote{Encoding image 1 of 2} meanwhile.

A bad image is a 400 that names the part: a scheme other than https or data, a failed download, broken base64, a format the server does not read (it reads five: PNG, JPEG, WEBP, GIF, BMP), more than 20~MB, a header that declares more than 64 million pixels, more than 16 images in the newest message, more than 65,536 image rows in the newest message, an image in a system message, a video part, or any image when the server runs with \code{--vision off}. \code{--image-act-gb} (6) refuses an image whose encode would need more memory than that, before it queues.

An https URL is a request the server makes for the client, so it is fenced. The host must resolve to public addresses only: loopback, private, link-local, shared, multicast and reserved ranges are refused, IPv4 inside IPv6 included. The connection goes to the address that was checked, with the certificate checked against the host name, so a second DNS answer cannot move it. A redirect is followed only to another https URL that passes the same check, at most 3 times. \code{--image-https-private on} lets URLs reach private addresses; \code{--image-https off} accepts data URLs only.

A long conversation is not refused for the images it has gathered. A coding agent that reads a screenshot every few turns sends every earlier image again on every turn, so once it has read 17, every later request would be over the cap of 16. When the images are over \code{--max-images}\index{max-images@\texttt{-\kern0pt-max-images}} or \code{--max-image-rows}\index{max-image-rows@\texttt{-\kern0pt-max-image-rows}} and the excess is in earlier messages, the oldest images are left out instead, each replaced in the prompt by a fixed text. They go in steps of half the cap (8 by default), counted from the conversation's start, so the prompt's head that the caches resume stays the same for several turns. The response header \code{X-Tandem-Images-Left-Out} says how many.

\section{Admission and limits}\index{admission}

A sequence serves one request at a time, so requests wait for a free one in a bounded queue. The shipped profile runs 8 of them (Section~\ref{sec:parallel-requests}).

\begin{table}[htbp]
\centering\small
\caption{Admission limits. Code defaults; the served values are in parentheses where they differ.}\label{tab:admission}
\begin{tabularx}{\linewidth}{@{}L{.24\linewidth}L{.16\linewidth}Y@{}}
\toprule
\hd{limit} & \hd{value} & \hd{what happens} \\
\midrule
\code{--max-queue} & 8 & the next request gets a 503 with \code{Retry-After: 5} \\
\code{--queue-timeout} & 120~s (240) & a request without a stream that waited that long gets a 429; a streamed chat request waits up to the request timeout \\
\code{--request-timeout} & 900~s (3,600) & a generation that runs that long ends with \code{finish_reason: "timeout"} \\
\code{--max-len} & 8,192 (262,144) & a prompt that does not fit gets a 400; \code{max_tokens} is clamped to what fits \\
\code{--default-max-tokens} & 8,192 (32,768; Ornith 131,072) & used when the request sends no \code{max_tokens}; a request's own \code{max_tokens} wins \\
\bottomrule
\end{tabularx}
\end{table}

At load the server allocates the KV buffer for the whole context, so the context length is a memory decision taken before the first request (Table~\ref{tab:memory-split}).

\section{The memory gate}\label{sec:memory-gate}\index{memory gate}

Three limits hold the engine's memory. \code{--mem-ceiling-gib}\index{mem-ceiling-gib@\texttt{-\kern0pt-mem-ceiling-gib}} (70) is what the engine may hold between requests, allocated plus the allocator's reserve; a cache grows only while the engine stays under it. \code{--mem-peak-gib}\index{mem-peak-gib@\texttt{-\kern0pt-mem-peak-gib}} (77) is the ceiling plus the workspace of a 262,144-token prefill, what the engine may hold during a prefill. \code{--reserve-slack-gib}\index{reserve-slack-gib@\texttt{-\kern0pt-reserve-slack-gib}} (2) is how many free blocks the allocator may keep between requests; over it they go back to the driver after the request, never inside one.

Before every prefill the gate brings the engine back under the ceiling and then checks that the prompt's estimated workspace fits under the peak. When it does not, it gives caches back, cheapest first (Figure~\ref{fig:memory-gate}). The last level releases the small models' own memory pools. Before a refusal the gate collects and empties the allocator once more, and a prompt whose workspace fits in the allocator's own free reserve is admitted. When nothing is enough, the request gets a 503 with \code{Retry-After} instead of running the board into its out-of-memory wedge, and the Dev view shows an error line with the allocated and reserved figures. With more than one lane, a request the admission cannot fit waits for memory (Section~\ref{sec:memory-wait}).

\begin{figure}[htbp]
\centering
\input{figures/memory-gate}
\caption[The memory gate]{The memory gate before every prefill. Each level is given back only while the prompt's workspace does not fit under the declared peak (77~GiB); a cache given back costs a later cold read, never a different token.}\label{fig:memory-gate}
\end{figure}

\section{The side lane}\label{sec:side-lane}\index{side lane}

A short request does not wait for a long turn to end. When a request finds the engine busy and is short (at most 2,048 prompt tokens and 64 output tokens by default, one choice, no images, thinking off) and is either a raw request or a client's side request such as a title call, it becomes a job. The generation that holds the engine serves it at its next safe point, between two decode blocks or between two prefill chunks (\code{--side-lane decode} leaves the prefill out, \code{off} turns the lane off). The guest is prefilled from position 0 on the prefill's grid and decoded one token a forward with no drafter, so its answer is plain greedy decoding of its prompt.

The generation gets the engine back as it was. \code{engine/cache.py}\index{engine/cache.py@\texttt{engine/cache.py}} \code{Borrow} copies what the guest can write: the KV rows the guest can reach, the recurrent state, the convolution tails, a walked chain state, and the record of a pending commit, copied unsettled so the next verify folds the same record into the same bytes. The drafter is never touched. A copy costs about 0.3~GB for a 2,000-token guest and is freed after it. \code{tests/test_side_lane.py} checks on the random model that the next block's logits, the state and every KV row are bit for bit an uninterrupted run's. The turn's text equals an uninterrupted run's only when its tree shapes are the same: the length router sizes blocks from timings measured inside the request (and, without \code{--no-learn-cost}, learned across requests), so a guest can change a later tree and a near tie can go the other way, as it can between two runs alone (Section~\ref{sec:two-runs}).

\begin{figure}[htbp]
\centering
\input{figures/side-lane}
\caption[The scheduler: side lane and queue]{The scheduler. A long turn holds the engine. Short guests (a title request, the router's own question) are served at the turn's next safe point and the turn gets its state back bit for bit. A second full turn waits in the queue; its stream opens after 10~s and carries a comment every 10~s.}\label{fig:side-lane}
\end{figure}

Measured on the engine (the 0.4.0 series, \code{tools/lane_check.py}, a 5,000-token prompt with 768 output tokens): with 10 short requests fired into it, the long request's reasoning and answer matched two runs alone in that hold. Nine of the ten were served in the lane, 1.2~s at the median and 5.1~s at most; the tenth asked for more output than the lane's cap and waited. The long request took 7.4~s longer, for 4.7~s the guests held the engine (\measured). A later measurement on the release's code parted at token 61 with guests, and at token 12 between two runs alone, both at near ties (Table~\ref{tab:two-runs}).

Three guards apply. A job waits for a later safe point when the engine's memory plus the guest's copy and workspace would pass the declared peak. A disk-tier write that has not copied the rows below the guest yet stops where it is. A job no holder takes runs on its own thread under the lock. Side-lane rows are no speed sample: a short plain decode beside a turn says nothing about the engine's decode rate.

\section{A turn behind a running turn}\index{waiting line}

The side lane serves short requests only. A second session's full turn still waits for the running turn to end. A streamed chat request waits as long as one request may run (\code{--request-timeout}\index{request-timeout@\texttt{-\kern0pt-request-timeout}}). After 10~s in the queue its stream opens and an SSE comment\index{SSE comment} goes out every 10~s:

\begin{plain}
: waiting, position 1, about 40 s, 30 s so far
\end{plain}

A comment is not an event, so no client reads it as a token, and a client or proxy with an idle-read timeout keeps the connection. With the profile's progress lines on, the waiting line carries the same estimate (Chapter~\ref{ch:lines}). The estimate (\code{server/queue_wait.py}) reads the work of the request that holds the engine: the prefill it still has to read at its live rate, plus the tokens it is likely to still write, the 75th percentile of what earlier requests of its kind (thinking or not) wrote past the tokens it has written, never past its \code{max_tokens}, at its own decode rate. Each waiter ahead adds its typical time. Replayed over 3,716 requests of the reference box, with the waiter arriving at 10 to 90\,\% of each hold, the estimate for waits of 10 to 60~s was a median 0.52 of the real wait when it used only the hold times of earlier requests, and is 1.06 now; for waits over a minute 0.17 before and 0.59 now (\measured). Without a view of the holder, or with too few finished requests of its kind, the earlier rule stays: the median remaining hold time of requests that ran as long, and with too few of those the position only. A request without a stream cannot show that it is alive, so it keeps \code{--queue-timeout}\index{queue-timeout@\texttt{-\kern0pt-queue-timeout}} and its 429. A full queue is still refused at once with a 503. A request that has waited less than 2~s gets no line. When TandemLLM's own work holds the engine (the shadow reader, a summary prepared ahead), the line says so in plain words, and since that work yields at its next chunk, its estimate is the median of the waits requests really had behind it; the ledger's \code{wait_own_ms} records each request's time behind it.

Two full turns are not interleaved on one sequence. Serving a second turn between the first one's blocks would mean prefilling its whole prompt and copying every KV row it can reach, each time it gets a slice; for a full turn that would cost more than it saves and hold two conversations' state at once. The client is told how long the wait is instead, unless the engine holds more than one sequence (Section~\ref{sec:parallel-requests}).

\section{Parallel requests}\label{sec:parallel-requests}\index{parallel requests}

With \code{--parallel N}\index{parallel@\texttt{-\kern0pt-parallel}} (\code{QSE_PARALLEL}; the flag's default is 1, and the shipped profile runs 8 lanes) the engine builds N~-~1 side slots beside itself. Each is one sequence's complete state over the one copy of the weights: its own KV, recurrent state, verify graphs and drafter stack, as long as slot~0's window (\code{--side-max-len} 0; a shorter value is kept). Slot~0 keeps the full window and every cache. A request takes a side slot when slot~0 is busy, its prompt and \code{max_tokens} fit, and it does not continue slot~0's resident conversation; everything else about it is served by the same code. Whether a request continues that conversation is read off the client's own tokens: the server keeps the last slot-0 prompt whose rows stayed resident as the client's messages render it (no router note, no generated row), and a request whose rendered prompt starts with all of it, within one prefill chunk, waits for slot~0 and resumes in place. The resident rows themselves would not do: under a conductor plan they hold the router's notes and the end of the last generation, so an agent's next turn would never match and would read its whole conversation again on a side slot. TandemLLM's own work on slot~0 (the shadow reader, a summary being prepared, the memory writer, the router's question) gives way to a waiting request, so it never sends one to a side slot (a request waits for it up to 15~s). Each side slot keeps the conversation it served last, as slot~0 does (FX16, SRV-528). The next turn of that conversation goes back to its slot and resumes there at its end; while another request holds that slot it waits up to \code{QWEN38_SIDE_HOME_WAIT_S} (10~s), then takes another slot, where its rows are carried in or read again. A conversation the KV pool's reclaim wrote to the NVMe from a side slot is read back on that slot when the slot holds no other conversation (SRV-779, Section~\ref{sec:spill}). Among free side slots of the same window, a new conversation takes the one served longest ago (\code{QWEN38_SIDE_LRU}). In 0.3.46-test, before side slots kept a conversation (0.4.8-test), two long agent sessions mostly took turns on slot~0: measured then with two opencode sessions started together (\measured, single runs), both were done in 301~s taking turns, against 490~s with one of them on the side slot and 193~s one after the other. Sharing a client's head is not continuing a conversation: Open WebUI's system prompt and tool block (about 6,700 tokens) and opencode's (about 7,400) are the same for every chat, so a request continues slot~0's conversation only when the part it shares reaches past the end of its own first user message (\code{QWEN38_SIDE_CONV}). Without that rule a second Open WebUI tab asking a one-line question beside a short essay chat shared 6,675 of slot~0's 7,096 rows and waited 35.8~s with the side slot idle; with it, it ran on the side slot at once (TTFT 2.6~s, \measured, n~=~1). A side slot reads slot~0's disk tier, so that tab started from 6,144 tokens of the client's head on the NVMe (Section~\ref{sec:disk-tier}). One kind of request takes an idle side slot even when slot~0 is free: a guest of slot~0's conversation, short (at most \code{--resident-guest-max} tokens and under a quarter of it), sharing less than one prefill chunk with it, while slot~0 holds two chunks or more. The router's own questions before a new user message (its class and its package, 150 to 450 tokens) and a client's title, tags and follow-up tasks are such requests: on a side slot they cost the same prefill and leave slot~0's rows alone (\code{QWEN38_SIDE_GUESTS}, on). The side slots' buffers stay in the engine's memory limits whatever value a residency document gives the ceiling and the peak (\code{/health} \code{cache.memory.fixed}); in the test builds before this release the first request's document set them back and the caches lost the slot's 7.65~GiB, and the router's questions, refused their stash, cut the conversation's rows: replayed with one build (\measured, n~=~1 each), two new user messages after a 32,875-token conversation read 0\,\% from the cache (67.5 and 70.7~s) with the bytes lost and 99.9 and 99.3\,\% (6.4 and 9.4~s) with them. The \code{[parallel]} log line names the slot and why; the usage ledger's \code{slot} column records it.

The GPU is shared in turns. Two or more decoding requests draft in one packed pass over the drafter and verify their trees in one joint pass: every projection runs once over all their rows, so the weights are read once, and each request's tree is cut to its share so the pass stays at or under 32 rows, the range where a row's result does not depend on its neighbours. A request alone keeps its own graphed verify. Admission by bytes keeps MemAvailable at or above \code{--parallel-floor-gb} (15): a request waits while its own prefill workspace, the other requests' and what the caches may still grow into, together at most the engine's room to \code{--mem-peak-gib}, would take it under the floor. Every cache and every prefill passes the memory policy, so the engine never takes more than that room. The wait is queue time (\code{wait_mem_ms} in the ledger) and the client gets a waiting line. A 26,236- and a 17,582-token prompt prefilled and decoded at once with MemAvailable at 28.7~GiB or more (\measured).

Measured on one DGX Spark (\measured, single runs, four long-answer prompts of 384 tokens, store off): one request at a time 48.8~tok/s; two at once 67.4 in total (1.38$\times$), four 69.2 (1.42$\times$; 76.5 with \code{--parallel-cut joint}). With the flag on, one request at a time costs nothing. A side slot of 8,192 tokens takes 6.9 to 7.5~GiB beyond its window, most of it the pools of its 121 verify graphs.

\paragraph{Beside the small models.} A box that also runs Laya and the narrator has less room. On a 128~GB DGX Spark at the 262,144-token window with Laya-english, the narrator (Qwen3.5-2B) and the shadow reader, MemAvailable at the worst moment with one sequence is about 22.7~GiB (\measured). Four 8,192-token slots (about 25~GiB) do not fit; two do, with one side slot of 65,536 tokens and \code{--side-graphs off}\index{side-graphs@\texttt{-\kern0pt-side-graphs}}: the slot then holds no verify graphs and verifies eager when it decodes alone (a joint pass is eager either way), about 1.5~GiB beyond its window instead of 7.5. The window matters more than the number of slots: a request takes a side slot only when its prompt, its \code{max_tokens} and the conductor's margin (8,192 tokens for a request it shapes) fit, and a chat client's \code{max_tokens} of 32,768 alone rules out an 8,192-token window. Side slots are built while the engine loads, before the residency planner loads Laya and the narrator, so the \code{residency} document's claims must be the models' measured sizes; with larger declared ones the planner refuses them. While two requests run, Laya answers only when slot~0 is idle (a decision forward takes the engine lock), and the narrator serves both requests, one line at a time.

A request beside others writes its alone text up to a near tie (Chapter~\ref{ch:exactness}): its tree is cut to its share, another shape than alone, and the shape can move a near tie. On 30 prompts 17 texts were identical to their alone text at two requests and 15 at four, every first difference at a top-2 gap under 0.12 nats. \code{tools/parallel_bench.py} measures the totals and the agreement and proves from \code{/health} that the requests decoded together.

\paragraph{Waiting for slot~0.} A request that continues slot~0's conversation prefers slot~0, but since 0.8.2 it waits there at most \code{QWEN38_SIDE_SLOT0_WAIT_S} seconds (3), then takes a free side slot its prompt fits in and prefills there what the disk tier does not hold. A request too long for every side slot still waits for slot~0. A match on the NVMe now counts as the request's conversation only when it reaches past the first user message; a match at the client's head, such as Open WebUI's system prompt and tool list, is logged as \code{spill_head=} and the request is a new conversation. On v0.8.0 such a head match held a new chat for 437~s while both side slots were free. \code{/health} \code{parallel} says how many requests may run at once (\code{max_running} with \code{max_running_reason}) and counts \code{spill_waits} and \code{slot0_gave_up}; the Lanes view's queued line says why a request waits, for example \enquote{Queued: all 3 slots are busy}.

\paragraph{Images on a side slot.} Since 0.12.1 a request with images takes a side slot the way a text request does (\code{QWEN38_SIDE_IMAGES}, on; \code{off} keeps images on slot~0, as 0.12.0). Each side slot is its own engine over the same weights, with its own image context and rotary offset. The router expands each image to its rows on the host and keys them by the image's digest, as the caches do, and the window, the output cut and the memory admission count those rows: a screenshot of 1,280 by 800 pixels is 1,000 rows. On the 27B, one image request wrote the same 300 tokens on slot~1 as twice on slot~0 (\measured). Super VL's tower stays on slot~0.

\section{Waiting for memory}\label{sec:memory-wait}\index{waiting for memory}\index{memory wait}

With more than one lane, a request whose bytes do not fit above the floor waits for memory before its prefill. The server pauses it and does not refuse it: the request keeps its place in the line, and its prefill starts by itself once memory frees. The client sends nothing again. The floors are those of Section~\ref{sec:parallel-requests}: \code{--parallel-floor-gb} (15) for the whole request, and while its prefill runs the KV pool's \code{--kv-hard-floor-gb} (10). Each floor holds the largest prefill chunk in flight or of this turn, one KV commit, a chat sandbox's last ask and \code{QSE_GUARD_MARGIN_GB} (5), at most the configured value.

A streamed chat request gets its stream's headers at once, with \code{X-Tandem-Wait: memory}, and an SSE comment every 15~s:

\begin{plain}
: waiting for memory, position <n>, <x> GiB short, <s> s so far, refused in <s> s unless memory frees
\end{plain}

The live view and the Lanes view show the request as \code{waiting_memory} with the server's reason and when the wait ends (TandemApp's dashboard); the live view's \code{activity.memory} has the same numbers as fields. The engine log prints \code{[parallel] <id> waits for memory} with the check's terms, again every 15~s, and \code{admitted after <s> s} at the end; the ledger books the wait as \code{wait_mem_ms}. The wait does not count toward \code{--request-timeout}.

\paragraph{The line.}\index{memory priority} Requests that wait stand in line by their chat's priority, then by arrival. The priority is \code{tandem.priority} in the request body, \code{"high"}, \code{"normal"} or \code{"low"}; any other value, or none, is normal. For every 30 minutes a request has waited it counts one step higher, so a low chat still gets its turn. A new request does not take the bytes a higher-priority waiter is held for, and the KV pool's reclaim takes the idle conversations of lower-priority chats first, after the older copies of a chat. A lane that gave its memory back in the middle of its answer goes to the head of the line. TandemApp's Chat tab sets the priority per chat.

\paragraph{How long.} A streamed request waits up to \code{QSE_MEM_WAIT_S}\index{qsememwaits@\texttt{QSE\_MEM\_WAIT\_S}} (10,800~s, 3~hours). A request without a stream waits up to \code{--queue-timeout} (120~s; the served profile 240), unless it asks for the long wait with \code{wait_for_memory: true} in the body or in its \code{tandem} object, or with the header \code{X-Tandem-Wait-For-Memory: 1}. At the bound the server checks once more, then answers 503 with \code{Retry-After: 30}, type \code{server_busy} and code \code{memory_wait_timeout}; the message says what the turn needed and what was free above the guard. A stream that opened while it waited ends with the same error as an event. A planned stop ends every wait at once with a 503 to send again.

\paragraph{What a turn is charged (0.15.4, SRV-801).} The admission reserves the workspace of the rows the prefill forwards and the KV pages its prompt and up to \code{QSE_KV_ADMIT_ROWS} (8,192) of its \code{max_tokens} may still commit. It counts the rows from where the prompt resumes on its slot: the rows in place, the end anchor or the conversation's entry on the NVMe. The \code{[req]} line has the charge as \code{ws_rows} and \code{ws_gib} and the rows forwarded as \code{fwd}. On 0.15.3 slot~0 had no share of its own conversation, so a turn of the chat it held was charged a cold read of the whole chat; on a user's box such a turn was charged 212,938 rows (10.75~GiB), waited 148.6~s and then forwarded 449. Since 0.15.4:

\begin{itemize}
\item slot~0 counts the share of the conversation it holds, as a side slot does;
\item the price is the served prompt's, with the router's notes and the replayed spans in it;
\item a regenerate or an edit asks the slot's rows where its prefill will resume and pays the rows from there (with images it cannot ask them, takes no share and pays more);
\item a next turn whose copy leaves out the last answer's reasoning is charged from the end anchor, so the last answer no longer counts twice;
\item the floor a prefill's pages commit under holds the chunk the turn forwards, not a whole 4,096-row chunk, unless the rows are not known;
\item before the prefill's first forward the server compares the rows it forwards with the price. When they need more workspace it raises the request's own reserve at the floor (\code{[memory] <id> forwards <n> rows, priced at <m>}). When the floor does not hold the raise, the lane gives way before the forward allocates anything: a streamed one waits and reads again, one not streamed that runs alone ends with \code{memory_full}.
\end{itemize}

A hand-back that frees nothing during a wait now prints one line, then a count once a minute and at the wait's end (with \code{--verbose}); before, each one-second retry printed eight.

Measured on one DGX Spark (\measured, 2026-11-10, Flash-Next 3.05, 8 lanes, a synthetic chat of about 203,000 tokens, then short turns): every request line had \code{ws_rows} equal to \code{fwd}. Follow-ups forwarded 931 to 1,059 rows, a regenerate 7,281 and an edit 2,737. No turn waited in that run, so a wait at the new price is not measured yet. A cold read took more host memory than its reserve, 0.37~GiB more for a 3,412-row prompt and 0.44~GiB for a 203,415-row one; the cause is not known (SRV-805, open).

\section{A second model behind the same server}\label{sec:serving-flashnext}\index{Flash-Next!serving}

\code{--model} with an EXL3 pack of Qwen3.8-Flash-Next serves that model through this server. The server reads the checkpoint's own configuration, finds its architecture in the registry (\code{engine/models}), and calls that entry's builder; it has no branch on the model's name (\code{tests/test_load_shared.py} fails on one). Everything in this chapter is the same code for every model: the request fields, reasoning formats and budgets, tool calls and grammar masks, penalties and \code{logit_bias}, logprobs, seeds, stops, streaming, the queue and the memory gate, the live view and the usage rows. Through the server, 19 of the 20 requests of \code{tools/fn_server_check.py} gave the same text, tool calls and logprobs with and without speculation; the 20th, a think budget at low effort, closes at a block boundary and differs after it (Section~\ref{sec:fn-exact}).

A model's capabilities say what the server may ask beyond the request loop (\code{caps}). Flash-Next holds every cache, the KV pool and side slots, so parallel requests run on it without changing any answer, and its profiles turn them on with three slots (Chapter~\ref{ch:flashnext}). Images go through the same path: Flash-Next's builder hands the server its own tower on the EXL3 kernels (VIS-97), and its profiles run with \code{VISION=on}. One thing is off on Flash-Next: the conductor's shadow reader, because its outputs were validated on the 27B's reasoning only; its status says why. \code{/health} names the pack, its draft layer, the load time and the self-check.

Since 0.8.2 the reference box serves Flash-Next with the conductor on, as it serves the 27B. Asked who it is, the model then answers that it is TandemLLM running the served model on the user's own machine: the conductor's identity note names the model from the pack's card. With the conductor off, v0.8.0 named another assistant in 20 of 20 answers. The Flash-Next profiles also load Laya, because the memory features ask a decision model whether a message means remember, forget or recall (about 20~ms); without it memory did not work on Flash-Next.

\section{The tool-argument loop guard}\label{sec:argloop}\index{tool-argument loop guard}

A model can loop inside a tool call, where the reasoning guards do not look. In one recorded case the model wrote the same 97-character row of a tile map 12,866 times inside a \code{write} call: 65,536 tokens, 327~s, nothing written, and the cut text made the next prompt 157,806 tokens. The lookup drafter copied the loop: all 31 draft tokens were accepted in 2,007 of its 2,209 blocks. The guard (\code{--argloop on|log|off}, \code{server/argloop.py}) reads the open call's arguments every 64 tokens once the block has 1,024 and fires on either of two rules. The periodic rule: the end of the block is one unit of at most 4,096 characters repeated exactly, 6 times over at least 1,024 tokens, or, for a unit of 256 characters or more, 3 times over at least 384 tokens. The size rule: the call passed 32,768 tokens and three times the request's longest message, and its last 16~KB compress 16 times or more.

On the tile-map incident it fires 2,176 tokens into the call, 1,055 after the loop began. On a recorded Open WebUI call that looped on a 405-character piece of a sprite, it fires 448 tokens after the loop began, at three copies; the six-copy rule alone needed 1,024. On 4,780 real writes and edits from an opencode database it fires on none, with the long-unit rule or without it (\measured, \code{tools/argloop_fp.py}). Three identical frames of a sprite written in a row do not fire; a literal block of one line repeated over more than 1,024 tokens (600 zeros written out) does, with either rule. When it fires, the generation stops, the client gets one sentence and \code{finish_reason: "stop"}, and the session's next request with tools carries a hidden note naming the tool and the repeated piece: write it without the repetition, generate repeated rows with code, write a large file in parts. A call that already streamed part of its arguments keeps what was sent, and the cut text is never sent again as content.

\paragraph{A stopped call sent back.} Open WebUI's tool loop answers a call whose arguments never closed with a tool error and sends the conversation again, the stopped arguments included. In the recorded case the model copied its own loop five times, byte for byte, about 55~s each, the prompt growing from 17,692 to 46,419 tokens. Now an assistant call in the history whose arguments do not parse (or a value of opencode's \code{invalid} stand-in) and end in a periodic stretch is cut to its first 1,500 characters and one marker before anything reads the messages; a call that parsed is never touched. The markers after the last user message are the retry streak: the retry's note says when the same loop came again, the router thinks one level higher on it (at least medium) unless the client turned thinking off, and after two identical stops in a row, or three of one tool, the engine does not generate again. The client gets one sentence and no call, and its tool loop ends.

\section{Logs and privacy}\index{privacy}

Each request leaves one log line (Chapter~\ref{ch:operations}), and no line carries the text of a prompt or an answer. The usage ledger (Chapter~\ref{ch:ref-metrics}) stores counts and timings, never text. The conductor's journal keeps the hidden texts the conductor itself wrote, for 30 days, and its transcripts keep a session's tool turns with secrets redacted, under the engine's home at mode 0600.

\input{chapters/measurement}
\chapter{Adding a model}\label{ch:adding-model}

The engine serves big models such as Qwen3.8-27B, a hybrid of Gated DeltaNet and attention layers, and Qwen3.8-Flash-Next, a mixture of experts on the same hybrid (Chapter~\ref{ch:flashnext}). The small models beside them (the Qwen3.5 narrator and shadow reader) run on the 27B's forward code. Flash-Next was the first model added through the adapter path below. This chapter lists the seams where a new model plugs in, what each one expects, and the tests a new model must pass before its numbers mean anything. Some of the engine is still specific to this family, and the last section says which parts.

\section{The adapter path}\label{sec:adapter-path}\index{model adapter}\index{ModelState}

A new architecture brings a package under \code{engine/models/<name>/}, which registers itself in its own \code{arch.py}, and data files, and no edit anywhere else. The methods read a model only through its adapter. The entry says, as data, the form the model writes tool calls in and how it reasons (\code{ServeSpec}), the files it shares with other families or builds on (\code{shares}, \code{uses}) and the kernel options it turns on (\code{EngineHooks}); its own knobs live in its \code{knobs.py}, and its own C++ or CUDA kernels in its \code{csrc/} folder (one subfolder a module that compiles them). \code{tests/test_adapter_guards.py} holds every family's fingerprint to its golden, checks that a change to one package moves only its own and its declared dependents', and fails when a shared module names a model, a family or a tool-call form; \code{tests/test_load_shared.py} fails if code under \code{server/} asks an engine what it is. Five layers depend only downward (Figure~\ref{fig:engine-layers}). The upper three exist, the formats as today's \code{Linear} protocol, and so does the server's one loader. Four pieces are still design: format plugins (ENG-253), a kernel registry (ENG-251), device profiles (ENG-255) and the mesh (ENG-259).

\begin{figure}[htbp]
\centering
\begin{tikzpicture}[node distance=4mm]
  \node[hot, minimum width=110mm] (m) {methods, written once\\[-1pt]{\scriptsize StairCut and its prices, the MTP drafter, the walk, every cache tier, the server, the governor}};
  \node[box, minimum width=110mm, below=of m] (a) {model adapters, one per architecture\\[-1pt]{\scriptsize \texttt{qwen3\_5}: the 27B \quad \texttt{qwen4\_exp}: Flash-Next \quad hot and cold verbs, \texttt{ModelState}}};
  \node[box, minimum width=110mm, below=of a] (f) {weight formats\\[-1pt]{\scriptsize the \texttt{Linear} protocol today; plugins with \texttt{read} and \texttt{bind} planned}};
  \node[quiet, minimum width=110mm, below=of f] (h) {hardware, planned\\[-1pt]{\scriptsize \texttt{KernelImpl} keyed by op, format and device; device profiles}};
  \node[quiet, minimum width=110mm, below=of h] (t) {topology, planned\\[-1pt]{\scriptsize \texttt{Mesh} at world 1: no collective on the code path}};
  \draw[flow] (m) -- (a); \draw[flow] (a) -- (f); \draw[flow] (f) -- (h); \draw[flow] (h) -- (t);
\end{tikzpicture}
\caption[The engine's layers]{The engine's layers. A layer reads only the layers under it, and per-model numbers are data the adapter returns or files under \texttt{models/} and \texttt{ops/}. Filled: built; outlined: planned.}\label{fig:engine-layers}
\end{figure}

An adapter is the engine object with two sets of verbs. The hot verbs run every round and are the engine's own bound methods, never a wrapper, so graphs capture the same launches: \code{forward}, \code{forward_tree}, \code{forward_block}, \code{commit_tree}, \code{rollback_to}, \code{picks}, \code{hidden_post_norm}, \code{kv.length} and \code{reset}. The cold verbs run around a cache copy or at load: \code{caps} (what the server may ask: caches, side slots, the cache grid, the shadow reader, joint verify), \code{cost()} (state bytes and the prefill workspace), \code{cache_state()}, \code{resume_at(n)}, \code{settle()}, \code{discard_pending()}, \code{trim_pending(n)}, \code{manifest()} (what its stored rows depend on, for the disk tier's fingerprint), \code{memory_parts()}, \code{resident_files()} and \code{precapture()}. The 27B's cold verbs are installed on its class by \code{engine/models/qwen3_5/adapter.py}, so \code{engine/model.py} keeps its bytes.

\code{cache_state()} returns a \code{ModelState}: \emph{sections}, the recurrent state by name and copied whole, and \emph{views}, the position-indexed buffers with their token axis and sliced by length. Every tier keeps a sequence in that form, and a signature over names, shapes and dtypes refuses a restore across models.

A new model then needs, in order: a config reader that refuses unknown fields; a reader for its checkpoint and a \code{Linear} object for each weight; a reference forward checked layer by layer against the published implementation; the engine with the hot verbs, every verified row equal to the decode step's row; the cold verbs; a draft head if the checkpoint ships one (the shared MTP drafter asks for \code{depth}, \code{reset}, \code{prime}, \code{capture_all}, \code{launch}, \code{read}, \code{views} and \code{resume}); a server builder \code{serve(a, ctx)}; the registry entry; and its data files: a StairCut price file per profile and device, the weight profiles (\code{ops/profiles/<name>.env}) and the Setup Guide's registry entry (\code{models/<id>/model.yaml}, described in \code{docs/setup.md}).

\section{The seams}

\begin{figure}[htbp]
\centering
\begin{tikzpicture}
  \node[quiet, minimum width=24mm] (env) at (0,1.5) {environment};
  \node[quiet, minimum width=24mm] (cfg) at (0,0) {\texttt{config.json}};
  \node[quiet, minimum width=24mm] (ck) at (0,-1.5) {checkpoint files};
  \node[box, minimum width=28mm] (set) at (3.6,1.5) {\texttt{EngineSettings}\\[-1pt]{\scriptsize engine/settings.py}};
  \node[box, minimum width=28mm] (tc) at (3.6,0) {\texttt{TextConfig}\\[-1pt]{\scriptsize engine/config.py}};
  \node[box, minimum width=28mm] (lay) at (3.6,-1.5) {\texttt{Layout}\\[-1pt]{\scriptsize engine/loader.py}};
  \node[box, minimum width=26mm] (lin) at (7.2,-1.5) {\texttt{Linear}\\[-1pt]{\scriptsize engine/linear.py}};
  \node[hot, minimum width=28mm] (model) at (10.6,0) {target model\\[-1pt]{\scriptsize engine/model.py}};
  \node[box, minimum width=28mm] (dr) at (10.6,-2.8) {drafters declare\\[-1pt]{\scriptsize what they read}};
  \draw[flow] (env) -- (set);
  \draw[flow] (cfg) -- (tc);
  \draw[flow] (ck) -- (lay);
  \draw[flow] (lay) -- (lin);
  \draw[flow] (set.east) -| ([xshift=-4mm]model.north);
  \draw[flow] (tc) -- (model);
  \draw[flow] (lin.east) -| ([xshift=-6mm]model.south);
  \draw[flow] ([xshift=6mm]model.south) -- ([xshift=6mm]dr.north);
\end{tikzpicture}
\caption{The seams a new model plugs into.}\label{fig:seams}
\end{figure}

\paragraph{Config.} \code{engine/config.py::load_config} reads \code{config.json} into a \code{TextConfig} with every size the engine needs, from the layer count to the rotary settings, and answers which layers are recurrent. It accepts a vision-language config (the text model under \code{text_config}) and a text-only one. A new family either fits \code{TextConfig} or extends it; the rest of the engine asks \code{TextConfig}, never the raw JSON.

\paragraph{Layout.} \code{engine/loader.py::Layout} says how a checkpoint names its tensors, in one place. The engine addresses every tensor by a canonical name, the language model's own naming with the checkpoint's wrapper prefix taken off. The layout lists the prefixes to strip, the tensors the language model never reads (such as a vision tower), where the checkpoint keeps its prediction head, and the projections that may be stored as FP8 pairs or as plain matrices. \code{tests/test_layout.py} pins the served checkpoint's key set: 1,250 language-model keys read, 333 vision keys skipped.

\paragraph{Weight formats.} Every weight the engine multiplies by implements one small interface (\code{engine/linear.py}): a shape, the bytes read per use, and \code{matmul}. \code{FP8Block}, \code{NVFP4Block}, \code{FP8Head} and \code{BF16Block} implement it. The model calls \code{w.matmul(x)} and never asks which format it holds. A new format is a new class with those three members, plus a kernel that keeps the row-invariance rule (Chapter~\ref{ch:kernels}).

\paragraph{Settings.} \code{engine/settings.py} is the registry of every \code{QWEN38_*} switch, read only through \code{SETTINGS.get(name)}. \code{EngineSettings(env, prefix)} gives the same reads over another mapping or another prefix.

\paragraph{What a drafter reads.} A drafter reads parts of the target: hidden states at some layers, the embedding, the head. Each drafter declares these through \code{Drafter.requires()}: the hidden size, the tap layers, the target tensors it reads by name, and the vocabulary size. At load the engine compares them with the target and refuses a mismatch by name, so a DFlash2 drafter trained for another model fails there instead of drafting garbage. A new model needs its own block drafter; the lookup drafter works for any model once its corpus is rebuilt with the new tokenizer (\code{tools/build_corpus.py}). The lookup stores record the tokenizer's sha256. Under another tokenizer a read skips the store with one log line, and a write refuses it.

\paragraph{The tokenizer and the chat template.} The server renders prompts with the model's own chat template and reads the model's own tool-call and reasoning formats. A model with another format needs its own parser behind the same OpenAI output.

\section{The tests a new model must pass}

In this order, each on the new model:

\begin{enumerate}
\item \code{tools/refcheck.py}: argmax agreement with the published reference implementation on teacher-forced positions, the reference in its own process because only one model fits on the board at a time. On Qwen3.8 the engine agrees on 99.25~\% of 400 positions and on all 242 confident ones (\measured).
\item \code{tools/quality_gate.py} for every quantised weight set (Chapter~\ref{ch:quantisation}).
\item \code{tools/verify_spec.py}\index{tools/verifyspec.py@\texttt{tools/verify\_spec.py}}, the agreement gate with its near-tie rule, with every drafter (Chapter~\ref{ch:exactness}).
\item The CPU suite, including \code{tests/test_refactor_gate.py}. That test records the exact bytes of a tiny model's run (every logit and every byte of state after a prefill, a decode, a block verify, a rollback, a tree verify and a commit) and fails if any byte moves. A change that is supposed to change arithmetic re-records the fixture and says why.
\item \code{ops/gate.sh}\index{ops/gate.sh@\texttt{ops/gate.sh}} and the client smoke set before any number is published (Chapter~\ref{ch:measurement}).
\end{enumerate}

Flash-Next took its own form of these gates: the one loader (\code{tests/test_load_shared.py}), every cache tier against a cold prefill (\code{tests/test_cache_any_model.py} and its server check), a fingerprint that moves with its kernels and tables and with nothing else (\code{tests/test_fp_v3.py}), the load-time self-check of every verify width (3,936 rows bit for bit on the 3.05 pack), and speculative against plain decoding in the engine and through the server (\code{tools/fn_spec_check.py}, \code{tools/fn_server_check.py --compare}).

\section{What is still specific to Qwen3.8}

\begin{itemize}
\item \code{engine/model.py} implements the 27B's forward and \code{engine/models/flashnext/} Flash-Next's. A new family needs its own layer code behind the same verify, commit and cache interfaces; Flash-Next routes its experts inside the verify, and every row reads its own experts.
\item The recurrent-layer kernels assume Gated DeltaNet's state shape, heads of 128 by 128.
\item Tile tables and router prices are measured for each model's shapes on the GB10. A new model or a new board needs its own sweep and its own price file; the kernel registry and device profiles that would make a board data are planned, not built.
\item Some modules read their switches into globals at import, so two engines with different settings cannot share one process yet.
\item Runtime names still carry the old project name, from the \code{QWEN38_} and \code{QSE_} prefixes to the \code{Qwen38Engine} class.
\end{itemize}

The Qwen3.5 small models already run on the engine's own forward: the narrator and the shadow reader load them through \code{engine/generative/qwen35.py}, the same layer code, because they share Qwen3.8's architecture family and its tokenizer (TandemApp's manual, \enquote{Status lines, the voice and the shadow reader}).

\input{chapters/flashnext}
\input{chapters/lfm25}
\input{chapters/nemotron}
\input{chapters/ornith}
\input{chapters/flint}

\partintro{%
The reference. Chapters~\ref{ch:ref-cli} to \ref{ch:ref-metrics} are generated from the code: every flag, variable, route, metric and ledger column, with the value the served profile sets. Then troubleshooting by symptom (\ref{ch:troubleshooting}), questions and answers (\ref{ch:faq}), the roadmap (\ref{ch:roadmap}) and the changelog (\ref{ch:changelog}). The glossary and the index follow.}
\part{Reference}
\chapter{Command-line reference}\label{ch:ref-cli}

The tables in this part are generated from the code by \code{manual/tools/gen_reference.py}, which parses the argument parsers, the settings registry, the metric registrations and the ledger's schema with Python's \code{ast} module and never imports the project. The manual's build runs it first, so the tables follow the code they were built from. A description is the first sentence of the flag's own help text; where the code gives none, the generator carries a short one. A default is the code's; where the served profile passes another value, the cell says \enquote{served:} and that value, as \code{ops/start.sh}\index{ops/start.sh@\texttt{ops/start.sh}} builds it from \code{ops/serve.env}\index{ops/serve.env@\texttt{ops/serve.env}}.

\section{The engine's server}\index{server flags}

\code{python server/app.py [flags]}, or \code{ops/start.sh}, which builds the flags from \code{serve.env}. \code{python server/app.py --help} prints the full help texts.

% Generated by manual/tools/gen_reference.py from the code. Do not edit by hand.

\begin{reftable}{Every flag of the engine's server, server/app.py, with the value the served profile passes where it differs}{tab:ref-server-flags}{\hd{Flag} & \hd{Default} & \hd{What it does}}
\texttt{-{}-model} & \texttt{none} & the checkpoint directory of the target model. \\
\texttt{-{}-served-\allowbreak{}model} & \texttt{qwen3.8-\allowbreak{}27b-\allowbreak{}spark-\allowbreak{}engine; served:\allowbreak{} qwen38-\allowbreak{}spark-\allowbreak{}engine} & the raw engine's model name; with the conductor on it is a hidden alias of TandemLLM. \\
\texttt{-{}-host} & \texttt{127.0.0.1; served:\allowbreak{} 0.0.0.0} & the address to bind. \\
\texttt{-{}-port} & \texttt{8000} & the port to listen on. \\
\texttt{-{}-max-\allowbreak{}len} & \texttt{8192; served:\allowbreak{} 262144} & context length. The KV buffer is allocated for all of it up front, so this is a memory decision as much as a capability one: 64 KiB a token of engine KV plus 20 KiB a token per drafter arm, 27.9 GB of KV and 53.4 GB of whole ... \\
\texttt{-{}-rep-\allowbreak{}penalty} & \texttt{1} & deterministic repetition penalty on the target's logits, before the argmax (1.0 = off, today's output). \\
\texttt{-{}-presence-\allowbreak{}penalty} & \texttt{0} & flat subtraction from tokens already in the history (0 = off). \\
\texttt{-{}-frequency-\allowbreak{}penalty} & \texttt{0} & subtraction per occurrence, on tokens already in the history (0 = off). \\
\texttt{-{}-temperature} & \texttt{0} & server default sampling temperature (0 = greedy, the exact path). \\
\texttt{-{}-top-\allowbreak{}p} & \texttt{none} & server default for top\_p. \\
\texttt{-{}-top-\allowbreak{}k} & \texttt{none} & server default for top\_k (0 is off). \\
\texttt{-{}-no-\allowbreak{}repeat-\allowbreak{}ngram} & \texttt{0} & forbid the token that would complete an n-gram already seen (HF's no\_repeat\_ngram\_size; 0 = off, must be \textgreater{}= 2). \\
\texttt{-{}-pattern-\allowbreak{}stop} & \texttt{none} & end a generation that repeats one short pattern forever, as MAX:MIN:COUNT (e.g. 64:1:8, a block of 1..64 tokens repeated 8 times ends the generation with finish\_reason stop). \\
\texttt{-{}-argloop} & \texttt{on} & the tool-argument loop guard (server/api/argloop.py): \texttt{on} stops a held tool call whose arguments loop and tells the model on its next turn. One of: on, log, off. \\
\texttt{-{}-default-\allowbreak{}max-\allowbreak{}tokens} & \texttt{8192; served:\allowbreak{} 32768} & what a request that does not send max\_tokens gets. \\
\texttt{-{}-reasoning-\allowbreak{}format} & \texttt{tags} & how the reasoning block reaches the client: inside \textless{}think\textgreater{}\textless{}/think\textgreater{} in content (default, and what Open WebUI folds on), as OpenAI-style reasoning\_content deltas, or both. One of: tags, reasoning\_content, both. \\
\texttt{-{}-drafter} & \texttt{mtp} & the drafter set. lenrouter, the served one, holds both DFlash2 widths beside the lookup drafter; none decodes one token a step. One of: mtp, router, merged, dflash2, lenrouter, none. \\
\texttt{-{}-tree} & \texttt{off (a switch)} & verify a draft TREE per step instead of a chain, where the drafter builds one; see docs/speculative-decoding.md. \\
\texttt{-{}-regenerate} & \texttt{sample (QWEN38\_\allowbreak{}REGENERATE)} & : a regenerate (the same conversation again after it was answered, under a TandemLLM plan; or X-Tandem-Regenerate: 1) samples with a fresh seed; every other request keeps its own sampling (greedy by default). One of: sample, off. \\
\texttt{-{}-regenerate-\allowbreak{}temperature} & \texttt{0.7 (QWEN38\_\allowbreak{}REGEN\_\allowbreak{}TEMPERATURE)} & a regenerate's temperature (QWEN38\_REGEN\_TEMPERATURE). \\
\texttt{-{}-regenerate-\allowbreak{}top-\allowbreak{}p} & \texttt{0.95 (QWEN38\_\allowbreak{}REGEN\_\allowbreak{}TOP\_\allowbreak{}P)} & a regenerate's top-p (QWEN38\_REGEN\_TOP\_P). \\
\texttt{-{}-regenerate-\allowbreak{}top-\allowbreak{}k} & \texttt{20 (QWEN38\_\allowbreak{}REGEN\_\allowbreak{}TOP\_\allowbreak{}K)} & a regenerate's top-k (QWEN38\_REGEN\_TOP\_K). \\
\texttt{-{}-sampled-\allowbreak{}tree} & \texttt{none; served:\allowbreak{} det} & let a SAMPLED request keep the tree verify. One of: det, mixed. \\
\texttt{-{}-budget} & \texttt{16} & nodes per tree, anchor included. \\
\texttt{-{}-df2-\allowbreak{}temp} & \texttt{1} & the block drafter's temperature for its own proposals. \\
\texttt{-{}-corpus} & \texttt{none (QWEN38\_\allowbreak{}CORPUS)} & the lookup drafter's corpus: a directory of token ids with a suffix array (tools/build\_corpus.py). \\
\texttt{-{}-depth} & \texttt{3} & the drafter's depth for the older chain drafters. \\
\texttt{-{}-k} & \texttt{0} & verify block size; 0 = the drafter's depth. \\
\texttt{-{}-draft-\allowbreak{}head} & \texttt{none} & an optional separate head file for the drafter. \\
\texttt{-{}-nvfp4} & \texttt{none} & the NVFP4 weight files to load, one per target (mlp, gdn, attn); empty means the checkpoint's FP8. \\
\texttt{-{}-price-\allowbreak{}table} & \texttt{none (QSE\_\allowbreak{}PRICE\_\allowbreak{}TABLE)} & FILE[:NAME], the routers' verify/draft/rollback prices for this weight set (engine/prices.py); empty = the NVFP4 constants in the code. \\
\texttt{-{}-fp8-\allowbreak{}head} & \texttt{none} & e4m3 lm\_head from tools/quant\_head.py; halves the head's 2.54 GB in both the verify pass and the block drafter's top-k read. \\
\texttt{-{}-dflash2-\allowbreak{}blocks} & \texttt{1} & chained 8-wide draft blocks; 1 proposes 7 tokens, 2 proposes 14. \\
\texttt{-{}-dflash2-\allowbreak{}path} & \texttt{greedy} & how the block drafter walks its lattice: the greedy path or Viterbi. One of: greedy, viterbi. \\
\texttt{-{}-dflash2-\allowbreak{}ckpt} & \texttt{none} & the narrow (block 8) DFlash2 drafter checkpoint. \\
\texttt{-{}-dflash2-\allowbreak{}ckpt16} & \texttt{none} & the sixteen-wide drafter, for -\kern0pt-drafter lenrouter. \\
\texttt{-{}-len-\allowbreak{}fixed} & \texttt{0} & pin -\kern0pt-drafter lenrouter to one block width (8 or 16). 0 lets the policy choose, which means -\kern0pt-len-latch: one decision a request. 8 and 16 measure the fixed baselines through the same code. \\
\texttt{-{}-len-\allowbreak{}latch} & \texttt{on (-{}-no-\allowbreak{}len-\allowbreak{}latch turns it off)} & decide the block width ONCE a request, four wide blocks, up to four narrow probes, then no more switching. \\
\texttt{-{}-learn-\allowbreak{}cost} & \texttt{on (-{}-no-\allowbreak{}learn-\allowbreak{}cost turns it off)} & the length router learns its verify and draft costs from the blocks it times, across requests (default). -\kern0pt-no-learn-cost keeps the served price table; timings inside a request still size its blocks, so two runs can part at a near tie. \\
\texttt{-{}-len-\allowbreak{}explore} & \texttt{32} & blocks between forced wide probes when nothing suggests one. \\
\texttt{-{}-deep} & \texttt{0 (QWEN38\_\allowbreak{}DEEP)} & with -\kern0pt-drafter lenrouter -\kern0pt-tree: after -\kern0pt-deep-after wide blocks in a row commit their whole width, propose the lookup drafter's long exact continuation of this request's text as one chain of up to this many rows. 0 = off. \\
\texttt{-{}-deep-\allowbreak{}after} & \texttt{2 (QWEN38\_\allowbreak{}DEEP\_\allowbreak{}AFTER)} & full wide blocks in a row before a deep chain (2 keeps it off new text). \\
\texttt{-{}-narrow-\allowbreak{}arm} & \texttt{on} & with -\kern0pt-drafter lenrouter: build the block-8 drafter arm (on, the served default) or run the block-16 arm alone. One of: on, off. \\
\texttt{-{}-drop-\allowbreak{}idle} & \texttt{off (-{}-drop-\allowbreak{}idle turns it on); served:\allowbreak{} on} & release the arm that LOSES the latch for the rest of the request: no tap, no sync, and nothing of it in the state snapshot. \\
\texttt{-{}-relax-\allowbreak{}tau} & \texttt{1} & Breaks the agreement rule: accept a drafted token whose probability is at least this fraction of the argmax's. 1.0, the default, keeps the rule. \\
\texttt{-{}-relax-\allowbreak{}rank} & \texttt{1} & Breaks the agreement rule: accept a drafted token among the target's top r. 1, the default, keeps the rule. \\
\texttt{-{}-think-\allowbreak{}stall} & \texttt{on} & with -\kern0pt-conductor shadow or on, run the reasoning-loop arms of the policy's reasoning\_loop task (a pattern repeated over 128+ tokens, the history-wide 16-gram repeat, novelty collapse); the monitor decides whether a hit closes the ... \\
\texttt{-{}-no-\allowbreak{}think-\allowbreak{}stall} & \texttt{sets think\_\allowbreak{}stall off} & never run the reasoning-loop arms (the budget still applies). \\
\texttt{-{}-conductor} & \texttt{raw} & the second-level harness (app/tandemapp/conductor/monitor/). raw, the default (off is the same): nothing of it runs and the output is the raw engine's. shadow: every rule runs and journals, nothing is done. on: each task acts in ... One of: raw, off, shadow, on. \\
\texttt{-{}-conductor-\allowbreak{}policy} & \texttt{none} & the monitor's policy file (default app/ops/supervisor.yaml), validated at start against app/ops/supervisor.schema.json. \\
\texttt{-{}-conductor-\allowbreak{}watch-\allowbreak{}s} & \texttt{30} & with -\kern0pt-conductor shadow or on, how often the stall watchdog's watch looks for sessions whose client has gone. \\
\texttt{-{}-think-\allowbreak{}budget} & \texttt{0} & default cap on reasoning tokens per request, 0 = uncapped. \\
\texttt{-{}-answer-\allowbreak{}reserve} & \texttt{auto (QSE\_\allowbreak{}ANSWER\_\allowbreak{}RESERVE)} & tokens of the output limit kept for the answer when a client asks for a thinking budget in any shape (server/api/thinking.py): its reasoning closes early enough to leave them. auto = min(2048, a quarter of the limit), 0 = none ... \\
\texttt{-{}-reasoning-\allowbreak{}effort} & \texttt{none; served:\allowbreak{} medium} & default \texttt{reasoning\_\allowbreak{}effort} for the chat template. One of: none, off, low, medium, xhigh. \\
\texttt{-{}-cache-\allowbreak{}budget-\allowbreak{}gb} & \texttt{2} & RAM budget for the state cache, in GiB (CACHE\_GB). \\
\texttt{-{}-store-\allowbreak{}entry-\allowbreak{}gb} & \texttt{1} & GiB one state-store snapshot may take (STORE\_ENTRY\_GB), whatever budget is free: 1 GiB is about 7k tokens. \\
\texttt{-{}-no-\allowbreak{}session-\allowbreak{}cache} & \texttt{off (a switch)} & do not keep a conversation's state after its turn. \\
\texttt{-{}-no-\allowbreak{}prefix-\allowbreak{}cache} & \texttt{off (a switch)} & do not checkpoint a prefill at chunk boundaries. \\
\texttt{-{}-resident-\allowbreak{}gb} & \texttt{2.5} & GiB for the resident prefix's anchors (RESIDENT\_GB), the recurrent state (146.8 MiB) at chunk boundaries of the prompt the KV buffer holds, so the next turn of a growing conversation resumes in place, KV not copied, at any length ... \\
\texttt{-{}-resident-\allowbreak{}stash-\allowbreak{}gb} & \texttt{1} & GiB for the rows a short unrelated request between two turns may overwrite, copied aside and back (RESIDENT\_STASH\_GB; \textasciitilde{}104 kB a row: 1 GiB holds a -\kern0pt-resident-guest-max guest and its 512-row slack). \\
\texttt{-{}-resident-\allowbreak{}guest-\allowbreak{}max} & \texttt{8192} & the longest request that is a guest (RESIDENT\_GUEST\_MAX): opencode's title call is 1-3k tokens; a longer one replaces the resident prompt with its head pinned. \\
\texttt{-{}-resident-\allowbreak{}tail} & \texttt{4} & the newest anchors never evicted (the next turn resumes near the end). \\
\texttt{-{}-resident-\allowbreak{}pin-\allowbreak{}gb} & \texttt{4} & GiB for the head of the resident prompt a long unrelated request would otherwise cut (RESIDENT\_PIN\_GB): its KV rows and the anchors in it are copied aside, and a later request that shares more with that head than with the rows in ... \\
\texttt{-{}-resident-\allowbreak{}park-\allowbreak{}gb} & \texttt{4} & GiB for one other session's rows (RESIDENT\_PARK\_GB): when a request of another session (x-session-id) replaces the resident prompt, the rows past their shared head are copied aside with its newest anchors, and that session's next ... \\
\texttt{-{}-end-\allowbreak{}anchor} & \texttt{on} & the resident prefix's end anchor (END\_ANCHOR): at the end of every generation the recurrent state at its last position (147 MiB, one per resident prompt), so the next request that starts with the prompt and every token the engine ... One of: on, off. \\
\texttt{-{}-compaction} & \texttt{continue (QSE\_\allowbreak{}COMPACTION)} & a client's compaction request (QSE\_COMPACTION, app/tandemapp/compaction.py): continue = served as the continuation of the conversation the engine holds for its session (the prefill reads only the instructions); client = as the ... One of: continue, client. \\
\texttt{-{}-compaction-\allowbreak{}prepare} & \texttt{on (QSE\_\allowbreak{}COMPACTION\_\allowbreak{}PREPARE)} & write an idle opencode session's summary ahead, so its /compact is answered without the engine (QSE\_COMPACTION\_PREPARE); a request that waits stops it at the next chunk or block. One of: on, off. \\
\texttt{-{}-compaction-\allowbreak{}idle-\allowbreak{}s} & \texttt{30 (QSE\_\allowbreak{}COMPACTION\_\allowbreak{}IDLE\_\allowbreak{}S)} & seconds without a request before a summary is prepared. \\
\texttt{-{}-compaction-\allowbreak{}min-\allowbreak{}tokens} & \texttt{32768 (QSE\_\allowbreak{}COMPACTION\_\allowbreak{}MIN\_\allowbreak{}TOKENS)} & the shortest conversation a summary is prepared for. \\
\texttt{-{}-end-\allowbreak{}align} & \texttt{on} & the end anchor for a copy whose whitespace is not the rows' (END\_ALIGN, server/endalign.py): a follow-up whose copy of an earlier generation differs from what the engine wrote in whitespace alone (a tool call's trimmed parameter ... One of: on, off. \\
\texttt{-{}-mem-\allowbreak{}ceiling-\allowbreak{}gib} & \texttt{70} & GiB the engine may hold between requests, allocated plus the allocator's reserve (MEM\_CEILING\_GIB): a cache grows only while the engine stays under it less -\kern0pt-reserve-slack-gib, and over it the caches are given back, cheapest ... \\
\texttt{-{}-mem-\allowbreak{}peak-\allowbreak{}gib} & \texttt{77} & GiB the engine may hold during a prefill (MEM\_PEAK\_GIB): the ceiling plus the workspace of a 262,144-token prompt. \\
\texttt{-{}-reserve-\allowbreak{}slack-\allowbreak{}gib} & \texttt{2} & GiB of free blocks the allocator may keep between requests (RESERVE\_SLACK\_GIB); over it they go back to the driver once a request has finished, never inside one. \\
\texttt{-{}-prompt-\allowbreak{}token-\allowbreak{}cache} & \texttt{4} & millions of tokens of tokenised chat messages kept, so a growing conversation's prompt is tokenised only where it is new (server/tokcache.py). 0 = tokenise every prompt whole. \\
\texttt{-{}-session-\allowbreak{}trim} & \texttt{off} & cut the engine back to the committed tokens before the session snapshot when the last verified block ran past them (end token, stop, budget) and its commit is still pending, so the put is not declined. One of: on, off. \\
\texttt{-{}-store-\allowbreak{}boundaries} & \texttt{1} & chunk boundaries of one prefill the state store keeps: its longest ones under the per-entry cap. 0 = every boundary, where one cold prefill's nested snapshots fill the budget and evict everything else. \\
\texttt{-{}-prefill-\allowbreak{}heartbeat-\allowbreak{}s} & \texttt{5} & a streamed request whose prefill has run this long gets an SSE comment (': prefill done/total') after every chunk, at most once a second, so a client and every proxy between see the stream is alive. \\
\texttt{-{}-live-\allowbreak{}activity} & \texttt{on} & the live activity on /v1/dashboard/live (states, prefill progress, timeline, the last 20 stops) at QSE\_LIVE\_HZ events a second while busy. off = contract 1.1 with those fields null at one event a second: the kill switch, and the ... One of: on, off. \\
\texttt{-{}-max-\allowbreak{}prefill-\allowbreak{}rows} & \texttt{8192} & with the prefix cache off, forward a long prompt in chunks of this many rows instead of one call; a 131k-token single forward exhausted the board. 0 = one call. \\
\texttt{-{}-prefix-\allowbreak{}chunk} & \texttt{1024} & tokens between prefill checkpoints, and the forward size of EVERY prefill while the prefix cache is on, the two have to agree or a warm prefill is not the same arithmetic as a cold one. \\
\texttt{-{}-response-\allowbreak{}cache} & \texttt{off (a switch)} & OPT-IN. Answer an identical (prompt, params) request from memory. \\
\texttt{-{}-response-\allowbreak{}cache-\allowbreak{}mb} & \texttt{256} & the response cache's budget in MiB. \\
\texttt{-{}-response-\allowbreak{}cache-\allowbreak{}ttl} & \texttt{3600} & seconds a response cache entry is kept. \\
\texttt{-{}-suffix-\allowbreak{}store} & \texttt{\textasciitilde{}/\allowbreak{}.qwen38-\allowbreak{}spark-\allowbreak{}engine/\allowbreak{}suffix (QWEN38\_\allowbreak{}SUFFIX\_\allowbreak{}STORE)} & directory for the persistent suffix store of what this engine has read and written, which the lookup drafter reads as a second corpus. \\
\texttt{-{}-suffix-\allowbreak{}store-\allowbreak{}readonly} & \texttt{off (a switch)} & read the suffix store and never append to it: a fixed benchmark measured against a store of real traffic must not write itself into it. \\
\texttt{-{}-suffix-\allowbreak{}store-\allowbreak{}mb} & \texttt{192} & cap on the store, in MiB of int32 token ids (192 MiB = 48 M tokens). \\
\texttt{-{}-disk-\allowbreak{}cache} & \texttt{off; served:\allowbreak{} (a path)} & the disk tier: a directory on the NVMe for prompt-prefix state, prefixes prefilled more than once, a session's first long request, under the caches in RAM, or off (the default). \\
\texttt{-{}-disk-\allowbreak{}cache-\allowbreak{}min-\allowbreak{}tokens} & \texttt{4096} & the shortest prefix the disk tier writes an anchor for. \\
\texttt{-{}-conductor-\allowbreak{}defaults} & \texttt{builtin; served:\allowbreak{} shipped} & shipped, the live documents are stored as version 1 for every document the conductor's directory has no version of; builtin, the minimal built-ins stay version 0. One of: builtin, shipped. \\
\texttt{-{}-session-\allowbreak{}spill} & \texttt{on} & a live conversation on the NVMe, under the disk tier's directory: written whole at a planned stop and read back at the next start, written before another client's request evicts it, read back at its next turn. One of: on, off. \\
\texttt{-{}-session-\allowbreak{}spill-\allowbreak{}gb} & \texttt{64} & the session spill's own budget on the NVMe (a 256k conversation is about 27 GB as it is in the engine, about 16 GB compressed); the disk tier's free-space floor holds for both; settings.memory.session\_spill overrides it live. \\
\texttt{-{}-session-\allowbreak{}spill-\allowbreak{}codec} & \texttt{zstd} & new spill blocks as byte planes compressed with zstd (exact under the agreement rule, 1.71x on a real entry; blocks on the disk are read either way). One of: zstd, off. \\
\texttt{-{}-disk-\allowbreak{}write-\allowbreak{}budget-\allowbreak{}gb} & \texttt{200} & the engine's NVMe writes a local day past which optional writes stop (the tier's, and the spill's but a drain and evictions of 32k+ tokens; past 1.5x only the drain); 0 = no budget; settings.memory.disk\_writes.budget\_gb overrides ... \\
\texttt{-{}-session-\allowbreak{}spill-\allowbreak{}idle} & \texttt{on} & write every slot's conversation of 4,096+ tokens to the NVMe once it is idle (its rows stay in RAM), so a later reclaim under memory pressure only trims it; optional under the day's write budget. One of: on, off. \\
\texttt{-{}-session-\allowbreak{}spill-\allowbreak{}idle-\allowbreak{}s} & \texttt{20} & how long a conversation is idle before the idle backfill writes it. \\
\texttt{-{}-session-\allowbreak{}spill-\allowbreak{}idle-\allowbreak{}every-\allowbreak{}s} & \texttt{60} & the least time between two idle writes of one conversation. \\
\texttt{-{}-session-\allowbreak{}spill-\allowbreak{}min-\allowbreak{}tokens} & \texttt{16384} & the shortest resident conversation an eviction writes to the NVMe (shorter ones are the RAM park's). \\
\texttt{-{}-disk-\allowbreak{}cache-\allowbreak{}head-\allowbreak{}anchors} & \texttt{4} & anchors a session's first long request gets, at the grid lines below where its last user message starts (where the next session parts). \\
\texttt{-{}-suffix-\allowbreak{}store-\allowbreak{}scope} & \texttt{all} & \texttt{all} remembers prompts and answers, \texttt{outputs} only what the engine wrote. One of: all, outputs. \\
\texttt{-{}-request-\allowbreak{}timeout} & \texttt{900; served:\allowbreak{} 3600} & wall-clock cap on one generation, checked between blocks; 0 disables. \\
\texttt{-{}-max-\allowbreak{}queue} & \texttt{8} & requests allowed to wait for the engine before new ones get a 503 with Retry-After. \\
\texttt{-{}-queue-\allowbreak{}timeout} & \texttt{120; served:\allowbreak{} 240} & how long a request waits for the engine before a 429. \\
\texttt{-{}-side-\allowbreak{}lane} & \texttt{on} & a short request, a client's side request or a raw one, thinking off, within the two caps below, is served at the next safe point of the generation that holds the engine (on: between decode blocks and between prefill chunks ... One of: on, decode, off. \\
\texttt{-{}-side-\allowbreak{}lane-\allowbreak{}prompt-\allowbreak{}max} & \texttt{2048} & the longest prompt the side lane serves, in tokens. \\
\texttt{-{}-side-\allowbreak{}lane-\allowbreak{}new-\allowbreak{}max} & \texttt{64} & the most output tokens the side lane serves. \\
\texttt{-{}-parallel} & \texttt{1 (QSE\_\allowbreak{}PARALLEL)} & sequence slots: 1 (the default) is one sequence, as always. \\
\texttt{-{}-side-\allowbreak{}max-\allowbreak{}len} & \texttt{0 (QSE\_\allowbreak{}SIDE\_\allowbreak{}MAX\_\allowbreak{}LEN)} & a side slot's window in tokens; 0: -\kern0pt-max-len, every slot the same window. \\
\texttt{-{}-parallel-\allowbreak{}floor-\allowbreak{}gb} & \texttt{15 (QSE\_\allowbreak{}PARALLEL\_\allowbreak{}FLOOR\_\allowbreak{}GB)} & admission: MemAvailable less the caches' growth, the requests in flight and this request's worst case must stay at or above this; a side slot is built only if it does with the slot's bytes added. \\
\texttt{-{}-parallel-\allowbreak{}batch} & \texttt{on (QSE\_\allowbreak{}PARALLEL\_\allowbreak{}BATCH)} & on: decoding slots verify their trees in one joint pass (lockstep rounds); off: they take turns, one verify each. One of: on, off. \\
\texttt{-{}-parallel-\allowbreak{}draft} & \texttt{on (QSE\_\allowbreak{}PARALLEL\_\allowbreak{}DRAFT)} & with the joint verify on: the decoding slots' draft calls in one packed pass over the drafter (on), or one after the other (off). One of: on, off. \\
\texttt{-{}-parallel-\allowbreak{}switch-\allowbreak{}ms} & \texttt{0 (QSE\_\allowbreak{}PARALLEL\_\allowbreak{}SWITCH\_\allowbreak{}MS)} & with side slots: the interpreter's thread switch interval in ms (0 keeps Python's 5 ms); a woken slot waits up to this long for the GIL. \\
\texttt{-{}-parallel-\allowbreak{}cut} & \texttt{none (QSE\_\allowbreak{}PARALLEL\_\allowbreak{}CUT)} & how a joint pass's rows are shared: share (each slot cuts its tree to a fixed share before it submits: 16 at two slots, 10 at three, 8 at four), joint (the round cuts every slot's whole tree together, best nodes first), or auto ... One of: share, joint, auto. \\
\texttt{-{}-parallel-\allowbreak{}max-\allowbreak{}rows} & \texttt{0 (QSE\_\allowbreak{}PARALLEL\_\allowbreak{}MAX\_\allowbreak{}ROWS)} & rows of a joint pass at most; 0 (the default): the model family's (\texttt{joint\_\allowbreak{}policy}, else 32, the row-invariant kernels' range).: the 27B on EXL3 weights takes 64, each engine's joint rows bit for bit its own verify. \\
\texttt{-{}-parallel-\allowbreak{}draft-\allowbreak{}rows} & \texttt{128 (QSE\_\allowbreak{}PARALLEL\_\allowbreak{}DRAFT\_\allowbreak{}ROWS)} & rows of a packed draft pass at most: drafting decides speed only, so the drafter's pass is not held to the verify's row-invariant range; at 128, eight slots' 16-row blocks draft in one pass (32 packed two). \\
\texttt{-{}-batch-\allowbreak{}timeout-\allowbreak{}ms} & \texttt{2} & a joint round runs without a missing slot once the GPU turn has been idle this long. \\
\texttt{-{}-turn-\allowbreak{}quantum-\allowbreak{}ms} & \texttt{50} & the GPU turn goes to a waiting slot this far behind in GPU time. \\
\texttt{-{}-parallel-\allowbreak{}settle} & \texttt{on} & synchronize the device before the turn changes hands (the routers' timings then hold only their own slot's work). One of: on, off. \\
\texttt{-{}-side-\allowbreak{}precapture} & \texttt{off} & capture a side slot's verify graphs at start (on, about a minute a slot) or at their first use (off: the build takes seconds; each width's first lone round on a side slot pays its capture, a fraction of a second). One of: on, off. \\
\texttt{-{}-side-\allowbreak{}graphs} & \texttt{on (QSE\_\allowbreak{}SIDE\_\allowbreak{}GRAPHS)} & a side slot's own verify graphs: on (captured as -\kern0pt-side-precapture says; most of a slot's constant bytes) or off (the slot verifies eager when it decodes alone and holds no graph pool; a joint pass is eager either way). One of: on, off. \\
\texttt{-{}-side-\allowbreak{}slot-\allowbreak{}const-\allowbreak{}gb} & \texttt{0 (QSE\_\allowbreak{}SIDE\_\allowbreak{}SLOT\_\allowbreak{}CONST\_\allowbreak{}GB)} & a side slot's bytes beyond its window's, in the build's estimate; 0 (the default) follows -\kern0pt-side-graphs: 7.5 GiB with it on (6.9 to 7.5 GiB allocated measured for 8,192-token slots, most of it the verify graphs' pools: 121 ... \\
\texttt{-{}-kv-\allowbreak{}pool} & \texttt{off (QSE\_\allowbreak{}KV\_\allowbreak{}POOL)} & on: every slot's KV window (engine attention KV and the drafters' context KV) is reserved as virtual addresses and its pages are committed as a conversation's rows are written and given back when it ends, from one pool for every ... One of: on, off. \\
\texttt{-{}-kv-\allowbreak{}hard-\allowbreak{}floor-\allowbreak{}gb} & \texttt{10 (QSE\_\allowbreak{}KV\_\allowbreak{}HARD\_\allowbreak{}FLOOR\_\allowbreak{}GB)} & with -\kern0pt-kv-pool on: a commit never takes MemAvailable under this (GiB); the pool gives back caches and idle conversations first, else the request is not run (admission keeps requests out at -\kern0pt-parallel-floor-gb). \\
\texttt{-{}-kv-\allowbreak{}margin-\allowbreak{}rows} & \texttt{256 (QSE\_\allowbreak{}KV\_\allowbreak{}MARGIN\_\allowbreak{}ROWS)} & with -\kern0pt-kv-pool on: rows committed past every write (a masked tile's over-read stays on mapped pages). \\
\texttt{-{}-vision} & \texttt{on; served:\allowbreak{} off} & load the checkpoint's vision tower (0.92 GB bf16) and accept image\_url parts in chat requests; off refuses them with a 400. One of: on, off. \\
\texttt{-{}-image-\allowbreak{}max-\allowbreak{}mb} & \texttt{20} & largest image a request may send or link, after base64 decoding. \\
\texttt{-{}-image-\allowbreak{}max-\allowbreak{}decode-\allowbreak{}pixels} & \texttt{64000000} & largest canvas an image may declare, checked before its pixels are decoded; the processor then resizes to its own limits. \\
\texttt{-{}-image-\allowbreak{}max-\allowbreak{}pixels} & \texttt{0} & the processor's resize target cap in pixels (0: the checkpoint's preprocessor\_config.json, e.g. 16,777,216 = 16,384 rows an image). \\
\texttt{-{}-image-\allowbreak{}backend} & \texttt{torchvision (QSE\_\allowbreak{}IMAGE\_\allowbreak{}BACKEND)} & the image processor's backend, pinned so the patches do not depend on what is installed (the two differ by up to two 8-bit levels after a resize). torchvision (the default, what every served measurement ran) needs torchvision ... One of: torchvision, pil. \\
\texttt{-{}-max-\allowbreak{}images} & \texttt{16} & images one request may send. \\
\texttt{-{}-max-\allowbreak{}image-\allowbreak{}rows} & \texttt{65536} & prompt rows all of a request's images may take together (one row per 2x2 block of 16-pixel patches; a 1024x1024 image is 1024 rows). \\
\texttt{-{}-image-\allowbreak{}https} & \texttt{on} & fetch https:// image URLs (off: data: URLs only). One of: on, off. \\
\texttt{-{}-image-\allowbreak{}https-\allowbreak{}private} & \texttt{off} & let https image URLs reach loopback, private and link-local addresses (off: public hosts only). One of: on, off. \\
\texttt{-{}-image-\allowbreak{}fetch-\allowbreak{}timeout} & \texttt{15} & seconds for a whole https image download, redirects included. \\
\texttt{-{}-image-\allowbreak{}cache-\allowbreak{}mb} & \texttt{512} & encoded images kept on the device by content, so a conversation that sends an image again does not encode it again. \\
\texttt{-{}-image-\allowbreak{}act-\allowbreak{}gb} & \texttt{6} & admission: an image whose encode would need more than this (the tower's byte math, engine/vision.py) is refused with a 400. \\
\texttt{-{}-usage-\allowbreak{}default} & \texttt{on (QSE\_\allowbreak{}USAGE\_\allowbreak{}DEFAULT)} & a streamed request that sends no stream\_options gets usage, timings and metrics on its finish chunk (on, the default: Open WebUI asks for include\_usage only when a model's Usage capability is ticked), or nothing, as before (off). ... One of: on, off. \\
\texttt{-{}-usage-\allowbreak{}context} & \texttt{client (QSE\_\allowbreak{}USAGE\_\allowbreak{}CONTEXT)} & usage.prompt\_tokens counts the client's prompt (client, the default) or the full context occupied after replay (engine). One of: client, engine. \\
\texttt{-{}-usage-\allowbreak{}ledger} & \texttt{off; served:\allowbreak{} (a path)} & the SQLite usage ledger's path, or off (the default). \\
\texttt{-{}-conductor-\allowbreak{}config} & \texttt{off} & the conductor's config directory (versioned profiles, router, roles, packages and its journal, app/tandemapp/conductor/), or off (the default). \\
\texttt{-{}-backends} & \texttt{off} & the proxy's config database (app/proxy/, app/docs/proxy.md), or off (the default). \\
\texttt{-{}-host-\allowbreak{}collector} & \texttt{on} & the System view's hosts (app/proxy/hostcollect.py): this machine's memory, GPU, CPU, disks and network, and the process behind each backend on it. One of: on, off. \\
\texttt{-{}-usage-\allowbreak{}retention-\allowbreak{}days} & \texttt{400} & days the usage ledger keeps its rows (at least 400). \\
\texttt{-{}-trust-\allowbreak{}loopback} & \texttt{on (-{}-no-\allowbreak{}trust-\allowbreak{}loopback turns it off)} & a request from loopback with no X-Forwarded-For / X-Real-IP header (the box's own watchdog, row3, gate, soak) needs no token for /metrics, the full /health and the cache routes once an admin password or metrics token is set. \\
\texttt{-{}-dashboard-\allowbreak{}login} & \texttt{off (QSE\_\allowbreak{}DASHBOARD\_\allowbreak{}LOGIN)} & on, with an admin password (QSE\_ADMIN\_TOKEN) set: /v1/dashboard/* needs the admin token or a session too (the page has no sign-in screen yet). off, the default: the dashboard's read API answers without a token. One of: on, off. \\
\texttt{-{}-dashboard-\allowbreak{}dir} & \texttt{none} & where the dashboard's built files are (default: app/dashboard/dist in this repository); /dashboard/ serves them, or a placeholder when missing. \\
\texttt{-{}-setup} & \texttt{off (a switch)} & setup mode (docs/setup.md, \texttt{ops/\allowbreak{}start.sh -{}-setup}): no served model is loaded; the dashboard, /v1/setup/* and Wren, the setup guide's assistant, alone. \\
\texttt{-{}-fake-\allowbreak{}engine} & \texttt{off (a switch)} & the e2e: no weights, no GPU, a deterministic CPU token source (server/fake\_engine.py) behind the real handler, auth, ledger, metrics and logs. \\
\texttt{-{}-fake-\allowbreak{}tps} & \texttt{200} & with -\kern0pt-fake-engine: tokens a second it writes. \\
\texttt{-{}-fake-\allowbreak{}prefill-\allowbreak{}tps} & \texttt{0} & with -\kern0pt-fake-engine: the prefill rate of a FAKE\_SLOW\_PREFILL prompt (0 = 2,000 tokens a second). \\
\texttt{-{}-fake-\allowbreak{}load-\allowbreak{}s} & \texttt{0} & with -\kern0pt-fake-engine: a load of this many seconds before the port opens, writing the load's phases as the real engine does. \\
\texttt{-{}-structured-\allowbreak{}outputs} & \texttt{on (-{}-no-\allowbreak{}structured-\allowbreak{}outputs turns it off)} & serve response\_format json\_object / json\_schema and structured\_outputs (regex, choice, json) as masks over the target's rows, exact under speculation. \\
\texttt{-{}-log-\allowbreak{}content} & \texttt{off (a switch)} & let exception messages that quote a request into the log. \\
\texttt{-{}-log-\allowbreak{}request-\allowbreak{}keys} & \texttt{off (a switch)} & one [body] line per request with the parameter names and the scalar values of the non-content ones (model, stream, max\_tokens, temperature,...); never messages, prompt, tools or stop strings. \\
\texttt{-{}-verbose} & \texttt{off (a switch); served:\allowbreak{} on} & more log lines. \\
\texttt{-{}-stage} & \texttt{none} & serve the layers [LO, HI) of -\kern0pt-model as one of two stage processes: 0:k is the front stage, k:layers the tail, the one a client talks to. \\
\texttt{-{}-link-\allowbreak{}listen} & \texttt{none} & with -\kern0pt-stage: wait for the other stage on this address (an IP literal). \\
\texttt{-{}-link-\allowbreak{}peer} & \texttt{none} & with -\kern0pt-stage: connect to the other stage at this address (an IP literal). \\
\texttt{-{}-stage-\allowbreak{}graphs} & \texttt{on} & with -\kern0pt-stage: a stage whose family captures graphs for its passes captures its own (on), or runs every pass eager (off). One of: on, off. \\
\texttt{-{}-stage-\allowbreak{}caches} & \texttt{off} & with -\kern0pt-stage: keep state on both stages under one key: the prefix, session and resident caches, the disk tier and side slots; off: one sequence, nothing kept. One of: on, off. \\
\texttt{-{}-stage-\allowbreak{}snapshots-\allowbreak{}gb} & \texttt{2} & with -\kern0pt-stage-caches on, the front stage: memory for the states it keeps, a slot. \\
\texttt{-{}-stage-\allowbreak{}cache-\allowbreak{}dir} & \texttt{none} & with -\kern0pt-stage-caches on, the front stage: the directory for its half of what the tail's disk tier keeps (none: a state read from the tail's disk is rebuilt or missed). \\
\texttt{-{}-stage-\allowbreak{}cache-\allowbreak{}gb} & \texttt{32} & with -\kern0pt-stage-cache-dir: the directory's size; the files unused longest go first. \\
\texttt{-{}-link-\allowbreak{}key-\allowbreak{}file} & \texttt{none} & with -\kern0pt-stage: the key both stages share, a file of mode 0600 with at least 32 bytes. \\
\end{reftable}


\section{TandemApp}\index{tandem-app}

\code{tandem-app}, \code{tandem-app launch} and \code{tandem-app hostcollect} are TandemApp's commands: TandemApp's manual lists their flags.

\section{The service scripts}

\begin{table}[htbp]
\centering\small
\caption{The scripts in \texttt{ops/} and the installed control script.}\label{tab:ref-scripts}
\begin{tabularx}{\linewidth}{@{}L{.36\linewidth}Y@{}}
\toprule
\hd{command} & \hd{what it does} \\
\midrule
\code{bin/tandem start|stop|restart|status|smoke|logs} & the installed engine's control script (Chapter~\ref{ch:install-engine}) \\
\code{ops/start.sh} & starts the engine from \code{serve.env}, waits for \code{/health}, drops the weights' page cache, loads the roles \\
\code{ops/stop.sh [seconds]} & SIGTERM, the drain, then a kill only after the wait \\
\code{ops/watchdog.sh} & one health check, for cron every minute \\
\code{ops/hold.sh <minutes> <command>} & a command with the board to itself, then the service back \\
\code{ops/gate.sh <label> [--profile F]} & the release gate \\
\code{ops/make-secrets.sh [--admin] [--rotate]} & the metrics token; \code{--admin} adds the admin password \\
\code{ops/roles.sh load|up|down|settle} & roles beside the engine, through the residency planner \\
\code{. ops/engines.sh; engine_pids} & every engine process on the board \\
\bottomrule
\end{tabularx}
\end{table}

\section{Measurement and maintenance tools}

The \code{tools/} directory holds about 180 scripts: kernels, quantisers, probes, checks and benches. The ones this manual refers to are listed here; each prints its own \code{--help}.

\begin{table}[htbp]
\centering\small
\caption{The tools this manual refers to.}\label{tab:ref-tools}
\begin{tabularx}{\linewidth}{@{}L{.32\linewidth}Y@{}}
\toprule
\hd{tool} & \hd{what it does} \\
\midrule
\code{tools/row3.py} & the benchmark row, N runs through one server, with \code{--compare} \\
\code{tools/block_ab.py} & the in-process A/B for changes that keep the bits \\
\code{tools/forced_bench.py} & the teacher-forced bench \\
\code{tools/verify_spec.py} & the agreement gate with its near-tie rule \\
\code{tools/flagoff_identity.py} & the identity gate with the new flags off \\
\code{tools/api_text_check.py} & 13 plain requests against the current release \\
\code{tools/client_smoke.py} & the request shapes real clients send \\
\code{tools/verify_curve.py} & the verify prices StairCut cuts on \\
\code{tools/quant_nvfp4.py}, \code{quant_head.py} & the quantisers \\
\code{tools/quality_gate.py} & the quality gate \\
\code{tools/calib_corpus.py}, \code{build_corpus.py} & the calibration text and the lookup corpus \\
\code{tools/mem_audit.py} & the memory a configuration holds \\
\code{tools/disk_cache_live.py}, \code{disk_restore_check.py} & the disk tier's live and bit-exact checks \\
\code{tools/session_spill_check.py} & the session spill's bit-exact check \\
\code{tools/lane_check.py} & the side lane on the real engine \\
\code{tools/tandem_eval.py} & the conductor's evaluation suites \\
\code{tools/conductor_policy.py} & the monitor's policy schema and its check \\
\code{tools/shadow_ab.py}, \code{shadow_measure.py} & the shadow reader's cost and read rate \\
\code{tools/voice_replay.py} & a recorded session's lines through the voice \\
\code{tools/parallel_load.py} & agent-like sessions over several backends \\
\code{tools/package_check.sh} & the wheel builds and installs without torch \\
\code{manual/tools/gen_reference.py} & this part's tables \\
\bottomrule
\end{tabularx}
\end{table}

\chapter{Environment and configuration reference}\label{ch:ref-env}

TandemLLM reads four families of settings. Engine switches are \code{QWEN38_*} variables, read through one registry. Server settings are command-line flags (Chapter~\ref{ch:ref-cli}); \code{ops/start.sh}\index{ops/start.sh@\texttt{ops/start.sh}} builds them from \code{ops/serve.env}\index{ops/serve.env@\texttt{ops/serve.env}}, and a few \code{QSE_*} variables set flags' defaults or serve the scripts. The installer and the service scripts read a few \code{TANDEM_*} variables. And the conductor's behaviour lives in its versioned documents, which the dashboard edits. The \code{QWEN38_*} knobs the extension reads, and the variables of TandemApp's proxy, collector and plugins, are in TandemApp's manual.

\section{The served profile}\index{serve.env}

\code{ops/serve.env} is the served configuration. On an installed engine it is \code{~/TandemLLM/run/ops/serve.env}, rewritten by every run of the installer, and \code{local.env} beside it is read last for your own values.

% Generated by manual/tools/gen_reference.py from the code. Do not edit by hand.

\begin{reftable}{The keys ops/serve.env sets, in file order (paths not shown)}{tab:ref-serve-env}{\hd{Key} & \hd{Value} & \hd{Kind}}
\texttt{HOME\_\allowbreak{}DIR} & \texttt{(a path)} & read by ops/start.sh \\
\texttt{REPO} & \texttt{(a path)} & read by ops/start.sh \\
\texttt{PY} & \texttt{(a path)} & read by ops/start.sh \\
\texttt{PYTHONPATH} & \texttt{(a path)} & read by ops/start.sh \\
\texttt{HOST} & \texttt{0.0.0.0} & read by ops/start.sh \\
\texttt{PORT} & \texttt{8000} & read by ops/start.sh \\
\texttt{SERVED\_\allowbreak{}MODEL} & \texttt{qwen38-\allowbreak{}spark-\allowbreak{}engine} & read by ops/start.sh \\
\texttt{NV} & \texttt{(a path)} & read by ops/start.sh \\
\texttt{HEAD} & \texttt{(a path)} & read by ops/start.sh \\
\texttt{CKPT8} & \texttt{(a path)} & read by ops/start.sh \\
\texttt{CKPT16} & \texttt{(a path)} & read by ops/start.sh \\
\texttt{CORPUS} & \texttt{(a path)} & read by ops/start.sh \\
\texttt{MAX\_\allowbreak{}LEN} & \texttt{262144} & read by ops/start.sh \\
\texttt{DEFAULT\_\allowbreak{}MAX\_\allowbreak{}TOKENS} & \texttt{32768} & read by ops/start.sh \\
\texttt{PATTERN\_\allowbreak{}STOP} & \texttt{(empty)} & read by ops/start.sh \\
\texttt{QSE\_\allowbreak{}MODE} & \texttt{(empty)} & server setting \\
\texttt{CONDUCTOR} & \texttt{on} & read by ops/start.sh \\
\texttt{REASONING\_\allowbreak{}FORMAT} & \texttt{tags} & read by ops/start.sh \\
\texttt{REASONING\_\allowbreak{}EFFORT} & \texttt{medium} & read by ops/start.sh \\
\texttt{LEN\_\allowbreak{}FIXED} & \texttt{0} & read by ops/start.sh \\
\texttt{LEN\_\allowbreak{}LATCH} & \texttt{1} & read by ops/start.sh \\
\texttt{BUDGET} & \texttt{16} & read by ops/start.sh \\
\texttt{MEM\_\allowbreak{}CEILING\_\allowbreak{}GIB} & \texttt{70} & read by ops/start.sh \\
\texttt{MEM\_\allowbreak{}PEAK\_\allowbreak{}GIB} & \texttt{77} & read by ops/start.sh \\
\texttt{RESERVE\_\allowbreak{}SLACK\_\allowbreak{}GIB} & \texttt{2} & read by ops/start.sh \\
\texttt{RESIDENT\_\allowbreak{}GB} & \texttt{2.5} & read by ops/start.sh \\
\texttt{RESIDENT\_\allowbreak{}STASH\_\allowbreak{}GB} & \texttt{1.0} & read by ops/start.sh \\
\texttt{RESIDENT\_\allowbreak{}GUEST\_\allowbreak{}MAX} & \texttt{8192} & read by ops/start.sh \\
\texttt{RESIDENT\_\allowbreak{}PIN\_\allowbreak{}GB} & \texttt{4.0} & read by ops/start.sh \\
\texttt{RESIDENT\_\allowbreak{}PARK\_\allowbreak{}GB} & \texttt{4} & read by ops/start.sh \\
\texttt{CACHE\_\allowbreak{}GB} & \texttt{2} & read by ops/start.sh \\
\texttt{STORE\_\allowbreak{}ENTRY\_\allowbreak{}GB} & \texttt{1} & read by ops/start.sh \\
\texttt{NARROW\_\allowbreak{}ARM} & \texttt{1} & read by ops/start.sh \\
\texttt{END\_\allowbreak{}ANCHOR} & \texttt{on} & read by ops/start.sh \\
\texttt{END\_\allowbreak{}ALIGN} & \texttt{on} & read by ops/start.sh \\
\texttt{DROP\_\allowbreak{}PAGE\_\allowbreak{}CACHE} & \texttt{1} & read by ops/start.sh \\
\texttt{REQUEST\_\allowbreak{}TIMEOUT} & \texttt{3600} & read by ops/start.sh \\
\texttt{MAX\_\allowbreak{}QUEUE} & \texttt{8} & read by ops/start.sh \\
\texttt{QUEUE\_\allowbreak{}TIMEOUT} & \texttt{240} & read by ops/start.sh \\
\texttt{FUSE\_\allowbreak{}PROJ} & \texttt{0} & read by ops/start.sh \\
\texttt{QWEN38\_\allowbreak{}DF2\_\allowbreak{}TREE\_\allowbreak{}MODE} & \texttt{nodes} & engine switch \\
\texttt{QWEN38\_\allowbreak{}TREE\_\allowbreak{}ALIAS\_\allowbreak{}STATE} & \texttt{0} & engine switch \\
\texttt{DROP\_\allowbreak{}IDLE} & \texttt{1} & read by ops/start.sh \\
\texttt{QWEN38\_\allowbreak{}DRAFT\_\allowbreak{}TEMP} & \texttt{0.3} & engine switch \\
\texttt{SAMPLED\_\allowbreak{}TREE} & \texttt{det} & read by ops/start.sh \\
\texttt{QWEN38\_\allowbreak{}FUSED\_\allowbreak{}GDNPREFILL} & \texttt{1} & engine switch \\
\texttt{QWEN38\_\allowbreak{}DECODE\_\allowbreak{}ATTN} & \texttt{1} & engine switch \\
\texttt{QWEN38\_\allowbreak{}DRAFT\_\allowbreak{}NVFP4} & \texttt{1} & engine switch \\
\texttt{QWEN38\_\allowbreak{}NVFP4\_\allowbreak{}PREFILL\_\allowbreak{}V2} & \texttt{1} & engine switch \\
\texttt{QWEN38\_\allowbreak{}NVFP4\_\allowbreak{}PREFILL\_\allowbreak{}V2\_\allowbreak{}UNTIL} & \texttt{1025} & engine switch \\
\texttt{QWEN38\_\allowbreak{}NVFP4\_\allowbreak{}SKINNY} & \texttt{1} & engine switch \\
\texttt{QWEN38\_\allowbreak{}SKINNY\_\allowbreak{}TILES} & \texttt{(a path)} & engine switch \\
\texttt{QWEN38\_\allowbreak{}FUSED\_\allowbreak{}COMMIT} & \texttt{1} & engine switch \\
\texttt{QWEN38\_\allowbreak{}FUSED\_\allowbreak{}GDNVERIFY} & \texttt{1} & engine switch \\
\texttt{QWEN38\_\allowbreak{}DRAFT\_\allowbreak{}FC\_\allowbreak{}NVFP4} & \texttt{1} & engine switch \\
\texttt{QWEN38\_\allowbreak{}FUSED\_\allowbreak{}ADDNORM} & \texttt{1} & engine switch \\
\texttt{QWEN38\_\allowbreak{}TREE\_\allowbreak{}HOST\_\allowbreak{}DEPTH} & \texttt{1} & engine switch \\
\texttt{QWEN38\_\allowbreak{}HOST\_\allowbreak{}WALK} & \texttt{1} & engine switch \\
\texttt{QWEN38\_\allowbreak{}VERIFY\_\allowbreak{}GRAPH} & \texttt{1} & engine switch \\
\texttt{QWEN38\_\allowbreak{}COMMIT\_\allowbreak{}IN\_\allowbreak{}VERIFY} & \texttt{1} & engine switch \\
\texttt{QWEN38\_\allowbreak{}DRAFT\_\allowbreak{}GRAPH} & \texttt{1} & engine switch \\
\texttt{QWEN38\_\allowbreak{}GDNV\_\allowbreak{}ONE\_\allowbreak{}WARP} & \texttt{1} & engine switch \\
\texttt{QWEN38\_\allowbreak{}TREE\_\allowbreak{}NODES\_\allowbreak{}NARROW} & \texttt{16} & engine switch \\
\texttt{QWEN38\_\allowbreak{}TREE\_\allowbreak{}MS} & \texttt{8:\allowbreak{}77.9,\allowbreak{}16:\allowbreak{}78.9,\allowbreak{}24:\allowbreak{}87.0,\allowbreak{}32:\allowbreak{}97.0} & engine switch \\
\texttt{QWEN38\_\allowbreak{}VERIFY\_\allowbreak{}ROWS} & \texttt{32} & engine switch \\
\texttt{QWEN38\_\allowbreak{}SKINNY\_\allowbreak{}TILES\_\allowbreak{}WIDE} & \texttt{(a path)} & engine switch \\
\texttt{QWEN38\_\allowbreak{}TREE\_\allowbreak{}NODES} & \texttt{24} & engine switch \\
\texttt{QWEN38\_\allowbreak{}TREE\_\allowbreak{}WIDE\_\allowbreak{}AFTER} & \texttt{32} & engine switch \\
\texttt{QWEN38\_\allowbreak{}DEEP} & \texttt{32} & engine switch \\
\texttt{QWEN38\_\allowbreak{}DEEP\_\allowbreak{}AFTER} & \texttt{2} & engine switch \\
\texttt{QWEN38\_\allowbreak{}LEN\_\allowbreak{}SWITCH} & \texttt{1} & engine switch \\
\texttt{QWEN38\_\allowbreak{}LEN\_\allowbreak{}MODE} & \texttt{wide} & engine switch \\
\texttt{QWEN38\_\allowbreak{}STAIR\_\allowbreak{}TABLES} & \texttt{(a path)} & engine switch \\
\texttt{QWEN38\_\allowbreak{}STAIR\_\allowbreak{}RHO} & \texttt{(a path)} & engine switch \\
\texttt{QWEN38\_\allowbreak{}STAIR\_\allowbreak{}SKIP} & \texttt{1} & engine switch \\
\texttt{QWEN38\_\allowbreak{}SKINNY\_\allowbreak{}SRUN} & \texttt{1} & engine switch \\
\texttt{QSE\_\allowbreak{}USAGE\_\allowbreak{}LEDGER} & \texttt{(a path)} & server setting \\
\texttt{QSE\_\allowbreak{}CONDUCTOR\_\allowbreak{}CONFIG} & \texttt{(a path)} & server setting \\
\texttt{QSE\_\allowbreak{}BACKENDS} & \texttt{(empty)} & server setting \\
\texttt{QSE\_\allowbreak{}PARALLEL} & \texttt{8} & server setting \\
\texttt{QSE\_\allowbreak{}SIDE\_\allowbreak{}MAX\_\allowbreak{}LEN} & \texttt{0} & server setting \\
\texttt{QSE\_\allowbreak{}PARALLEL\_\allowbreak{}FLOOR\_\allowbreak{}GB} & \texttt{15} & server setting \\
\texttt{QSE\_\allowbreak{}SIDE\_\allowbreak{}GRAPHS} & \texttt{off} & server setting \\
\texttt{QSE\_\allowbreak{}SIDE\_\allowbreak{}SLOT\_\allowbreak{}CONST\_\allowbreak{}GB} & \texttt{1.5} & server setting \\
\texttt{QSE\_\allowbreak{}KV\_\allowbreak{}POOL} & \texttt{on} & server setting \\
\texttt{QSE\_\allowbreak{}KV\_\allowbreak{}HARD\_\allowbreak{}FLOOR\_\allowbreak{}GB} & \texttt{10} & server setting \\
\texttt{QSE\_\allowbreak{}KV\_\allowbreak{}RECOMPUTE} & \texttt{on} & server setting \\
\texttt{QSE\_\allowbreak{}KV\_\allowbreak{}ADMIT\_\allowbreak{}ROWS} & \texttt{8192} & server setting \\
\texttt{QSE\_\allowbreak{}BARRIER\_\allowbreak{}IDLE\_\allowbreak{}OWNER} & \texttt{on} & server setting \\
\texttt{QWEN38\_\allowbreak{}GDNV\_\allowbreak{}CONV\_\allowbreak{}BLOCK} & \texttt{32} & engine switch \\
\texttt{QWEN38\_\allowbreak{}GDNV\_\allowbreak{}CONV\_\allowbreak{}WARPS} & \texttt{1} & engine switch \\
\texttt{QWEN38\_\allowbreak{}GDNV\_\allowbreak{}WY} & \texttt{1} & engine switch \\
\texttt{QWEN38\_\allowbreak{}GDNV\_\allowbreak{}WY\_\allowbreak{}FUSED} & \texttt{1} & engine switch \\
\texttt{QWEN38\_\allowbreak{}GDNV\_\allowbreak{}WY\_\allowbreak{}MAXT} & \texttt{16} & engine switch \\
\texttt{QWEN38\_\allowbreak{}GDNV\_\allowbreak{}WY\_\allowbreak{}CHAIN\_\allowbreak{}MAXT} & \texttt{16} & engine switch \\
\texttt{QWEN38\_\allowbreak{}GDNV\_\allowbreak{}WY\_\allowbreak{}FUSED\_\allowbreak{}MAXT} & \texttt{16} & engine switch \\
\texttt{LIVE\_\allowbreak{}ACTIVITY} & \texttt{on} & read by ops/start.sh \\
\texttt{QSE\_\allowbreak{}DASHBOARD\_\allowbreak{}LOGIN} & \texttt{off} & server setting \\
\texttt{ROLES\_\allowbreak{}AUTOLOAD} & \texttt{(empty)} & read by ops/start.sh \\
\texttt{QWEN38\_\allowbreak{}DECIDER\_\allowbreak{}MODELS} & \texttt{(empty)} & engine switch \\
\texttt{QWEN38\_\allowbreak{}DECIDER\_\allowbreak{}DIR} & \texttt{(a path)} & engine switch \\
\texttt{QWEN38\_\allowbreak{}NARRATOR\_\allowbreak{}MODEL} & \texttt{(empty)} & engine switch \\
\texttt{QWEN38\_\allowbreak{}NARRATOR\_\allowbreak{}DIR} & \texttt{(empty)} & engine switch \\
\texttt{QWEN38\_\allowbreak{}SHADOW\_\allowbreak{}MODEL} & \texttt{(empty)} & engine switch \\
\texttt{QWEN38\_\allowbreak{}SHADOW\_\allowbreak{}DIR} & \texttt{(empty)} & engine switch \\
\texttt{QWEN38\_\allowbreak{}SHADOW\_\allowbreak{}GB} & \texttt{4} & engine switch \\
\texttt{QWEN38\_\allowbreak{}SHADOW\_\allowbreak{}GRID} & \texttt{512} & engine switch \\
\texttt{QWEN38\_\allowbreak{}SHADOW\_\allowbreak{}DISK} & \texttt{auto} & engine switch \\
\texttt{QWEN38\_\allowbreak{}SHADOW\_\allowbreak{}DISK\_\allowbreak{}GB} & \texttt{20} & engine switch \\
\texttt{QWEN38\_\allowbreak{}PROBE\_\allowbreak{}SNAPSHOTS} & \texttt{off} & engine switch \\
\texttt{QWEN38\_\allowbreak{}PROBE\_\allowbreak{}DIR} & \texttt{(a path)} & engine switch \\
\texttt{QSE\_\allowbreak{}DISK\_\allowbreak{}CACHE} & \texttt{(a path)} & server setting \\
\texttt{SESSION\_\allowbreak{}SPILL} & \texttt{on} & read by ops/start.sh \\
\texttt{SESSION\_\allowbreak{}SPILL\_\allowbreak{}GB} & \texttt{64} & read by ops/start.sh \\
\texttt{SESSION\_\allowbreak{}SPILL\_\allowbreak{}CODEC} & \texttt{zstd} & read by ops/start.sh \\
\texttt{DISK\_\allowbreak{}WRITE\_\allowbreak{}BUDGET\_\allowbreak{}GB} & \texttt{200} & read by ops/start.sh \\
\texttt{SESSION\_\allowbreak{}SPILL\_\allowbreak{}IDLE} & \texttt{on} & read by ops/start.sh \\
\texttt{VISION} & \texttt{on} & read by ops/start.sh \\
\texttt{RESPONSE\_\allowbreak{}CACHE} & \texttt{off} & read by ops/start.sh \\
\texttt{QSE\_\allowbreak{}STATUS\_\allowbreak{}PORT} & \texttt{(empty)} & server setting \\
\texttt{QSE\_\allowbreak{}PROFILE\_\allowbreak{}NAME} & \texttt{(a path)} & server setting \\
\texttt{QSE\_\allowbreak{}PACK\_\allowbreak{}VISION} & \texttt{(empty)} & server setting \\
\texttt{QSE\_\allowbreak{}SETUP\_\allowbreak{}ENV\_\allowbreak{}FILE} & \texttt{(a path)} & server setting \\
\texttt{QSE\_\allowbreak{}VISION\_\allowbreak{}ASKED} & \texttt{(a path)} & server setting \\
\texttt{VISION} & \texttt{off} & read by ops/start.sh \\
\texttt{QSE\_\allowbreak{}MODE} & \texttt{engine} & server setting \\
\texttt{QSE\_\allowbreak{}MODE} & \texttt{chat} & server setting \\
\texttt{QSE\_\allowbreak{}CONDUCTOR\_\allowbreak{}CONFIG} & \texttt{(a path)} & server setting \\
\texttt{QSE\_\allowbreak{}EXT} & \texttt{(a path)} & server setting \\
\texttt{QSE\_\allowbreak{}MODE} & \texttt{engine} & server setting \\
\texttt{CONDUCTOR} & \texttt{raw} & read by ops/start.sh \\
\texttt{QSE\_\allowbreak{}CONDUCTOR\_\allowbreak{}CONFIG} & \texttt{off} & server setting \\
\texttt{QWEN38\_\allowbreak{}DECIDER\_\allowbreak{}MODELS} & \texttt{(empty)} & engine switch \\
\texttt{QWEN38\_\allowbreak{}NARRATOR\_\allowbreak{}MODEL} & \texttt{(empty)} & engine switch \\
\texttt{QWEN38\_\allowbreak{}SHADOW\_\allowbreak{}MODEL} & \texttt{(empty)} & engine switch \\
\texttt{ROLES\_\allowbreak{}AUTOLOAD} & \texttt{(empty)} & read by ops/start.sh \\
\texttt{QSE\_\allowbreak{}EXT} & \texttt{(empty)} & server setting \\
\texttt{QSE\_\allowbreak{}CHATD\_\allowbreak{}ASKED} & \texttt{(a path)} & server setting \\
\texttt{QSE\_\allowbreak{}CHATD} & \texttt{off ;;} & server setting \\
\end{reftable}


\section{Engine switches}\index{QWEN38 variables}

\code{engine/settings.py} lists every engine switch with its default and the module that reads it; \code{tests/test_settings.py} fails when a module reads one another way. The table shows each switch, its code default, and the value the served profile sets when it sets one.

% Generated by manual/tools/gen_reference.py from the code. Do not edit by hand.

\begin{reftable}{Every QWEN38\_* switch in engine/settings.py, with the value ops/serve.env sets}{tab:ref-knobs}{\hd{Variable} & \hd{Code default} & \hd{Served value \textbullet{} read in}}
\texttt{QWEN38\_\allowbreak{}ATTN\_\allowbreak{}INV} & \texttt{0} & \texttt{tools/\allowbreak{}attn\_\allowbreak{}kernels.py} \\
\texttt{QWEN38\_\allowbreak{}ATTN\_\allowbreak{}INV\_\allowbreak{}NT} & \texttt{4} & \texttt{tools/\allowbreak{}attn\_\allowbreak{}kernels.py} \\
\texttt{QWEN38\_\allowbreak{}ATTN\_\allowbreak{}INV\_\allowbreak{}SPLIT} & \texttt{1024} & \texttt{tools/\allowbreak{}attn\_\allowbreak{}kernels.py} \\
\texttt{QWEN38\_\allowbreak{}ATTN\_\allowbreak{}STAGES3} & \texttt{1} & \texttt{tools/\allowbreak{}attn\_\allowbreak{}kernels.py} \\
\texttt{QWEN38\_\allowbreak{}BLOCKING\_\allowbreak{}SYNC} & \texttt{0} & \texttt{server/\allowbreak{}app.py} \\
\texttt{QWEN38\_\allowbreak{}CLIENT\_\allowbreak{}CWD} & \texttt{note} & \texttt{app/\allowbreak{}tandemapp/\allowbreak{}conductor/\allowbreak{}router.py} \\
\texttt{QWEN38\_\allowbreak{}CLIENT\_\allowbreak{}REASONING} & \texttt{*=\allowbreak{}reasoning\_\allowbreak{}content,\allowbreak{}hermes=\allowbreak{}tags} & \texttt{server/\allowbreak{}compat.py} \\
\texttt{QWEN38\_\allowbreak{}COMMIT\_\allowbreak{}IN\_\allowbreak{}VERIFY} & \texttt{0} & \texttt{1} \textbullet{} \texttt{engine/\allowbreak{}model.py} \\
\texttt{QWEN38\_\allowbreak{}CONDUCTOR\_\allowbreak{}FIT} & \texttt{model} & \texttt{app/\allowbreak{}tandemapp/\allowbreak{}conductor/\allowbreak{}modelfit.py} \\
\texttt{QWEN38\_\allowbreak{}CORPUS} & \texttt{none} & \texttt{engine/\allowbreak{}drafters/\allowbreak{}ngram.py} \\
\texttt{QWEN38\_\allowbreak{}DECIDER\_\allowbreak{}COLLATE} & \texttt{call} & \texttt{server/\allowbreak{}decider\_\allowbreak{}api.py} \\
\texttt{QWEN38\_\allowbreak{}DECIDER\_\allowbreak{}DEVICE} & \texttt{auto} & \texttt{server/\allowbreak{}decider\_\allowbreak{}api.py} \\
\texttt{QWEN38\_\allowbreak{}DECIDER\_\allowbreak{}DIR} & \texttt{\textasciitilde{}/\allowbreak{}.qwen38-\allowbreak{}spark-\allowbreak{}engine/\allowbreak{}deciders/\allowbreak{}laya} & \texttt{(a path)} \textbullet{} \texttt{server/\allowbreak{}decider\_\allowbreak{}api.py} \\
\texttt{QWEN38\_\allowbreak{}DECIDER\_\allowbreak{}MAX\_\allowbreak{}BATCH} & \texttt{64} & \texttt{server/\allowbreak{}decider\_\allowbreak{}api.py} \\
\texttt{QWEN38\_\allowbreak{}DECIDER\_\allowbreak{}MAX\_\allowbreak{}QUEUE} & \texttt{256} & \texttt{server/\allowbreak{}decider\_\allowbreak{}api.py} \\
\texttt{QWEN38\_\allowbreak{}DECIDER\_\allowbreak{}MODELS} & \texttt{none} & \texttt{server/\allowbreak{}decider\_\allowbreak{}api.py} \\
\texttt{QWEN38\_\allowbreak{}DECIDER\_\allowbreak{}WAIT\_\allowbreak{}MS} & \texttt{2} & \texttt{server/\allowbreak{}decider\_\allowbreak{}api.py} \\
\texttt{QWEN38\_\allowbreak{}DECODE\_\allowbreak{}ATTN} & \texttt{0} & \texttt{1} \textbullet{} \texttt{engine/\allowbreak{}model.py} \\
\texttt{QWEN38\_\allowbreak{}DEEP} & \texttt{0} & \texttt{32} \textbullet{} \texttt{server/\allowbreak{}app.py} \\
\texttt{QWEN38\_\allowbreak{}DEEP\_\allowbreak{}AFTER} & \texttt{2} & \texttt{2} \textbullet{} \texttt{server/\allowbreak{}app.py} \\
\texttt{QWEN38\_\allowbreak{}DF2\_\allowbreak{}DEFER\_\allowbreak{}SYNC} & \texttt{1} & \texttt{engine/\allowbreak{}drafters/\allowbreak{}dflash2.py} \\
\texttt{QWEN38\_\allowbreak{}DF2\_\allowbreak{}KV\_\allowbreak{}WINDOW} & \texttt{0} & \texttt{engine/\allowbreak{}drafters/\allowbreak{}dflash2.py} \\
\texttt{QWEN38\_\allowbreak{}DF2\_\allowbreak{}PACK\_\allowbreak{}BATCH} & \texttt{1} & \texttt{engine/\allowbreak{}drafters/\allowbreak{}dflash2.py} \\
\texttt{QWEN38\_\allowbreak{}DF2\_\allowbreak{}PACK\_\allowbreak{}FOLD} & \texttt{0} & \texttt{engine/\allowbreak{}drafters/\allowbreak{}dflash2.py} \\
\texttt{QWEN38\_\allowbreak{}DF2\_\allowbreak{}PACK\_\allowbreak{}GQA} & \texttt{0} & \texttt{engine/\allowbreak{}drafters/\allowbreak{}dflash2.py} \\
\texttt{QWEN38\_\allowbreak{}DF2\_\allowbreak{}PACK\_\allowbreak{}HOST1} & \texttt{1} & \texttt{engine/\allowbreak{}drafters/\allowbreak{}draft\_\allowbreak{}graph.py} \\
\texttt{QWEN38\_\allowbreak{}DF2\_\allowbreak{}PACK\_\allowbreak{}KP} & \texttt{1} & \texttt{engine/\allowbreak{}drafters/\allowbreak{}dflash2.py} \\
\texttt{QWEN38\_\allowbreak{}DF2\_\allowbreak{}PACK\_\allowbreak{}TAIL} & \texttt{0} & \texttt{engine/\allowbreak{}drafters/\allowbreak{}dflash2.py} \\
\texttt{QWEN38\_\allowbreak{}DF2\_\allowbreak{}PACK\_\allowbreak{}TOPK} & \texttt{0} & \texttt{engine/\allowbreak{}drafters/\allowbreak{}dflash2.py} \\
\texttt{QWEN38\_\allowbreak{}DF2\_\allowbreak{}SPAN\_\allowbreak{}STEP} & \texttt{256} & \texttt{engine/\allowbreak{}drafters/\allowbreak{}draft\_\allowbreak{}graph.py} \\
\texttt{QWEN38\_\allowbreak{}DF2\_\allowbreak{}TEMP} & \texttt{1.0} & \texttt{engine/\allowbreak{}drafters/\allowbreak{}dflash2.py} \\
\texttt{QWEN38\_\allowbreak{}DF2\_\allowbreak{}TOPK\_\allowbreak{}FAST} & \texttt{1} & \texttt{engine/\allowbreak{}drafters/\allowbreak{}topk\_\allowbreak{}small.py} \\
\texttt{QWEN38\_\allowbreak{}DF2\_\allowbreak{}TREE\_\allowbreak{}MODE} & \texttt{paths} & \texttt{nodes} \textbullet{} \texttt{engine/\allowbreak{}drafters/\allowbreak{}dflash2.py} \\
\texttt{QWEN38\_\allowbreak{}DFLASH2} & \texttt{none} & \texttt{engine/\allowbreak{}drafters/\allowbreak{}dflash2.py} \\
\texttt{QWEN38\_\allowbreak{}DRAFT\_\allowbreak{}EXL3} & \texttt{none} & \texttt{engine/\allowbreak{}drafters/\allowbreak{}dflash2.py} \\
\texttt{QWEN38\_\allowbreak{}DRAFT\_\allowbreak{}FC\_\allowbreak{}NVFP4} & \texttt{0} & \texttt{1} \textbullet{} \texttt{engine/\allowbreak{}drafters/\allowbreak{}dflash2.py} \\
\texttt{QWEN38\_\allowbreak{}DRAFT\_\allowbreak{}GRAPH} & \texttt{0} & \texttt{1} \textbullet{} \texttt{engine/\allowbreak{}drafters/\allowbreak{}dflash2.py} \\
\texttt{QWEN38\_\allowbreak{}DRAFT\_\allowbreak{}HEAD} & \texttt{none} & \texttt{engine/\allowbreak{}drafters/\allowbreak{}dflash2.py} \\
\texttt{QWEN38\_\allowbreak{}DRAFT\_\allowbreak{}HEAD\_\allowbreak{}EXL3} & \texttt{none} & \texttt{engine/\allowbreak{}drafters/\allowbreak{}dflash2.py} \\
\texttt{QWEN38\_\allowbreak{}DRAFT\_\allowbreak{}HEAD\_\allowbreak{}NVFP4} & \texttt{0} & \texttt{engine/\allowbreak{}drafters/\allowbreak{}dflash2.py} \\
\texttt{QWEN38\_\allowbreak{}DRAFT\_\allowbreak{}NVFP4} & \texttt{0} & \texttt{1} \textbullet{} \texttt{engine/\allowbreak{}drafters/\allowbreak{}dflash2.py} \\
\texttt{QWEN38\_\allowbreak{}DRAFT\_\allowbreak{}TEMP} & \texttt{1.0} & \texttt{0.3} \textbullet{} \texttt{engine/\allowbreak{}drafters/\allowbreak{}dflash2.py} \\
\texttt{QWEN38\_\allowbreak{}DSPARK} & \texttt{none} & \texttt{engine/\allowbreak{}drafters/\allowbreak{}dspark.py} \\
\texttt{QWEN38\_\allowbreak{}DSPARK\_\allowbreak{}NO\_\allowbreak{}YARN} & \texttt{none} & \texttt{engine/\allowbreak{}drafters/\allowbreak{}dspark.py} \\
\texttt{QWEN38\_\allowbreak{}END\_\allowbreak{}ALIGN\_\allowbreak{}CALLS} & \texttt{on} & \texttt{server/\allowbreak{}endalign.py} \\
\texttt{QWEN38\_\allowbreak{}END\_\allowbreak{}ALIGN\_\allowbreak{}MOVED} & \texttt{on} & \texttt{server/\allowbreak{}endalign.py} \\
\texttt{QWEN38\_\allowbreak{}END\_\allowbreak{}ALIGN\_\allowbreak{}REASONING} & \texttt{on} & \texttt{server/\allowbreak{}endalign.py} \\
\texttt{QWEN38\_\allowbreak{}END\_\allowbreak{}RESERVE} & \texttt{on} & \texttt{server/\allowbreak{}app.py} \\
\texttt{QWEN38\_\allowbreak{}EXL3} & \texttt{none} & \texttt{engine/\allowbreak{}loader.py} \\
\texttt{QWEN38\_\allowbreak{}EXL3\_\allowbreak{}CHUNK\_\allowbreak{}UNTIL} & \texttt{128} & \texttt{tools/\allowbreak{}exl3\_\allowbreak{}linear.py} \\
\texttt{QWEN38\_\allowbreak{}EXL3\_\allowbreak{}FRAG} & \texttt{1} & \texttt{engine/\allowbreak{}models/\allowbreak{}qwen3\_\allowbreak{}5/\allowbreak{}frag.py} \\
\texttt{QWEN38\_\allowbreak{}EXL3\_\allowbreak{}FUSE} & \texttt{1} & \texttt{engine/\allowbreak{}loader.py} \\
\texttt{QWEN38\_\allowbreak{}EXL3\_\allowbreak{}MLP\_\allowbreak{}BRIDGE} & \texttt{1} & \texttt{tools/\allowbreak{}exl3\_\allowbreak{}linear.py} \\
\texttt{QWEN38\_\allowbreak{}EXL3\_\allowbreak{}PDL} & \texttt{0} & \texttt{tools/\allowbreak{}exl3\_\allowbreak{}linear.py} \\
\texttt{QWEN38\_\allowbreak{}EXL3\_\allowbreak{}ROWGRID} & \texttt{1} & \texttt{tools/\allowbreak{}exl3\_\allowbreak{}linear.py} \\
\texttt{QWEN38\_\allowbreak{}EXL3\_\allowbreak{}SILU\_\allowbreak{}FUSED} & \texttt{1} & \texttt{tools/\allowbreak{}exl3\_\allowbreak{}linear.py} \\
\texttt{QWEN38\_\allowbreak{}EXL3\_\allowbreak{}STAGES} & \texttt{3} & \texttt{tools/\allowbreak{}exl3\_\allowbreak{}linear.py} \\
\texttt{QWEN38\_\allowbreak{}EXL3\_\allowbreak{}TALL} & \texttt{none} & \texttt{engine/\allowbreak{}models/\allowbreak{}qwen3\_\allowbreak{}5/\allowbreak{}tall.py} \\
\texttt{QWEN38\_\allowbreak{}EXL3\_\allowbreak{}TILES} & \texttt{none} & \texttt{tools/\allowbreak{}exl3\_\allowbreak{}linear.py} \\
\texttt{QWEN38\_\allowbreak{}FAIL\_\allowbreak{}OPEN} & \texttt{general} & \texttt{app/\allowbreak{}tandemapp/\allowbreak{}conductor/\allowbreak{}router.py} \\
\texttt{QWEN38\_\allowbreak{}FLINT\_\allowbreak{}ANCHOR\_\allowbreak{}BACKOFF} & \texttt{0} & \texttt{engine/\allowbreak{}models/\allowbreak{}flint/\allowbreak{}serving.py} \\
\texttt{QWEN38\_\allowbreak{}FLINT\_\allowbreak{}KV\_\allowbreak{}FP8} & \texttt{0} & \texttt{engine/\allowbreak{}models/\allowbreak{}flint/\allowbreak{}engine.py} \\
\texttt{QWEN38\_\allowbreak{}FLINT\_\allowbreak{}LOOKUP} & \texttt{1} & \texttt{engine/\allowbreak{}models/\allowbreak{}flint/\allowbreak{}serving.py} \\
\texttt{QWEN38\_\allowbreak{}FLINT\_\allowbreak{}LOOK\_\allowbreak{}BACKOFF} & \texttt{0} & \texttt{engine/\allowbreak{}models/\allowbreak{}flint/\allowbreak{}serving.py} \\
\texttt{QWEN38\_\allowbreak{}FLINT\_\allowbreak{}PDL} & \texttt{1} & \texttt{engine/\allowbreak{}models/\allowbreak{}flint/\allowbreak{}engine.py} \\
\texttt{QWEN38\_\allowbreak{}FLINT\_\allowbreak{}SIDE\_\allowbreak{}STREAM} & \texttt{1} & \texttt{engine/\allowbreak{}models/\allowbreak{}flint/\allowbreak{}engine.py} \\
\texttt{QWEN38\_\allowbreak{}FN\_\allowbreak{}ATTN\_\allowbreak{}MULTI} & \texttt{0} & \texttt{engine/\allowbreak{}models/\allowbreak{}flashnext/\allowbreak{}engine.py} \\
\texttt{QWEN38\_\allowbreak{}FN\_\allowbreak{}DENSE} & \texttt{0} & \texttt{engine/\allowbreak{}models/\allowbreak{}flashnext/\allowbreak{}engine.py} \\
\texttt{QWEN38\_\allowbreak{}FN\_\allowbreak{}DRAFT\_\allowbreak{}FREEZE} & \texttt{0} & \texttt{engine/\allowbreak{}models/\allowbreak{}flashnext/\allowbreak{}serving.py} \\
\texttt{QWEN38\_\allowbreak{}FN\_\allowbreak{}DRAFT\_\allowbreak{}HOT} & \texttt{32768} & \texttt{engine/\allowbreak{}models/\allowbreak{}flashnext/\allowbreak{}serving.py} \\
\texttt{QWEN38\_\allowbreak{}FN\_\allowbreak{}DROP\_\allowbreak{}SHARD\_\allowbreak{}CACHE} & \texttt{1} & \texttt{engine/\allowbreak{}models/\allowbreak{}flashnext/\allowbreak{}load.py} \\
\texttt{QWEN38\_\allowbreak{}FN\_\allowbreak{}FOLD\_\allowbreak{}COMMIT} & \texttt{1} & \texttt{engine/\allowbreak{}models/\allowbreak{}flashnext/\allowbreak{}engine.py} \\
\texttt{QWEN38\_\allowbreak{}FN\_\allowbreak{}GDN\_\allowbreak{}DEFER} & \texttt{0} & \texttt{engine/\allowbreak{}models/\allowbreak{}flashnext/\allowbreak{}engine.py} \\
\texttt{QWEN38\_\allowbreak{}FN\_\allowbreak{}GDN\_\allowbreak{}MULTI} & \texttt{1} & \texttt{engine/\allowbreak{}models/\allowbreak{}flashnext/\allowbreak{}engine.py} \\
\texttt{QWEN38\_\allowbreak{}FN\_\allowbreak{}GDN\_\allowbreak{}PRE\_\allowbreak{}MULTI} & \texttt{0} & \texttt{engine/\allowbreak{}models/\allowbreak{}flashnext/\allowbreak{}engine.py} \\
\texttt{QWEN38\_\allowbreak{}FN\_\allowbreak{}GEMV\_\allowbreak{}MC} & \texttt{0} & \texttt{tools/\allowbreak{}fn\_\allowbreak{}kernels.py} \\
\texttt{QWEN38\_\allowbreak{}FN\_\allowbreak{}GEMV\_\allowbreak{}NB} & \texttt{0} & \texttt{tools/\allowbreak{}fn\_\allowbreak{}kernels.py} \\
\texttt{QWEN38\_\allowbreak{}FN\_\allowbreak{}GEMV\_\allowbreak{}SPLIT} & \texttt{1} & \texttt{tools/\allowbreak{}fn\_\allowbreak{}kernels.py} \\
\texttt{QWEN38\_\allowbreak{}FN\_\allowbreak{}GROUP\_\allowbreak{}ROWS} & \texttt{64} & \texttt{tools/\allowbreak{}fn\_\allowbreak{}prefill\_\allowbreak{}kernels.py} \\
\texttt{QWEN38\_\allowbreak{}FN\_\allowbreak{}JG\_\allowbreak{}PLE\_\allowbreak{}LATE} & \texttt{0} & \texttt{engine/\allowbreak{}models/\allowbreak{}flashnext/\allowbreak{}engine.py} \\
\texttt{QWEN38\_\allowbreak{}FN\_\allowbreak{}JOINT\_\allowbreak{}GRAPH} & \texttt{0} & \texttt{engine/\allowbreak{}models/\allowbreak{}flashnext/\allowbreak{}engine.py} \\
\texttt{QWEN38\_\allowbreak{}FN\_\allowbreak{}JOINT\_\allowbreak{}GRAPH\_\allowbreak{}AFTER} & \texttt{2} & \texttt{engine/\allowbreak{}models/\allowbreak{}flashnext/\allowbreak{}engine.py} \\
\texttt{QWEN38\_\allowbreak{}FN\_\allowbreak{}JOINT\_\allowbreak{}GRAPH\_\allowbreak{}MAX} & \texttt{12} & \texttt{engine/\allowbreak{}models/\allowbreak{}flashnext/\allowbreak{}engine.py} \\
\texttt{QWEN38\_\allowbreak{}FN\_\allowbreak{}JOINT\_\allowbreak{}MOE\_\allowbreak{}GROUP} & \texttt{1} & \texttt{engine/\allowbreak{}models/\allowbreak{}flashnext/\allowbreak{}engine.py} \\
\texttt{QWEN38\_\allowbreak{}FN\_\allowbreak{}JOINT\_\allowbreak{}PREFILL} & \texttt{0} & \texttt{engine/\allowbreak{}models/\allowbreak{}flashnext/\allowbreak{}joint\_\allowbreak{}prefill.py} \\
\texttt{QWEN38\_\allowbreak{}FN\_\allowbreak{}JOINT\_\allowbreak{}ROWS8} & \texttt{24} & \texttt{engine/\allowbreak{}models/\allowbreak{}flashnext/\allowbreak{}serving.py} \\
\texttt{QWEN38\_\allowbreak{}FN\_\allowbreak{}JOINT\_\allowbreak{}ROWWISE} & \texttt{1} & \texttt{engine/\allowbreak{}models/\allowbreak{}flashnext/\allowbreak{}engine.py} \\
\texttt{QWEN38\_\allowbreak{}FN\_\allowbreak{}JOINT\_\allowbreak{}TIMING} & \texttt{0} & \texttt{engine/\allowbreak{}models/\allowbreak{}flashnext/\allowbreak{}engine.py} \\
\texttt{QWEN38\_\allowbreak{}FN\_\allowbreak{}KV\_\allowbreak{}FP8} & \texttt{0} & \texttt{engine/\allowbreak{}models/\allowbreak{}flashnext/\allowbreak{}engine.py} \\
\texttt{QWEN38\_\allowbreak{}FN\_\allowbreak{}LOOKUP} & \texttt{0} & \texttt{engine/\allowbreak{}models/\allowbreak{}flashnext/\allowbreak{}serving.py} \\
\texttt{QWEN38\_\allowbreak{}FN\_\allowbreak{}LOOK\_\allowbreak{}BACKOFF} & \texttt{0} & \texttt{engine/\allowbreak{}models/\allowbreak{}flashnext/\allowbreak{}serving.py} \\
\texttt{QWEN38\_\allowbreak{}FN\_\allowbreak{}MIX} & \texttt{int8} & \texttt{engine/\allowbreak{}models/\allowbreak{}flashnext/\allowbreak{}pack.py} \\
\texttt{QWEN38\_\allowbreak{}FN\_\allowbreak{}MIXER\_\allowbreak{}MMA} & \texttt{1} & \texttt{tools/\allowbreak{}fn\_\allowbreak{}kernels.py} \\
\texttt{QWEN38\_\allowbreak{}FN\_\allowbreak{}MIX\_\allowbreak{}FUSED} & \texttt{0} & \texttt{engine/\allowbreak{}models/\allowbreak{}flashnext/\allowbreak{}engine.py} \\
\texttt{QWEN38\_\allowbreak{}FN\_\allowbreak{}MOE\_\allowbreak{}AS} & \texttt{0} & \texttt{tools/\allowbreak{}fn\_\allowbreak{}kernels.py} \\
\texttt{QWEN38\_\allowbreak{}FN\_\allowbreak{}MOE\_\allowbreak{}GPAR} & \texttt{0} & \texttt{tools/\allowbreak{}fn\_\allowbreak{}kernels.py} \\
\texttt{QWEN38\_\allowbreak{}FN\_\allowbreak{}MOE\_\allowbreak{}GR} & \texttt{8} & \texttt{tools/\allowbreak{}fn\_\allowbreak{}kernels.py} \\
\texttt{QWEN38\_\allowbreak{}FN\_\allowbreak{}MOE\_\allowbreak{}GY} & \texttt{0} & \texttt{tools/\allowbreak{}fn\_\allowbreak{}kernels.py} \\
\texttt{QWEN38\_\allowbreak{}FN\_\allowbreak{}MOE\_\allowbreak{}ST} & \texttt{3} & \texttt{tools/\allowbreak{}gpu/\allowbreak{}gb10.py} \\
\texttt{QWEN38\_\allowbreak{}FN\_\allowbreak{}MOE\_\allowbreak{}STAGE} & \texttt{0} & \texttt{tools/\allowbreak{}fn\_\allowbreak{}kernels.py} \\
\texttt{QWEN38\_\allowbreak{}FN\_\allowbreak{}MOE\_\allowbreak{}WARPS} & \texttt{4} & \texttt{tools/\allowbreak{}gpu/\allowbreak{}gb10.py} \\
\texttt{QWEN38\_\allowbreak{}FN\_\allowbreak{}MTP\_\allowbreak{}CORRECTOR} & \texttt{0} & \texttt{engine/\allowbreak{}models/\allowbreak{}flashnext/\allowbreak{}corrector.py} \\
\texttt{QWEN38\_\allowbreak{}FN\_\allowbreak{}MTP\_\allowbreak{}CORRECTOR\_\allowbreak{}FILE} & \texttt{auto} & \texttt{engine/\allowbreak{}models/\allowbreak{}flashnext/\allowbreak{}corrector.py} \\
\texttt{QWEN38\_\allowbreak{}FN\_\allowbreak{}MTP\_\allowbreak{}COUPLED} & \texttt{0} & \texttt{engine/\allowbreak{}models/\allowbreak{}flashnext/\allowbreak{}coupled.py} \\
\texttt{QWEN38\_\allowbreak{}FN\_\allowbreak{}MTP\_\allowbreak{}COUPLED\_\allowbreak{}CAL} & \texttt{0} & \texttt{engine/\allowbreak{}models/\allowbreak{}flashnext/\allowbreak{}coupled.py} \\
\texttt{QWEN38\_\allowbreak{}FN\_\allowbreak{}MTP\_\allowbreak{}COUPLED\_\allowbreak{}FILTER} & \texttt{1} & \texttt{engine/\allowbreak{}models/\allowbreak{}flashnext/\allowbreak{}coupled.py} \\
\texttt{QWEN38\_\allowbreak{}FN\_\allowbreak{}MTP\_\allowbreak{}COUPLED\_\allowbreak{}HOT} & \texttt{0} & \texttt{engine/\allowbreak{}models/\allowbreak{}flashnext/\allowbreak{}coupled.py} \\
\texttt{QWEN38\_\allowbreak{}FN\_\allowbreak{}MTP\_\allowbreak{}COUPLED\_\allowbreak{}HOT\_\allowbreak{}LONE} & \texttt{0} & \texttt{engine/\allowbreak{}models/\allowbreak{}flashnext/\allowbreak{}coupled.py} \\
\texttt{QWEN38\_\allowbreak{}FN\_\allowbreak{}MTP\_\allowbreak{}COUPLED\_\allowbreak{}SIB} & \texttt{0} & \texttt{engine/\allowbreak{}models/\allowbreak{}flashnext/\allowbreak{}coupled.py} \\
\texttt{QWEN38\_\allowbreak{}FN\_\allowbreak{}MTP\_\allowbreak{}COUPLED\_\allowbreak{}SIB\_\allowbreak{}MARGIN} & \texttt{1.3} & \texttt{engine/\allowbreak{}models/\allowbreak{}flashnext/\allowbreak{}coupled.py} \\
\texttt{QWEN38\_\allowbreak{}FN\_\allowbreak{}MTP\_\allowbreak{}COUPLED\_\allowbreak{}TAU} & \texttt{0.85} & \texttt{engine/\allowbreak{}models/\allowbreak{}flashnext/\allowbreak{}coupled.py} \\
\texttt{QWEN38\_\allowbreak{}FN\_\allowbreak{}PACK\_\allowbreak{}GRAPH} & \texttt{1} & \texttt{engine/\allowbreak{}models/\allowbreak{}flashnext/\allowbreak{}engine.py} \\
\texttt{QWEN38\_\allowbreak{}FN\_\allowbreak{}PACK\_\allowbreak{}GRAPH\_\allowbreak{}AFTER} & \texttt{2} & \texttt{engine/\allowbreak{}models/\allowbreak{}flashnext/\allowbreak{}engine.py} \\
\texttt{QWEN38\_\allowbreak{}FN\_\allowbreak{}PDL} & \texttt{1} & \texttt{tools/\allowbreak{}fn\_\allowbreak{}kernels.py} \\
\texttt{QWEN38\_\allowbreak{}FN\_\allowbreak{}PF\_\allowbreak{}ATTN\_\allowbreak{}STAGES} & \texttt{1} & \texttt{engine/\allowbreak{}models/\allowbreak{}flashnext/\allowbreak{}engine.py} \\
\texttt{QWEN38\_\allowbreak{}FN\_\allowbreak{}PF\_\allowbreak{}HC\_\allowbreak{}FUSED} & \texttt{1} & \texttt{engine/\allowbreak{}models/\allowbreak{}flashnext/\allowbreak{}engine.py} \\
\texttt{QWEN38\_\allowbreak{}FN\_\allowbreak{}PF\_\allowbreak{}INJ\_\allowbreak{}FUSED} & \texttt{1} & \texttt{engine/\allowbreak{}models/\allowbreak{}flashnext/\allowbreak{}engine.py} \\
\texttt{QWEN38\_\allowbreak{}FN\_\allowbreak{}PF\_\allowbreak{}NOCOPY} & \texttt{1} & \texttt{engine/\allowbreak{}models/\allowbreak{}flashnext/\allowbreak{}engine.py} \\
\texttt{QWEN38\_\allowbreak{}FN\_\allowbreak{}PF\_\allowbreak{}ROT\_\allowbreak{}ROWS} & \texttt{0} & \texttt{tools/\allowbreak{}fn\_\allowbreak{}prefill\_\allowbreak{}kernels.py} \\
\texttt{QWEN38\_\allowbreak{}FN\_\allowbreak{}PF\_\allowbreak{}WINDOW} & \texttt{8192} & \texttt{engine/\allowbreak{}models/\allowbreak{}flashnext/\allowbreak{}layer\_\allowbreak{}major.py} \\
\texttt{QWEN38\_\allowbreak{}FN\_\allowbreak{}PLE\_\allowbreak{}BATCH} & \texttt{1} & \texttt{engine/\allowbreak{}models/\allowbreak{}flashnext/\allowbreak{}engine.py} \\
\texttt{QWEN38\_\allowbreak{}FN\_\allowbreak{}PLE\_\allowbreak{}LATE} & \texttt{1} & \texttt{engine/\allowbreak{}models/\allowbreak{}flashnext/\allowbreak{}engine.py} \\
\texttt{QWEN38\_\allowbreak{}FN\_\allowbreak{}PLE\_\allowbreak{}MULTI} & \texttt{0} & \texttt{engine/\allowbreak{}models/\allowbreak{}flashnext/\allowbreak{}engine.py} \\
\texttt{QWEN38\_\allowbreak{}FN\_\allowbreak{}PLE\_\allowbreak{}PREFETCH} & \texttt{0} & \texttt{engine/\allowbreak{}models/\allowbreak{}flashnext/\allowbreak{}engine.py} \\
\texttt{QWEN38\_\allowbreak{}FN\_\allowbreak{}PLE\_\allowbreak{}THREADS} & \texttt{16} & \texttt{engine/\allowbreak{}models/\allowbreak{}flashnext/\allowbreak{}ple.py} \\
\texttt{QWEN38\_\allowbreak{}FN\_\allowbreak{}POOL\_\allowbreak{}MULTI} & \texttt{0} & \texttt{engine/\allowbreak{}models/\allowbreak{}flashnext/\allowbreak{}engine.py} \\
\texttt{QWEN38\_\allowbreak{}FN\_\allowbreak{}PREFILL} & \texttt{kernels} & \texttt{engine/\allowbreak{}models/\allowbreak{}flashnext/\allowbreak{}engine.py} \\
\texttt{QWEN38\_\allowbreak{}FN\_\allowbreak{}PREFILL\_\allowbreak{}CHUNK} & \texttt{4096} & \texttt{engine/\allowbreak{}models/\allowbreak{}flashnext/\allowbreak{}engine.py} \\
\texttt{QWEN38\_\allowbreak{}FN\_\allowbreak{}PREFILL\_\allowbreak{}MM} & \texttt{fp32} & \texttt{tools/\allowbreak{}fn\_\allowbreak{}prefill\_\allowbreak{}kernels.py} \\
\texttt{QWEN38\_\allowbreak{}FN\_\allowbreak{}QSA\_\allowbreak{}EPI\_\allowbreak{}TORCH} & \texttt{0} & \texttt{tools/\allowbreak{}fn\_\allowbreak{}prefill\_\allowbreak{}kernels.py} \\
\texttt{QWEN38\_\allowbreak{}FN\_\allowbreak{}QSA\_\allowbreak{}LIVE} & \texttt{1} & \texttt{engine/\allowbreak{}models/\allowbreak{}flashnext/\allowbreak{}engine.py} \\
\texttt{QWEN38\_\allowbreak{}FN\_\allowbreak{}QSA\_\allowbreak{}TORCH\_\allowbreak{}BIG} & \texttt{0} & \texttt{tools/\allowbreak{}fn\_\allowbreak{}seq\_\allowbreak{}kernels.py} \\
\texttt{QWEN38\_\allowbreak{}FN\_\allowbreak{}ROUTE\_\allowbreak{}FUSED} & \texttt{1} & \texttt{engine/\allowbreak{}models/\allowbreak{}flashnext/\allowbreak{}engine.py} \\
\texttt{QWEN38\_\allowbreak{}FN\_\allowbreak{}TABLE} & \texttt{warm} & \texttt{engine/\allowbreak{}models/\allowbreak{}flashnext/\allowbreak{}load.py} \\
\texttt{QWEN38\_\allowbreak{}FN\_\allowbreak{}TABLE\_\allowbreak{}CACHE\_\allowbreak{}GB} & \texttt{2} & \texttt{engine/\allowbreak{}models/\allowbreak{}flashnext/\allowbreak{}load.py} \\
\texttt{QWEN38\_\allowbreak{}FN\_\allowbreak{}TREE\_\allowbreak{}PLAN} & \texttt{0} & \texttt{engine/\allowbreak{}models/\allowbreak{}flashnext/\allowbreak{}engine.py} \\
\texttt{QWEN38\_\allowbreak{}FN\_\allowbreak{}UNION\_\allowbreak{}EVERY} & \texttt{0} & \texttt{engine/\allowbreak{}models/\allowbreak{}flashnext/\allowbreak{}engine.py} \\
\texttt{QWEN38\_\allowbreak{}FN\_\allowbreak{}WARM\_\allowbreak{}TABLE} & \texttt{1} & \texttt{engine/\allowbreak{}models/\allowbreak{}flashnext/\allowbreak{}load.py} \\
\texttt{QWEN38\_\allowbreak{}FP8\_\allowbreak{}HEAD} & \texttt{none} & \texttt{engine/\allowbreak{}loader.py} \\
\texttt{QWEN38\_\allowbreak{}FP8\_\allowbreak{}TILES} & \texttt{none} & \texttt{tools/\allowbreak{}fp8\_\allowbreak{}linear.py} \\
\texttt{QWEN38\_\allowbreak{}FRESH\_\allowbreak{}SEED} & \texttt{auto} & \texttt{server/\allowbreak{}app.py} \\
\texttt{QWEN38\_\allowbreak{}FUSED\_\allowbreak{}ADDNORM} & \texttt{0} & \texttt{1} \textbullet{} \texttt{engine/\allowbreak{}model.py} \\
\texttt{QWEN38\_\allowbreak{}FUSED\_\allowbreak{}ATTN} & \texttt{0} & \texttt{engine/\allowbreak{}model.py} \\
\texttt{QWEN38\_\allowbreak{}FUSED\_\allowbreak{}ATTN\_\allowbreak{}PREP} & \texttt{1} & \texttt{engine/\allowbreak{}model.py} \\
\texttt{QWEN38\_\allowbreak{}FUSED\_\allowbreak{}COMMIT} & \texttt{0} & \texttt{1} \textbullet{} \texttt{engine/\allowbreak{}model.py} \\
\texttt{QWEN38\_\allowbreak{}FUSED\_\allowbreak{}GDN} & \texttt{1} & \texttt{engine/\allowbreak{}model.py} \\
\texttt{QWEN38\_\allowbreak{}FUSED\_\allowbreak{}GDNBLOCK} & \texttt{1} & \texttt{engine/\allowbreak{}model.py} \\
\texttt{QWEN38\_\allowbreak{}FUSED\_\allowbreak{}GDNPRE} & \texttt{1} & \texttt{engine/\allowbreak{}model.py} \\
\texttt{QWEN38\_\allowbreak{}FUSED\_\allowbreak{}GDNPREFILL} & \texttt{0} & \texttt{1} \textbullet{} \texttt{engine/\allowbreak{}model.py} \\
\texttt{QWEN38\_\allowbreak{}FUSED\_\allowbreak{}GDNTREE} & \texttt{1} & \texttt{engine/\allowbreak{}model.py} \\
\texttt{QWEN38\_\allowbreak{}FUSED\_\allowbreak{}GDNVERIFY} & \texttt{0} & \texttt{1} \textbullet{} \texttt{engine/\allowbreak{}model.py} \\
\texttt{QWEN38\_\allowbreak{}FUSED\_\allowbreak{}HEAD} & \texttt{1} & \texttt{engine/\allowbreak{}model.py} \\
\texttt{QWEN38\_\allowbreak{}FUSED\_\allowbreak{}NORM} & \texttt{1} & \texttt{engine/\allowbreak{}model.py} \\
\texttt{QWEN38\_\allowbreak{}FUSE\_\allowbreak{}PROJ} & \texttt{0} & \texttt{tools/\allowbreak{}nvfp4\_\allowbreak{}linear\_\allowbreak{}v2.py} \\
\texttt{QWEN38\_\allowbreak{}GDNV\_\allowbreak{}BV} & \texttt{16} & \texttt{tools/\allowbreak{}gdn\_\allowbreak{}verify\_\allowbreak{}kernels.py} \\
\texttt{QWEN38\_\allowbreak{}GDNV\_\allowbreak{}CONV\_\allowbreak{}BLOCK} & \texttt{256} & \texttt{32} \textbullet{} \texttt{tools/\allowbreak{}gdn\_\allowbreak{}verify\_\allowbreak{}kernels.py} \\
\texttt{QWEN38\_\allowbreak{}GDNV\_\allowbreak{}CONV\_\allowbreak{}WARPS} & \texttt{4} & \texttt{1} \textbullet{} \texttt{tools/\allowbreak{}gdn\_\allowbreak{}verify\_\allowbreak{}kernels.py} \\
\texttt{QWEN38\_\allowbreak{}GDNV\_\allowbreak{}ONE\_\allowbreak{}WARP} & \texttt{0} & \texttt{1} \textbullet{} \texttt{tools/\allowbreak{}gdn\_\allowbreak{}verify\_\allowbreak{}kernels.py} \\
\texttt{QWEN38\_\allowbreak{}GDNV\_\allowbreak{}TREE\_\allowbreak{}PF} & \texttt{0} & \texttt{tools/\allowbreak{}gdn\_\allowbreak{}verify\_\allowbreak{}kernels.py} \\
\texttt{QWEN38\_\allowbreak{}GDNV\_\allowbreak{}WARPS} & \texttt{4} & \texttt{tools/\allowbreak{}gdn\_\allowbreak{}verify\_\allowbreak{}kernels.py} \\
\texttt{QWEN38\_\allowbreak{}GDNV\_\allowbreak{}WY} & \texttt{0} & \texttt{1} \textbullet{} \texttt{tools/\allowbreak{}gdn\_\allowbreak{}verify\_\allowbreak{}kernels.py} \\
\texttt{QWEN38\_\allowbreak{}GDNV\_\allowbreak{}WY\_\allowbreak{}CHAIN\_\allowbreak{}MAXT} & \texttt{32} & \texttt{16} \textbullet{} \texttt{tools/\allowbreak{}gdn\_\allowbreak{}verify\_\allowbreak{}kernels.py} \\
\texttt{QWEN38\_\allowbreak{}GDNV\_\allowbreak{}WY\_\allowbreak{}FUSED} & \texttt{0} & \texttt{1} \textbullet{} \texttt{tools/\allowbreak{}gdn\_\allowbreak{}verify\_\allowbreak{}kernels.py} \\
\texttt{QWEN38\_\allowbreak{}GDNV\_\allowbreak{}WY\_\allowbreak{}FUSED\_\allowbreak{}MAXT} & \texttt{32} & \texttt{16} \textbullet{} \texttt{tools/\allowbreak{}gdn\_\allowbreak{}verify\_\allowbreak{}kernels.py} \\
\texttt{QWEN38\_\allowbreak{}GDNV\_\allowbreak{}WY\_\allowbreak{}KC} & \texttt{32} & \texttt{tools/\allowbreak{}gdn\_\allowbreak{}verify\_\allowbreak{}kernels.py} \\
\texttt{QWEN38\_\allowbreak{}GDNV\_\allowbreak{}WY\_\allowbreak{}MAXT} & \texttt{32} & \texttt{16} \textbullet{} \texttt{tools/\allowbreak{}gdn\_\allowbreak{}verify\_\allowbreak{}kernels.py} \\
\texttt{QWEN38\_\allowbreak{}GDN\_\allowbreak{}AB} & \texttt{1} & \texttt{engine/\allowbreak{}model.py} \\
\texttt{QWEN38\_\allowbreak{}GDN\_\allowbreak{}CHUNK} & \texttt{64} & \texttt{engine/\allowbreak{}model.py} \\
\texttt{QWEN38\_\allowbreak{}GDN\_\allowbreak{}CHUNKED} & \texttt{1} & \texttt{engine/\allowbreak{}model.py} \\
\texttt{QWEN38\_\allowbreak{}GDN\_\allowbreak{}CHUNKED\_\allowbreak{}BV} & \texttt{16} & \texttt{tools/\allowbreak{}gdn\_\allowbreak{}chunked\_\allowbreak{}kernels.py} \\
\texttt{QWEN38\_\allowbreak{}GDN\_\allowbreak{}CHUNKED\_\allowbreak{}C} & \texttt{64} & \texttt{tools/\allowbreak{}gdn\_\allowbreak{}chunked\_\allowbreak{}kernels.py} \\
\texttt{QWEN38\_\allowbreak{}GDN\_\allowbreak{}CHUNKED\_\allowbreak{}WARPS} & \texttt{4} & \texttt{tools/\allowbreak{}gdn\_\allowbreak{}chunked\_\allowbreak{}kernels.py} \\
\texttt{QWEN38\_\allowbreak{}GDN\_\allowbreak{}CHUNKED\_\allowbreak{}WARPS\_\allowbreak{}INTRA} & \texttt{8} & \texttt{tools/\allowbreak{}gdn\_\allowbreak{}chunked\_\allowbreak{}kernels.py} \\
\texttt{QWEN38\_\allowbreak{}GDN\_\allowbreak{}MM} & \texttt{fp32} & \texttt{engine/\allowbreak{}gdn.py} \\
\texttt{QWEN38\_\allowbreak{}GDN\_\allowbreak{}MM\_\allowbreak{}FROM} & \texttt{64} & \texttt{engine/\allowbreak{}gdn.py} \\
\texttt{QWEN38\_\allowbreak{}GDN\_\allowbreak{}PREFILL\_\allowbreak{}BV} & \texttt{32} & \texttt{tools/\allowbreak{}gdn\_\allowbreak{}prefill\_\allowbreak{}kernels.py} \\
\texttt{QWEN38\_\allowbreak{}GDN\_\allowbreak{}PREFILL\_\allowbreak{}NOREP} & \texttt{0} & \texttt{engine/\allowbreak{}model.py} \\
\texttt{QWEN38\_\allowbreak{}GDN\_\allowbreak{}PREFILL\_\allowbreak{}PREC} & \texttt{bf16x3} & \texttt{tools/\allowbreak{}gdn\_\allowbreak{}prefill\_\allowbreak{}kernels.py} \\
\texttt{QWEN38\_\allowbreak{}GDN\_\allowbreak{}PREFILL\_\allowbreak{}PREC\_\allowbreak{}A} & \texttt{tf32} & \texttt{tools/\allowbreak{}gdn\_\allowbreak{}prefill\_\allowbreak{}kernels.py} \\
\texttt{QWEN38\_\allowbreak{}GDN\_\allowbreak{}PREFILL\_\allowbreak{}PREC\_\allowbreak{}SCAN} & \texttt{none} & \texttt{tools/\allowbreak{}gdn\_\allowbreak{}prefill\_\allowbreak{}kernels.py} \\
\texttt{QWEN38\_\allowbreak{}GDN\_\allowbreak{}PREFILL\_\allowbreak{}SERIES} & \texttt{0} & \texttt{tools/\allowbreak{}gdn\_\allowbreak{}prefill\_\allowbreak{}kernels.py} \\
\texttt{QWEN38\_\allowbreak{}GDN\_\allowbreak{}PREFILL\_\allowbreak{}STAGES} & \texttt{1} & \texttt{tools/\allowbreak{}gdn\_\allowbreak{}prefill\_\allowbreak{}kernels.py} \\
\texttt{QWEN38\_\allowbreak{}GDN\_\allowbreak{}PREFILL\_\allowbreak{}WARPS\_\allowbreak{}INTRA} & \texttt{8} & \texttt{tools/\allowbreak{}gdn\_\allowbreak{}prefill\_\allowbreak{}kernels.py} \\
\texttt{QWEN38\_\allowbreak{}GDN\_\allowbreak{}PREFILL\_\allowbreak{}WARPS\_\allowbreak{}SCAN} & \texttt{8} & \texttt{tools/\allowbreak{}gdn\_\allowbreak{}prefill\_\allowbreak{}kernels.py} \\
\texttt{QWEN38\_\allowbreak{}GDN\_\allowbreak{}PREP} & \texttt{1} & \texttt{engine/\allowbreak{}model.py} \\
\texttt{QWEN38\_\allowbreak{}GDN\_\allowbreak{}PREP\_\allowbreak{}BV} & \texttt{16} & \texttt{engine/\allowbreak{}model.py} \\
\texttt{QWEN38\_\allowbreak{}GDN\_\allowbreak{}PREP\_\allowbreak{}WARPS} & \texttt{4} & \texttt{engine/\allowbreak{}model.py} \\
\texttt{QWEN38\_\allowbreak{}GDN\_\allowbreak{}SEQ} & \texttt{0} & \texttt{tools/\allowbreak{}gdn\_\allowbreak{}seq\_\allowbreak{}kernels.py} \\
\texttt{QWEN38\_\allowbreak{}GDN\_\allowbreak{}SEQ\_\allowbreak{}BV} & \texttt{16} & \texttt{tools/\allowbreak{}gdn\_\allowbreak{}seq\_\allowbreak{}kernels.py} \\
\texttt{QWEN38\_\allowbreak{}GDN\_\allowbreak{}SEQ\_\allowbreak{}BYNODE} & \texttt{0} & \texttt{tools/\allowbreak{}gdn\_\allowbreak{}seq\_\allowbreak{}kernels.py} \\
\texttt{QWEN38\_\allowbreak{}GDN\_\allowbreak{}SEQ\_\allowbreak{}WARPS} & \texttt{1} & \texttt{tools/\allowbreak{}gdn\_\allowbreak{}seq\_\allowbreak{}kernels.py} \\
\texttt{QWEN38\_\allowbreak{}GQA\_\allowbreak{}FROM} & \texttt{64} & \texttt{engine/\allowbreak{}model.py} \\
\texttt{QWEN38\_\allowbreak{}GRAPH\_\allowbreak{}MAX\_\allowbreak{}CTX} & \texttt{32768} & \texttt{engine/\allowbreak{}verify\_\allowbreak{}graph.py} \\
\texttt{QWEN38\_\allowbreak{}GRAPH\_\allowbreak{}SHARE\_\allowbreak{}FROM} & \texttt{2} & \texttt{engine/\allowbreak{}verify\_\allowbreak{}graph.py} \\
\texttt{QWEN38\_\allowbreak{}GRAPH\_\allowbreak{}SHARE\_\allowbreak{}OUT} & \texttt{0} & \texttt{engine/\allowbreak{}verify\_\allowbreak{}graph.py} \\
\texttt{QWEN38\_\allowbreak{}HEAD\_\allowbreak{}GEMM} & \texttt{none} & \texttt{tools/\allowbreak{}head\_\allowbreak{}gemv.py} \\
\texttt{QWEN38\_\allowbreak{}HOST\_\allowbreak{}ASYNC} & \texttt{0} & \texttt{engine/\allowbreak{}model.py} \\
\texttt{QWEN38\_\allowbreak{}HOST\_\allowbreak{}WALK} & \texttt{0} & \texttt{1} \textbullet{} \texttt{engine/\allowbreak{}drafters/\allowbreak{}dflash2.py} \\
\texttt{QWEN38\_\allowbreak{}JOINT\_\allowbreak{}ATTN\_\allowbreak{}GROUP} & \texttt{1} & \texttt{engine/\allowbreak{}models/\allowbreak{}qwen3\_\allowbreak{}5/\allowbreak{}attn\_\allowbreak{}group.py} \\
\texttt{QWEN38\_\allowbreak{}JOINT\_\allowbreak{}CHAINS} & \texttt{1} & \texttt{engine/\allowbreak{}models/\allowbreak{}qwen3\_\allowbreak{}5/\allowbreak{}adapter.py} \\
\texttt{QWEN38\_\allowbreak{}JOINT\_\allowbreak{}GDN} & \texttt{1} & \texttt{engine/\allowbreak{}model.py} \\
\texttt{QWEN38\_\allowbreak{}JOINT\_\allowbreak{}GRAPHS} & \texttt{1} & \texttt{engine/\allowbreak{}models/\allowbreak{}qwen3\_\allowbreak{}5/\allowbreak{}joint\_\allowbreak{}graphs.py} \\
\texttt{QWEN38\_\allowbreak{}JOINT\_\allowbreak{}GRAPHS\_\allowbreak{}AFTER} & \texttt{4} & \texttt{engine/\allowbreak{}models/\allowbreak{}qwen3\_\allowbreak{}5/\allowbreak{}joint\_\allowbreak{}graphs.py} \\
\texttt{QWEN38\_\allowbreak{}JOINT\_\allowbreak{}GRAPHS\_\allowbreak{}MAX} & \texttt{24} & \texttt{engine/\allowbreak{}models/\allowbreak{}qwen3\_\allowbreak{}5/\allowbreak{}joint\_\allowbreak{}graphs.py} \\
\texttt{QWEN38\_\allowbreak{}JOINT\_\allowbreak{}MAX\_\allowbreak{}ROWS} & \texttt{none} & \texttt{engine/\allowbreak{}models/\allowbreak{}qwen3\_\allowbreak{}5/\allowbreak{}adapter.py} \\
\texttt{QWEN38\_\allowbreak{}JOINT\_\allowbreak{}ROWS\_\allowbreak{}AT\_\allowbreak{}LANES} & \texttt{4:\allowbreak{}8} & \texttt{engine/\allowbreak{}models/\allowbreak{}qwen3\_\allowbreak{}5/\allowbreak{}joint\_\allowbreak{}graphs.py} \\
\texttt{QWEN38\_\allowbreak{}JOINT\_\allowbreak{}ROWS\_\allowbreak{}A\_\allowbreak{}SLOT} & \texttt{0} & \texttt{engine/\allowbreak{}models/\allowbreak{}qwen3\_\allowbreak{}5/\allowbreak{}joint\_\allowbreak{}graphs.py} \\
\texttt{QWEN38\_\allowbreak{}JOINT\_\allowbreak{}ROWWISE} & \texttt{1} & \texttt{engine/\allowbreak{}model.py} \\
\texttt{QWEN38\_\allowbreak{}JOIN\_\allowbreak{}VIEWS} & \texttt{1} & \texttt{engine/\allowbreak{}model.py} \\
\texttt{QWEN38\_\allowbreak{}KV\_\allowbreak{}FP8} & \texttt{0} & \texttt{engine/\allowbreak{}model.py} \\
\texttt{QWEN38\_\allowbreak{}LATCH\_\allowbreak{}PRICE} & \texttt{0} & \texttt{engine/\allowbreak{}lenrouter.py} \\
\texttt{QWEN38\_\allowbreak{}LAUNCH\_\allowbreak{}FIRST} & \texttt{0} & \texttt{server/\allowbreak{}app.py} \\
\texttt{QWEN38\_\allowbreak{}LEN\_\allowbreak{}MODE} & \texttt{none} & \texttt{wide} \textbullet{} \texttt{engine/\allowbreak{}lenrouter.py} \\
\texttt{QWEN38\_\allowbreak{}LEN\_\allowbreak{}SWITCH} & \texttt{0} & \texttt{1} \textbullet{} \texttt{engine/\allowbreak{}lenrouter.py} \\
\texttt{QWEN38\_\allowbreak{}LFM2\_\allowbreak{}KV\_\allowbreak{}FP8} & \texttt{0} & \texttt{engine/\allowbreak{}models/\allowbreak{}lfm2\_\allowbreak{}moe/\allowbreak{}engine.py} \\
\texttt{QWEN38\_\allowbreak{}LIBRARY\_\allowbreak{}HOLD} & \texttt{message} & \texttt{app/\allowbreak{}tandemapp/\allowbreak{}conductor/\allowbreak{}libsearch.py} \\
\texttt{QWEN38\_\allowbreak{}LOWER\_\allowbreak{}RIGHT\_\allowbreak{}FROM} & \texttt{64} & \texttt{engine/\allowbreak{}model.py} \\
\texttt{QWEN38\_\allowbreak{}MEMORY\_\allowbreak{}PROXIES} & \texttt{none} & \texttt{server/\allowbreak{}app.py} \\
\texttt{QWEN38\_\allowbreak{}MEMORY\_\allowbreak{}ROUTES} & \texttt{auto} & \texttt{app/\allowbreak{}tandemapp/\allowbreak{}conductor/\allowbreak{}memory.py} \\
\texttt{QWEN38\_\allowbreak{}MEMORY\_\allowbreak{}WRITER\_\allowbreak{}TEXT} & \texttt{none} & \texttt{server/\allowbreak{}app.py} \\
\texttt{QWEN38\_\allowbreak{}MM\_\allowbreak{}END} & \texttt{on} & \texttt{server/\allowbreak{}app.py} \\
\texttt{QWEN38\_\allowbreak{}MODEL} & \texttt{none} & \texttt{engine/\allowbreak{}config.py} \\
\texttt{QWEN38\_\allowbreak{}MODEL\_\allowbreak{}NAME} & \texttt{none} & \texttt{app/\allowbreak{}tandemapp/\allowbreak{}conductor/\allowbreak{}schema.py} \\
\texttt{QWEN38\_\allowbreak{}MTP\_\allowbreak{}BACKLOG\_\allowbreak{}CAP} & \texttt{2048} & \texttt{engine/\allowbreak{}methods/\allowbreak{}mtp\_\allowbreak{}drafter.py} \\
\texttt{QWEN38\_\allowbreak{}NARRATOR\_\allowbreak{}DEVICE} & \texttt{auto} & \texttt{server/\allowbreak{}conductor/\allowbreak{}inproc.py} \\
\texttt{QWEN38\_\allowbreak{}NARRATOR\_\allowbreak{}DIR} & \texttt{none} & \texttt{server/\allowbreak{}conductor/\allowbreak{}inproc.py} \\
\texttt{QWEN38\_\allowbreak{}NARRATOR\_\allowbreak{}MODEL} & \texttt{none} & \texttt{app/\allowbreak{}tandemapp/\allowbreak{}conductor/\allowbreak{}narrator.py} \\
\texttt{QWEN38\_\allowbreak{}NH\_\allowbreak{}KV\_\allowbreak{}FP8} & \texttt{0} & \texttt{engine/\allowbreak{}models/\allowbreak{}nemotron\_\allowbreak{}h/\allowbreak{}engine.py} \\
\texttt{QWEN38\_\allowbreak{}NOTES\_\allowbreak{}HEAD} & \texttt{off} & \texttt{app/\allowbreak{}tandemapp/\allowbreak{}conductor/\allowbreak{}router.py} \\
\texttt{QWEN38\_\allowbreak{}NOTES\_\allowbreak{}HEAD\_\allowbreak{}ORDER} & \texttt{memory} & \texttt{app/\allowbreak{}tandemapp/\allowbreak{}conductor/\allowbreak{}router.py} \\
\texttt{QWEN38\_\allowbreak{}NOTE\_\allowbreak{}ANCHOR} & \texttt{on} & \texttt{app/\allowbreak{}tandemapp/\allowbreak{}conductor/\allowbreak{}router.py} \\
\texttt{QWEN38\_\allowbreak{}NOTE\_\allowbreak{}DEDUP} & \texttt{parts} & \texttt{app/\allowbreak{}tandemapp/\allowbreak{}conductor/\allowbreak{}router.py} \\
\texttt{QWEN38\_\allowbreak{}NVFP4} & \texttt{none} & \texttt{engine/\allowbreak{}loader.py} \\
\texttt{QWEN38\_\allowbreak{}NVFP4\_\allowbreak{}BLOCK\_\allowbreak{}TILES} & \texttt{0} & \texttt{tools/\allowbreak{}nvfp4\_\allowbreak{}linear.py} \\
\texttt{QWEN38\_\allowbreak{}NVFP4\_\allowbreak{}DEQUANT\_\allowbreak{}FROM} & \texttt{512} & \texttt{tools/\allowbreak{}nvfp4\_\allowbreak{}linear.py} \\
\texttt{QWEN38\_\allowbreak{}NVFP4\_\allowbreak{}PREFILL\_\allowbreak{}V2} & \texttt{0} & \texttt{1} \textbullet{} \texttt{tools/\allowbreak{}nvfp4\_\allowbreak{}linear.py} \\
\texttt{QWEN38\_\allowbreak{}NVFP4\_\allowbreak{}PREFILL\_\allowbreak{}V2\_\allowbreak{}UNTIL} & \texttt{1024} & \texttt{1025} \textbullet{} \texttt{tools/\allowbreak{}nvfp4\_\allowbreak{}linear.py} \\
\texttt{QWEN38\_\allowbreak{}NVFP4\_\allowbreak{}SKINNY} & \texttt{0} & \texttt{1} \textbullet{} \texttt{tools/\allowbreak{}nvfp4\_\allowbreak{}skinny.py} \\
\texttt{QWEN38\_\allowbreak{}NVFP4\_\allowbreak{}SPLITK} & \texttt{0} & \texttt{tools/\allowbreak{}nvfp4\_\allowbreak{}linear.py} \\
\texttt{QWEN38\_\allowbreak{}NVFP4\_\allowbreak{}V2} & \texttt{1} & \texttt{tools/\allowbreak{}nvfp4\_\allowbreak{}linear\_\allowbreak{}v2.py} \\
\texttt{QWEN38\_\allowbreak{}NVFP4\_\allowbreak{}V2\_\allowbreak{}BN} & \texttt{0} & \texttt{tools/\allowbreak{}nvfp4\_\allowbreak{}linear\_\allowbreak{}v2.py} \\
\texttt{QWEN38\_\allowbreak{}NVFP4\_\allowbreak{}V2\_\allowbreak{}DOTS} & \texttt{1} & \texttt{tools/\allowbreak{}nvfp4\_\allowbreak{}linear\_\allowbreak{}v2.py} \\
\texttt{QWEN38\_\allowbreak{}NVFP4\_\allowbreak{}V2\_\allowbreak{}MAX} & \texttt{32} & \texttt{tools/\allowbreak{}nvfp4\_\allowbreak{}linear\_\allowbreak{}v2.py} \\
\texttt{QWEN38\_\allowbreak{}NVFP4\_\allowbreak{}V2\_\allowbreak{}MIN} & \texttt{1} & \texttt{tools/\allowbreak{}nvfp4\_\allowbreak{}linear\_\allowbreak{}v2.py} \\
\texttt{QWEN38\_\allowbreak{}NVFP4\_\allowbreak{}V2\_\allowbreak{}PREFETCH} & \texttt{0} & \texttt{tools/\allowbreak{}nvfp4\_\allowbreak{}linear\_\allowbreak{}v2.py} \\
\texttt{QWEN38\_\allowbreak{}NVFP4\_\allowbreak{}V2\_\allowbreak{}W} & \texttt{0} & \texttt{tools/\allowbreak{}nvfp4\_\allowbreak{}linear\_\allowbreak{}v2.py} \\
\texttt{QWEN38\_\allowbreak{}OWN\_\allowbreak{}BURST\_\allowbreak{}RATE} & \texttt{peek} & \texttt{server/\allowbreak{}app.py} \\
\texttt{QWEN38\_\allowbreak{}OWN\_\allowbreak{}ETA} & \texttt{on} & \texttt{server/\allowbreak{}app.py} \\
\texttt{QWEN38\_\allowbreak{}OWN\_\allowbreak{}RATE} & \texttt{solo} & \texttt{server/\allowbreak{}app.py} \\
\texttt{QWEN38\_\allowbreak{}OWN\_\allowbreak{}SIDE} & \texttt{on} & \texttt{server/\allowbreak{}app.py} \\
\texttt{QWEN38\_\allowbreak{}OWN\_\allowbreak{}UNITS} & \texttt{prefill} & \texttt{server/\allowbreak{}app.py} \\
\texttt{QWEN38\_\allowbreak{}PACKAGE\_\allowbreak{}HOLD} & \texttt{pick} & \texttt{app/\allowbreak{}tandemapp/\allowbreak{}conductor/\allowbreak{}router.py} \\
\texttt{QWEN38\_\allowbreak{}POOL\_\allowbreak{}SNAPSHOT\_\allowbreak{}TTL\_\allowbreak{}S} & \texttt{5.0} & \texttt{engine/\allowbreak{}decision/\allowbreak{}device.py} \\
\texttt{QWEN38\_\allowbreak{}PREFILL\_\allowbreak{}CAUSAL} & \texttt{1} & \texttt{engine/\allowbreak{}model.py} \\
\texttt{QWEN38\_\allowbreak{}PREFILL\_\allowbreak{}WINDOW} & \texttt{8192} & \texttt{engine/\allowbreak{}cache.py} \\
\texttt{QWEN38\_\allowbreak{}PREFILL\_\allowbreak{}WINDOW\_\allowbreak{}FREE} & \texttt{32768} & \texttt{engine/\allowbreak{}cache.py} \\
\texttt{QWEN38\_\allowbreak{}PROBE\_\allowbreak{}DAY\_\allowbreak{}MB} & \texttt{100} & \texttt{server/\allowbreak{}probe.py} \\
\texttt{QWEN38\_\allowbreak{}PROBE\_\allowbreak{}DIR} & \texttt{\textasciitilde{}/\allowbreak{}.qwen38-\allowbreak{}spark-\allowbreak{}engine/\allowbreak{}probe} & \texttt{(a path)} \textbullet{} \texttt{server/\allowbreak{}probe.py} \\
\texttt{QWEN38\_\allowbreak{}PROBE\_\allowbreak{}EVERY} & \texttt{1024} & \texttt{server/\allowbreak{}probe.py} \\
\texttt{QWEN38\_\allowbreak{}PROBE\_\allowbreak{}SNAPSHOTS} & \texttt{off} & \texttt{off} \textbullet{} \texttt{server/\allowbreak{}probe.py} \\
\texttt{QWEN38\_\allowbreak{}PROMPT\_\allowbreak{}END} & \texttt{1} & \texttt{engine/\allowbreak{}cache.py} \\
\texttt{QWEN38\_\allowbreak{}Q35M\_\allowbreak{}ANCHOR\_\allowbreak{}BACKOFF} & \texttt{0} & \texttt{engine/\allowbreak{}models/\allowbreak{}qwen3\_\allowbreak{}5\_\allowbreak{}moe/\allowbreak{}serving.py} \\
\texttt{QWEN38\_\allowbreak{}Q35M\_\allowbreak{}DFLASH2\_\allowbreak{}HOT} & \texttt{0} & \texttt{engine/\allowbreak{}models/\allowbreak{}qwen3\_\allowbreak{}5\_\allowbreak{}moe/\allowbreak{}dflash2.py} \\
\texttt{QWEN38\_\allowbreak{}Q35M\_\allowbreak{}DFLASH2\_\allowbreak{}NVFP4} & \texttt{0} & \texttt{engine/\allowbreak{}models/\allowbreak{}qwen3\_\allowbreak{}5\_\allowbreak{}moe/\allowbreak{}dflash2.py} \\
\texttt{QWEN38\_\allowbreak{}Q35M\_\allowbreak{}KV\_\allowbreak{}FP8} & \texttt{0} & \texttt{engine/\allowbreak{}models/\allowbreak{}qwen3\_\allowbreak{}5\_\allowbreak{}moe/\allowbreak{}engine.py} \\
\texttt{QWEN38\_\allowbreak{}Q35M\_\allowbreak{}LOOKUP} & \texttt{1} & \texttt{engine/\allowbreak{}models/\allowbreak{}qwen3\_\allowbreak{}5\_\allowbreak{}moe/\allowbreak{}serving.py} \\
\texttt{QWEN38\_\allowbreak{}Q35M\_\allowbreak{}SIBLINGS} & \texttt{0} & \texttt{engine/\allowbreak{}models/\allowbreak{}qwen3\_\allowbreak{}5\_\allowbreak{}moe/\allowbreak{}serving.py} \\
\texttt{QWEN38\_\allowbreak{}RANKK} & \texttt{1} & \texttt{engine/\allowbreak{}model.py} \\
\texttt{QWEN38\_\allowbreak{}REGENERATE} & \texttt{sample} & \texttt{server/\allowbreak{}app.py} \\
\texttt{QWEN38\_\allowbreak{}REGEN\_\allowbreak{}TEMPERATURE} & \texttt{0.7} & \texttt{server/\allowbreak{}app.py} \\
\texttt{QWEN38\_\allowbreak{}REGEN\_\allowbreak{}TOP\_\allowbreak{}K} & \texttt{20} & \texttt{server/\allowbreak{}app.py} \\
\texttt{QWEN38\_\allowbreak{}REGEN\_\allowbreak{}TOP\_\allowbreak{}P} & \texttt{0.95} & \texttt{server/\allowbreak{}app.py} \\
\texttt{QWEN38\_\allowbreak{}REREAD\_\allowbreak{}SAY\_\allowbreak{}TOKENS} & \texttt{8192} & \texttt{server/\allowbreak{}app.py} \\
\texttt{QWEN38\_\allowbreak{}ROUND\_\allowbreak{}LOG} & \texttt{none} & \texttt{server/\allowbreak{}roundlog.py} \\
\texttt{QWEN38\_\allowbreak{}ROUND\_\allowbreak{}LOG\_\allowbreak{}HIDDEN} & \texttt{0} & \texttt{server/\allowbreak{}roundlog.py} \\
\texttt{QWEN38\_\allowbreak{}ROUND\_\allowbreak{}LOG\_\allowbreak{}LATTICE} & \texttt{0} & \texttt{server/\allowbreak{}roundlog.py} \\
\texttt{QWEN38\_\allowbreak{}ROUND\_\allowbreak{}PLAN} & \texttt{off} & \texttt{engine/\allowbreak{}models/\allowbreak{}qwen3\_\allowbreak{}5/\allowbreak{}roundplan.py} \\
\texttt{QWEN38\_\allowbreak{}ROUND\_\allowbreak{}PLAN\_\allowbreak{}BUCKETS} & \texttt{4,\allowbreak{}8,\allowbreak{}16} & \texttt{engine/\allowbreak{}models/\allowbreak{}qwen3\_\allowbreak{}5/\allowbreak{}roundplan.py} \\
\texttt{QWEN38\_\allowbreak{}ROUND\_\allowbreak{}PLAN\_\allowbreak{}PRICES} & \texttt{none} & \texttt{engine/\allowbreak{}models/\allowbreak{}qwen3\_\allowbreak{}5/\allowbreak{}roundplan.py} \\
\texttt{QWEN38\_\allowbreak{}ROUND\_\allowbreak{}PLAN\_\allowbreak{}SAVE\_\allowbreak{}EVERY} & \texttt{512} & \texttt{engine/\allowbreak{}models/\allowbreak{}qwen3\_\allowbreak{}5/\allowbreak{}roundplan.py} \\
\texttt{QWEN38\_\allowbreak{}ROUND\_\allowbreak{}PLAN\_\allowbreak{}STATE} & \texttt{none} & \texttt{engine/\allowbreak{}models/\allowbreak{}qwen3\_\allowbreak{}5/\allowbreak{}roundplan.py} \\
\texttt{QWEN38\_\allowbreak{}SAMPLED\_\allowbreak{}CUT} & \texttt{0} & \texttt{engine/\allowbreak{}router.py} \\
\texttt{QWEN38\_\allowbreak{}SAMPLE\_\allowbreak{}BATCH} & \texttt{1} & \texttt{engine/\allowbreak{}sample.py} \\
\texttt{QWEN38\_\allowbreak{}SAMPLE\_\allowbreak{}KEYED\_\allowbreak{}KERNEL} & \texttt{0} & \texttt{engine/\allowbreak{}sample.py} \\
\texttt{QWEN38\_\allowbreak{}SAMPLE\_\allowbreak{}NUCLEUS} & \texttt{1} & \texttt{engine/\allowbreak{}sample.py} \\
\texttt{QWEN38\_\allowbreak{}SAMPLE\_\allowbreak{}SPARSE\_\allowbreak{}KEYED} & \texttt{1} & \texttt{engine/\allowbreak{}sample.py} \\
\texttt{QWEN38\_\allowbreak{}SAMPLE\_\allowbreak{}TOPP\_\allowbreak{}K} & \texttt{1024} & \texttt{engine/\allowbreak{}sample.py} \\
\texttt{QWEN38\_\allowbreak{}SCALE\_\allowbreak{}ON} & \texttt{weight} & \texttt{tools/\allowbreak{}fp8\_\allowbreak{}linear.py} \\
\texttt{QWEN38\_\allowbreak{}SESSION\_\allowbreak{}GENERATIONS} & \texttt{on} & \texttt{app/\allowbreak{}tandemapp/\allowbreak{}conductor/\allowbreak{}router.py} \\
\texttt{QWEN38\_\allowbreak{}SHADOW\_\allowbreak{}ASK} & \texttt{idle} & \texttt{app/\allowbreak{}tandemapp/\allowbreak{}conductor/\allowbreak{}deferred.py} \\
\texttt{QWEN38\_\allowbreak{}SHADOW\_\allowbreak{}ASK\_\allowbreak{}IDLE\_\allowbreak{}S} & \texttt{5} & \texttt{app/\allowbreak{}tandemapp/\allowbreak{}conductor/\allowbreak{}deferred.py} \\
\texttt{QWEN38\_\allowbreak{}SHADOW\_\allowbreak{}DIR} & \texttt{none} & \texttt{server/\allowbreak{}conductor/\allowbreak{}shadow.py} \\
\texttt{QWEN38\_\allowbreak{}SHADOW\_\allowbreak{}DISK} & \texttt{none} & \texttt{auto} \textbullet{} \texttt{server/\allowbreak{}conductor/\allowbreak{}shadow.py} \\
\texttt{QWEN38\_\allowbreak{}SHADOW\_\allowbreak{}DISK\_\allowbreak{}GB} & \texttt{20} & \texttt{20} \textbullet{} \texttt{server/\allowbreak{}conductor/\allowbreak{}shadow.py} \\
\texttt{QWEN38\_\allowbreak{}SHADOW\_\allowbreak{}GB} & \texttt{4} & \texttt{4} \textbullet{} \texttt{server/\allowbreak{}conductor/\allowbreak{}shadow.py} \\
\texttt{QWEN38\_\allowbreak{}SHADOW\_\allowbreak{}GRID} & \texttt{512} & \texttt{512} \textbullet{} \texttt{server/\allowbreak{}conductor/\allowbreak{}shadow.py} \\
\texttt{QWEN38\_\allowbreak{}SHADOW\_\allowbreak{}MODEL} & \texttt{none} & \texttt{server/\allowbreak{}conductor/\allowbreak{}shadow.py} \\
\texttt{QWEN38\_\allowbreak{}SHADOW\_\allowbreak{}POOL\_\allowbreak{}SLACK\_\allowbreak{}GIB} & \texttt{2} & \texttt{server/\allowbreak{}conductor/\allowbreak{}shadow.py} \\
\texttt{QWEN38\_\allowbreak{}SIDE\_\allowbreak{}CARRY} & \texttt{on} & \texttt{server/\allowbreak{}app.py} \\
\texttt{QWEN38\_\allowbreak{}SIDE\_\allowbreak{}CONTINUE} & \texttt{view} & \texttt{server/\allowbreak{}app.py} \\
\texttt{QWEN38\_\allowbreak{}SIDE\_\allowbreak{}CONV} & \texttt{on} & \texttt{server/\allowbreak{}app.py} \\
\texttt{QWEN38\_\allowbreak{}SIDE\_\allowbreak{}DISK} & \texttt{read} & \texttt{server/\allowbreak{}app.py} \\
\texttt{QWEN38\_\allowbreak{}SIDE\_\allowbreak{}GUESTS} & \texttt{on} & \texttt{server/\allowbreak{}app.py} \\
\texttt{QWEN38\_\allowbreak{}SIDE\_\allowbreak{}GUEST\_\allowbreak{}END} & \texttt{promote} & \texttt{server/\allowbreak{}app.py} \\
\texttt{QWEN38\_\allowbreak{}SIDE\_\allowbreak{}GUEST\_\allowbreak{}IDLE\_\allowbreak{}S} & \texttt{300} & \texttt{server/\allowbreak{}app.py} \\
\texttt{QWEN38\_\allowbreak{}SIDE\_\allowbreak{}HOME\_\allowbreak{}WAIT\_\allowbreak{}S} & \texttt{10} & \texttt{server/\allowbreak{}app.py} \\
\texttt{QWEN38\_\allowbreak{}SIDE\_\allowbreak{}IMAGES} & \texttt{on} & \texttt{server/\allowbreak{}app.py} \\
\texttt{QWEN38\_\allowbreak{}SIDE\_\allowbreak{}LRU} & \texttt{on} & \texttt{server/\allowbreak{}app.py} \\
\texttt{QWEN38\_\allowbreak{}SIDE\_\allowbreak{}MIN\_\allowbreak{}ROOM} & \texttt{16384} & \texttt{server/\allowbreak{}app.py} \\
\texttt{QWEN38\_\allowbreak{}SIDE\_\allowbreak{}NEW\_\allowbreak{}FIRST} & \texttt{on} & \texttt{server/\allowbreak{}app.py} \\
\texttt{QWEN38\_\allowbreak{}SIDE\_\allowbreak{}NEW\_\allowbreak{}WAIT\_\allowbreak{}S} & \texttt{20} & \texttt{server/\allowbreak{}app.py} \\
\texttt{QWEN38\_\allowbreak{}SIDE\_\allowbreak{}RESIDENT\_\allowbreak{}GB} & \texttt{0.5} & \texttt{server/\allowbreak{}app.py} \\
\texttt{QWEN38\_\allowbreak{}SIDE\_\allowbreak{}SLOT0\_\allowbreak{}WAIT\_\allowbreak{}S} & \texttt{3} & \texttt{server/\allowbreak{}app.py} \\
\texttt{QWEN38\_\allowbreak{}SIDE\_\allowbreak{}STASH\_\allowbreak{}GB} & \texttt{1.0} & \texttt{server/\allowbreak{}app.py} \\
\texttt{QWEN38\_\allowbreak{}SKINNY\_\allowbreak{}LDW} & \texttt{0} & \texttt{tools/\allowbreak{}nvfp4\_\allowbreak{}skinny.py} \\
\texttt{QWEN38\_\allowbreak{}SKINNY\_\allowbreak{}PDL} & \texttt{0} & \texttt{tools/\allowbreak{}nvfp4\_\allowbreak{}skinny.py} \\
\texttt{QWEN38\_\allowbreak{}SKINNY\_\allowbreak{}SRUN} & \texttt{0} & \texttt{1} \textbullet{} \texttt{tools/\allowbreak{}nvfp4\_\allowbreak{}skinny.py} \\
\texttt{QWEN38\_\allowbreak{}SKINNY\_\allowbreak{}SSTUB} & \texttt{0} & \texttt{tools/\allowbreak{}nvfp4\_\allowbreak{}skinny.py} \\
\texttt{QWEN38\_\allowbreak{}SKINNY\_\allowbreak{}TILES} & \texttt{none} & \texttt{(a path)} \textbullet{} \texttt{tools/\allowbreak{}nvfp4\_\allowbreak{}skinny.py} \\
\texttt{QWEN38\_\allowbreak{}SKINNY\_\allowbreak{}TILES\_\allowbreak{}B} & \texttt{none} & \texttt{tools/\allowbreak{}nvfp4\_\allowbreak{}skinny.py} \\
\texttt{QWEN38\_\allowbreak{}SKINNY\_\allowbreak{}TILES\_\allowbreak{}C} & \texttt{none} & \texttt{tools/\allowbreak{}nvfp4\_\allowbreak{}skinny.py} \\
\texttt{QWEN38\_\allowbreak{}SKINNY\_\allowbreak{}TILES\_\allowbreak{}WIDE} & \texttt{none} & \texttt{(a path)} \textbullet{} \texttt{tools/\allowbreak{}nvfp4\_\allowbreak{}skinny.py} \\
\texttt{QWEN38\_\allowbreak{}SKINNY\_\allowbreak{}TILES\_\allowbreak{}WIDE\_\allowbreak{}B} & \texttt{none} & \texttt{tools/\allowbreak{}nvfp4\_\allowbreak{}skinny.py} \\
\texttt{QWEN38\_\allowbreak{}SKINNY\_\allowbreak{}XSTUB} & \texttt{0} & \texttt{tools/\allowbreak{}nvfp4\_\allowbreak{}skinny.py} \\
\texttt{QWEN38\_\allowbreak{}SLAB\_\allowbreak{}TRITON} & \texttt{1} & \texttt{engine/\allowbreak{}model.py} \\
\texttt{QWEN38\_\allowbreak{}SLOT0\_\allowbreak{}KEEP\_\allowbreak{}S} & \texttt{1800} & \texttt{server/\allowbreak{}app.py} \\
\texttt{QWEN38\_\allowbreak{}SPILL\_\allowbreak{}PARK} & \texttt{on} & \texttt{engine/\allowbreak{}cache.py} \\
\texttt{QWEN38\_\allowbreak{}STAIR\_\allowbreak{}MAX\_\allowbreak{}NODES} & \texttt{none} & \texttt{engine/\allowbreak{}lenrouter.py} \\
\texttt{QWEN38\_\allowbreak{}STAIR\_\allowbreak{}RHO} & \texttt{none} & \texttt{(a path)} \textbullet{} \texttt{engine/\allowbreak{}lenrouter.py} \\
\texttt{QWEN38\_\allowbreak{}STAIR\_\allowbreak{}SKIP} & \texttt{0} & \texttt{1} \textbullet{} \texttt{engine/\allowbreak{}lenrouter.py} \\
\texttt{QWEN38\_\allowbreak{}STAIR\_\allowbreak{}TABLES} & \texttt{none} & \texttt{(a path)} \textbullet{} \texttt{engine/\allowbreak{}lenrouter.py} \\
\texttt{QWEN38\_\allowbreak{}STAIR\_\allowbreak{}TEMP\_\allowbreak{}SAMPLED} & \texttt{none} & \texttt{engine/\allowbreak{}lenrouter.py} \\
\texttt{QWEN38\_\allowbreak{}STATUS\_\allowbreak{}CLIENTS} & \texttt{open-\allowbreak{}webui} & \texttt{app/\allowbreak{}tandemapp/\allowbreak{}conductor/\allowbreak{}states.py} \\
\texttt{QWEN38\_\allowbreak{}SUFFIX\_\allowbreak{}STORE} & \texttt{\textasciitilde{}/\allowbreak{}.qwen38-\allowbreak{}spark-\allowbreak{}engine/\allowbreak{}suffix} & \texttt{server/\allowbreak{}app.py} \\
\texttt{QWEN38\_\allowbreak{}TANDEM\_\allowbreak{}SCRATCH} & \texttt{none} & \texttt{app/\allowbreak{}tandemapp/\allowbreak{}conductor/\allowbreak{}packages.py} \\
\texttt{QWEN38\_\allowbreak{}TANDEM\_\allowbreak{}URL} & \texttt{none} & \texttt{app/\allowbreak{}tandemapp/\allowbreak{}conductor/\allowbreak{}packages.py} \\
\texttt{QWEN38\_\allowbreak{}TOOL\_\allowbreak{}FIT} & \texttt{on} & \texttt{app/\allowbreak{}tandemapp/\allowbreak{}conductor/\allowbreak{}router.py} \\
\texttt{QWEN38\_\allowbreak{}TREE\_\allowbreak{}ALIAS\_\allowbreak{}STATE} & \texttt{0} & \texttt{0} \textbullet{} \texttt{engine/\allowbreak{}model.py} \\
\texttt{QWEN38\_\allowbreak{}TREE\_\allowbreak{}CHAIN\_\allowbreak{}DELEGATE} & \texttt{1} & \texttt{engine/\allowbreak{}model.py} \\
\texttt{QWEN38\_\allowbreak{}TREE\_\allowbreak{}HOST\_\allowbreak{}DEPTH} & \texttt{0} & \texttt{1} \textbullet{} \texttt{engine/\allowbreak{}model.py} \\
\texttt{QWEN38\_\allowbreak{}TREE\_\allowbreak{}MS} & \texttt{none} & \texttt{8:\allowbreak{}77.9,\allowbreak{}16:\allowbreak{}78.9,\allowbreak{}24:\allowbreak{}87.0,\allowbreak{}32:\allowbreak{}97.0} \textbullet{} \texttt{engine/\allowbreak{}router.py} \\
\texttt{QWEN38\_\allowbreak{}TREE\_\allowbreak{}NATIVE} & \texttt{1} & \texttt{engine/\allowbreak{}native\_\allowbreak{}tree.py} \\
\texttt{QWEN38\_\allowbreak{}TREE\_\allowbreak{}NODES} & \texttt{0} & \texttt{24} \textbullet{} \texttt{engine/\allowbreak{}router.py} \\
\texttt{QWEN38\_\allowbreak{}TREE\_\allowbreak{}NODES\_\allowbreak{}NARROW} & \texttt{0} & \texttt{16} \textbullet{} \texttt{engine/\allowbreak{}router.py} \\
\texttt{QWEN38\_\allowbreak{}TREE\_\allowbreak{}WIDE\_\allowbreak{}AFTER} & \texttt{0} & \texttt{32} \textbullet{} \texttt{engine/\allowbreak{}lenrouter.py} \\
\texttt{QWEN38\_\allowbreak{}TWO\_\allowbreak{}STREAM} & \texttt{0} & \texttt{engine/\allowbreak{}model.py} \\
\texttt{QWEN38\_\allowbreak{}UT\_\allowbreak{}INVERSE} & \texttt{1} & \texttt{engine/\allowbreak{}gdn.py} \\
\texttt{QWEN38\_\allowbreak{}VERIFY\_\allowbreak{}GRAPH} & \texttt{0} & \texttt{1} \textbullet{} \texttt{engine/\allowbreak{}model.py} \\
\texttt{QWEN38\_\allowbreak{}VERIFY\_\allowbreak{}ROWS} & \texttt{16} & \texttt{32} \textbullet{} \texttt{engine/\allowbreak{}model.py} \\
\texttt{QWEN38\_\allowbreak{}VERIFY\_\allowbreak{}TILES} & \texttt{1} & \texttt{tools/\allowbreak{}nvfp4\_\allowbreak{}verify\_\allowbreak{}tiles.py} \\
\texttt{QWEN38\_\allowbreak{}WARM\_\allowbreak{}PREFILL\_\allowbreak{}TOKENS} & \texttt{2600} & \texttt{server/\allowbreak{}app.py} \\
\texttt{QWEN38\_\allowbreak{}WINDOW\_\allowbreak{}MLP} & \texttt{0} & \texttt{engine/\allowbreak{}model.py} \\
\texttt{QWEN38\_\allowbreak{}WINDOW\_\allowbreak{}NORM\_\allowbreak{}ROT} & \texttt{1} & \texttt{engine/\allowbreak{}model.py} \\
\texttt{QWEN38\_\allowbreak{}WINDOW\_\allowbreak{}ROWWISE} & \texttt{1} & \texttt{engine/\allowbreak{}model.py} \\
\texttt{QWEN38\_\allowbreak{}WY\_\allowbreak{}BLOCK} & \texttt{8} & \texttt{tools/\allowbreak{}gdn\_\allowbreak{}wy\_\allowbreak{}kernels.py} \\
\texttt{QWEN38\_\allowbreak{}WY\_\allowbreak{}BV} & \texttt{16} & \texttt{tools/\allowbreak{}gdn\_\allowbreak{}wy\_\allowbreak{}kernels.py} \\
\texttt{QWEN38\_\allowbreak{}WY\_\allowbreak{}DBG} & \texttt{0} & \texttt{tools/\allowbreak{}gdn\_\allowbreak{}wy\_\allowbreak{}kernels.py} \\
\texttt{QWEN38\_\allowbreak{}WY\_\allowbreak{}PREC} & \texttt{ieee} & \texttt{tools/\allowbreak{}gdn\_\allowbreak{}wy\_\allowbreak{}kernels.py} \\
\texttt{QWEN38\_\allowbreak{}WY\_\allowbreak{}WARPS} & \texttt{4} & \texttt{tools/\allowbreak{}gdn\_\allowbreak{}wy\_\allowbreak{}kernels.py} \\
\texttt{QWEN38\_\allowbreak{}WY\_\allowbreak{}WARPS\_\allowbreak{}PREP} & \texttt{4} & \texttt{tools/\allowbreak{}gdn\_\allowbreak{}wy\_\allowbreak{}kernels.py} \\
\end{reftable}


\section{Server variables}\index{QSE variables}

% Generated by manual/tools/gen_reference.py from the code. Do not edit by hand.

\begin{reftabletwo}{QSE\_* variables: the server's settings from the environment or serve.env}{tab:ref-qse}{\hd{Variable} & \hd{Read in}}
\texttt{QSE\_\allowbreak{}ADMIN\_\allowbreak{}TOKEN} & \texttt{ops/\allowbreak{}make-\allowbreak{}secrets.sh,\allowbreak{} ops/\allowbreak{}serve.env,\allowbreak{} ops/\allowbreak{}start.sh and 5 more} \\
\texttt{QSE\_\allowbreak{}ANSWER\_\allowbreak{}RESERVE} & \texttt{server/\allowbreak{}app.py} \\
\texttt{QSE\_\allowbreak{}ASSETS} & \texttt{engine/\allowbreak{}models/\allowbreak{}qwen3\_\allowbreak{}5/\allowbreak{}community.py,\allowbreak{} ops/\allowbreak{}profiles/\allowbreak{}\_\allowbreak{}flashnext.env,\allowbreak{} ops/\allowbreak{}profiles/\allowbreak{}\_\allowbreak{}flint.env and 16 more} \\
\texttt{QSE\_\allowbreak{}BACKENDS} & \texttt{ops/\allowbreak{}engines.sh,\allowbreak{} ops/\allowbreak{}serve.env,\allowbreak{} ops/\allowbreak{}start.sh and 2 more} \\
\texttt{QSE\_\allowbreak{}BARRIER\_\allowbreak{}IDLE\_\allowbreak{}OWNER} & \texttt{ops/\allowbreak{}serve.env,\allowbreak{} server/\allowbreak{}slots.py} \\
\texttt{QSE\_\allowbreak{}CAPTURE\_\allowbreak{}GUARD\_\allowbreak{}DEBUG} & \texttt{engine/\allowbreak{}capture\_\allowbreak{}lock.py} \\
\texttt{QSE\_\allowbreak{}CHATD} & \texttt{install.sh,\allowbreak{} ops/\allowbreak{}engines.sh,\allowbreak{} ops/\allowbreak{}serve.env and 4 more} \\
\texttt{QSE\_\allowbreak{}CHATD\_\allowbreak{}ASKED} & \texttt{install.sh,\allowbreak{} ops/\allowbreak{}serve.env,\allowbreak{} ops/\allowbreak{}start.sh} \\
\texttt{QSE\_\allowbreak{}CHATD\_\allowbreak{}ENV} & \texttt{install.sh} \\
\texttt{QSE\_\allowbreak{}CHAT\_\allowbreak{}DATA} & \texttt{install.sh,\allowbreak{} server/\allowbreak{}setup/\allowbreak{}doctor.py} \\
\texttt{QSE\_\allowbreak{}CHAT\_\allowbreak{}FORWARD\_\allowbreak{}KEY} & \texttt{ops/\allowbreak{}make-\allowbreak{}secrets.sh} \\
\texttt{QSE\_\allowbreak{}CHAT\_\allowbreak{}URL} & \texttt{server/\allowbreak{}api/\allowbreak{}\_\allowbreak{}\_\allowbreak{}init\_\allowbreak{}\_\allowbreak{}.py,\allowbreak{} server/\allowbreak{}api/\allowbreak{}chat\_\allowbreak{}url.py,\allowbreak{} server/\allowbreak{}mode.py} \\
\texttt{QSE\_\allowbreak{}COMMUNITY\_\allowbreak{}DIR} & \texttt{server/\allowbreak{}setup/\allowbreak{}community.py} \\
\texttt{QSE\_\allowbreak{}COMMUNITY\_\allowbreak{}ENV\_\allowbreak{}DIR} & \texttt{ops/\allowbreak{}serve.env,\allowbreak{} ops/\allowbreak{}use-\allowbreak{}profile.sh} \\
\texttt{QSE\_\allowbreak{}COMMUNITY\_\allowbreak{}URL} & \texttt{ops/\allowbreak{}serve.env,\allowbreak{} server/\allowbreak{}setup/\allowbreak{}community.py} \\
\texttt{QSE\_\allowbreak{}COMPACTION} & \texttt{server/\allowbreak{}app.py} \\
\texttt{QSE\_\allowbreak{}COMPACTION\_\allowbreak{}IDLE\_\allowbreak{}S} & \texttt{server/\allowbreak{}app.py} \\
\texttt{QSE\_\allowbreak{}COMPACTION\_\allowbreak{}MIN\_\allowbreak{}TOKENS} & \texttt{server/\allowbreak{}app.py} \\
\texttt{QSE\_\allowbreak{}COMPACTION\_\allowbreak{}PREPARE} & \texttt{server/\allowbreak{}app.py} \\
\texttt{QSE\_\allowbreak{}CONDUCTOR\_\allowbreak{}CONFIG} & \texttt{ops/\allowbreak{}engines.sh,\allowbreak{} ops/\allowbreak{}serve.env,\allowbreak{} ops/\allowbreak{}start.sh and 4 more} \\
\texttt{QSE\_\allowbreak{}CONDUCTOR\_\allowbreak{}ROOT} & \texttt{ops/\allowbreak{}roles.sh,\allowbreak{} server/\allowbreak{}conductor/\allowbreak{}residency.py} \\
\texttt{QSE\_\allowbreak{}CORPUS\_\allowbreak{}BUILD} & \texttt{ops/\allowbreak{}use-\allowbreak{}profile.sh} \\
\texttt{QSE\_\allowbreak{}CORRECTOR\_\allowbreak{}HF} & \texttt{engine/\allowbreak{}models/\allowbreak{}flashnext/\allowbreak{}corrector.py,\allowbreak{} ops/\allowbreak{}fetch\_\allowbreak{}weights.py,\allowbreak{} ops/\allowbreak{}profiles/\allowbreak{}\_\allowbreak{}flashnext.env and 2 more} \\
\texttt{QSE\_\allowbreak{}CORRECTOR\_\allowbreak{}PATH} & \texttt{engine/\allowbreak{}models/\allowbreak{}flashnext/\allowbreak{}corrector.py,\allowbreak{} ops/\allowbreak{}fetch\_\allowbreak{}weights.py,\allowbreak{} ops/\allowbreak{}profiles/\allowbreak{}flashnext-\allowbreak{}3.05.env and 1 more} \\
\texttt{QSE\_\allowbreak{}CORRECTOR\_\allowbreak{}REVISION} & \texttt{engine/\allowbreak{}models/\allowbreak{}flashnext/\allowbreak{}corrector.py,\allowbreak{} ops/\allowbreak{}fetch\_\allowbreak{}weights.py,\allowbreak{} ops/\allowbreak{}profiles/\allowbreak{}flashnext-\allowbreak{}3.05.env and 1 more} \\
\texttt{QSE\_\allowbreak{}CORRECTOR\_\allowbreak{}SHA256} & \texttt{engine/\allowbreak{}models/\allowbreak{}flashnext/\allowbreak{}corrector.py,\allowbreak{} ops/\allowbreak{}fetch\_\allowbreak{}weights.py,\allowbreak{} ops/\allowbreak{}profiles/\allowbreak{}flashnext-\allowbreak{}3.05.env and 1 more} \\
\texttt{QSE\_\allowbreak{}DASHBOARD\_\allowbreak{}LOGIN} & \texttt{ops/\allowbreak{}serve.env,\allowbreak{} ops/\allowbreak{}start.sh,\allowbreak{} server/\allowbreak{}api/\allowbreak{}auth.py and 1 more} \\
\texttt{QSE\_\allowbreak{}DISK\_\allowbreak{}ADOPT\_\allowbreak{}FROM} & \texttt{ops/\allowbreak{}serve.env,\allowbreak{} server/\allowbreak{}diskcache.py} \\
\texttt{QSE\_\allowbreak{}DISK\_\allowbreak{}CACHE} & \texttt{engine/\allowbreak{}models/\allowbreak{}nemotron\_\allowbreak{}h/\allowbreak{}community.py,\allowbreak{} engine/\allowbreak{}models/\allowbreak{}qwen3\_\allowbreak{}5/\allowbreak{}community.py,\allowbreak{} ops/\allowbreak{}engines.sh and 11 more} \\
\texttt{QSE\_\allowbreak{}ENGINE\_\allowbreak{}ALLOC\_\allowbreak{}CONF} & \texttt{engine/\allowbreak{}cache.py,\allowbreak{} ops/\allowbreak{}profiles/\allowbreak{}\_\allowbreak{}flashnext.env,\allowbreak{} ops/\allowbreak{}start.sh} \\
\texttt{QSE\_\allowbreak{}ENGINE\_\allowbreak{}FLAGS} & \texttt{engine/\allowbreak{}models/\allowbreak{}nemotron\_\allowbreak{}h/\allowbreak{}community.py,\allowbreak{} engine/\allowbreak{}models/\allowbreak{}qwen3\_\allowbreak{}5/\allowbreak{}community.py,\allowbreak{} ops/\allowbreak{}engines.sh and 9 more} \\
\texttt{QSE\_\allowbreak{}ENGINE\_\allowbreak{}PORT} & \texttt{ops/\allowbreak{}engine\_\allowbreak{}status.py} \\
\texttt{QSE\_\allowbreak{}EXL3\_\allowbreak{}BPW} & \texttt{ops/\allowbreak{}profiles/\allowbreak{}\_\allowbreak{}exl3.env,\allowbreak{} ops/\allowbreak{}profiles/\allowbreak{}exl3-\allowbreak{}2.5.env,\allowbreak{} ops/\allowbreak{}profiles/\allowbreak{}exl3-\allowbreak{}3.0.env and 4 more} \\
\texttt{QSE\_\allowbreak{}EXT} & \texttt{ops/\allowbreak{}serve.env,\allowbreak{} server/\allowbreak{}app.py,\allowbreak{} server/\allowbreak{}conductor/\allowbreak{}residency.py and 2 more} \\
\texttt{QSE\_\allowbreak{}EXT\_\allowbreak{}ARGS} & \texttt{server/\allowbreak{}app.py,\allowbreak{} server/\allowbreak{}ext.py} \\
\texttt{QSE\_\allowbreak{}FAKE\_\allowbreak{}CAPTURE\_\allowbreak{}HOLD\_\allowbreak{}S} & \texttt{server/\allowbreak{}fake\_\allowbreak{}engine.py} \\
\texttt{QSE\_\allowbreak{}FAKE\_\allowbreak{}CAPTURE\_\allowbreak{}MARK} & \texttt{server/\allowbreak{}fake\_\allowbreak{}engine.py} \\
\texttt{QSE\_\allowbreak{}FAKE\_\allowbreak{}CARD\_\allowbreak{}DIR} & \texttt{server/\allowbreak{}fake\_\allowbreak{}engine.py} \\
\texttt{QSE\_\allowbreak{}FAKE\_\allowbreak{}PROMPTS} & \texttt{server/\allowbreak{}fake\_\allowbreak{}engine.py} \\
\texttt{QSE\_\allowbreak{}FETCH\_\allowbreak{}ASSETS} & \texttt{ops/\allowbreak{}use-\allowbreak{}profile.sh,\allowbreak{} server/\allowbreak{}setup/\allowbreak{}apply.py} \\
\texttt{QSE\_\allowbreak{}FLINT\_\allowbreak{}GROUPED} & \texttt{engine/\allowbreak{}models/\allowbreak{}flint/\allowbreak{}engine.py} \\
\texttt{QSE\_\allowbreak{}FLINT\_\allowbreak{}VARIANT} & \texttt{ops/\allowbreak{}profiles/\allowbreak{}\_\allowbreak{}flint.env,\allowbreak{} ops/\allowbreak{}profiles/\allowbreak{}flint-\allowbreak{}bf16.env,\allowbreak{} ops/\allowbreak{}profiles/\allowbreak{}flint-\allowbreak{}exl3-\allowbreak{}4.0.env and 1 more} \\
\texttt{QSE\_\allowbreak{}FN\_\allowbreak{}BPW} & \texttt{ops/\allowbreak{}profiles/\allowbreak{}\_\allowbreak{}flashnext.env,\allowbreak{} ops/\allowbreak{}profiles/\allowbreak{}flashnext-\allowbreak{}2.05.env,\allowbreak{} ops/\allowbreak{}profiles/\allowbreak{}flashnext-\allowbreak{}3.05.env and 3 more} \\
\texttt{QSE\_\allowbreak{}FN\_\allowbreak{}TABLE} & \texttt{ops/\allowbreak{}profiles/\allowbreak{}\_\allowbreak{}flashnext.env,\allowbreak{} ops/\allowbreak{}profiles/\allowbreak{}flashnext-\allowbreak{}2.05.env,\allowbreak{} ops/\allowbreak{}profiles/\allowbreak{}flashnext-\allowbreak{}3.05.env and 3 more} \\
\texttt{QSE\_\allowbreak{}FN\_\allowbreak{}VISION\_\allowbreak{}FILE} & \texttt{ops/\allowbreak{}profiles/\allowbreak{}\_\allowbreak{}flashnext.env,\allowbreak{} ops/\allowbreak{}profiles/\allowbreak{}flashnext-\allowbreak{}4.05.env} \\
\texttt{QSE\_\allowbreak{}GRAPH\_\allowbreak{}POOL\_\allowbreak{}GIB} & \texttt{server/\allowbreak{}app.py} \\
\texttt{QSE\_\allowbreak{}GRAPH\_\allowbreak{}SET\_\allowbreak{}MIN\_\allowbreak{}IDLE\_\allowbreak{}S} & \texttt{server/\allowbreak{}app.py} \\
\texttt{QSE\_\allowbreak{}GUARD\_\allowbreak{}MARGIN\_\allowbreak{}GB} & \texttt{server/\allowbreak{}app.py} \\
\texttt{QSE\_\allowbreak{}HANDLER\_\allowbreak{}TIMEOUT\_\allowbreak{}S} & \texttt{server/\allowbreak{}app.py} \\
\texttt{QSE\_\allowbreak{}HEALTH\_\allowbreak{}TTL\_\allowbreak{}MS} & \texttt{ops/\allowbreak{}profiles/\allowbreak{}\_\allowbreak{}flashnext.env,\allowbreak{} server/\allowbreak{}app.py,\allowbreak{} server/\allowbreak{}health\_\allowbreak{}cache.py} \\
\texttt{QSE\_\allowbreak{}IMAGE\_\allowbreak{}BACKEND} & \texttt{server/\allowbreak{}app.py} \\
\texttt{QSE\_\allowbreak{}KV\_\allowbreak{}ADMIT\_\allowbreak{}ROWS} & \texttt{ops/\allowbreak{}serve.env,\allowbreak{} server/\allowbreak{}app.py} \\
\texttt{QSE\_\allowbreak{}KV\_\allowbreak{}HARD\_\allowbreak{}FLOOR\_\allowbreak{}GB} & \texttt{ops/\allowbreak{}profiles/\allowbreak{}\_\allowbreak{}flashnext.env,\allowbreak{} ops/\allowbreak{}serve.env,\allowbreak{} server/\allowbreak{}app.py} \\
\texttt{QSE\_\allowbreak{}KV\_\allowbreak{}MARGIN\_\allowbreak{}ROWS} & \texttt{server/\allowbreak{}app.py} \\
\texttt{QSE\_\allowbreak{}KV\_\allowbreak{}POOL} & \texttt{engine/\allowbreak{}kvpool.py,\allowbreak{} ops/\allowbreak{}engines.sh,\allowbreak{} ops/\allowbreak{}profiles/\allowbreak{}\_\allowbreak{}flashnext.env and 7 more} \\
\texttt{QSE\_\allowbreak{}KV\_\allowbreak{}POOL\_\allowbreak{}MIN\_\allowbreak{}ROWS} & \texttt{engine/\allowbreak{}kvpool.py} \\
\texttt{QSE\_\allowbreak{}KV\_\allowbreak{}RECOMPUTE} & \texttt{ops/\allowbreak{}serve.env,\allowbreak{} server/\allowbreak{}app.py} \\
\texttt{QSE\_\allowbreak{}KV\_\allowbreak{}TEST\_\allowbreak{}REFUSE\_\allowbreak{}AFTER} & \texttt{server/\allowbreak{}round\_\allowbreak{}runner.py} \\
\texttt{QSE\_\allowbreak{}KV\_\allowbreak{}TEST\_\allowbreak{}REFUSE\_\allowbreak{}SLOTS} & \texttt{server/\allowbreak{}round\_\allowbreak{}runner.py} \\
\texttt{QSE\_\allowbreak{}LFM\_\allowbreak{}VARIANT} & \texttt{ops/\allowbreak{}profiles/\allowbreak{}\_\allowbreak{}lfm25.env,\allowbreak{} ops/\allowbreak{}profiles/\allowbreak{}lfm25-\allowbreak{}bf16.env,\allowbreak{} ops/\allowbreak{}profiles/\allowbreak{}lfm25-\allowbreak{}exl3-\allowbreak{}2.5.env and 5 more} \\
\texttt{QSE\_\allowbreak{}LIVE\_\allowbreak{}DECODE\_\allowbreak{}WARMUP\_\allowbreak{}TOKENS} & \texttt{ops/\allowbreak{}serve.env,\allowbreak{} server/\allowbreak{}api/\allowbreak{}live.py} \\
\texttt{QSE\_\allowbreak{}LIVE\_\allowbreak{}HZ} & \texttt{ops/\allowbreak{}serve.env,\allowbreak{} ops/\allowbreak{}start.sh,\allowbreak{} server/\allowbreak{}api/\allowbreak{}live.py and 1 more} \\
\texttt{QSE\_\allowbreak{}LIVE\_\allowbreak{}MOSTLY\_\allowbreak{}CACHED\_\allowbreak{}PCT} & \texttt{ops/\allowbreak{}serve.env,\allowbreak{} server/\allowbreak{}api/\allowbreak{}live.py} \\
\texttt{QSE\_\allowbreak{}LIVE\_\allowbreak{}MOSTLY\_\allowbreak{}CACHED\_\allowbreak{}TOKENS} & \texttt{ops/\allowbreak{}serve.env,\allowbreak{} server/\allowbreak{}api/\allowbreak{}live.py} \\
\texttt{QSE\_\allowbreak{}LIVE\_\allowbreak{}NOW\_\allowbreak{}WINDOW\_\allowbreak{}S} & \texttt{ops/\allowbreak{}serve.env,\allowbreak{} server/\allowbreak{}api/\allowbreak{}live.py} \\
\texttt{QSE\_\allowbreak{}LIVE\_\allowbreak{}PREFILL\_\allowbreak{}WARMUP\_\allowbreak{}CHUNKS} & \texttt{ops/\allowbreak{}serve.env,\allowbreak{} server/\allowbreak{}api/\allowbreak{}live.py} \\
\texttt{QSE\_\allowbreak{}MAX\_\allowbreak{}BODY\_\allowbreak{}MB} & \texttt{server/\allowbreak{}api/\allowbreak{}http\_\allowbreak{}body.py,\allowbreak{} server/\allowbreak{}app.py} \\
\texttt{QSE\_\allowbreak{}MEMINFO} & \texttt{ops/\allowbreak{}engines.sh} \\
\texttt{QSE\_\allowbreak{}MEM\_\allowbreak{}NEED\_\allowbreak{}GIB} & \texttt{ops/\allowbreak{}profiles/\allowbreak{}flashnext-\allowbreak{}2.05.env,\allowbreak{} ops/\allowbreak{}profiles/\allowbreak{}flashnext-\allowbreak{}3.05.env,\allowbreak{} ops/\allowbreak{}profiles/\allowbreak{}flashnext-\allowbreak{}4.05.env and 27 more} \\
\texttt{QSE\_\allowbreak{}MEM\_\allowbreak{}WAIT\_\allowbreak{}S} & \texttt{server/\allowbreak{}app.py} \\
\texttt{QSE\_\allowbreak{}METRICS\_\allowbreak{}TOKEN} & \texttt{ops/\allowbreak{}make-\allowbreak{}secrets.sh,\allowbreak{} ops/\allowbreak{}start.sh,\allowbreak{} server/\allowbreak{}api/\allowbreak{}auth.py} \\
\texttt{QSE\_\allowbreak{}MODE} & \texttt{install.sh,\allowbreak{} ops/\allowbreak{}engines.sh,\allowbreak{} ops/\allowbreak{}serve.env and 2 more} \\
\texttt{QSE\_\allowbreak{}NH\_\allowbreak{}DENSE\_\allowbreak{}ROWS} & \texttt{engine/\allowbreak{}models/\allowbreak{}nemotron\_\allowbreak{}h/\allowbreak{}engine.py} \\
\texttt{QSE\_\allowbreak{}NH\_\allowbreak{}DIR} & \texttt{engine/\allowbreak{}models/\allowbreak{}nemotron\_\allowbreak{}h/\allowbreak{}community.py,\allowbreak{} ops/\allowbreak{}profiles/\allowbreak{}\_\allowbreak{}nemotron.env,\allowbreak{} ops/\allowbreak{}profiles/\allowbreak{}nemotron-\allowbreak{}super-\allowbreak{}vl-\allowbreak{}exl3-\allowbreak{}2.0.env and 2 more} \\
\texttt{QSE\_\allowbreak{}NH\_\allowbreak{}GGUF} & \texttt{engine/\allowbreak{}models/\allowbreak{}nemotron\_\allowbreak{}h/\allowbreak{}load.py,\allowbreak{} ops/\allowbreak{}profiles/\allowbreak{}nemotron-\allowbreak{}bolt-\allowbreak{}iq2.env,\allowbreak{} server/\allowbreak{}diskcache.py} \\
\texttt{QSE\_\allowbreak{}NH\_\allowbreak{}HEAD} & \texttt{engine/\allowbreak{}models/\allowbreak{}nemotron\_\allowbreak{}h/\allowbreak{}community.py,\allowbreak{} ops/\allowbreak{}profiles/\allowbreak{}\_\allowbreak{}nemotron.env,\allowbreak{} ops/\allowbreak{}profiles/\allowbreak{}nemotron-\allowbreak{}bolt-\allowbreak{}exl3-\allowbreak{}2.0.env and 10 more} \\
\texttt{QSE\_\allowbreak{}NH\_\allowbreak{}JOINT\_\allowbreak{}ROWS8} & \texttt{engine/\allowbreak{}models/\allowbreak{}nemotron\_\allowbreak{}h/\allowbreak{}serving.py} \\
\texttt{QSE\_\allowbreak{}NH\_\allowbreak{}LILICORR} & \texttt{engine/\allowbreak{}models/\allowbreak{}nemotron\_\allowbreak{}h/\allowbreak{}serving.py,\allowbreak{} ops/\allowbreak{}profiles/\allowbreak{}nemotron-\allowbreak{}bolt-\allowbreak{}exl3-\allowbreak{}2.0.env,\allowbreak{} ops/\allowbreak{}profiles/\allowbreak{}nemotron-\allowbreak{}bolt-\allowbreak{}exl3-\allowbreak{}2.5.env and 8 more} \\
\texttt{QSE\_\allowbreak{}NH\_\allowbreak{}LILICORR\_\allowbreak{}PACK} & \texttt{engine/\allowbreak{}drafters/\allowbreak{}lilicorr.py} \\
\texttt{QSE\_\allowbreak{}NH\_\allowbreak{}LILICORR\_\allowbreak{}PACK\_\allowbreak{}GRAPHS} & \texttt{engine/\allowbreak{}drafters/\allowbreak{}lilicorr.py} \\
\texttt{QSE\_\allowbreak{}NH\_\allowbreak{}LOOKUP} & \texttt{engine/\allowbreak{}models/\allowbreak{}nemotron\_\allowbreak{}h/\allowbreak{}serving.py} \\
\texttt{QSE\_\allowbreak{}NH\_\allowbreak{}MODEL} & \texttt{engine/\allowbreak{}models/\allowbreak{}nemotron\_\allowbreak{}h/\allowbreak{}serving.py,\allowbreak{} ops/\allowbreak{}profiles/\allowbreak{}\_\allowbreak{}nemotron.env,\allowbreak{} ops/\allowbreak{}profiles/\allowbreak{}nemotron-\allowbreak{}bolt-\allowbreak{}exl3-\allowbreak{}2.0.env and 11 more} \\
\texttt{QSE\_\allowbreak{}NH\_\allowbreak{}VARIANT} & \texttt{engine/\allowbreak{}models/\allowbreak{}nemotron\_\allowbreak{}h/\allowbreak{}community.py,\allowbreak{} ops/\allowbreak{}profiles/\allowbreak{}\_\allowbreak{}nemotron.env,\allowbreak{} ops/\allowbreak{}profiles/\allowbreak{}nemotron-\allowbreak{}bolt-\allowbreak{}exl3-\allowbreak{}2.0.env and 11 more} \\
\texttt{QSE\_\allowbreak{}NO\_\allowbreak{}ROLLBACK} & \texttt{ops/\allowbreak{}use-\allowbreak{}profile.sh,\allowbreak{} server/\allowbreak{}setup/\allowbreak{}apply.py} \\
\texttt{QSE\_\allowbreak{}OPS\_\allowbreak{}DIR} & \texttt{ops/\allowbreak{}start.sh,\allowbreak{} server/\allowbreak{}setup/\allowbreak{}apply.py,\allowbreak{} server/\allowbreak{}setup/\allowbreak{}device.py} \\
\texttt{QSE\_\allowbreak{}ORN\_\allowbreak{}DFLASH2} & \texttt{engine/\allowbreak{}models/\allowbreak{}qwen3\_\allowbreak{}5\_\allowbreak{}moe/\allowbreak{}serving.py,\allowbreak{} ops/\allowbreak{}profiles/\allowbreak{}\_\allowbreak{}ornith.env,\allowbreak{} ops/\allowbreak{}profiles/\allowbreak{}ornith-\allowbreak{}1.5-\allowbreak{}122b-\allowbreak{}exl3-\allowbreak{}4.0.env} \\
\texttt{QSE\_\allowbreak{}ORN\_\allowbreak{}DRAFTER} & \texttt{ops/\allowbreak{}profiles/\allowbreak{}\_\allowbreak{}ornith.env} \\
\texttt{QSE\_\allowbreak{}ORN\_\allowbreak{}HEAD} & \texttt{ops/\allowbreak{}profiles/\allowbreak{}\_\allowbreak{}ornith.env,\allowbreak{} ops/\allowbreak{}profiles/\allowbreak{}ornith-\allowbreak{}1.5-\allowbreak{}122b-\allowbreak{}exl3-\allowbreak{}4.0.env} \\
\texttt{QSE\_\allowbreak{}ORN\_\allowbreak{}VARIANT} & \texttt{ops/\allowbreak{}profiles/\allowbreak{}\_\allowbreak{}ornith.env,\allowbreak{} ops/\allowbreak{}profiles/\allowbreak{}ornith-\allowbreak{}1.5-\allowbreak{}122b-\allowbreak{}exl3-\allowbreak{}4.0.env} \\
\texttt{QSE\_\allowbreak{}PACK\_\allowbreak{}VISION} & \texttt{engine/\allowbreak{}models/\allowbreak{}nemotron\_\allowbreak{}h/\allowbreak{}community.py,\allowbreak{} engine/\allowbreak{}models/\allowbreak{}qwen3\_\allowbreak{}5/\allowbreak{}community.py,\allowbreak{} ops/\allowbreak{}serve.env and 1 more} \\
\texttt{QSE\_\allowbreak{}PAGE\_\allowbreak{}CACHE\_\allowbreak{}DROP} & \texttt{engine/\allowbreak{}models/\allowbreak{}nemotron\_\allowbreak{}h/\allowbreak{}community.py,\allowbreak{} ops/\allowbreak{}profiles/\allowbreak{}\_\allowbreak{}flashnext.env,\allowbreak{} ops/\allowbreak{}profiles/\allowbreak{}\_\allowbreak{}flint.env and 5 more} \\
\texttt{QSE\_\allowbreak{}PAGE\_\allowbreak{}CACHE\_\allowbreak{}KEEP} & \texttt{ops/\allowbreak{}profiles/\allowbreak{}\_\allowbreak{}flashnext.env,\allowbreak{} ops/\allowbreak{}serve.env} \\
\texttt{QSE\_\allowbreak{}PARALLEL} & \texttt{ops/\allowbreak{}engines.sh,\allowbreak{} ops/\allowbreak{}profiles/\allowbreak{}\_\allowbreak{}flashnext.env,\allowbreak{} ops/\allowbreak{}profiles/\allowbreak{}\_\allowbreak{}flint.env and 6 more} \\
\texttt{QSE\_\allowbreak{}PARALLEL\_\allowbreak{}AUTO\_\allowbreak{}N} & \texttt{server/\allowbreak{}app.py,\allowbreak{} server/\allowbreak{}slots.py} \\
\texttt{QSE\_\allowbreak{}PARALLEL\_\allowbreak{}BATCH} & \texttt{server/\allowbreak{}app.py} \\
\texttt{QSE\_\allowbreak{}PARALLEL\_\allowbreak{}BUSY\_\allowbreak{}WAIT\_\allowbreak{}MS} & \texttt{server/\allowbreak{}slots.py} \\
\texttt{QSE\_\allowbreak{}PARALLEL\_\allowbreak{}CUT} & \texttt{server/\allowbreak{}app.py} \\
\texttt{QSE\_\allowbreak{}PARALLEL\_\allowbreak{}DRAFT} & \texttt{server/\allowbreak{}app.py} \\
\texttt{QSE\_\allowbreak{}PARALLEL\_\allowbreak{}DRAFT\_\allowbreak{}ROWS} & \texttt{server/\allowbreak{}app.py} \\
\texttt{QSE\_\allowbreak{}PARALLEL\_\allowbreak{}FLOOR\_\allowbreak{}GB} & \texttt{engine/\allowbreak{}kvpool.py,\allowbreak{} ops/\allowbreak{}profiles/\allowbreak{}\_\allowbreak{}flashnext.env,\allowbreak{} ops/\allowbreak{}serve.env and 1 more} \\
\texttt{QSE\_\allowbreak{}PARALLEL\_\allowbreak{}LAYER\_\allowbreak{}YIELD} & \texttt{server/\allowbreak{}app.py} \\
\texttt{QSE\_\allowbreak{}PARALLEL\_\allowbreak{}MAX\_\allowbreak{}ROWS} & \texttt{server/\allowbreak{}app.py} \\
\texttt{QSE\_\allowbreak{}PARALLEL\_\allowbreak{}ROW\_\allowbreak{}CAP} & \texttt{server/\allowbreak{}app.py} \\
\texttt{QSE\_\allowbreak{}PARALLEL\_\allowbreak{}SWITCH\_\allowbreak{}MS} & \texttt{server/\allowbreak{}app.py} \\
\texttt{QSE\_\allowbreak{}PARALLEL\_\allowbreak{}TIMING} & \texttt{server/\allowbreak{}app.py} \\
\texttt{QSE\_\allowbreak{}PARALLEL\_\allowbreak{}WINDOW\_\allowbreak{}ROWS} & \texttt{server/\allowbreak{}app.py} \\
\texttt{QSE\_\allowbreak{}PEAK\_\allowbreak{}STRICT} & \texttt{engine/\allowbreak{}cache.py,\allowbreak{} server/\allowbreak{}app.py} \\
\texttt{QSE\_\allowbreak{}PEER} & \texttt{ops/\allowbreak{}engines.sh,\allowbreak{} ops/\allowbreak{}start.sh,\allowbreak{} ops/\allowbreak{}stop.sh} \\
\texttt{QSE\_\allowbreak{}PEER\_\allowbreak{}DISK\_\allowbreak{}CACHE} & \texttt{ops/\allowbreak{}engines.sh} \\
\texttt{QSE\_\allowbreak{}PEER\_\allowbreak{}ENGINE\_\allowbreak{}FLAGS} & \texttt{ops/\allowbreak{}engines.sh} \\
\texttt{QSE\_\allowbreak{}PEER\_\allowbreak{}FLOOR\_\allowbreak{}GB} & \texttt{ops/\allowbreak{}engines.sh} \\
\texttt{QSE\_\allowbreak{}PEER\_\allowbreak{}HOST} & \texttt{ops/\allowbreak{}engines.sh} \\
\texttt{QSE\_\allowbreak{}PEER\_\allowbreak{}NEED\_\allowbreak{}GB} & \texttt{ops/\allowbreak{}engines.sh,\allowbreak{} ops/\allowbreak{}start.sh} \\
\texttt{QSE\_\allowbreak{}PEER\_\allowbreak{}PORT} & \texttt{ops/\allowbreak{}engines.sh,\allowbreak{} ops/\allowbreak{}start.sh} \\
\texttt{QSE\_\allowbreak{}PREFILL\_\allowbreak{}BURST} & \texttt{server/\allowbreak{}app.py,\allowbreak{} server/\allowbreak{}prefill\_\allowbreak{}burst.py} \\
\texttt{QSE\_\allowbreak{}PREFILL\_\allowbreak{}BURST\_\allowbreak{}ASKING} & \texttt{server/\allowbreak{}app.py} \\
\texttt{QSE\_\allowbreak{}PREFILL\_\allowbreak{}BURST\_\allowbreak{}WAIT\_\allowbreak{}S} & \texttt{server/\allowbreak{}app.py} \\
\texttt{QSE\_\allowbreak{}PRICE\_\allowbreak{}TABLE} & \texttt{server/\allowbreak{}app.py} \\
\texttt{QSE\_\allowbreak{}PROFILE\_\allowbreak{}NAME} & \texttt{ops/\allowbreak{}serve.env,\allowbreak{} ops/\allowbreak{}start.sh} \\
\texttt{QSE\_\allowbreak{}PROFILE\_\allowbreak{}OVERRIDE} & \texttt{ops/\allowbreak{}serve.env,\allowbreak{} ops/\allowbreak{}use-\allowbreak{}profile.sh,\allowbreak{} server/\allowbreak{}setup/\allowbreak{}apply.py and 1 more} \\
\texttt{QSE\_\allowbreak{}RECLAIM\_\allowbreak{}GRAPHS} & \texttt{server/\allowbreak{}app.py} \\
\texttt{QSE\_\allowbreak{}REQ\_\allowbreak{}PEAK} & \texttt{server/\allowbreak{}app.py} \\
\texttt{QSE\_\allowbreak{}ROLES\_\allowbreak{}DYNAMIC} & \texttt{ops/\allowbreak{}serve.env,\allowbreak{} server/\allowbreak{}rolemem.py} \\
\texttt{QSE\_\allowbreak{}ROLE\_\allowbreak{}MAX\_\allowbreak{}HOLD\_\allowbreak{}S} & \texttt{ops/\allowbreak{}serve.env,\allowbreak{} server/\allowbreak{}rolemem.py} \\
\texttt{QSE\_\allowbreak{}ROLE\_\allowbreak{}RELOAD\_\allowbreak{}MARGIN\_\allowbreak{}GB} & \texttt{ops/\allowbreak{}serve.env,\allowbreak{} server/\allowbreak{}rolemem.py} \\
\texttt{QSE\_\allowbreak{}ROLE\_\allowbreak{}RELOAD\_\allowbreak{}S} & \texttt{ops/\allowbreak{}serve.env,\allowbreak{} server/\allowbreak{}rolemem.py} \\
\texttt{QSE\_\allowbreak{}ROLE\_\allowbreak{}THRASH\_\allowbreak{}S} & \texttt{ops/\allowbreak{}serve.env,\allowbreak{} server/\allowbreak{}rolemem.py} \\
\texttt{QSE\_\allowbreak{}ROLE\_\allowbreak{}TICK\_\allowbreak{}FAKE} & \texttt{server/\allowbreak{}app.py} \\
\texttt{QSE\_\allowbreak{}ROLE\_\allowbreak{}TICK\_\allowbreak{}S} & \texttt{server/\allowbreak{}rolemem.py} \\
\texttt{QSE\_\allowbreak{}ROLE\_\allowbreak{}UNLOAD\_\allowbreak{}LAYA} & \texttt{ops/\allowbreak{}serve.env,\allowbreak{} server/\allowbreak{}app.py,\allowbreak{} server/\allowbreak{}rolemem.py} \\
\texttt{QSE\_\allowbreak{}ROUND\_\allowbreak{}RUNNER} & \texttt{server/\allowbreak{}round\_\allowbreak{}runner.py} \\
\texttt{QSE\_\allowbreak{}ROUND\_\allowbreak{}WALK} & \texttt{server/\allowbreak{}app.py} \\
\texttt{QSE\_\allowbreak{}RUNNER\_\allowbreak{}BUSY\_\allowbreak{}WORK} & \texttt{ops/\allowbreak{}profiles/\allowbreak{}\_\allowbreak{}flashnext.env,\allowbreak{} server/\allowbreak{}api/\allowbreak{}live.py,\allowbreak{} server/\allowbreak{}health\_\allowbreak{}cache.py and 1 more} \\
\texttt{QSE\_\allowbreak{}RUNNER\_\allowbreak{}COMMIT\_\allowbreak{}MULTI} & \texttt{server/\allowbreak{}round\_\allowbreak{}runner.py} \\
\texttt{QSE\_\allowbreak{}RUNNER\_\allowbreak{}COUNTERS} & \texttt{engine/\allowbreak{}models/\allowbreak{}flashnext/\allowbreak{}serving.py,\allowbreak{} server/\allowbreak{}round\_\allowbreak{}runner.py} \\
\texttt{QSE\_\allowbreak{}RUNNER\_\allowbreak{}GC} & \texttt{engine/\allowbreak{}cache.py,\allowbreak{} ops/\allowbreak{}profiles/\allowbreak{}\_\allowbreak{}flashnext.env,\allowbreak{} server/\allowbreak{}app.py and 1 more} \\
\texttt{QSE\_\allowbreak{}RUNNER\_\allowbreak{}GC\_\allowbreak{}FULL\_\allowbreak{}ROUNDS} & \texttt{ops/\allowbreak{}profiles/\allowbreak{}\_\allowbreak{}flashnext.env,\allowbreak{} server/\allowbreak{}round\_\allowbreak{}runner.py} \\
\texttt{QSE\_\allowbreak{}RUNNER\_\allowbreak{}HOST\_\allowbreak{}BLOCKS} & \texttt{engine/\allowbreak{}models/\allowbreak{}flashnext/\allowbreak{}serving.py,\allowbreak{} server/\allowbreak{}round\_\allowbreak{}runner.py} \\
\texttt{QSE\_\allowbreak{}RUNNER\_\allowbreak{}JOINT1} & \texttt{engine/\allowbreak{}model.py,\allowbreak{} engine/\allowbreak{}models/\allowbreak{}qwen3\_\allowbreak{}5/\allowbreak{}joint\_\allowbreak{}graphs.py,\allowbreak{} server/\allowbreak{}round\_\allowbreak{}runner.py} \\
\texttt{QSE\_\allowbreak{}RUNNER\_\allowbreak{}JOINT\_\allowbreak{}PRICE} & \texttt{server/\allowbreak{}round\_\allowbreak{}runner.py} \\
\texttt{QSE\_\allowbreak{}RUNNER\_\allowbreak{}LANE\_\allowbreak{}MARGIN} & \texttt{server/\allowbreak{}round\_\allowbreak{}runner.py} \\
\texttt{QSE\_\allowbreak{}RUNNER\_\allowbreak{}LANE\_\allowbreak{}PLAIN} & \texttt{engine/\allowbreak{}models/\allowbreak{}nemotron\_\allowbreak{}h/\allowbreak{}serving.py,\allowbreak{} server/\allowbreak{}round\_\allowbreak{}runner.py} \\
\texttt{QSE\_\allowbreak{}RUNNER\_\allowbreak{}LANE\_\allowbreak{}PROBE} & \texttt{server/\allowbreak{}round\_\allowbreak{}runner.py} \\
\texttt{QSE\_\allowbreak{}RUNNER\_\allowbreak{}MTP\_\allowbreak{}PACK} & \texttt{server/\allowbreak{}round\_\allowbreak{}runner.py} \\
\texttt{QSE\_\allowbreak{}RUNNER\_\allowbreak{}PIPE} & \texttt{engine/\allowbreak{}models/\allowbreak{}flashnext/\allowbreak{}engine.py,\allowbreak{} server/\allowbreak{}round\_\allowbreak{}runner.py} \\
\texttt{QSE\_\allowbreak{}RUNNER\_\allowbreak{}PRICED\_\allowbreak{}CUT} & \texttt{server/\allowbreak{}round\_\allowbreak{}runner.py} \\
\texttt{QSE\_\allowbreak{}RUNNER\_\allowbreak{}PRICED\_\allowbreak{}EXPLORE} & \texttt{server/\allowbreak{}round\_\allowbreak{}runner.py} \\
\texttt{QSE\_\allowbreak{}RUNNER\_\allowbreak{}ROWCAP} & \texttt{server/\allowbreak{}app.py} \\
\texttt{QSE\_\allowbreak{}RUNNER\_\allowbreak{}ROWCOPY} & \texttt{engine/\allowbreak{}rowcopy.py,\allowbreak{} server/\allowbreak{}round\_\allowbreak{}runner.py} \\
\texttt{QSE\_\allowbreak{}RUNNER\_\allowbreak{}SUPERVISED} & \texttt{server/\allowbreak{}round\_\allowbreak{}runner.py} \\
\texttt{QSE\_\allowbreak{}RUNNER\_\allowbreak{}VERDICT\_\allowbreak{}WAIT\_\allowbreak{}MS} & \texttt{server/\allowbreak{}round\_\allowbreak{}runner.py} \\
\texttt{QSE\_\allowbreak{}SECRETS} & \texttt{install.sh,\allowbreak{} ops/\allowbreak{}make-\allowbreak{}secrets.sh,\allowbreak{} ops/\allowbreak{}start.sh and 3 more} \\
\texttt{QSE\_\allowbreak{}SESSION\_\allowbreak{}S} & \texttt{server/\allowbreak{}api/\allowbreak{}auth.py} \\
\texttt{QSE\_\allowbreak{}SETUP\_\allowbreak{}ENV} & \texttt{ops/\allowbreak{}serve.env,\allowbreak{} server/\allowbreak{}setup/\allowbreak{}apply.py} \\
\texttt{QSE\_\allowbreak{}SETUP\_\allowbreak{}ENV\_\allowbreak{}FILE} & \texttt{ops/\allowbreak{}serve.env} \\
\texttt{QSE\_\allowbreak{}SETUP\_\allowbreak{}MODE} & \texttt{ops/\allowbreak{}start.sh} \\
\texttt{QSE\_\allowbreak{}SIDE\_\allowbreak{}GRAPHS} & \texttt{ops/\allowbreak{}serve.env,\allowbreak{} server/\allowbreak{}app.py} \\
\texttt{QSE\_\allowbreak{}SIDE\_\allowbreak{}MAX\_\allowbreak{}LEN} & \texttt{ops/\allowbreak{}profiles/\allowbreak{}\_\allowbreak{}flashnext.env,\allowbreak{} ops/\allowbreak{}serve.env,\allowbreak{} server/\allowbreak{}app.py and 2 more} \\
\texttt{QSE\_\allowbreak{}SIDE\_\allowbreak{}RESIDENT\_\allowbreak{}GB} & \texttt{server/\allowbreak{}app.py} \\
\texttt{QSE\_\allowbreak{}SIDE\_\allowbreak{}SLOT\_\allowbreak{}CONST\_\allowbreak{}GB} & \texttt{ops/\allowbreak{}serve.env,\allowbreak{} server/\allowbreak{}app.py} \\
\texttt{QSE\_\allowbreak{}STACK\_\allowbreak{}SAMPLE} & \texttt{server/\allowbreak{}app.py,\allowbreak{} server/\allowbreak{}stack\_\allowbreak{}sampler.py} \\
\texttt{QSE\_\allowbreak{}STATE\_\allowbreak{}DIR} & \texttt{ops/\allowbreak{}make-\allowbreak{}secrets.sh,\allowbreak{} ops/\allowbreak{}start.sh,\allowbreak{} server/\allowbreak{}setup/\allowbreak{}apply.py and 2 more} \\
\texttt{QSE\_\allowbreak{}STATUS\_\allowbreak{}FILE} & \texttt{ops/\allowbreak{}engine\_\allowbreak{}status.py} \\
\texttt{QSE\_\allowbreak{}STATUS\_\allowbreak{}HOST} & \texttt{ops/\allowbreak{}engine\_\allowbreak{}status.py,\allowbreak{} ops/\allowbreak{}hold.sh,\allowbreak{} ops/\allowbreak{}serve.env and 1 more} \\
\texttt{QSE\_\allowbreak{}STATUS\_\allowbreak{}KEEP} & \texttt{ops/\allowbreak{}hold.sh,\allowbreak{} ops/\allowbreak{}stop.sh} \\
\texttt{QSE\_\allowbreak{}STATUS\_\allowbreak{}PORT} & \texttt{install.sh,\allowbreak{} ops/\allowbreak{}engine\_\allowbreak{}status.py,\allowbreak{} ops/\allowbreak{}engines.sh and 4 more} \\
\texttt{QSE\_\allowbreak{}STATUS\_\allowbreak{}TREE} & \texttt{install.sh,\allowbreak{} ops/\allowbreak{}engine\_\allowbreak{}status.py,\allowbreak{} ops/\allowbreak{}hold.sh} \\
\texttt{QSE\_\allowbreak{}SUFFIX\_\allowbreak{}RAW} & \texttt{server/\allowbreak{}app.py} \\
\texttt{QSE\_\allowbreak{}TEST} & \texttt{server/\allowbreak{}api/\allowbreak{}ledger.py,\allowbreak{} server/\allowbreak{}app.py} \\
\texttt{QSE\_\allowbreak{}USAGE\_\allowbreak{}CONTEXT} & \texttt{server/\allowbreak{}app.py} \\
\texttt{QSE\_\allowbreak{}USAGE\_\allowbreak{}DEFAULT} & \texttt{server/\allowbreak{}app.py} \\
\texttt{QSE\_\allowbreak{}USAGE\_\allowbreak{}LEDGER} & \texttt{install.sh,\allowbreak{} ops/\allowbreak{}engines.sh,\allowbreak{} ops/\allowbreak{}serve.env and 4 more} \\
\texttt{QSE\_\allowbreak{}VISION\_\allowbreak{}ASKED} & \texttt{ops/\allowbreak{}serve.env,\allowbreak{} ops/\allowbreak{}start.sh} \\
\texttt{QSE\_\allowbreak{}WEIGHTS} & \texttt{engine/\allowbreak{}models/\allowbreak{}flashnext/\allowbreak{}corrector.py,\allowbreak{} engine/\allowbreak{}models/\allowbreak{}flashnext/\allowbreak{}coupled.py,\allowbreak{} engine/\allowbreak{}models/\allowbreak{}flashnext/\allowbreak{}serving.py and 20 more} \\
\texttt{QSE\_\allowbreak{}WEIGHTS\_\allowbreak{}ALSO} & \texttt{ops/\allowbreak{}profiles/\allowbreak{}\_\allowbreak{}flashnext.env,\allowbreak{} ops/\allowbreak{}use-\allowbreak{}profile.sh} \\
\texttt{QSE\_\allowbreak{}WEIGHTS\_\allowbreak{}DIR} & \texttt{engine/\allowbreak{}models/\allowbreak{}nemotron\_\allowbreak{}h/\allowbreak{}community.py,\allowbreak{} engine/\allowbreak{}models/\allowbreak{}qwen3\_\allowbreak{}5/\allowbreak{}community.py,\allowbreak{} ops/\allowbreak{}profiles/\allowbreak{}\_\allowbreak{}exl3.env and 12 more} \\
\texttt{QSE\_\allowbreak{}WEIGHTS\_\allowbreak{}FILES} & \texttt{engine/\allowbreak{}models/\allowbreak{}nemotron\_\allowbreak{}h/\allowbreak{}community.py,\allowbreak{} engine/\allowbreak{}models/\allowbreak{}qwen3\_\allowbreak{}5/\allowbreak{}community.py,\allowbreak{} ops/\allowbreak{}profiles/\allowbreak{}\_\allowbreak{}exl3.env and 9 more} \\
\texttt{QSE\_\allowbreak{}WEIGHTS\_\allowbreak{}GB} & \texttt{engine/\allowbreak{}models/\allowbreak{}nemotron\_\allowbreak{}h/\allowbreak{}community.py,\allowbreak{} engine/\allowbreak{}models/\allowbreak{}qwen3\_\allowbreak{}5/\allowbreak{}community.py,\allowbreak{} ops/\allowbreak{}profiles/\allowbreak{}exl3-\allowbreak{}2.5.env and 35 more} \\
\texttt{QSE\_\allowbreak{}WEIGHTS\_\allowbreak{}HF} & \texttt{engine/\allowbreak{}models/\allowbreak{}nemotron\_\allowbreak{}h/\allowbreak{}community.py,\allowbreak{} engine/\allowbreak{}models/\allowbreak{}qwen3\_\allowbreak{}5/\allowbreak{}community.py,\allowbreak{} ops/\allowbreak{}profiles/\allowbreak{}\_\allowbreak{}exl3.env and 27 more} \\
\texttt{QSE\_\allowbreak{}WEIGHTS\_\allowbreak{}REVISION} & \texttt{engine/\allowbreak{}models/\allowbreak{}nemotron\_\allowbreak{}h/\allowbreak{}community.py,\allowbreak{} engine/\allowbreak{}models/\allowbreak{}qwen3\_\allowbreak{}5/\allowbreak{}community.py,\allowbreak{} ops/\allowbreak{}profiles/\allowbreak{}\_\allowbreak{}flashnext.env and 30 more} \\
\texttt{QSE\_\allowbreak{}WREN\_\allowbreak{}MODEL} & \texttt{server/\allowbreak{}setup/\allowbreak{}wren.py} \\
\end{reftabletwo}


The most used: \code{QSE_CONDUCTOR_CONFIG}\index{qseconductorconfig@\texttt{QSE\_CONDUCTOR\_CONFIG}}, \code{QSE_USAGE_LEDGER}\index{qseusageledger@\texttt{QSE\_USAGE\_LEDGER}}, \code{QSE_DISK_CACHE}\index{qsediskcache@\texttt{QSE\_DISK\_CACHE}} and \code{QSE_BACKENDS}\index{qsebackends@\texttt{QSE\_BACKENDS}} name the conductor's directory, the ledger, the disk tier's directory and the proxy's config database; \code{ops/start.sh} passes them to the service only, so a benchmark server started with \code{serve.env} exported never writes them. \code{QSE_DASHBOARD_LOGIN}\index{qsedashboardlogin@\texttt{QSE\_DASHBOARD\_LOGIN}} turns the dashboard login on. \code{QSE_ADMIN_TOKEN} and \code{QSE_METRICS_TOKEN} live in \code{secrets.env}, which \code{ops/make-secrets.sh} writes.

The wait for memory (Section~\ref{sec:memory-wait}) reads these, with their code defaults: \code{QSE_MEM_WAIT_S} (10,800) is the longest wait in seconds of a streamed request, or of one that sends \code{wait_for_memory: true}; \code{QSE_PARALLEL_FLOOR_GB} (15) is the floor the admission keeps MemAvailable at, the same as \code{--parallel-floor-gb}; \code{QSE_KV_HARD_FLOOR_GB} (10) is the floor the KV pool's commits keep, the same as \code{--kv-hard-floor-gb}; \code{QSE_GUARD_MARGIN_GB} (5) is the margin each floor holds under the largest step the engine takes between two checks; \code{QSE_KV_ADMIT_ROWS} (8,192) is the part of \code{max_tokens} the admission reserves KV pages for, 0 for all of it; \code{QSE_KV_RECOMPUTE} (\code{on}) lets a lane the pool refuses give way and go on later; \code{QSE_PEAK_STRICT} (0) set to 1 makes a lone request that the memory gate refuses wait too.

\section{TANDEM variables}\index{TANDEM variables}

% Generated by manual/tools/gen_reference.py from the code. Do not edit by hand.

\begin{reftable}{TANDEM\_* variables: the proxy, the collector, the plugins and the installer}{tab:ref-tandem}{\hd{Variable} & \hd{Example} & \hd{What it sets}}
\texttt{TANDEM\_\allowbreak{}BACKENDS} & \texttt{deepseek,\allowbreak{}vllm1} & the backends the environment declares \\
\texttt{TANDEM\_\allowbreak{}BACKEND\_\allowbreak{}\textless{}ID\textgreater{}\_\allowbreak{}KIND} & \texttt{deepseek} & generic \textbar{} tandemllm \textbar{} vllm \textbar{} sglang \textbar{} openai \textbar{} deepseek \\
\texttt{TANDEM\_\allowbreak{}BACKEND\_\allowbreak{}\textless{}ID\textgreater{}\_\allowbreak{}URL} & \texttt{https:\allowbreak{}/\allowbreak{}/\allowbreak{}api.deepseek.com/\allowbreak{}v1} &  \\
\texttt{TANDEM\_\allowbreak{}BACKEND\_\allowbreak{}\textless{}ID\textgreater{}\_\allowbreak{}KEY} & \texttt{(a key)} & held in memory, never stored \\
\texttt{TANDEM\_\allowbreak{}BACKEND\_\allowbreak{}\textless{}ID\textgreater{}\_\allowbreak{}LABEL} & \texttt{DeepSeek} &  \\
\texttt{TANDEM\_\allowbreak{}BACKEND\_\allowbreak{}\textless{}ID\textgreater{}\_\allowbreak{}MODELS} & \texttt{deepseek-\allowbreak{}flash,\allowbreak{}deepseek-\allowbreak{}v4-\allowbreak{}pro} &  \\
\texttt{TANDEM\_\allowbreak{}BACKEND\_\allowbreak{}\textless{}ID\textgreater{}\_\allowbreak{}MODELS\_\allowbreak{}MODE} & \texttt{discover} & discover (GET /models, plus MODELS) \textbar{} list \\
\texttt{TANDEM\_\allowbreak{}BACKEND\_\allowbreak{}\textless{}ID\textgreater{}\_\allowbreak{}FALLBACK} & \texttt{openrouter} & backends tried after this one, in order \\
\texttt{TANDEM\_\allowbreak{}BACKEND\_\allowbreak{}\textless{}ID\textgreater{}\_\allowbreak{}MAX\_\allowbreak{}INFLIGHT} & \texttt{none} & requests at once, past it a fallback or a 503; none (the default): no limit from the proxy, the backend takes what it accepts \\
\texttt{TANDEM\_\allowbreak{}BACKEND\_\allowbreak{}\textless{}ID\textgreater{}\_\allowbreak{}ENABLED} & \texttt{1} &  \\
\texttt{TANDEM\_\allowbreak{}BACKEND\_\allowbreak{}\textless{}ID\textgreater{}\_\allowbreak{}INJECT\_\allowbreak{}USAGE} & \texttt{1} & ask a stream for its usage frame \\
\texttt{TANDEM\_\allowbreak{}BACKEND\_\allowbreak{}\textless{}ID\textgreater{}\_\allowbreak{}HARNESS} & \texttt{off} & off \textbar{} shadow \textbar{} on \\
\texttt{TANDEM\_\allowbreak{}BACKEND\_\allowbreak{}\textless{}ID\textgreater{}\_\allowbreak{}HARNESS\_\allowbreak{}PROFILE} & \texttt{coding} & the conductor profile of its requests \\
\texttt{TANDEM\_\allowbreak{}BACKEND\_\allowbreak{}\textless{}ID\textgreater{}\_\allowbreak{}WRITER} & \texttt{auto} & the background memory writer on this backend (proxy/memwriter.py): auto (on, off for a paid cloud model) \textbar{} off \textbar{} shadow \textbar{} on \\
\texttt{TANDEM\_\allowbreak{}BACKEND\_\allowbreak{}\textless{}ID\textgreater{}\_\allowbreak{}WRITER\_\allowbreak{}USD\_\allowbreak{}DAY} & \texttt{0.25} & its spend cap a day on a paid model \\
\texttt{TANDEM\_\allowbreak{}BACKEND\_\allowbreak{}\textless{}ID\textgreater{}\_\allowbreak{}WRITER\_\allowbreak{}RUNS\_\allowbreak{}DAY} & \texttt{48} & its runs a day \\
\texttt{TANDEM\_\allowbreak{}BACKEND\_\allowbreak{}\textless{}ID\textgreater{}\_\allowbreak{}WRITER\_\allowbreak{}PROMPT\_\allowbreak{}TOKENS} & \texttt{32000} & its prompt at most (older messages dropped) \\
\texttt{TANDEM\_\allowbreak{}BACKEND\_\allowbreak{}\textless{}ID\textgreater{}\_\allowbreak{}TIMEOUT\_\allowbreak{}CONNECT} & \texttt{10} & seconds; also \_TIMEOUT\_FIRST\_BYTE, \_TIMEOUT\_IDLE \\
\texttt{TANDEM\_\allowbreak{}BACKEND\_\allowbreak{}\textless{}ID\textgreater{}\_\allowbreak{}HOST} & \texttt{auto} & where the System view reads the backend's machine: auto (this machine when the URL is), same, collector, exporter, off \\
\texttt{TANDEM\_\allowbreak{}BACKEND\_\allowbreak{}\textless{}ID\textgreater{}\_\allowbreak{}HOST\_\allowbreak{}URL} & \texttt{(a URL)} & the remote collector (python -m proxy.hostcollect), or a node\_exporter page \\
\texttt{TANDEM\_\allowbreak{}BACKEND\_\allowbreak{}\textless{}ID\textgreater{}\_\allowbreak{}HOST\_\allowbreak{}DCGM\_\allowbreak{}URL} & \texttt{(a URL)} & a DCGM exporter page (mode exporter) \\
\texttt{TANDEM\_\allowbreak{}COLLECTOR\_\allowbreak{}TOKEN} & \texttt{...} & the remote collectors' bearer token, in memory only (proxy/hostcollect.py) \\
\texttt{TANDEM\_\allowbreak{}ROUTE\_\allowbreak{}\textless{}n\textgreater{}} & \texttt{gpt-\allowbreak{}*=\allowbreak{}openai} & routing rules, ordered by n \\
\texttt{TANDEM\_\allowbreak{}PRICE\_\allowbreak{}\textless{}n\textgreater{}} & \texttt{deepseek/\allowbreak{}deepseek-\allowbreak{}flash=\allowbreak{}0.30,\allowbreak{}0.006,\allowbreak{}1.2} & USD per million: input, cached input, output, and optionally cache write \\
\texttt{TANDEM\_\allowbreak{}PRICE\_\allowbreak{}\textless{}n\textgreater{}\_\allowbreak{}OFFPEAK} & \texttt{0.15,\allowbreak{}0.003,\allowbreak{}0.6} & the same model outside the provider's peak hours (DeepSeek) \\
\texttt{TANDEM\_\allowbreak{}API\_\allowbreak{}KEYS} & \texttt{(a key)} & the proxy's own keys, in memory only \\
\texttt{TANDEM\_\allowbreak{}API\_\allowbreak{}USERS\_\allowbreak{}FROM} & \texttt{\textasciitilde{}/\allowbreak{}.engine/\allowbreak{}config.sqlite3} & : the keys another config database's users hold count here too, read from that file at each check (read only: the label, the key's sha256, enabled); a controller in front of an engine takes the ... \\
\texttt{TANDEM\_\allowbreak{}DEFAULT\_\allowbreak{}BACKEND} & \texttt{engine} & : a model id that no deployment, rule or backend serves goes to this backend as asked (a TandemLLM engine answers any id with its own model); unset: OpenAI's 404 \\
\texttt{TANDEM\_\allowbreak{}SECRETS\_\allowbreak{}KEY} & \texttt{\textless{}base64} & of 32 bytes\textgreater{} the master key (proxy/secrets.py) \\
\texttt{TANDEM\_\allowbreak{}ADMIN\_\allowbreak{}TOKEN} & & read in \texttt{app/\allowbreak{}packaging/\allowbreak{}app.env.example,\allowbreak{} app/\allowbreak{}proxy/\allowbreak{}\_\allowbreak{}\_\allowbreak{}main\_\allowbreak{}\_\allowbreak{}.py} \\
\texttt{TANDEM\_\allowbreak{}BACKEND\_\allowbreak{}DEEPSEEK\_\allowbreak{}KEY} & & read in \texttt{app/\allowbreak{}packaging/\allowbreak{}app.env.example} \\
\texttt{TANDEM\_\allowbreak{}BACKEND\_\allowbreak{}DEEPSEEK\_\allowbreak{}KIND} & & read in \texttt{app/\allowbreak{}packaging/\allowbreak{}app.env.example} \\
\texttt{TANDEM\_\allowbreak{}BACKEND\_\allowbreak{}DEEPSEEK\_\allowbreak{}LABEL} & & read in \texttt{app/\allowbreak{}packaging/\allowbreak{}app.env.example} \\
\texttt{TANDEM\_\allowbreak{}BACKEND\_\allowbreak{}VLLM1\_\allowbreak{}FALLBACK} & & read in \texttt{app/\allowbreak{}packaging/\allowbreak{}app.env.example} \\
\texttt{TANDEM\_\allowbreak{}BACKEND\_\allowbreak{}VLLM1\_\allowbreak{}HOST} & & read in \texttt{app/\allowbreak{}packaging/\allowbreak{}app.env.example} \\
\texttt{TANDEM\_\allowbreak{}BACKEND\_\allowbreak{}VLLM1\_\allowbreak{}HOST\_\allowbreak{}DCGM\_\allowbreak{}URL} & & read in \texttt{app/\allowbreak{}packaging/\allowbreak{}app.env.example} \\
\texttt{TANDEM\_\allowbreak{}BACKEND\_\allowbreak{}VLLM1\_\allowbreak{}HOST\_\allowbreak{}URL} & & read in \texttt{app/\allowbreak{}packaging/\allowbreak{}app.env.example} \\
\texttt{TANDEM\_\allowbreak{}BACKEND\_\allowbreak{}VLLM1\_\allowbreak{}KIND} & & read in \texttt{app/\allowbreak{}packaging/\allowbreak{}app.env.example} \\
\texttt{TANDEM\_\allowbreak{}BACKEND\_\allowbreak{}VLLM1\_\allowbreak{}MAX\_\allowbreak{}INFLIGHT} & & read in \texttt{app/\allowbreak{}packaging/\allowbreak{}app.env.example} \\
\texttt{TANDEM\_\allowbreak{}BACKEND\_\allowbreak{}VLLM1\_\allowbreak{}TIMEOUT\_\allowbreak{}FIRST\_\allowbreak{}BYTE} & & read in \texttt{app/\allowbreak{}packaging/\allowbreak{}app.env.example} \\
\texttt{TANDEM\_\allowbreak{}BACKEND\_\allowbreak{}VLLM1\_\allowbreak{}URL} & & read in \texttt{app/\allowbreak{}packaging/\allowbreak{}app.env.example} \\
\texttt{TANDEM\_\allowbreak{}COLLECT} & & read in \texttt{app/\allowbreak{}proxy/\allowbreak{}launch.py} \\
\texttt{TANDEM\_\allowbreak{}COLLECT\_\allowbreak{}BASE} & & read in \texttt{app/\allowbreak{}proxy/\allowbreak{}launch.py} \\
\texttt{TANDEM\_\allowbreak{}COLLECT\_\allowbreak{}TOKEN} & & read in \texttt{app/\allowbreak{}plugins/\allowbreak{}sglang/\allowbreak{}tandem\_\allowbreak{}sglang\_\allowbreak{}plugin/\allowbreak{}\_\allowbreak{}push.py,\allowbreak{} app/\allowbreak{}plugins/\allowbreak{}vllm/\allowbreak{}tandem\_\allowbreak{}vllm\_\allowbreak{}plugin/\allowbreak{}\_\allowbreak{}\_\allowbreak{}init\_\allowbreak{}\_\allowbreak{}.py} \\
\texttt{TANDEM\_\allowbreak{}COLLECT\_\allowbreak{}URL} & & read in \texttt{app/\allowbreak{}plugins/\allowbreak{}sglang/\allowbreak{}tandem\_\allowbreak{}sglang\_\allowbreak{}plugin/\allowbreak{}\_\allowbreak{}push.py,\allowbreak{} app/\allowbreak{}plugins/\allowbreak{}vllm/\allowbreak{}tandem\_\allowbreak{}vllm\_\allowbreak{}plugin/\allowbreak{}\_\allowbreak{}\_\allowbreak{}init\_\allowbreak{}\_\allowbreak{}.py} \\
\texttt{TANDEM\_\allowbreak{}CONDUCTOR} & & read in \texttt{app/\allowbreak{}proxy/\allowbreak{}\_\allowbreak{}\_\allowbreak{}main\_\allowbreak{}\_\allowbreak{}.py} \\
\texttt{TANDEM\_\allowbreak{}CONFIG\_\allowbreak{}DB} & & read in \texttt{app/\allowbreak{}packaging/\allowbreak{}app.env.example,\allowbreak{} app/\allowbreak{}proxy/\allowbreak{}\_\allowbreak{}\_\allowbreak{}main\_\allowbreak{}\_\allowbreak{}.py} \\
\texttt{TANDEM\_\allowbreak{}DASHBOARD\_\allowbreak{}LOGIN} & & read in \texttt{app/\allowbreak{}packaging/\allowbreak{}app.env.example,\allowbreak{} app/\allowbreak{}proxy/\allowbreak{}\_\allowbreak{}\_\allowbreak{}main\_\allowbreak{}\_\allowbreak{}.py} \\
\texttt{TANDEM\_\allowbreak{}DIR} & & read in \texttt{app/\allowbreak{}proxy/\allowbreak{}nodeagent.py,\allowbreak{} install.sh} \\
\texttt{TANDEM\_\allowbreak{}DMI\_\allowbreak{}DIR} & & read in \texttt{app/\allowbreak{}proxy/\allowbreak{}fleettrust.py} \\
\texttt{TANDEM\_\allowbreak{}ENGINE\_\allowbreak{}INSTALL\_\allowbreak{}SH} & & read in \texttt{app/\allowbreak{}packaging/\allowbreak{}app.env.example,\allowbreak{} app/\allowbreak{}proxy/\allowbreak{}fleetupdate.py} \\
\texttt{TANDEM\_\allowbreak{}ENGINE\_\allowbreak{}REF} & & read in \texttt{app/\allowbreak{}packaging/\allowbreak{}app.env.example,\allowbreak{} app/\allowbreak{}proxy/\allowbreak{}fleetupdate.py} \\
\texttt{TANDEM\_\allowbreak{}FLEET\_\allowbreak{}RETENTION\_\allowbreak{}DAYS} & & read in \texttt{app/\allowbreak{}proxy/\allowbreak{}\_\allowbreak{}\_\allowbreak{}main\_\allowbreak{}\_\allowbreak{}.py} \\
\texttt{TANDEM\_\allowbreak{}FLEET\_\allowbreak{}STORE} & & read in \texttt{app/\allowbreak{}proxy/\allowbreak{}\_\allowbreak{}\_\allowbreak{}main\_\allowbreak{}\_\allowbreak{}.py} \\
\texttt{TANDEM\_\allowbreak{}HEALTH\_\allowbreak{}TIMEOUT} & & read in \texttt{install.sh} \\
\texttt{TANDEM\_\allowbreak{}HOST} & & read in \texttt{app/\allowbreak{}packaging/\allowbreak{}app.env.example,\allowbreak{} app/\allowbreak{}proxy/\allowbreak{}\_\allowbreak{}\_\allowbreak{}main\_\allowbreak{}\_\allowbreak{}.py} \\
\texttt{TANDEM\_\allowbreak{}IDLE\_\allowbreak{}TIMEOUT\_\allowbreak{}S} & & read in \texttt{app/\allowbreak{}proxy/\allowbreak{}\_\allowbreak{}\_\allowbreak{}main\_\allowbreak{}\_\allowbreak{}.py,\allowbreak{} app/\allowbreak{}proxy/\allowbreak{}app.py} \\
\texttt{TANDEM\_\allowbreak{}LOCAL\_\allowbreak{}NODE} & & read in \texttt{app/\allowbreak{}proxy/\allowbreak{}app.py,\allowbreak{} app/\allowbreak{}proxy/\allowbreak{}fleet.py} \\
\texttt{TANDEM\_\allowbreak{}MACHINE\_\allowbreak{}ID\_\allowbreak{}FILE} & & read in \texttt{app/\allowbreak{}proxy/\allowbreak{}fleettrust.py} \\
\texttt{TANDEM\_\allowbreak{}MIN\_\allowbreak{}FREE\_\allowbreak{}GIB} & & read in \texttt{install.sh} \\
\texttt{TANDEM\_\allowbreak{}MODEL\_\allowbreak{}NAME} & & read in \texttt{install.sh} \\
\texttt{TANDEM\_\allowbreak{}NODE\_\allowbreak{}AGENT} & & read in \texttt{app/\allowbreak{}proxy/\allowbreak{}app.py,\allowbreak{} app/\allowbreak{}proxy/\allowbreak{}nodeagent.py} \\
\texttt{TANDEM\_\allowbreak{}NODE\_\allowbreak{}DIR} & & read in \texttt{app/\allowbreak{}proxy/\allowbreak{}nodeagent.py} \\
\texttt{TANDEM\_\allowbreak{}NODE\_\allowbreak{}HOST} & & read in \texttt{app/\allowbreak{}proxy/\allowbreak{}nodeagent.py} \\
\texttt{TANDEM\_\allowbreak{}NODE\_\allowbreak{}MDNS} & & read in \texttt{app/\allowbreak{}proxy/\allowbreak{}nodeagent.py} \\
\texttt{TANDEM\_\allowbreak{}NODE\_\allowbreak{}PORT} & & read in \texttt{app/\allowbreak{}proxy/\allowbreak{}nodeagent.py} \\
\texttt{TANDEM\_\allowbreak{}NODE\_\allowbreak{}SECRET} & & read in \texttt{app/\allowbreak{}proxy/\allowbreak{}app.py,\allowbreak{} app/\allowbreak{}proxy/\allowbreak{}nodeagent.py} \\
\texttt{TANDEM\_\allowbreak{}NODE\_\allowbreak{}VIRTUAL} & & read in \texttt{app/\allowbreak{}proxy/\allowbreak{}nodeagent.py} \\
\texttt{TANDEM\_\allowbreak{}NODE\_\allowbreak{}WEIGHTS} & & read in \texttt{app/\allowbreak{}proxy/\allowbreak{}nodeagent.py} \\
\texttt{TANDEM\_\allowbreak{}NODE\_\allowbreak{}WHEEL} & & read in \texttt{app/\allowbreak{}proxy/\allowbreak{}fleetupdate.py,\allowbreak{} app/\allowbreak{}proxy/\allowbreak{}nodeinstall.py} \\
\texttt{TANDEM\_\allowbreak{}NOTIFY\_\allowbreak{}URL} & & read in \texttt{app/\allowbreak{}tandemapp/\allowbreak{}conductor/\allowbreak{}monitor/\allowbreak{}notify.py,\allowbreak{} app/\allowbreak{}tandemapp/\allowbreak{}conductor/\allowbreak{}monitor/\allowbreak{}sessiontasks.py} \\
\texttt{TANDEM\_\allowbreak{}PF\_\allowbreak{}PREFIX} & & read in \texttt{app/\allowbreak{}proxy/\allowbreak{}nodeinstall.py,\allowbreak{} app/\allowbreak{}proxy/\allowbreak{}preflight.py} \\
\texttt{TANDEM\_\allowbreak{}PF\_\allowbreak{}SELF} & & read in \texttt{app/\allowbreak{}proxy/\allowbreak{}nodeagent.py,\allowbreak{} app/\allowbreak{}proxy/\allowbreak{}preflight.py} \\
\texttt{TANDEM\_\allowbreak{}PORT} & & read in \texttt{app/\allowbreak{}packaging/\allowbreak{}app.env.example,\allowbreak{} app/\allowbreak{}proxy/\allowbreak{}\_\allowbreak{}\_\allowbreak{}main\_\allowbreak{}\_\allowbreak{}.py} \\
\texttt{TANDEM\_\allowbreak{}PRICE\_\allowbreak{}1} & & read in \texttt{app/\allowbreak{}packaging/\allowbreak{}app.env.example} \\
\texttt{TANDEM\_\allowbreak{}PRICE\_\allowbreak{}1\_\allowbreak{}OFFPEAK} & & read in \texttt{app/\allowbreak{}packaging/\allowbreak{}app.env.example} \\
\texttt{TANDEM\_\allowbreak{}REF} & & read in \texttt{install.sh} \\
\texttt{TANDEM\_\allowbreak{}REPO} & & read in \texttt{install.sh} \\
\texttt{TANDEM\_\allowbreak{}ROUTE\_\allowbreak{}1} & & read in \texttt{app/\allowbreak{}packaging/\allowbreak{}app.env.example} \\
\texttt{TANDEM\_\allowbreak{}SCRATCH} & & read in \texttt{app/\allowbreak{}tandemapp/\allowbreak{}conductor/\allowbreak{}packages.py,\allowbreak{} app/\allowbreak{}tandemapp/\allowbreak{}knobs.py} \\
\texttt{TANDEM\_\allowbreak{}SKIP\_\allowbreak{}MEMCHECK} & & read in \texttt{install.sh} \\
\texttt{TANDEM\_\allowbreak{}SSH\_\allowbreak{}SECRET} & & read in \texttt{app/\allowbreak{}proxy/\allowbreak{}sshonce.py} \\
\texttt{TANDEM\_\allowbreak{}TRUSTED\_\allowbreak{}PROXY} & & read in \texttt{app/\allowbreak{}proxy/\allowbreak{}app.py,\allowbreak{} app/\allowbreak{}tandemapp/\allowbreak{}auth.py} \\
\texttt{TANDEM\_\allowbreak{}URL} & & read in \texttt{app/\allowbreak{}tandemapp/\allowbreak{}conductor/\allowbreak{}packages.py,\allowbreak{} app/\allowbreak{}tandemapp/\allowbreak{}knobs.py} \\
\texttt{TANDEM\_\allowbreak{}USAGE\_\allowbreak{}LEDGER} & & read in \texttt{app/\allowbreak{}packaging/\allowbreak{}app.env.example,\allowbreak{} app/\allowbreak{}proxy/\allowbreak{}\_\allowbreak{}\_\allowbreak{}main\_\allowbreak{}\_\allowbreak{}.py} \\
\end{reftable}


\section{The conductor's documents}\index{conductor documents}

\begin{table}[htbp]
\centering\small
\caption{The conductor's versioned documents.}\label{tab:ref-documents}
\begin{tabularx}{\linewidth}{@{}L{.16\linewidth}YL{.22\linewidth}@{}}
\toprule
\hd{document} & \hd{what it holds} & \hd{edited on} \\
\midrule
\code{profiles} & one document per profile: mode, decisions, task modes and caps, thinking, \code{ui}, outcomes, library, memory, shadow & Profiles \\
\code{router} & the classes and their policy, the classifier stages and thresholds, the structural rules, side shapes, outcomes, library & Routing \\
\code{packages} & every package: source, switch, limits, flags & Packages \\
\code{roles} & every role's contract & Profiles, Roles \\
\code{settings} & the default profile, transcripts, the disk tier, memory notes and the writer & Memory, Profiles \\
\code{deciders} & the decision models, their questions, Laya for scripts & Deciders \\
\code{residency} & the floor, the margin, the settle time, the load order, the claims & Profiles, Residency \\
\code{monitor} & the session monitor's policy (version 0 is \code{ops/supervisor.yaml}) & Monitor \\
\code{harness} & the harness's settings per backend, behind the proxy & Backends \\
\bottomrule
\end{tabularx}
\end{table}

A document is JSON, checked against the conductor's schema (\code{GET /v1/conductor/meta} serves it, and \code{docs/contract/conductor-v1/} holds a copy). A version stored by an older release loads with defaults for every field it lacks, and the journal says so once.

\chapter{HTTP API reference}\label{ch:ref-http}

This chapter lists every route, every header and every error of the engine's server. TandemApp's proxy speaks the same API for clients and has its own routes, headers and errors, and so does its extension inside the engine: TandemApp's manual lists them. The tables are written in \code{manual/tools/http_ref.py}, and \code{gen_reference.py --check} compares them with the code on every build: a route a handler tests, a header name or an error type the code uses that is missing here, or a row the code no longer has, fails the check.

\section{Who may call what}\index{auth}

The admin password is the admin token, \code{QSE_ADMIN_TOKEN}\index{qseadmintoken@\texttt{QSE\_ADMIN\_TOKEN}} in \code{secrets.env} It is not set by default. Then every route below answers anyone who reaches the port, and only \code{/metrics} keeps the metrics token when one is set (Section~\ref{sec:dash-login}). With a password set, an admin route needs it as a bearer header or a dashboard session (a cookie the token buys). A request from the machine itself (loopback, no proxy headers) needs neither, because the watchdog and the measurement tools rely on that (\code{--trust-loopback}\index{trust-loopback@\texttt{-\kern0pt-trust-loopback}}).

The OpenAI routes are open on the engine, with or without a password. Table~\ref{tab:ref-http-who} names the labels the route tables use.

\input{generated/http-who}

\section{The engine's routes}

Table~\ref{tab:ref-engine-routes} lists the engine server's routes. A route the server does not know answers 404. A route TandemApp serves answers 404 \code{not_on_engine} on a headless engine, and the page 404 \code{no_ui}; in chat mode the extension serves them. The proxy's routes without \code{--backends} answer 404 with a message that names the flag, and the decision routes without decision models as an unknown route. Table~\ref{tab:ref-side-routes} adds the status responder, which answers while the engine is stopped.

% Generated by manual/tools/gen_reference.py from the code. Do not edit by hand.

\begin{httptable}{The engine server's routes}{tab:ref-engine-routes}
\texttt{POST}\newline \texttt{/\allowbreak{}v1/\allowbreak{}chat/\allowbreak{}completions} & open & A chat completion, streamed as server-sent events or not. With \texttt{-{}-backends} a model of another backend is relayed there. \emph{Body:} OpenAI's chat body: model, messages, tools, stream, \texttt{max\_\allowbreak{}tokens} and the other fields of Section~\ref{sec:api-fields}. \emph{Answer:} chat.completion, or chat.completion.chunk events and \texttt{[DONE]}; usage carries \texttt{prompt\_\allowbreak{}tokens\_\allowbreak{}details.hidden\_\allowbreak{}tokens} for the conductor's notes. \\
\texttt{POST}\newline \texttt{/\allowbreak{}v1/\allowbreak{}completions} & open & A plain text completion. \emph{Body:} prompt, \texttt{max\_\allowbreak{}tokens} and the sampling fields. \emph{Answer:} \texttt{text\_\allowbreak{}completion}, or its chunks and \texttt{[DONE]}. \\
\texttt{POST}\newline \texttt{/\allowbreak{}v1/\allowbreak{}embeddings} & key & Relayed to the backend that serves the model; only with \texttt{-{}-backends}. \emph{Body:} OpenAI's embeddings body. \emph{Answer:} the backend's answer, unchanged. \\
\texttt{POST}\newline \texttt{/\allowbreak{}v1/\allowbreak{}responses} & key & OpenAI's Responses API, relayed to an openai backend; only with \texttt{-{}-backends}. \emph{Body:} the client's body, sent on with the backend's model id. \emph{Answer:} the backend's answer, unchanged. \\
\texttt{GET}\newline \texttt{/\allowbreak{}v1/\allowbreak{}models} & open & The model list: TandemLLM alone with the conductor on, else the served name; with \texttt{-{}-backends} every backend's models after it, behind the users' keys. \emph{Answer:} \texttt{\{object:\allowbreak{} list,\allowbreak{} data:\allowbreak{} [...]\}}, and \texttt{decision\_\allowbreak{}models} beside data when decision models are served. \\
\texttt{GET}\newline \texttt{/\allowbreak{}v1/\allowbreak{}models/\allowbreak{}\{id\}} & open & One model card: the engine's, a decision model's, or a backend's. \emph{Answer:} the card, or 404 \texttt{model\_\allowbreak{}not\_\allowbreak{}found}. \\
\texttt{GET}\newline \texttt{/\allowbreak{}v1/\allowbreak{}replica/\allowbreak{}identity} & open & What two engines compare before they serve one model as replicas: the code, the family, the weights, drafters, rows and tokenizer by content, and the decode settings (SRV-741). Follows the caller rule of \texttt{/\allowbreak{}v1/\allowbreak{}models.} \emph{Answer:} \texttt{\{identity:\allowbreak{} 1,\allowbreak{} served\_\allowbreak{}model,\allowbreak{} base\_\allowbreak{}model,\allowbreak{} version,\allowbreak{} code,\allowbreak{} family,\allowbreak{} layers,\allowbreak{} profile,\allowbreak{} weights,\allowbreak{} drafters,\allowbreak{} rows,\allowbreak{} tokenizer,\allowbreak{} decode,\allowbreak{} max\_\allowbreak{}len,\allowbreak{} stage\}}; 503 \texttt{identity\_\allowbreak{}pending} until the first call's computation ends. \\
\texttt{GET}\newline \texttt{/\allowbreak{}health} & open & Whether the server takes work. The full body follows the admin rule; anyone else gets the status alone. Also \texttt{/\allowbreak{}healthz} and \texttt{/\allowbreak{}v1/\allowbreak{}health.} \emph{Answer:} \texttt{\{status:\allowbreak{} ok\textbar{}draining\}}; the full body adds model, uptime, inflight, the queue, limits, sampling, cache, memory and parallel; 503 while draining. \\
\texttt{GET}\newline \texttt{/\allowbreak{}metrics} & metrics & The Prometheus families of Chapter~\ref{ch:ref-metrics}. \emph{Answer:} Prometheus text format. \\
\texttt{GET}\newline \texttt{/\allowbreak{}metrics/\allowbreak{}up} & open & One gauge, \texttt{qse\_\allowbreak{}up}, so a scraper can tell an engine that is down from a wrong token. \emph{Answer:} \texttt{qse\_\allowbreak{}up} 1, or 0 while draining. \\
\texttt{GET}\newline \texttt{/\allowbreak{}v1/\allowbreak{}cache/\allowbreak{}stats} & admin & Every cache's counters. \emph{Answer:} the cache report \texttt{/\allowbreak{}health} carries. \\
\texttt{GET}\newline \texttt{/\allowbreak{}v1/\allowbreak{}cache/\allowbreak{}memory?pools=\allowbreak{}1} & admin & The engine's memory by owner against what it held once loaded: each slot's anchors, end anchor, stash, pin and park, the store, the KV pool's pages and what no cache holds (SPD-257). pools=1 adds the allocator's segments by pool; it walks every block. \emph{Answer:} \texttt{\{allocated,\allowbreak{} reserved,\allowbreak{} load,\allowbreak{} caches,\allowbreak{} by\_\allowbreak{}owner,\allowbreak{} unattributed,\allowbreak{} limits,\allowbreak{} kv\_\allowbreak{}pool,\allowbreak{} allocator\_\allowbreak{}pools?\}}. \\
\texttt{POST}\newline \texttt{/\allowbreak{}v1/\allowbreak{}cache/\allowbreak{}peek} & admin & Where a chat's rows are before it is sent: the body is rendered as a request would be, the session's router notes are replayed, and the cache tiers are asked without running the model (SRV-349). \emph{Body:} a chat body (messages, tools, the effort fields) and session. \emph{Answer:} \texttt{\{tier:\allowbreak{} resident\textbar{}ram\_\allowbreak{}park\textbar{}ram\_\allowbreak{}snapshot\textbar{}nvme\_\allowbreak{}spill\textbar{}disk\_\allowbreak{}tier\textbar{}none,\allowbreak{} cached\_\allowbreak{}tokens,\allowbreak{} reread\_\allowbreak{}tokens,\allowbreak{} reread\_\allowbreak{}eta\_\allowbreak{}ms,\allowbreak{} tokens,\allowbreak{} restore\_\allowbreak{}ms,\allowbreak{} session\}}; 409 while a request holds the engine. \\
\texttt{POST}\newline \texttt{/\allowbreak{}v1/\allowbreak{}cache/\allowbreak{}prewarm} & admin & Reads a conversation the peek finds on the NVMe back into the engine while nothing runs or waits, so its next request resumes at the end anchor. \emph{Body:} as \texttt{/\allowbreak{}v1/\allowbreak{}cache/\allowbreak{}peek}. \emph{Answer:} \texttt{\{prewarmed,\allowbreak{} why?,\allowbreak{} restore?\}} and the peek's fields; 409 while busy. \\
\texttt{POST}\newline \texttt{/\allowbreak{}v1/\allowbreak{}tokenize} & open & Exact token counts with the model's own chat template, split by part; nothing runs on the GPU (SRV-347). With ids: true also the ids, the chat prompt's or the text's as \texttt{/\allowbreak{}v1/\allowbreak{}completions} reads it (SPD-249). 240 a minute from one address. \emph{Body:} \texttt{\{messages,\allowbreak{} tools?,\allowbreak{} ids?\}} or \texttt{\{text,\allowbreak{} ids?\}}. \emph{Answer:} \texttt{\{tokens,\allowbreak{} parts:\allowbreak{} \{system,\allowbreak{} tools,\allowbreak{} messages\},\allowbreak{} ids?\}}. \\
\texttt{POST}\newline \texttt{/\allowbreak{}v1/\allowbreak{}cache/\allowbreak{}clear?tier=\allowbreak{}ram\textbar{}disk\textbar{}all} & admin-bearer & Empties the caches in memory (ram, the default), the disk tier and the session spill (disk), or both (all). Nothing else is reset. \emph{Answer:} the cache report after the clear; 400 for another tier. \\
\texttt{POST}\newline \texttt{/\allowbreak{}v1/\allowbreak{}cache/\allowbreak{}release} & admin-bearer & Gives back memory the engine holds idle, cheapest first: free allocator blocks and caches while slot 0 is idle, the small models, idle side conversations, slot 0's idle conversation (to the NVMe first). chatd asks it before a sandbox waits (SRV-483). \emph{Body:} \texttt{\{bytes\}}. \emph{Answer:} what each step gave back and the total. \\
\texttt{GET}\newline \texttt{/\allowbreak{}dashboard/\allowbreak{}} & open & The dashboard's static files. They hold no data. \emph{Answer:} the page. \\
\texttt{GET}\newline \texttt{/\allowbreak{}v1/\allowbreak{}dashboard/\allowbreak{}summary?from=\allowbreak{}\&to=\allowbreak{}\&tz=\allowbreak{}\&model=\allowbreak{}\&client=\allowbreak{}\&backend=\allowbreak{}\&user=\allowbreak{}\&internal=\allowbreak{}} & dash & The Usage view's totals for a range. \emph{Answer:} \texttt{summary.schema.json} in app/contracts/dashboard-v1. \\
\texttt{GET}\newline \texttt{/\allowbreak{}v1/\allowbreak{}dashboard/\allowbreak{}usage?bucket=\allowbreak{}hour\textbar{}day\&from=\allowbreak{}\&to=\allowbreak{}\&tz=\allowbreak{}\&model=\allowbreak{}\&client=\allowbreak{}\&backend=\allowbreak{}\&user=\allowbreak{}\&internal=\allowbreak{}} & dash & Tokens, requests and cost per hour or day. \emph{Answer:} \texttt{usage.schema.json}. \\
\texttt{GET}\newline \texttt{/\allowbreak{}v1/\allowbreak{}dashboard/\allowbreak{}requests?limit=\allowbreak{}\&before=\allowbreak{}\&finish=\allowbreak{}\&model=\allowbreak{}\&client=\allowbreak{}\&backend=\allowbreak{}\&user=\allowbreak{}\&internal=\allowbreak{}} & dash & Recent request rows of the usage ledger, newest first. \emph{Answer:} \texttt{requests.schema.json}. \\
\texttt{GET}\newline \texttt{/\allowbreak{}v1/\allowbreak{}dashboard/\allowbreak{}system} & dash & The machine, the engine, the memory plan, the slots and the proxy's mode. \emph{Answer:} \texttt{system.schema.json}. \\
\texttt{GET}\newline \texttt{/\allowbreak{}v1/\allowbreak{}dashboard/\allowbreak{}live} & dash & What each request is doing now, as server-sent events. \emph{Answer:} \texttt{live.schema.json} events. \\
\texttt{GET}\newline \texttt{/\allowbreak{}v1/\allowbreak{}dashboard/\allowbreak{}logs?level=\allowbreak{}\&grep=\allowbreak{}\&backlog=\allowbreak{}\&since=\allowbreak{}\&follow=\allowbreak{}} & dash & The log: a backlog of up to 5,000 lines, then live lines as server-sent events; Last-Event-ID resumes; follow=0 answers JSON once. \emph{Answer:} \texttt{logs-\allowbreak{}line.schema.json} events, or \texttt{logs-\allowbreak{}json.schema.json}; 429 past four streams at once. \\
\texttt{GET}\newline \texttt{/\allowbreak{}v1/\allowbreak{}dashboard/\allowbreak{}logs/\allowbreak{}export} & dash & The whole log ring as one text file: a header with the version, the commit, the uptime, the time of export and the line count, then every line the ring holds, at all levels, oldest first. It takes no stream place. \emph{Answer:} text/plain as an attachment named tandemllm-\textless{}version\textgreater{}-\textless{}time\textgreater{}.log. \\
\texttt{GET}\newline \texttt{/\allowbreak{}v1/\allowbreak{}dashboard/\allowbreak{}metrics} & dash & The \texttt{/\allowbreak{}metrics} page under the dashboard's rule, for the Performance view. \emph{Answer:} Prometheus text format. \\
\texttt{GET}\newline \texttt{/\allowbreak{}v1/\allowbreak{}dashboard/\allowbreak{}parallel?since=\allowbreak{}} & dash & Every request in flight, the engine's and the relayed ones. \emph{Answer:} \texttt{parallel.schema.json}. \\
\texttt{GET}\newline \texttt{/\allowbreak{}v1/\allowbreak{}dashboard/\allowbreak{}session} & open & Whether a sign-in exists and is needed. \emph{Answer:} \texttt{\{authenticated,\allowbreak{} login,\allowbreak{} password\_\allowbreak{}set,\allowbreak{} expires\_\allowbreak{}at\}}. \\
\texttt{POST}\newline \texttt{/\allowbreak{}v1/\allowbreak{}dashboard/\allowbreak{}session} & open & The sign-in, with a password set and the login on: sets the \texttt{qse\_\allowbreak{}dash} cookie. Five failures a minute from one address and the sixth gets a 429. \emph{Body:} \texttt{\{token\}}. \emph{Answer:} the session; 401 for a wrong token. \\
\texttt{DELETE}\newline \texttt{/\allowbreak{}v1/\allowbreak{}dashboard/\allowbreak{}session} & open & Signs this session out. \emph{Answer:} the session, signed out. \\
\texttt{*}\newline \texttt{/\allowbreak{}v1/\allowbreak{}conductor/\allowbreak{}*} & admin & The conductor's documents, journal, sessions and memory (Table~\ref{tab:ref-conductor-routes}). Reads follow the dashboard rule, writes and every memory route the admin rule. 404 without \texttt{-{}-conductor-\allowbreak{}config}. \\
\texttt{*}\newline \texttt{/\allowbreak{}v1/\allowbreak{}setup/\allowbreak{}*} & admin & The setup guide (docs/setup.md): the model registry, the box, the memory estimate, the review plan, the apply and reset jobs and Wren. Every POST must also be same-origin JSON: 415 for another content type, 403 for another origin. In setup mode (\texttt{-{}-setup}) these, the dashboard and \texttt{/\allowbreak{}health} are the only routes that answer. \\
\texttt{*}\newline \texttt{/\allowbreak{}v1/\allowbreak{}config/\allowbreak{}*} & admin & The Backends and Setup screens (Table~\ref{tab:ref-config-routes}). Reads follow the dashboard rule, writes the admin rule. 404 without \texttt{-{}-backends}. \\
\texttt{POST}\newline \texttt{/\allowbreak{}v1/\allowbreak{}collect/\allowbreak{}\{backend\}} & collect & A batch from the vLLM or SGLang plugin, joined to the ledger rows; only with \texttt{-{}-backends}. \emph{Body:} the plugin's batch: plugin, engine, requests, steps. \emph{Answer:} \texttt{\{accepted:\allowbreak{} \{requests,\allowbreak{} steps\},\allowbreak{} joined\}}. \\
\texttt{POST}\newline \texttt{/\allowbreak{}v1/\allowbreak{}systemone} & admin & Typed decisions from the decision models, in the engine's idle window. \texttt{/\allowbreak{}v1/\allowbreak{}predict} is the same route. Only with QWEN38\_DECIDER\_MODELS. \emph{Body:} \{state, model?, questions: \texttt{\{id:\allowbreak{} \{type,\allowbreak{} instructions,\allowbreak{} criteria?\}\}}, task?, lang?\}. \emph{Answer:} \texttt{\{model,\allowbreak{} answers,\allowbreak{} usage,\allowbreak{} routing,\allowbreak{} latency\_\allowbreak{}ms,\allowbreak{} path\}}. \\
\texttt{POST}\newline \texttt{/\allowbreak{}v1/\allowbreak{}predict} & admin & The same as \texttt{/\allowbreak{}v1/\allowbreak{}systemone.} \emph{Body:} as \texttt{/\allowbreak{}v1/\allowbreak{}systemone}. \emph{Answer:} as \texttt{/\allowbreak{}v1/\allowbreak{}systemone}. \\
\texttt{GET}\newline \texttt{/\allowbreak{}readyz} & open & Whether a decision model is loaded and warm. \emph{Answer:} 200, or 503 with the reason. \\
\texttt{GET}\newline \texttt{/\allowbreak{}v1/\allowbreak{}laya} & open & What the Laya route takes and answers, its limits, and whether it is on. \emph{Answer:} a description of the route. \\
\texttt{POST}\newline \texttt{/\allowbreak{}v1/\allowbreak{}laya/\allowbreak{}ask} & open & A typed answer from Laya on one text or many, for the model's own scripts; rate-limited. \emph{Body:} \texttt{\{question,\allowbreak{} choices \textbar{} scale \textbar{} neither,\allowbreak{} text \textbar{} texts,\allowbreak{} threshold?,\allowbreak{} model?,\allowbreak{} deadline\_\allowbreak{}ms?\}}. \emph{Answer:} \texttt{\{label,\allowbreak{} p,\allowbreak{} sure,\allowbreak{} probs,\allowbreak{} model,\allowbreak{} ms,\allowbreak{} ...\}}, or \texttt{\{answers:\allowbreak{} [...]\}} for texts. \\
\texttt{POST}\newline \texttt{/\allowbreak{}v1/\allowbreak{}laya/\allowbreak{}drive/\allowbreak{}end} & open & Ends a playtest driver session and writes its journal row. \emph{Body:} the session's summary. \emph{Answer:} the row. \\
\texttt{GET}\newline \texttt{/\allowbreak{}v1/\allowbreak{}laya/\allowbreak{}kit/\allowbreak{}\{file\}} & open & A file of the playtest kit the model copies into a project. \emph{Answer:} the file. \\
\texttt{OPTIONS}\newline \texttt{/\allowbreak{}v1/\allowbreak{}laya/\allowbreak{}*} & open & A browser's CORS preflight for the Laya routes. Every other path answers OPTIONS with 501. \emph{Answer:} 204 with the CORS headers. \\
\texttt{GET}\newline \texttt{/\allowbreak{}v1/\allowbreak{}library} & open & The package library: every enabled skill's name, description and classes. \emph{Answer:} the list and the two routes below. \\
\texttt{GET}\newline \texttt{/\allowbreak{}v1/\allowbreak{}library/\allowbreak{}search?q=\allowbreak{}\&k=\allowbreak{}\&class=\allowbreak{}} & open & BM25 search over the skills, best first (k from 1 to 20, default 5). \emph{Answer:} [\texttt{\{name,\allowbreak{} description,\allowbreak{} score\}}]. \\
\texttt{GET}\newline \texttt{/\allowbreak{}v1/\allowbreak{}library/\allowbreak{}\{name\}} & open & One skill's text and its reference. \emph{Answer:} \texttt{\{name,\allowbreak{} description,\allowbreak{} classes,\allowbreak{} body,\allowbreak{} reference\}}. \\
\texttt{GET}\newline \texttt{/\allowbreak{}s/\allowbreak{}\{token\}/\allowbreak{}v1/\allowbreak{}memory} & session & The memory index this session sees. The \texttt{/\allowbreak{}s/\allowbreak{}\{token\}/\allowbreak{}} prefix also serves the library and Laya routes, filed to the session. \emph{Answer:} the notes' index lines. \\
\texttt{GET}\newline \texttt{/\allowbreak{}s/\allowbreak{}\{token\}/\allowbreak{}v1/\allowbreak{}memory/\allowbreak{}search?q=\allowbreak{}\&k=\allowbreak{}} & session & Memory notes that match, best first (k from 1 to 10). \emph{Answer:} [\texttt{\{id,\allowbreak{} kind,\allowbreak{} scope,\allowbreak{} status,\allowbreak{} title,\allowbreak{} score\}}]. \\
\texttt{GET}\newline \texttt{/\allowbreak{}s/\allowbreak{}\{token\}/\allowbreak{}v1/\allowbreak{}memory/\allowbreak{}\{id\}} & session & One note. \emph{Answer:} \texttt{\{id,\allowbreak{} title,\allowbreak{} kind,\allowbreak{} scope,\allowbreak{} status,\allowbreak{} body,\allowbreak{} ...\}}. \\
\texttt{POST}\newline \texttt{/\allowbreak{}s/\allowbreak{}\{token\}/\allowbreak{}v1/\allowbreak{}memory} & session & Writes, changes or archives a note, with the fields of the hidden memory block. \emph{Body:} \texttt{\{op,\allowbreak{} scope,\allowbreak{} kind,\allowbreak{} title,\allowbreak{} tags,\allowbreak{} body,\allowbreak{} id,\allowbreak{} reason\}}. \emph{Answer:} \texttt{\{op,\allowbreak{} id,\allowbreak{} title,\allowbreak{} ...\}}, or 422 with the reason as the type. \\
\end{httptable}


\section{Request fields of TandemLLM's own}\label{sec:ref-tandem-fields}\index{tandem object}

Two fields of a chat or completion request concern the wait for memory (Section~\ref{sec:memory-wait}). The engine reads them; Section~\ref{sec:api-fields} lists the OpenAI fields.

\begin{itemize}
\item \code{tandem.priority}\index{memory priority}: \code{"high"}, \code{"normal"} or \code{"low"}, the chat's place in the line for memory and in the order the KV pool's reclaim takes idle conversations. Any other value, or none, is normal (0.15.3).
\item \code{wait_for_memory: true}, in the body or in its \code{tandem} object: a request without a stream waits for memory up to \code{QSE_MEM_WAIT_S} instead of \code{--queue-timeout}, as the header \code{X-Tandem-Wait-For-Memory: 1} asks (0.15.3).
\end{itemize}

\section{Headers}\index{headers}

Table~\ref{tab:ref-headers} lists every header the engine reads or writes beyond the ones HTTP itself needs. \emph{in} is a header the server reads, \emph{out} one it writes.

% Generated by manual/tools/gen_reference.py from the code. Do not edit by hand.

\begin{reftable}{Every header the engine reads or writes (both: TandemApp's controller too)}{tab:ref-headers}{\hd{header} & \hd{direction, where} & \hd{meaning}}
\texttt{X-\allowbreak{}Tandem} & both, both & raw, shadow or on: lowers the conductor's mode for one request; the response says which mode ran \\
\texttt{X-\allowbreak{}Tandem-\allowbreak{}Profile} & both, both & picks a conductor profile; the response names the profile that ran; an unknown name is a 400 \\
\texttt{X-\allowbreak{}Tandem-\allowbreak{}Memory} & in, engine & off: no memory note is read or written; read: reads only, a block the model writes is hidden and dropped \\
\texttt{X-\allowbreak{}Tandem-\allowbreak{}Probe} & in, engine & off takes the request out of the probe snapshots \\
\texttt{X-\allowbreak{}Tandem-\allowbreak{}User} & in, engine & a proxy's user name for the memory notes, honoured from loopback or QWEN38\_MEMORY\_PROXIES \\
\texttt{X-\allowbreak{}Tandem-\allowbreak{}Relay-\allowbreak{}Id} & in, engine & a proxy's id for a relayed request, for the live counts; a relayed request without X-Tandem-User touches no memory \\
\texttt{X-\allowbreak{}Tandem-\allowbreak{}Parent} & in, engine & the request a call of TandemLLM to itself serves; the row becomes client kind internal \\
\texttt{X-\allowbreak{}Tandem-\allowbreak{}Worker} & in, engine & 1 on an agent's worker request: its session comes back, so it is never a guest, and its rows go to the NVMe rather than away while it waits \\
\texttt{X-\allowbreak{}Tandem-\allowbreak{}Parent-\allowbreak{}Session} & in, engine & a worker's parent session (the chat's): for 600 s after the worker's last request the parent's rows go to the NVMe rather than away, and its return tries slot 0 first \\
\texttt{X-\allowbreak{}Tandem-\allowbreak{}Regenerate} & in, engine & 1 serves the request as a regenerate (sampled, a fresh seed); 0 turns the detection off for it \\
\texttt{X-\allowbreak{}Tandem-\allowbreak{}Images-\allowbreak{}Left-\allowbreak{}Out} & out, engine & how many earlier images were left out of the prompt \\
\texttt{X-\allowbreak{}Tandem-\allowbreak{}Wait} & out, engine & memory, on a response that waited for memory before its prefill; its stream says so every 15 s \\
\texttt{X-\allowbreak{}Tandem-\allowbreak{}Wait-\allowbreak{}For-\allowbreak{}Memory} & in, engine & 1: a request not streamed waits for memory up to \texttt{QSE\_\allowbreak{}MEM\_\allowbreak{}WAIT\_\allowbreak{}S}, as \texttt{wait\_\allowbreak{}for\_\allowbreak{}memory}: true in the body does, instead of a 503 after the queue timeout \\
\texttt{X-\allowbreak{}Engine-\allowbreak{}Stop} & out, engine & pattern-stop(\textless{}label\textgreater{}) on a response without a stream that a stop pattern ended \\
\texttt{x-\allowbreak{}session-\allowbreak{}id} & in, both & opencode's session id: the router's session and the caches' session key \\
\texttt{X-\allowbreak{}Conversation-\allowbreak{}Id} & in, both & a conversation id for clients that send no session field; read after the body's own ids \\
\texttt{X-\allowbreak{}OpenWebUI-\allowbreak{}Chat-\allowbreak{}Id} & in, both & Open WebUI's chat id (a connection header with \texttt{\{\{CHAT\_\allowbreak{}ID\}\}}): the router's session for that chat \\
\texttt{X-\allowbreak{}QSE-\allowbreak{}Client} & in, engine & dashboard marks the dashboard's own requests (client kind dashboard) \\
\texttt{X-\allowbreak{}Real-\allowbreak{}IP} & in, both & set by a reverse proxy; a request that carries it is never trusted as local; the sign-in's rate limit counts per address \\
\texttt{X-\allowbreak{}Forwarded-\allowbreak{}For} & in, both & as X-Real-IP: a request that carries it is never trusted as local \\
\texttt{X-\allowbreak{}Forwarded-\allowbreak{}Proto} & in, both & https makes the session cookie Secure, and the notes' \texttt{\{tandem\_\allowbreak{}url\}} https \\
\texttt{X-\allowbreak{}Forwarded-\allowbreak{}Host} & in, engine & the host the notes' \texttt{\{tandem\_\allowbreak{}url\}} names, unless QWEN38\_TANDEM\_URL sets it \\
\texttt{X-\allowbreak{}Decider-\allowbreak{}Deadline-\allowbreak{}Ms} & in, engine & how long a decision request waits for the idle window (default 30,000, at most 120,000) \\
\texttt{X-\allowbreak{}Decider-\allowbreak{}Window} & in, engine & \texttt{in\_\allowbreak{}decode} is refused: decisions run in the idle window only \\
\texttt{X-\allowbreak{}Accel-\allowbreak{}Buffering} & out, both & no, on every stream, so nginx sends each event at once \\
\texttt{X-\allowbreak{}Content-\allowbreak{}Type-\allowbreak{}Options} & out, both & nosniff on the dashboard's static files \\
\texttt{Authorization} & in, both & Bearer: the admin token, the metrics token, the collect token or a proxy user key, by route; the engine's OpenAI routes ignore it and file the request under a hash of it \\
\texttt{User-\allowbreak{}Agent} & in, both & the client kind in the ledger: open-webui, opencode, openai-sdk, curl, other (a Hermes-shaped body: hermes) \\
\texttt{Last-\allowbreak{}Event-\allowbreak{}ID} & in, both & resumes the log stream after that line \\
\texttt{Retry-\allowbreak{}After} & out, both & with a 503 (queue full, memory gate, draining, a full backend), a 429, or a decision model that is loading \\
\texttt{WWW-\allowbreak{}Authenticate} & out, both & Bearer, with a 401 \\
\texttt{Set-\allowbreak{}Cookie} & out, both & \texttt{qse\_\allowbreak{}dash}, the dashboard session: HttpOnly, SameSite=Strict, Secure behind https \\
\end{reftable}


\section{Errors}\index{errors}

Errors are OpenAI-shaped:

\begin{jsonc}
{"error": {"message": "max_tokens must be at least 1", "type": "invalid_request_error",
           "param": "max_tokens", "code": null}}
\end{jsonc}

A 400 names the field or the message part it refuses in \code{param}. A refusal the client can retry (a full queue, the memory gate, a server that drains) has the type \code{server_busy}, the status in \code{code}, and a \code{Retry-After} header. A refusal after a stream opened (a queue timeout, a prompt too long, the memory gate) ends the stream with a chunk whose \code{finish_reason} is \code{"error"} and which carries the \code{error} object, then \code{[DONE]}. The decision routes keep their own body, \code{{"type": "error", "error": {"type", "message"}}}. Table~\ref{tab:ref-errors} lists every status and type.

% Generated by manual/tools/gen_reference.py from the code. Do not edit by hand.

\begin{reftable}{Every error status and type}{tab:ref-errors}{\hd{status, Retry-After} & \hd{error.type or code} & \hd{when}}
\texttt{400} & \texttt{invalid\_\allowbreak{}request\_\allowbreak{}error} & a request the engine will not serve as asked: error.param names the field (messages, a role, \texttt{max\_\allowbreak{}tokens}, tools, an image part, a prompt too long, a body that is not JSON or not an object) \\
\texttt{400} & \texttt{bad\_\allowbreak{}request} & a dashboard or conductor query that does not parse (a level, a range, a number) \\
\texttt{400} & \texttt{stage\_\allowbreak{}text\_\allowbreak{}only} & a request with images on the tail stage of a split model (\texttt{-{}-stage}), as error.code beside \texttt{invalid\_\allowbreak{}request\_\allowbreak{}error} \\
\texttt{400} & \texttt{confirm\_\allowbreak{}required} & a setup reset without the confirmation its level asks for (true, or the typed word) \\
\texttt{401} & \texttt{unauthorized} & a guarded route without the token or the session it needs, once an admin password is set; a wrong collect token \\
\texttt{403} & \texttt{local\_\allowbreak{}only}, \texttt{window\_\allowbreak{}key} & opening or closing a pairing window (POST \texttt{/\allowbreak{}v1/\allowbreak{}engine/\allowbreak{}pair/\allowbreak{}open}, \texttt{/\allowbreak{}close}) from another machine, or without the window key \\
\texttt{404} & \texttt{not\_\allowbreak{}found} & no such route, or a route whose part is off on this server (the conductor, the proxy, the decision models) \\
\texttt{404} & \texttt{not\_\allowbreak{}on\_\allowbreak{}engine} & a route TandemApp serves (the conductor, its library, memory and Laya, the backends) on an engine where no extension serves it \\
\texttt{404} & \texttt{no\_\allowbreak{}pairing} & a pairing route on a process that serves no pairing \\
\texttt{404} & \texttt{no\_\allowbreak{}ui} & the page, on every engine: TandemApp's controller serves it \\
\texttt{404} & \texttt{model\_\allowbreak{}not\_\allowbreak{}found} & a model id no backend serves, or one two backends serve (the message names the prefixed forms) (code \texttt{model\_\allowbreak{}not\_\allowbreak{}found}) \\
\texttt{409} & \texttt{busy}, \texttt{not\_\allowbreak{}loaded}, \texttt{role\_\allowbreak{}off} & an action that needs a role or a decision model in another state, or a role the roles document switches off (\texttt{role\_\allowbreak{}off}) \\
\texttt{409} & \texttt{stage\_\allowbreak{}front} & any route but \texttt{/\allowbreak{}health} on the front stage of a split model (\texttt{-{}-stage}): requests go to the tail stage \\
\texttt{422} & \texttt{invalid\_\allowbreak{}request\_\allowbreak{}error} & a value the schema refuses: a config field (error.field), a field an environment variable locks (error.variable), a decision request's fields \\
\texttt{429}\newline with \texttt{Retry-After} & \texttt{server\_\allowbreak{}busy} & a request without a stream that waited \texttt{-{}-queue-\allowbreak{}timeout} for the engine \\
\texttt{429}\newline with \texttt{Retry-After} & \texttt{too\_\allowbreak{}many} & a fifth log or live stream at once; a sixth failed sign-in in a minute \\
\texttt{429}\newline with \texttt{Retry-After} & \texttt{rate\_\allowbreak{}limited} & \texttt{/\allowbreak{}v1/\allowbreak{}tokenize} past its 240 calls a minute from one address; a pairing route past 20 a minute from one address \\
\texttt{500} & \texttt{internal\_\allowbreak{}error} & a bug: the message is the exception's; the request is in the log \\
\texttt{500} & \texttt{internal} & a PATCH \texttt{/\allowbreak{}v1/\allowbreak{}engine/\allowbreak{}settings} that raised: the message is the exception's; the call is journaled \\
\texttt{500} & \texttt{extension\_\allowbreak{}error} & a route of the extension that raised: the extension's bug, in the engine's log \\
\texttt{501} & none & OPTIONS on any path but the Laya routes \\
\texttt{503}\newline with \texttt{Retry-After} & \texttt{server\_\allowbreak{}busy} & the queue is full, the server is draining, or the memory gate or the admission by bytes cannot make room \\
\texttt{503}\newline with \texttt{Retry-After} & \texttt{memory\_\allowbreak{}wait\_\allowbreak{}timeout} & a request that waited for memory to its bound: \texttt{-{}-queue-\allowbreak{}timeout} without a stream, \texttt{QSE\_\allowbreak{}MEM\_\allowbreak{}WAIT\_\allowbreak{}S} streamed or with \texttt{wait\_\allowbreak{}for\_\allowbreak{}memory}; a stream that opened while it waited ends with it as an error event (type \texttt{server\_\allowbreak{}busy}, Retry-After: 30) \\
\texttt{503}\newline with \texttt{Retry-After} & \texttt{setup\_\allowbreak{}mode} & any route but the dashboard, \texttt{/\allowbreak{}health}, \texttt{/\allowbreak{}v1/\allowbreak{}models} and \texttt{/\allowbreak{}v1/\allowbreak{}setup/\allowbreak{}*} on a box in setup mode (Retry-After: 60) \\
\texttt{503}\newline with \texttt{Retry-After} & \texttt{warming}, \texttt{loading}, \texttt{refused}, \texttt{failed}, \texttt{not\_\allowbreak{}loaded}, \texttt{busy\_\allowbreak{}decoding} & a decision request while the model loads, was refused memory, failed, or the engine kept decoding past the deadline \\
\texttt{503}\newline with \texttt{Retry-After} & \texttt{identity\_\allowbreak{}pending} & the replica identity while the engine computes it, once after the start (Retry-After: 1) \\
\texttt{503}\newline with \texttt{Retry-After} & \texttt{stage\_\allowbreak{}down} & a completion on the tail stage of a split model while its link to the front stage is down (Retry-After: 5); a request in flight when the link fails ends with 500 and the same word in its message \\
\texttt{503}\newline with \texttt{Retry-After} & \texttt{engine\_\allowbreak{}failed} & a completion after the engine's CUDA context was lost: the sentence says so, the watchdog restarts the engine (Retry-After: 60) \\
\texttt{503}\newline with \texttt{Retry-After} & \texttt{provider\_\allowbreak{}unavailable} & the status responder while the engine is stopped, starting or held (code engine\_\textless{}state\textgreater{}) \\
\texttt{529} & \texttt{overloaded\_\allowbreak{}error} & a decision request the runtime could not take \\
\end{reftable}


\chapter{Metrics and the ledger}\label{ch:ref-metrics}

\section{Prometheus metrics}\index{metrics}\index{Prometheus}

\code{GET /metrics} serves the families below in the Prometheus text format, to anyone while no token is set, else behind the metrics or admin token, or to a request from the machine itself. \code{ops/monitoring/} has a scrape configuration, alert rules with their tests, and a Grafana dashboard generated by \code{tools/grafana_dashboard.py}. With the metrics token set, a scrape job needs it as a bearer token:

\begin{plain}
- job_name: tandemllm
  authorization: {type: Bearer, credentials_file: /etc/prometheus/tandem-metrics.token}
  static_configs: [{targets: ['gpu-host:8000']}]
\end{plain}

The engine's families never count a relayed request. The table lists every family the engine registers with a literal name in the code; families built at scrape time from a table (the cache families by \code{cache} label) are listed by their prefix in the help text of the page itself. The extension's families (\code{qse_conductor_*}, marked \code{"owner": "ext"} in \code{contracts/engine-api-v1/metrics-families.json}), the proxy's and the host collector's are in TandemApp's manual.

% Generated by manual/tools/gen_reference.py from the code. Do not edit by hand.

\begin{reftable}{Prometheus families the engine registers (the extension's are in TandemApp's manual)}{tab:ref-metrics}{\hd{Metric} & \hd{Type, labels} & \hd{What it counts}}
\texttt{qse\_\allowbreak{}requests\_\allowbreak{}total} & \texttt{counter,\allowbreak{} finish\_\allowbreak{}reason} & generations that finished, by the reason they finished \\
\texttt{qse\_\allowbreak{}request\_\allowbreak{}success\_\allowbreak{}total} & \texttt{counter} & generations that finished on the model's own stop token or the caller's token budget \\
\texttt{qse\_\allowbreak{}prompt\_\allowbreak{}tokens\_\allowbreak{}total} & \texttt{counter} & prompt tokens accepted, including tokens served from a cache \\
\texttt{qse\_\allowbreak{}generation\_\allowbreak{}tokens\_\allowbreak{}total} & \texttt{counter} & tokens written by the engine \\
\texttt{qse\_\allowbreak{}errors\_\allowbreak{}total} & \texttt{counter,\allowbreak{} type} & exceptions that ended a generation, by exception class \\
\texttt{qse\_\allowbreak{}images\_\allowbreak{}left\_\allowbreak{}out\_\allowbreak{}total} & \texttt{counter} & a conversation's earlier images left out of a request over the image caps (server/images.py) \\
\texttt{qse\_\allowbreak{}time\_\allowbreak{}to\_\allowbreak{}first\_\allowbreak{}token\_\allowbreak{}seconds} & \texttt{histogram} & arrival of the request to its first token, queue wait included \\
\texttt{qse\_\allowbreak{}time\_\allowbreak{}per\_\allowbreak{}output\_\allowbreak{}token\_\allowbreak{}seconds} & \texttt{histogram} & gap between two tokens of one generation. \\
\texttt{qse\_\allowbreak{}e2e\_\allowbreak{}request\_\allowbreak{}latency\_\allowbreak{}seconds} & \texttt{histogram} & arrival of the request to the end of its stream, queue wait included \\
\texttt{qse\_\allowbreak{}spec\_\allowbreak{}decode\_\allowbreak{}num\_\allowbreak{}accepted\_\allowbreak{}tokens\_\allowbreak{}total} & \texttt{counter} & drafted tokens the target agreed with. The token the target writes itself is not counted here, so accepted over drafted is the draft acceptance rate \\
\texttt{qse\_\allowbreak{}spec\_\allowbreak{}decode\_\allowbreak{}num\_\allowbreak{}draft\_\allowbreak{}tokens\_\allowbreak{}total} & \texttt{counter} & tokens proposed by the drafter, anchor excluded \\
\texttt{qse\_\allowbreak{}spec\_\allowbreak{}decode\_\allowbreak{}num\_\allowbreak{}drafts\_\allowbreak{}total} & \texttt{counter} & blocks the drafter proposed \\
\texttt{qse\_\allowbreak{}spec\_\allowbreak{}accept\_\allowbreak{}per\_\allowbreak{}block} & \texttt{histogram} & tokens committed by one verified block, the target's own token included \\
\texttt{qse\_\allowbreak{}draft\_\allowbreak{}width\_\allowbreak{}chosen\_\allowbreak{}total} & \texttt{counter,\allowbreak{} width} & blocks proposed at each width, anchor included. \\
\texttt{qse\_\allowbreak{}verify\_\allowbreak{}seconds} & \texttt{histogram} & one verify pass over a draft block, as the decode loop paid for it \\
\texttt{qse\_\allowbreak{}draft\_\allowbreak{}seconds} & \texttt{histogram} & one call to the drafter, lookup and block drafter together \\
\texttt{qse\_\allowbreak{}ext\_\allowbreak{}errors\_\allowbreak{}total} & \texttt{counter,\allowbreak{} hook} & extension hooks that failed (raised, a wrong shape, an act out of place); the request went on without them, by hook \\
\texttt{qse\_\allowbreak{}requests\_\allowbreak{}running} & \texttt{gauge} & generations the engine is working on \\
\texttt{qse\_\allowbreak{}requests\_\allowbreak{}waiting} & \texttt{gauge} & requests holding a socket and waiting for the engine \\
\texttt{qse\_\allowbreak{}requests\_\allowbreak{}refused\_\allowbreak{}total} & \texttt{counter,\allowbreak{} reason} & requests turned away with a Retry-After before the engine saw them \\
\texttt{qse\_\allowbreak{}cache\_\allowbreak{}misses\_\allowbreak{}total} & \texttt{counter,\allowbreak{} cache} & prefix lookups that found nothing \\
\texttt{qse\_\allowbreak{}cache\_\allowbreak{}evictions\_\allowbreak{}total} & \texttt{counter,\allowbreak{} cache} & entries dropped to stay inside the byte budget \\
\texttt{qse\_\allowbreak{}cache\_\allowbreak{}bytes} & \texttt{gauge,\allowbreak{} cache} & bytes held by each cache \\
\texttt{qse\_\allowbreak{}ngram\_\allowbreak{}table\_\allowbreak{}cache\_\allowbreak{}hits\_\allowbreak{}total} & \texttt{counter} & n-gram table rows served from the hot-row cache \\
\texttt{qse\_\allowbreak{}ngram\_\allowbreak{}table\_\allowbreak{}cache\_\allowbreak{}misses\_\allowbreak{}total} & \texttt{counter} & n-gram table rows read from disk (O\_DIRECT) for a hot-row cache miss \\
\texttt{qse\_\allowbreak{}ngram\_\allowbreak{}table\_\allowbreak{}blocks\_\allowbreak{}read\_\allowbreak{}total} & \texttt{counter} & 4 KiB blocks read from disk for the n-gram table \\
\texttt{qse\_\allowbreak{}ngram\_\allowbreak{}table\_\allowbreak{}cache\_\allowbreak{}bytes} & \texttt{gauge} & the n-gram table's hot-row cache, in bytes \\
\texttt{qse\_\allowbreak{}ngram\_\allowbreak{}table\_\allowbreak{}file\_\allowbreak{}bytes} & \texttt{gauge} & the n-gram table's file size, in bytes \\
\texttt{qse\_\allowbreak{}prepare\_\allowbreak{}runs\_\allowbreak{}total} & \texttt{counter} & summaries the engine prepared ahead of a request \\
\texttt{qse\_\allowbreak{}prepare\_\allowbreak{}served\_\allowbreak{}total} & \texttt{counter} & prepared summaries a request went on from \\
\texttt{qse\_\allowbreak{}prepare\_\allowbreak{}yielded\_\allowbreak{}total} & \texttt{counter} & prepare runs that gave the engine to a client request \\
\texttt{qse\_\allowbreak{}prepare\_\allowbreak{}too\_\allowbreak{}long\_\allowbreak{}total} & \texttt{counter} & prepare runs skipped: the conversation and the summary would not fit the context \\
\texttt{qse\_\allowbreak{}memory\_\allowbreak{}writer\_\allowbreak{}runs\_\allowbreak{}total} & \texttt{counter} & memory writer turns the engine ran \\
\texttt{qse\_\allowbreak{}memory\_\allowbreak{}writer\_\allowbreak{}notes\_\allowbreak{}total} & \texttt{counter} & memory notes the writer stored \\
\texttt{qse\_\allowbreak{}memory\_\allowbreak{}bytes} & \texttt{gauge,\allowbreak{} part} & the engine's memory against its limits, by part \\
\texttt{qse\_\allowbreak{}memory\_\allowbreak{}shed\_\allowbreak{}total} & \texttt{counter,\allowbreak{} level} & caches given back to keep the engine under its ceiling or its peak, by level \\
\texttt{qse\_\allowbreak{}cache\_\allowbreak{}entries} & \texttt{gauge,\allowbreak{} cache} & entries held by each cache \\
\texttt{qse\_\allowbreak{}disk\_\allowbreak{}cache\_\allowbreak{}bytes} & \texttt{gauge,\allowbreak{} part} & the disk tier: what its entries hold, its budget, the file system's free space and the floor it never crosses, the host staging buffer \\
\texttt{qse\_\allowbreak{}disk\_\allowbreak{}cache\_\allowbreak{}entries} & \texttt{gauge,\allowbreak{} kind} & the disk tier's files: anchors (one state each) and chunks (1,024 tokens of rows each) \\
\texttt{qse\_\allowbreak{}disk\_\allowbreak{}cache\_\allowbreak{}lookups\_\allowbreak{}total} & \texttt{counter} & prefills the disk tier could have shortened by a chunk or more \\
\texttt{qse\_\allowbreak{}disk\_\allowbreak{}cache\_\allowbreak{}hits\_\allowbreak{}total} & \texttt{counter} & prefills that started from the disk tier \\
\texttt{qse\_\allowbreak{}disk\_\allowbreak{}cache\_\allowbreak{}misses\_\allowbreak{}total} & \texttt{counter} & lookups the disk tier did not answer \\
\texttt{qse\_\allowbreak{}disk\_\allowbreak{}cache\_\allowbreak{}tokens\_\allowbreak{}restored\_\allowbreak{}total} & \texttt{counter} & prompt tokens read back from the disk instead of prefilled \\
\texttt{qse\_\allowbreak{}disk\_\allowbreak{}cache\_\allowbreak{}read\_\allowbreak{}bytes\_\allowbreak{}total} & \texttt{counter} & bytes a restore read \\
\texttt{qse\_\allowbreak{}disk\_\allowbreak{}cache\_\allowbreak{}written\_\allowbreak{}bytes\_\allowbreak{}total} & \texttt{counter} & bytes the writer wrote \\
\texttt{qse\_\allowbreak{}disk\_\allowbreak{}cache\_\allowbreak{}writes\_\allowbreak{}total} & \texttt{counter,\allowbreak{} kind} & files the writer wrote, by kind \\
\texttt{qse\_\allowbreak{}disk\_\allowbreak{}cache\_\allowbreak{}evictions\_\allowbreak{}total} & \texttt{counter,\allowbreak{} reason} & entries the disk tier gave up, by reason (budget, floor, ttl, orphan) \\
\texttt{qse\_\allowbreak{}disk\_\allowbreak{}cache\_\allowbreak{}dropped\_\allowbreak{}total} & \texttt{counter,\allowbreak{} reason} & writes not done, by reason (window: the next prefill came first; budget, floor, staging, aborted, too\_big) \\
\texttt{qse\_\allowbreak{}disk\_\allowbreak{}cache\_\allowbreak{}failed\_\allowbreak{}restores\_\allowbreak{}total} & \texttt{counter} & restores that failed a check (checksum, layout) and took the cold path \\
\texttt{qse\_\allowbreak{}disk\_\allowbreak{}cache\_\allowbreak{}last\_\allowbreak{}restore\_\allowbreak{}seconds} & \texttt{gauge} & the last restore's wall time \\
\texttt{qse\_\allowbreak{}cache\_\allowbreak{}collisions\_\allowbreak{}total} & \texttt{counter,\allowbreak{} cache} & stored prefixes whose 64-bit hash matched and whose tokens did not \\
\texttt{qse\_\allowbreak{}prefill\_\allowbreak{}tokens\_\allowbreak{}reused\_\allowbreak{}total} & \texttt{counter} & prompt tokens restored from a snapshot instead of forwarded \\
\texttt{qse\_\allowbreak{}prefill\_\allowbreak{}tokens\_\allowbreak{}forwarded\_\allowbreak{}total} & \texttt{counter} & prompt tokens run through the 64 layers \\
\texttt{qse\_\allowbreak{}gpu\_\allowbreak{}memory\_\allowbreak{}used\_\allowbreak{}bytes} & \texttt{gauge} & allocator bytes held by the engine \\
\texttt{qse\_\allowbreak{}gpu\_\allowbreak{}memory\_\allowbreak{}reserved\_\allowbreak{}bytes} & \texttt{gauge} & bytes the allocator has taken from the driver and not given back \\
\texttt{qse\_\allowbreak{}unified\_\allowbreak{}memory\_\allowbreak{}free\_\allowbreak{}bytes} & \texttt{gauge} & free bytes in the board's one memory pool \\
\texttt{qse\_\allowbreak{}engine\_\allowbreak{}uptime\_\allowbreak{}seconds} & \texttt{gauge} & seconds since the server started \\
\texttt{qse\_\allowbreak{}engine\_\allowbreak{}info} & \texttt{gauge,\allowbreak{} version,\allowbreak{} model,\allowbreak{} drafter,\allowbreak{} width,\allowbreak{} tree,\allowbreak{} nvfp4,\allowbreak{} fp8\_\allowbreak{}head,\allowbreak{} max\_\allowbreak{}len,\allowbreak{} caches} & always 1. The labels are the configuration the numbers above were produced by \\
\texttt{qse\_\allowbreak{}http\_\allowbreak{}requests\_\allowbreak{}total} & \texttt{counter,\allowbreak{} route,\allowbreak{} code} & HTTP responses by route (a fixed set) and status code \\
\texttt{qse\_\allowbreak{}prompt\_\allowbreak{}tokens\_\allowbreak{}cached\_\allowbreak{}total} & \texttt{counter} & prompt tokens restored from the state store (all of them on a response-cache replay) \\
\texttt{qse\_\allowbreak{}reasoning\_\allowbreak{}tokens\_\allowbreak{}total} & \texttt{counter} & committed tokens inside the reasoning block \\
\texttt{qse\_\allowbreak{}suffix\_\allowbreak{}store\_\allowbreak{}drafts\_\allowbreak{}total} & \texttt{counter,\allowbreak{} outcome} & proposed blocks for which the persistent suffix store matched the context (hit) or did not (miss); counted once a block, never once a lookup \\
\texttt{qse\_\allowbreak{}request\_\allowbreak{}queue\_\allowbreak{}seconds} & \texttt{histogram} & arrival to the engine lock, per request \\
\texttt{qse\_\allowbreak{}request\_\allowbreak{}prefill\_\allowbreak{}seconds} & \texttt{histogram} & the engine lock to the first token (template, tokenise, prefill), per request \\
\texttt{qse\_\allowbreak{}request\_\allowbreak{}decode\_\allowbreak{}tokens\_\allowbreak{}per\_\allowbreak{}second} & \texttt{histogram} & per decoded request: (completion tokens - 1) / (last token - first token), the atlas row's convention, and the \texttt{predicted\_\allowbreak{}per\_\allowbreak{}second} a response carries \\
\texttt{qse\_\allowbreak{}request\_\allowbreak{}prefill\_\allowbreak{}tokens\_\allowbreak{}per\_\allowbreak{}second} & \texttt{histogram} & per request with a forwarded prompt: forwarded tokens / prefill seconds \\
\texttt{qse\_\allowbreak{}request\_\allowbreak{}prompt\_\allowbreak{}tokens} & \texttt{histogram} & prompt tokens per request, cached ones included \\
\texttt{qse\_\allowbreak{}request\_\allowbreak{}completion\_\allowbreak{}tokens} & \texttt{histogram} & completion tokens per request \\
\texttt{qse\_\allowbreak{}suffix\_\allowbreak{}store\_\allowbreak{}tokens} & \texttt{gauge} & token ids in the persistent suffix store \\
\texttt{qse\_\allowbreak{}state\_\allowbreak{}store\_\allowbreak{}budget\_\allowbreak{}bytes} & \texttt{gauge} & the state cache's byte budget \\
\texttt{qse\_\allowbreak{}process\_\allowbreak{}resident\_\allowbreak{}memory\_\allowbreak{}bytes} & \texttt{gauge} & resident set of the server process \\
\texttt{qse\_\allowbreak{}engine\_\allowbreak{}start\_\allowbreak{}time\_\allowbreak{}seconds} & \texttt{gauge} & when the server started, unix seconds \\
\texttt{qse\_\allowbreak{}log\_\allowbreak{}subscribers} & \texttt{gauge} & open /v1/dashboard/logs streams \\
\texttt{qse\_\allowbreak{}usage\_\allowbreak{}ledger\_\allowbreak{}rows\_\allowbreak{}total} & \texttt{counter} & rows the usage ledger has written \\
\texttt{qse\_\allowbreak{}usage\_\allowbreak{}ledger\_\allowbreak{}dropped\_\allowbreak{}total} & \texttt{counter} & rows the usage ledger lost: a full queue or a failed write \\
\texttt{qse\_\allowbreak{}build\_\allowbreak{}info} & \texttt{gauge} & always 1. version, the git commit, the code hash row3 records (tools/row3.py::code\_hash) and a hash of the effective flags \\
\end{reftable}


\section{The usage ledger}\label{sec:ledger}\index{usage ledger}

\code{server/api/ledger.py} writes one SQLite row per request: counts, timings, the finish reason, the cache source, a client id. Never the prompt, the answer, a tool name or argument, a conversation id, the \code{user} field, an IP address or a raw token. Refused requests (503, 429) and rejected ones (400) are rows too, with null token fields, so the request count is honest.

It never slows or fails a request. A row is put on a bounded queue (10,000), a writer thread commits in batches at most every second, and a disk error drops that batch, counts it and logs it at most once a minute. Rows older than \code{--usage-retention-days} (at least 400) are deleted once a day, and once an ISO week an online backup goes beside the file, the newest 8 kept. The server takes the path from \code{--usage-ledger}\index{usage-ledger@\texttt{-\kern0pt-usage-ledger}} only, never from the environment, so a benchmark server started with \code{serve.env} exported never writes the operator's ledger. Migrations are forward only and in place, and keep every row.

% Generated by manual/tools/gen_reference.py from the code. Do not edit by hand.

\begin{reftable}{The usage ledger's requests table (schema version 12)}{tab:ref-ledger}{\hd{Column} & \hd{Type, since} & \hd{What it holds}}
\texttt{ts\_\allowbreak{}ms} & \texttt{INTEGER,\allowbreak{} v1} & when the request arrived, ms since the epoch \\
\texttt{request\_\allowbreak{}id} & \texttt{TEXT,\allowbreak{} v1} & the response id \\
\texttt{model} & \texttt{TEXT,\allowbreak{} v1} & the model name the client sent \\
\texttt{client\_\allowbreak{}id} & \texttt{TEXT,\allowbreak{} v1} & k: and 12 hex digits of the bearer token's sha256, or anon \\
\texttt{client\_\allowbreak{}kind} & \texttt{TEXT,\allowbreak{} v1} & read from the User-Agent: open-webui, openai-sdk, opencode, hermes, curl, dashboard, other \\
\texttt{endpoint} & \texttt{TEXT,\allowbreak{} v1} & the route \\
\texttt{stream} & \texttt{INTEGER,\allowbreak{} v1} & 1 for a streamed request \\
\texttt{status} & \texttt{INTEGER,\allowbreak{} v1} & the HTTP status \\
\texttt{finish\_\allowbreak{}reason} & \texttt{TEXT,\allowbreak{} v1} & stop, length, tool\_calls, timeout, error, abandoned \\
\texttt{prompt\_\allowbreak{}tokens} & \texttt{INTEGER,\allowbreak{} v1} & prompt tokens as the client sent them \\
\texttt{cached\_\allowbreak{}tokens} & \texttt{INTEGER,\allowbreak{} v1} & prompt tokens served from a cache \\
\texttt{completion\_\allowbreak{}tokens} & \texttt{INTEGER,\allowbreak{} v1} & tokens streamed to the client \\
\texttt{reasoning\_\allowbreak{}tokens} & \texttt{INTEGER,\allowbreak{} v1} & visible reasoning tokens \\
\texttt{queue\_\allowbreak{}ms} & \texttt{REAL,\allowbreak{} v1} & time waiting for the engine \\
\texttt{prompt\_\allowbreak{}ms} & \texttt{REAL,\allowbreak{} v1} & prefill time \\
\texttt{ttft\_\allowbreak{}ms} & \texttt{REAL,\allowbreak{} v1} & time to the first token \\
\texttt{decode\_\allowbreak{}ms} & \texttt{REAL,\allowbreak{} v1} & time after the first token \\
\texttt{total\_\allowbreak{}ms} & \texttt{REAL,\allowbreak{} v1} & arrival to the end \\
\texttt{decode\_\allowbreak{}tps} & \texttt{REAL,\allowbreak{} v1} & (completion tokens - 1) / decode time; null for a replay \\
\texttt{prefill\_\allowbreak{}tps} & \texttt{REAL,\allowbreak{} v1} & forwarded prompt tokens a second \\
\texttt{blocks} & \texttt{INTEGER,\allowbreak{} v1} & verify rounds \\
\texttt{draft\_\allowbreak{}tokens} & \texttt{INTEGER,\allowbreak{} v1} & tokens the drafters proposed \\
\texttt{draft\_\allowbreak{}accepted} & \texttt{INTEGER,\allowbreak{} v1} & drafted tokens the model agreed with \\
\texttt{tool\_\allowbreak{}calls} & \texttt{INTEGER,\allowbreak{} v1} & tool calls in the answer \\
\texttt{thinking} & \texttt{INTEGER,\allowbreak{} v1} & 1 when the model thought \\
\texttt{cache\_\allowbreak{}source} & \texttt{TEXT,\allowbreak{} v1} & where the prompt resumed: resident, end, store, disk, spill, continuation, prepared, response, none \\
\texttt{max\_\allowbreak{}tokens} & \texttt{INTEGER,\allowbreak{} v1} & the output limit \\
\texttt{error\_\allowbreak{}type} & \texttt{TEXT,\allowbreak{} v1} & the exception class or the refusal \\
\texttt{engine\_\allowbreak{}version} & \texttt{TEXT,\allowbreak{} v1} & the VERSION file \\
\texttt{code\_\allowbreak{}sha} & \texttt{TEXT,\allowbreak{} v1} & a hash over the code folders \\
\texttt{model\_\allowbreak{}served} & \texttt{TEXT,\allowbreak{} v2} & TandemLLM with the conductor on \\
\texttt{profile} & \texttt{TEXT,\allowbreak{} v2} & the conductor profile \\
\texttt{profile\_\allowbreak{}version} & \texttt{INTEGER,\allowbreak{} v2} & its version \\
\texttt{conductor\_\allowbreak{}mode} & \texttt{TEXT,\allowbreak{} v2} & raw, shadow or on \\
\texttt{conductor\_\allowbreak{}acts} & \texttt{INTEGER,\allowbreak{} v2} & the session monitor's executed acts \\
\texttt{conductor\_\allowbreak{}hidden\_\allowbreak{}tokens} & \texttt{INTEGER,\allowbreak{} v2} & tokens the model saw that the client did not send \\
\texttt{conductor\_\allowbreak{}acted} & \texttt{INTEGER,\allowbreak{} v2} & 1 when the conductor changed what the model saw or wrote \\
\texttt{lane} & \texttt{TEXT,\allowbreak{} v2} & decode, prefill or self for a side-lane request \\
\texttt{backend} & \texttt{TEXT,\allowbreak{} v3} & the backend id of a relayed request \\
\texttt{backend\_\allowbreak{}kind} & \texttt{TEXT,\allowbreak{} v3} & its kind \\
\texttt{backend\_\allowbreak{}model} & \texttt{TEXT,\allowbreak{} v3} & the id sent upstream \\
\texttt{public\_\allowbreak{}model} & \texttt{TEXT,\allowbreak{} v3} & the id the client sent \\
\texttt{usage\_\allowbreak{}source} & \texttt{TEXT,\allowbreak{} v3} & backend or none \\
\texttt{cache\_\allowbreak{}write\_\allowbreak{}tokens} & \texttt{INTEGER,\allowbreak{} v3} & tokens a provider charged as a cache write \\
\texttt{cost\_\allowbreak{}micros} & \texttt{INTEGER,\allowbreak{} v3} & cost in micro-dollars at the model's price \\
\texttt{proxy\_\allowbreak{}ms} & \texttt{REAL,\allowbreak{} v3} & the proxy's own time \\
\texttt{retries} & \texttt{INTEGER,\allowbreak{} v3} & tries again before the first byte \\
\texttt{fallback\_\allowbreak{}from} & \texttt{TEXT,\allowbreak{} v3} & the backend tried first \\
\texttt{upstream\_\allowbreak{}status} & \texttt{INTEGER,\allowbreak{} v3} & the backend's HTTP status \\
\texttt{user} & \texttt{TEXT,\allowbreak{} v3} & the label of the proxy key \\
\texttt{preemptions} & \texttt{INTEGER,\allowbreak{} v4} & preemptions (vLLM) or retractions (SGLang), from the plugin \\
\texttt{engine\_\allowbreak{}finish} & \texttt{TEXT,\allowbreak{} v4} & the backend engine's own finish reason \\
\texttt{corrupted} & \texttt{INTEGER,\allowbreak{} v4} & vLLM's NaN check \\
\texttt{l2\_\allowbreak{}join} & \texttt{TEXT,\allowbreak{} v4} & exact or approximate join of the plugin's numbers \\
\texttt{wait\_\allowbreak{}own\_\allowbreak{}ms} & \texttt{REAL,\allowbreak{} v5} & ms the request waited behind the engine's own work (the shadow reader, a summary written ahead, the conductor's question) \\
\texttt{parent\_\allowbreak{}request} & \texttt{TEXT,\allowbreak{} v6} & the turn an internal row served (the router's question names the turn it routes) \\
\texttt{slot} & \texttt{INTEGER,\allowbreak{} v7} & the engine's sequence slot with -\kern0pt-parallel N (0 the main slot, 1 and up a side slot); null without \\
\texttt{engine\_\allowbreak{}cached} & \texttt{INTEGER,\allowbreak{} v8} & the prompt tokens the engine read from a cache, the hidden notes included (cached\_tokens counts the client's among them) \\
\texttt{forwarded\_\allowbreak{}tokens} & \texttt{INTEGER,\allowbreak{} v8} & what the prefill forwarded, hidden tokens included \\
\texttt{wait\_\allowbreak{}mem\_\allowbreak{}ms} & \texttt{REAL,\allowbreak{} v8} & ms the parallel memory admission held the request (inside queue\_ms) \\
\texttt{sampling} & \texttt{TEXT,\allowbreak{} v9} & how the request sampled: null greedy, client (its own temperature), regenerate (TandemLLM sampled a regenerate) \\
\texttt{system\_\allowbreak{}fingerprint} & \texttt{TEXT,\allowbreak{} v9} &  \\
\texttt{node} & \texttt{TEXT,\allowbreak{} v9} &  \\
\texttt{deployment} & \texttt{TEXT,\allowbreak{} v9} &  \\
\texttt{upstream\_\allowbreak{}request} & \texttt{TEXT,\allowbreak{} v9} &  \\
\end{reftable}
{\hbadness=3000 \noindent The conductor's journal shares the file: \texttt{conductor\_events} has 53 columns (\texttt{id}, \texttt{ts\_\allowbreak{}ms}, \texttt{request\_\allowbreak{}id}, \texttt{session\_\allowbreak{}key}, \texttt{kind}, \texttt{task}, \texttt{tier}, \texttt{mode}, \texttt{sp}, \texttt{rule}, \texttt{trigger}, \texttt{score}, \texttt{decision}, \texttt{action}, \texttt{cap}, \texttt{tier1\_\allowbreak{}model}, \texttt{question\_\allowbreak{}id}, \texttt{probabilities}, \texttt{tier1\_\allowbreak{}ms}, \texttt{hidden\_\allowbreak{}text\_\allowbreak{}id}, \texttt{hidden\_\allowbreak{}tokens}, \texttt{tokens\_\allowbreak{}at\_\allowbreak{}cut}, \texttt{position}, \texttt{outcome}, \texttt{next\_\allowbreak{}checker}, \texttt{policy\_\allowbreak{}version}, \texttt{client\_\allowbreak{}kind}, \texttt{model}, \texttt{profile}, \texttt{profile\_\allowbreak{}version}, \texttt{router\_\allowbreak{}version}, \texttt{class\_\allowbreak{}id}, \texttt{confidence}, \texttt{source}, \texttt{kept}, \texttt{stages}, \texttt{classifier\_\allowbreak{}ms}, \texttt{effort}, \texttt{budget}, \texttt{budget\_\allowbreak{}source}, \texttt{thinking}, \texttt{decisions}, \texttt{injected}, \texttt{replayed}, \texttt{shadow}, \texttt{digest}, \texttt{label}, \texttt{label\_\allowbreak{}note}, \texttt{labelled\_\allowbreak{}at}, \texttt{acted}, \texttt{changed}, \texttt{detail}, \texttt{extra}), and \texttt{conductor\_texts} keeps every hidden text once, by its hash.\par}


A row's speed figures follow the benchmark row's formula, and a replay or a side-lane request has no decode rate. A relayed request has null queue, prefill, blocks, draft counts and cache source unless a plugin reported them, in which case \code{l2_join} says how the row was joined. A quick look:

\begin{shell}
sqlite3 ~/.qwen38-spark-engine/usage/ledger.sqlite3 \
  "select datetime(ts_ms/1000,'unixepoch','localtime'), finish_reason, prompt_tokens,
          completion_tokens, round(decode_tps,1), cache_source from requests order by id desc limit 10"
\end{shell}

\chapter{Troubleshooting}\label{ch:troubleshooting}

Each entry is a symptom, what causes it, and what to do. \code{~/TandemLLM/bin/tandem logs} follows the server log on an installed engine, and \code{GET /v1/dashboard/logs} streams it. TandemApp's manual has the dashboard's, the Chat tab's and the proxy's entries.

\section{Installing}\label{sec:trouble-install}

\paragraph{The installer stops: no tag, branch or commit 'v0.15.4'.} The repository it clones has no such tag. The installer clones the repository \code{--repo} names, the public one by default, at the tag \code{--ref} names, \code{v0.15.4} by default. To install from another repository, such as a fork, pass its address with \code{--repo}. The same missing tag makes \code{curl} answer 404 for the installer's own URL.

\paragraph{\texttt{install.sh -{}-setup} refuses: the directory already has a configured engine.} A box that switched to a profile counts as configured, and \code{--setup} would stop that engine. To update it, add \code{--force-setup}, then run \code{use-profile.sh} with the profile again; its files are already on disk.

\paragraph{\texttt{use-profile.sh} refuses the download: not enough disk.} The fetcher keeps 60~GB free after every download, per filesystem, and counts the profile's weights and the files beside them first. \code{--dry-run} prints each file with its size. Free the space, or pick a smaller profile: \code{exl3-2.5} needs 12.0~GB of weights beside the 27B's shared 40~GB.

\paragraph{A download stopped halfway.} Run the same \code{use-profile.sh} command again. It keeps every file that finished and checked, pins the same commit, and continues; a folder appears only when all its files are there.

\paragraph{The setup guide's community section says \enquote{no published list yet}.} GitHub answered 404 for \texttt{models/community-index.json} on the public repository's \texttt{main}: no release there carries the list yet. That is a note, not an error. The page shows the list from the box's last fetch, or the one that came with its release, and Refresh asks again (Section~\ref{sec:community-sync}). A real failure shows one short line instead, such as \enquote{GitHub could not be reached}, and the server's log has the URL.

\section{The engine}

\paragraph{The start refuses: an engine is already alive.} One engine per board. Two engines loading side by side run a DGX Spark out of memory, and a board in that state stops answering until a power cycle. A loading engine counts although it answers nothing yet, because it binds its port only after the weights are in. \code{. ops/engines.sh; engine_pids} lists every engine process; stop the other one before starting. The same rule holds for a vLLM or SGLang on the same board: check the free memory first and keep at least 15~GiB of \code{MemAvailable} at the worst case.

\paragraph{A request gets a 503 with Retry-After, and the log names the peak.} The memory gate (Section~\ref{sec:memory-gate}). Before every prefill the engine gives memory back cheapest first until the prompt's workspace fits under the declared peak (77~GiB); when it still does not fit, it refuses rather than wedge the board. The System view's memory plan shows what holds memory: a small model beside the engine holds its claim, and a 262,144-token context holds about 13~GiB more than a 131,072-token one. If the refusals go on after a role loaded, unload it in the Conductor view or lower \code{QWEN38_SHADOW_GB}.

\paragraph{A role does not load.} The residency planner refused it: \code{refused_memory}, \code{refused_hold} or \code{refused_load_in_progress}. It loads one model at a time, never during a hold, and only when the board keeps its floor and margin (15 + 6~GiB) after the engine's growth to its peak and the role's claim. The Deciders screen and the System view say why, and \code{residency.jsonl} in the conductor's directory journals each decision. The narrator's keeper tries again in the next idle window.

\paragraph{The engine does not start: the image processor resolved to the pil backend.} torchvision is missing or too old. Install the pair that matches torch (\code{pillow==11.3.0}, \code{torchvision==0.28.0+cu130} from the CUDA 13.0 index), or start with \code{--image-backend pil}, or \code{--vision off} for text only. The installer's venv has both.

\paragraph{A benchmark is slower than the paper, or its text differs.} With the conductor on, every request runs through the router: a class, an effort cap and a hidden note. Send \code{X-Tandem: raw}, or the profile \code{bench}, for the raw path. Then check the lookup corpus: the installer builds a public one, and the paper's corpus is not published. And check the store view: a store that has seen the benchmark's prompts inflates it (Chapter~\ref{ch:measurement}).

\paragraph{A client asks for a model name the engine does not list.} With the conductor on, \code{/v1/models} lists \code{TandemLLM} alone, and any model name is served as \code{TandemLLM}. Nothing to fix.

\paragraph{A turn waits a long time before its first token.} With one sequence (\code{QSE_PARALLEL=1}; the shipped profile runs 8 lanes) the engine serves one turn at a time. A streamed request behind a running turn gets an SSE comment every 10~s with its position and an estimate, and with the progress lines on a waiting line in the client; it waits as long as a request may run. A request without a stream gets a 429 after \code{--queue-timeout}. Short requests go through the side lane. A full queue (8) is a 503 at once.

\paragraph{A turn waits for memory.} The chat or the client says it waits for memory, the response has \code{X-Tandem-Wait: memory}, or the Lanes view shows \code{waiting_memory}. The turn is paused and starts by itself once memory frees (Section~\ref{sec:memory-wait}). Its line in the dashboard says what it needs, what holds the rest and when the wait ends; the engine log repeats it every 15~s. A streamed turn gives up only after \code{QSE_MEM_WAIT_S} (3~hours) with a 503 \code{memory_wait_timeout}; one without a stream after \code{--queue-timeout}, unless it sends \code{wait_for_memory: true}. A shorter conversation or fewer chats at once needs less memory, and a high priority lets the chat that matters go first. Since 0.15.4 a short turn on a long chat is charged the rows it forwards, so a wait of minutes for such a turn is worth a report with the exported log (\code{GET /v1/dashboard/logs/export}).

\paragraph{The next turn re-reads the whole conversation.} Look at the request's cache source in the Performance view or the Dev view's detail: \code{end}, \code{resident}, \code{prefix}, \code{session}, \code{disk}, \code{spill}, \code{continuation}, or \code{none} for a cold prefill. A client that rewrites its history (another system text, a changed earlier message) parts from the cached rows where it changed. A short conversation started while a longer one is resident is a guest of it and re-reads its first answer at its second turn. Chapter~\ref{ch:caches} lists each case. With \code{--parallel N}, each slot keeps the conversation it served last, and a turn goes back to the slot that holds its conversation (the request row's \code{slot}, the \code{[parallel]} log line). When memory ran short, the KV pool may have given that conversation's rows back: with the session spill on it wrote them to the NVMe drive (the SSD) first, and the turn reads them back with the cache source \code{spill} (Section~\ref{sec:spill}). With the spill off it reads its prompt again. A new user message in Open WebUI after the agent changed directory used to part at the tool block (its working directory is in \code{run_command}'s description); under a conductor plan it no longer does (\code{QWEN38_CLIENT_CWD}). If the turn's source is \code{none} right after short requests of TandemLLM's own (\code{json} rows of the router), check \code{refused_stash} in \code{/health}'s resident report and \code{cache.memory} against the ceiling.

\paragraph{An agent turn ends with no answer and no tool call.} The model wrote its tool call inside its thinking and stopped there. With the session monitor acting and E1 on in the request's profile (a profile stored before this release may have it off: set it to act on the Monitor screen), the call goes out as the answer when it parses and names an offered tool, under the line \emph{tool call written in the thinking, sending it}; otherwise the turn is continued, and after two tries it ends with one sentence. On the raw path the engine returns what the model wrote.

\paragraph{A long conversation shows X-Tandem-Images-Left-Out.} A request carries at most 16 images and 65,536 image rows. When the excess is in earlier messages, the oldest images are left out in steps of 8 and the request is served; only the newest message's own images over a cap are a 400.

\paragraph{The dashboard says the engine is stopped for maintenance or starting.} The engine is away and the reverse proxy shows the status responder's page (Section~\ref{sec:status-page}): the reason, since when, about how long. It opens the dashboard by itself when \code{/health} answers.

\paragraph{A request answers 400 and \texttt{param} names a field.} The request's shape is wrong: no user message, an unknown role, a late system message, \code{max_tokens} below 1, an empty prompt (Section~\ref{sec:api-fields}).

\paragraph{Open WebUI shows a literal \texttt{</think>} while an answer streams.} Open WebUI 0.11.3 shows the stream piece that carries \code{</think>} as message text until the round ends. Under a conductor plan the engine sends Open WebUI its reasoning in \code{reasoning_content}, which has no tag; a request with \code{reasoning_format: tags}, the raw path, or \code{QWEN38_CLIENT_REASONING} set empty bring the tags back. The saved chat is clean either way.

\paragraph{The model writes ``Laya unavailable''.} Its scripts could not reach the engine from where they run. Set \code{QWEN38_TANDEM_URL} to an address that machine reaches.

\section{Memory notes}

\paragraph{The model never writes a note.} Every shipped profile has \code{memory: off}; set it to \code{shadow} or \code{on}. The background writer has its own switch, \code{settings.memory.notes.writer.mode}. The log says why a due session was not written: \code{hold}, \code{gpu}, \code{nothing}, \code{moved} or \code{no session}.

\paragraph{Behind the proxy the writer never runs on a cloud model.} \code{auto} is off on a paid cloud model. Switch it on for that backend, give the model a price, and check the day's caps on the backend's page.

\paragraph{A note was refused.} Caps refuse, never cut; a body that is mostly a secret is refused; a rule about the user needs the user's own message in the turn. The Notes tab lists the refusals with their reasons.

\chapter{Questions and answers}\label{ch:faq}

\paragraph{Does speculation change the text?} Under greedy decoding the engine writes what plain decoding on the same weights writes, up to the first near tie of at most one ulp between the two best tokens, where the batched verify can round the other way. The paper measured 17 of 30 outputs identical to the end of the answer and 13 that diverged at such a near tie. Chapter~\ref{ch:exactness} states it exactly.

\paragraph{Does the conductor change the text?} It changes what the model reads (a hidden note, a thinking cap, a hidden close) and so what it writes, by design; the arithmetic is the same greedy decoding. \code{X-Tandem: raw} gives the release's raw path. Even there, two runs of one request can part at a near tie, because the length router sizes blocks from timings (Section~\ref{sec:two-runs}).

\paragraph{Does it serve several requests at once?} Yes. With \code{--parallel N} the engine decodes up to N requests in the same rounds: on 0.4.0 two at once gave 1.38$\times$ the total of one, four 1.42$\times$, in single measured runs (Section~\ref{sec:parallel-requests}). The shipped profile runs 8 lanes since 0.15.0 (SRV-707): on 0.15.0 the 27B ThinkingCap EXL3 4.0 pack decoded 1, 2, 4 and 8 requests of 512 tokens at 42.5, 73.5, 108.9 and 155.9~tok/s in total, greedy (Section~\ref{sec:changelog-0150}). The headline numbers in this manual are single-request numbers. With one sequence a short request still goes through the side lane, and the proxy serves other backends in parallel either way.

\paragraph{Will it run on another GPU?} Best effort. The installer turns off the skinny kernel that is built for the GB10 only, and StairCut's prices were measured on the Spark, so the router starts from the wrong staircase and corrects its level online. Nothing about another GPU's speed is measured.

\paragraph{Why NVFP4 weights but 16-bit arithmetic?} On this board the FP4 tensor-core path still has to move operands through shared memory, and a 4-bit activation path would need its own quantiser and gate to save at most the 8~\% of traffic the activations are. NVFP4 saves bytes, and bytes are time on a memory-bound decode (Chapter~\ref{ch:kernels}).

\paragraph{Why are my speeds below the paper's?} Check four things: the raw path (\code{X-Tandem: raw}), the lookup corpus (the paper's is not published), the store view (a store with the benchmark's text inflates it, a clean one does not), and the board's lock (another engine on the board halves the bandwidth).

\paragraph{Can I use the dashboard without the engine?} Yes. \code{pip install '.[host]'} and \code{tandem-app} run the dashboard and the proxy on any machine with Python 3.11, in front of vLLM, SGLang, llama.cpp, LM Studio, Ollama, OpenAI, DeepSeek or another TandemLLM (Chapter~\ref{ch:install-app}).

\paragraph{Does the proxy add latency?} 0.33~ms to the time to first token at the median against a fake backend on one machine, and a decode rate within 0.8~\% of the engine's own over the local network (\measured, the 0.4.0 series).

\paragraph{Are my prompts stored anywhere?} No log line, ledger row or metric holds a prompt or an answer. The suffix store holds token ids, never text. The conductor's transcripts keep a session's tool turns under the engine's home, mode 0600, secrets redacted, 30 days; memory notes are files you can read, edit and delete. Probe snapshots, off by default, hold activations, never text.

\paragraph{Is the dashboard safe to expose?} No. It is meant for a trusted local network. No admin password is set by default, so anyone who reaches the port can read and change the configuration and the memory notes. This release has no sign-in screen: with a password set, the dashboard's edits and its Memory screen answer 401. If other people can reach the machine, keep the port on loopback and put a reverse proxy with its own sign-in, or a VPN, in front (Section~\ref{sec:access}).

\paragraph{Where are the docs this book draws on?} In \code{docs/} of the repository, one page per part, under CC BY 4.0. This book adds the diagrams, the charts and the generated reference.

\input{chapters/roadmap}
\chapter{Changelog summary}\label{ch:changelog}

\code{CHANGELOG.md} in the repository lists every release with its measurements. The paragraphs on the Chat tab, the dashboard and the proxy are in TandemApp's manual. This chapter summarises it. Every release is held to the one before it on the raw path by the identity gate (the same logits and state, bit for bit, on a fixed scenario), except where an entry says otherwise, and each one's speed is checked on the benchmark row against the release before it, in the same hour. The gate fixes the arithmetic, not the tree shapes the length router picks from timings, so the served text can part at a near tie between two runs of either release (Section~\ref{sec:two-runs}).

\begin{table}[htbp]
\centering\small
\caption{Releases of TandemLLM.}\label{tab:releases}
\begin{tabularx}{\linewidth}{@{}L{.2\linewidth}L{.15\linewidth}Y@{}}
\toprule
\hd{release} & \hd{date} & \hd{what it brought} \\
\midrule
0.1.0-rc1 to rc4 & 2026-09-18 to 09-25 & the engine, the length router, release of the losing drafter, serving defaults for long work; rc4 is the first release measured in two store views \\
0.1.0-rc5 to rc7 & 09-26 to 09-27 & sampled requests on the draft tree, the in-process block A/B, tool calls read in three forms and streamed as argument deltas \\
0.1.0-rc8 to rc10 & 09-27 to 09-28 & the resident prefix, live activity on the dashboard, the FP8 and balanced profiles, the first seams for other models \\
\code{v0.2.0-staircut} & 09-29 & StairCut, on in the served profile: the paper's release \\
\code{v0.2.1} & 09-29 & the dashboard login off by default; no decode path changed \\
\code{v0.3.0} & 09-30 & image input; the text path unchanged \\
0.4.0 & 10-06 & the conductor, memory notes, the caches for long agent sessions, parallel requests (off by default), and the dashboard as a standalone proxy in front of other engines and cloud APIs; not published, 0.4.7 replaces it \\
\code{v0.4.7} & 10-08 & 0.4.0 plus the package library, Regenerate, long Open WebUI chats back from the NVMe, the compressed session spill and a daily NVMe write budget \\
\code{v0.4.8} & 10-08 & two Open WebUI chats beside each other (a side slot keeps its conversation), a chat whose retrieval block moved resumes at its end, the park keeps the end anchor, the router's notes stay live by their message \\
0.5.0 & 10-09 & weight profiles: NVFP4 or EXL3 at 2.5 to 6.0 bits a weight, one at a time, switched with \code{ops/use-profile.sh} \\
0.5.1 & 10-10 & two dashboard numbers corrected; four parallel requests on the reference box \\
0.6.0 & 10-10 & one shared KV pool that maps pages as a conversation grows; a request that does not fit waits instead of failing \\
0.6.1 & 10-10 & the small models give memory back when the pool runs short; the memory split counts committed pages \\
0.6.2 & 10-10 & decoding no longer stops while another slot prefills; the Perf tab's live feed fixed \\
0.7.0 & 10-11 & the setup guide at \texttt{\#/setup} with Wren, its assistant; reset levels; setup mode \\
0.7.1 & 10-11 & Wren drawn as the DGX Spark box; dashboard only \\
0.7.2 & 10-11 & Wren answers a question in words again; setup guide only \\
0.8.0 & 10-13 & a second model, Qwen3.8-Flash-Next, from turboderp's EXL3 packs; the Chat tab with chatd; the cache peek and prewarm \\
0.8.1 & 10-13 & a new box downloads every file a weight profile needs, with no token; the EXL3 weights public; the engine unchanged \\
0.8.2 & 10-13 & community weight packs; permission modes in the Chat tab; standalone EXL3 packs of the 27B family; a chat no longer waits minutes for slot 0 while a side slot is free; the model says it is TandemLLM when asked \\
0.8.3 & 10-13 & repository hygiene for the public release; no code change \\
0.9.0 & 10-14 & a third model, LFM2.5-24B-A2B: six EXL3 sizes and BF16, the default 3.5 by a stated rule, the lookup drafting without a draft head; Wren beside it \\
0.9.1 & 10-14 & web search and media playback in the Chat tab; the loop guard, the loop escape and memory notes on in the shipped conductor profiles; Setup Guide fixes; the 27B's EXL3 profiles fetch only the base tensors they read \\
0.9.2 & 10-15 & install fixes: the checksum step reads a download's globs; after \code{install.sh --setup} the 27B's EXL3 profiles take \code{base-plain} \\
0.10.0 & 10-15 & a fourth model, Nemotron 3.5 Bolt 74B-A9B: seven EXL3 sizes, NVIDIA's IQ2 GGUF and its LiLiCorr drafter; a 1~GiB limit for each Chat tab sandbox; attachments never inlined past the window \\
0.11.0 & 10-16 & a fifth model, Nemotron 3.5 Super VL 120B-A12B: EXL3 packs at 4.0 and 2.0, NVIDIA's NVFP4 release read as stored, its LiLiCorr drafter \\
0.11.1 & 10-16 & chatd's access to the sandbox's folders follows no symlink; StairCut price files for Bolt's EXL3 2.5 to 6.0; each weight pack's maker in the setup guide; the install through the setup guide \\
0.11.2 & 10-16 & chatd starts with the engine and the Chat tab says when it is down; install script and secrets fixes; sandbox containers scoped to their chatd; the preview opens the running app; Context panel figures; no horizontal overflow on a phone; a model picker with one row a model; Setup Guide steppers for chats and windows; Super VL window up to 1,048,576 \\
0.11.3 & 10-17 & Setup Guide hardening: \code{tandem doctor}, a check before the restart, Retry and Copy diagnostics, the restart card, removing weights; Super VL reads images; a cloud model's figures and cost; slot affinity \\
0.11.4 & 10-18 & Agents mode in the Chat tab: a main agent with workers, the Task card, workers in the Context panel; a worker's and its parent's cache rows kept; file texts out of a chat's requests; a memory refusal compacts; cloud runs take no run slot \\
0.12.0 & 10-19 & the \code{tools} profile for API clients that offer tools; structured outputs with tools; C5, one mutation per turn; the secret echo mask; the raw path's thinking in \code{reasoning_content} for plain API clients; tool-eval-bench results \\
0.12.1 & 10-20 & the first community pack on Super VL, styles01's EXL3 3.5 bpw, text only; a request with images on a side slot; redirects kept in the port preview; eleven outside pull requests; \code{SECURITY.md} and a threat model \\
0.12.2 & 10-20 & a Chat-tab run's tool results fit the room left and \code{read_file} reads by whole lines; the port preview streams; LFM2.5 adopts its 0.9.x to 0.11.1 conversations again \\
0.13.0 & 10-20 & Ornith-1.5-122B on the Qwen3.5 MoE trunk, from our EXL3 4.0 pack, drafting with its MTP layer; the conductor's third tool-eval-bench round; room for Ornith's caches; \code{chatd} fixes for long sessions \\
0.14.0 & 10-22 & Recursive's Flint from our EXL3 4.0 and 6.0 packs and the BF16 checkpoint, drafting with the lookup; a resumed session keeps its end anchor; a re-rendered tool call served as written; prompt attention over the KV heads as they are; a lost CUDA context fails \code{/health}; \code{tools/release.sh} \\
0.14.1 & 10-22 & a call's argument names matched to the offered tool's schema; a call to a tool not offered mapped to one that is; Flint answers directly with thinking off; Flint's e4m3 KV cache, off by default; model families behind adapters \\
0.14.2 & 10-23 & with thinking off, reasoning written into the answer goes to its field; Ornith answers up to its card's 131,072 tokens and decodes plain at long context; \code{cache_prompt: false} as a cold read; a lookup corpus for each tokenizer; no conductor notes on a model without a place for them; \code{tandem-dash serve} without the engine \\
0.15.0 & 10-27 & 8 lanes by default, every lane the full window; a lane short of memory gives way and goes on; one round runner with a packed draft and a joint verify pass, from CUDA graphs on the 27B; Flash-Next and Ornith in joint passes; a burst prefill; the 27B's faster GDN prefill; the card's sampling filters for a request that names none; a speed row on every weight card \\
0.15.1 & 10-31 & Flash-Next's joint pass in fewer launches: every lane's attention, PLE and GDN prep at once, streamed expert rows, a dense pass from CUDA graphs; a burst's prompts in one forward; coupled drafts for seeded sampled requests, a learned corrector for the 3.05 pack, the lookup's backoff; the runner's host work and collector in the GPU wait; a fresh server's first raw request no longer waits behind a generation; no idle capture below its state's context class; weight cards measured on \code{/v1/completions} \\
0.15.2 & 11-04 & a prefill admitted over the engine's peak charged only its new memory at the box's floor; idle graph captures pause at the peak; chatd's Stop ends a request that waits for memory, and the engine books it abandoned; a \code{tandem_memory} tool call taken as the memory block; the allocator's counters read without cyclic garbage; 7-bit EXL3 tensors on the dense kernels; the Setup Guide's simple mode; the quality gate, the benchmark report and one adapter a GPU \\
0.15.3 & 11-09 & Flash-Next's CUDA graphs on each slot's own stream, dropped together and capped; a request short of memory waits at the KV pool's floor and goes on; long chats side by side: 3-hour waits, a chat's priority, idle conversations on the NVMe before a reclaim, turns charged the rows they forward; a compaction as one more turn of the chat; Flash-Next's long prefill by a radix select and its drafts past 32k; Ada kernels, every GGUF quant type on NemotronH, strict HTTP framing \\
0.15.4 & 11-10 & a turn charged the rows it forwards, on slot 0 and for a regenerate or an edit too; a short reserve raised at the first forward; the pool's floor at the chunk a turn forwards; a quiet log during a wait; the log's Export button; the wait for memory, a chat's priority and the swap to the SSD in the docs \\
\bottomrule
\end{tabularx}
\end{table}

\section{0.15.4, a turn charged what it reads}\label{sec:changelog-0154}

0.15.4 fixes the charge that made short turns on long chats wait minutes for memory (SRV-801). Slot~0 had no share of its own conversation, so the admission priced a turn of the chat it held on its whole prompt: on a user's box a turn charged 212,938 rows (10.75~GiB) waited 148.6~s and then forwarded 449. Slot~0 now counts the share it holds, the gate prices the served prompt with the server's notes in it, a regenerate or an edit asks the slot's rows where its prefill will resume, and a next turn whose copy leaves out the last answer's reasoning is charged from the end anchor, so the last answer no longer counts twice. Before the first forward the server raises a request's own reserve when the rows it forwards need more, or the lane gives way before the forward allocates. The floor a prefill's pages commit under holds the chunk the turn forwards, not a whole one, and a retry that frees nothing no longer floods the log. On one DGX Spark every request line of a check on a synthetic 203,000-token chat had its charge equal to its forwarded rows (Section~\ref{sec:memory-wait}); a cold read still takes more host memory than its reserve, 0.37 to 0.44~GiB, cause not known (SRV-805).

The Dev view's log has an Export button that saves the whole log as one text file with a build header (SRV-803, TandemApp's manual). The manual gains sections on the wait for memory, a chat's memory priority, the lanes and a request held for memory, and the spill section calls it the swap of long chats to the SSD. Text that main no longer matched is fixed: a side slot keeps its conversation (Section~\ref{sec:parallel-requests}), the shipped profile runs 8 lanes, and the dashboard map shows all nine views. No file under \code{engine/} or \code{tools/} changed, so every family's fingerprint stays.

\section{0.15.3, memory that comes back}\label{sec:changelog-0153}

0.15.3 fixes a server that refused requests with \code{memory_full} after a day of use. Flash-Next's CUDA graphs kept one private allocator pool each on torch's shared capture stream, and nothing reset one: after 29 hours a server held 3,905 pools and 33.0~GB, 21.2~GB of it free and unreleasable. A slot now leases one capture stream for its life and drops all its graphs at once; the KV pool's reclaim drops idle sets after the caches and before any conversation's rows, and \code{QSE_GRAPH_POOL_GIB} (8) caps their growth over the load-time base (SRV-795).

A request short of memory waits at one floor, the KV pool's, and goes on (SRV-797). The admission checks the prompt's pages against the pool's floor while the prefill runs, a refused lane gives way alone and during its prefill and backs off, and a streamed request waits up to \code{QSE_MEM_WAIT_S} (3~h). A request not streamed asks for the same wait with \code{wait_for_memory: true} or the header \code{X-Tandem-Wait-For-Memory: 1}; without it, it gets a 503 with \code{Retry-After}.

Long chats share the box. A chat has a priority for memory (\#169), an idle conversation is written to the NVMe while its rows stay and before a reclaim takes its pages (SRV-779, SRV-E6), a chat's older copy after a compaction goes first and unwritten (\#170), and a turn is charged the rows it forwards (\#167). On one DGX Spark five chats of 221k to 225k tokens, with room for four, took a turn at once in 32 to 68~s each with the session spill on and 659 to 894~s with it off. A compaction's summary is one more turn of the chat's own request: 36 to 38~s at 212k (\#168).

Flash-Next selects its long-context prefill blocks with a radix select and fused scores (\#158): a cold 185k prompt took 198.4~s, 312.2~s on 0.15.2. It drafts past 32k again by priming the MTP head with the last 2,048 unread rows (\#159): 62.8 instead of 37.6~tok/s at 60,907 tokens of context. Ada (\code{sm_89}) gets its own Triton kernels (\#144), NemotronH's GGUF loader reads every ggml quant type (\#143), and the HTTP server and the relay gain socket timeouts and strict body framing (\#37, \#42). Flash-Next 3.05 writes the same text as 0.15.2. Every family's fingerprint moves once.

\section{0.15.2, the memory wait}\label{sec:changelog-0152}

0.15.2 fixes a long chat that waited 15 minutes for memory on a server that had run for an hour. A prefill the memory gate admits over the engine's peak takes the allocator's free default-pool blocks before any new memory, and MemAvailable already leaves those out, so the admission now charges it its workspace less those blocks (SRV-775). A lone chat of 31,053 rows was charged 3.9~GiB with 2.3~GiB of it free in the reserve and waited below its 10~GiB floor; its line now names the new part.

The idle window captures no graphs while the reserve plus the smallest prefill's workspace is over the peak (SPD-273): each capture keeps a private pool that nothing returns, about 20~MB a graph. A request's own class still captures at its first use.

chatd's Stop shuts a model call whose response has not started, so it ends a request that waits for memory within a second instead of after chatd's 15-minute read timeout, and the engine books that request abandoned with a \code{[req]} line (SRV-777).

With a chat's tools in view the model can write its memory note as a \code{tandem_memory} tool call, which no client offers. The server takes the call as the block it means: stored or refused, hidden, never sent to the client, and an answer of notes alone goes on after a hidden result. A rule about the user written in a tool turn is kept for that conversation only, at session scope.

The engine reads the allocator's counters off torch's nested dict, which leaves no reference cycle: with \code{QSE_RUNNER_GC=1} the collection at a round's end took 40 to 80~ms while the GPU waited (\#145). The native dense kernels take 7-bit EXL3 tensors (\#136), so the text of a pack with 7-bit tensors can change, and NemotronH's packed IQ2, Q8 and FP8 weights run a 33-row verify block on their row kernels (\#142). Every family's fingerprint moves once.

A pull request that changes the engine carries a quality gate report (\code{python -m tools.gate run}), each release gets a benchmark report page (\code{python -m tools.release_report}), and every per-GPU choice goes through one adapter a GPU in \code{tools/gpu}.

\section{0.15.1, Flash-Next's joint pass}\label{sec:changelog-0151}

0.15.1 cuts the launches of Flash-Next's joint pass. Its profiles run every lane's attention, PLE and GDN prep in one launch a layer, stream the grouped experts' rows beside their tiles, so a gate and up CTA needs about 14~KB of shared memory instead of 42~KB, and run a dense joint pass from CUDA graphs. Each row keeps its arithmetic and the tokens are the same; every family's fingerprint moves once (ENG-411, then SPD-249). Seeded texts equal 0.15.0's; greedy prose can part from 0.15.0's at a near tie, a word choice 157 to 646 characters in, and the cause is not isolated. The profiles also prefill the prompts that arrive together in one forward (\code{QWEN38_FN_JOINT_PREFILL}), each prompt's rows its own prefill's bits: more tokens per second in total, and a later first token for the long prompts of a burst, since every first token waits for the one forward.

Flash-Next drafts a seeded sampled request under the target's own keyed noise, so a draft lands on the target's draw whenever the MTP head and the target put the same token on top of that noise. A learned corrector re-ranks the head's top 8 tokens, so more drafts land on the target's draw. \code{ops/use-profile.sh flashnext-3.05} and the Setup Guide download its file (34.9~MB, \code{0xBakeer/TandemLLM-Qwen3.8-Flash-Next-EXL3-3.05-corrector}) at a pinned commit and check its sha256. Without it, or with a file that fails the load check, the drafts stay plain coupled and \code{/health} says why. The 3.05 profile keeps the drafters' cut on its price file's prices (\code{QWEN38_FN_DRAFT_FREEZE}), so a request drafts the same rounds whatever ran before it. On rounds that other lanes share, an MTP lane's lookup backs off after a round it gave nothing. The round runner pushes tokens, writes its report line and rebuilds \code{/health} while the GPU runs the next joint pass, and keeps the collector out of its rounds. A seeded text is the same with each of these on or off (SPD-249).

A fresh server's first raw request no longer waits for slot 0 behind a whole generation: its prompt is counted on the host, so no copy breaks the idle window's graph capture. An idle capture takes no context class below its state's, and CUDA work off the engine's thread waits for an open capture. The engine skips a lookup corpus of another tokenizer at load instead of stopping before it binds its port (\#105), and NemotronH runs every projection and the head of a 33-row verify block on the row kernels (SPD-255).

A Flash-Next block just under a context-class edge verifies and commits eagerly, and a verify, commit or draft graph whose capture fails keeps no key, puts the state back and runs its call eagerly (SPD-261). Each resident prefix holds its end anchor's buffers from the load, so a chat's follow-up resumes at its end at the memory ceiling; the Flash-Next profile starts the engine with expandable allocator segments, and graph captures run on plain segments (SPD-257). In the chat path, in beta, a new chat the router cannot classify in time is a plain chat at the profile's default effort, not \code{code_edit} at xhigh (SRV-763), new chats of one user share the notes head (SRV-765), a shadow library question waits until the server is idle (VIS-119), and with joint prefill on the router's class question no longer waits in a burst for its own turn (SRV-769).

Every speed number of this release is measured on \code{/v1/completions}, the engine: on one DGX Spark Flash-Next 3.05 decodes 8 requests at once at 146.8~tok/s in total, and at 147.2~tok/s with the chat template's thinking on and a new prompt in every round (156.7~tok/s when the same 8 prompts repeat; \measured). \code{/v1/chat/completions} works and is in beta; it is not measured.

\section{0.15.0, eight lanes}\label{sec:changelog-0150}

0.15.0 runs 8 lanes by default on the KV pool, every lane with the window the setup names, and the Setup Guide counts how many lanes can all hold it at once. A lane whose next pages the pool refuses while others run gives its pages back, waits at the head of the line and goes on in the same answer; alone it stops with \code{memory_full}. Ornith builds side slots, and the 27B, Flash-Next, Ornith and Flint cards default to 8 lanes (\measured).

One thread runs every lane's round: the proposals, one packed draft, one joint verify pass and the walks. The 27B's joint pass takes 64 rows on EXL3 weights and runs from CUDA graphs once its shape has come four times, every lane's rows bit for bit its own verify; each graph owner captures on a stream of its own, so no graph's reset frees the cuBLAS workspace another graph reads. Flash-Next and Ornith run their row-wise parts once over every lane's rows and pack their MTP drafts. NemotronH packs its lanes' LiLiCorr drafts and takes the joint plain step where drafts do not pay: Nemotron 3.5 Bolt at 8 requests at once runs at 92.2~tok/s in total on its card (\measured).

Requests that arrive on an idle engine prefill before any of them decodes, and on the 27B's EXL3 packs their prompts of up to 512 tokens run in one forward, each row its own prefill's bits. The 27B's GDN prefill runs at a 64-row chunk; with the prep kernel's scan at 16 by 4 this moves the summation order once, held by a gate (top-1 64 or 65 of 65), and Flash-Next's tensor-core mixers move theirs (Chapter~\ref{ch:exactness}). A sampled request that names no filter takes the card's top\_k and top\_p, and one without a seed draws its own and takes the keyed walk. Every family's fingerprint moves once. Every weight card carries a speed row of this build: decode alone, the first token of a cold 2K, 8K and 32K prompt, and 1, 2, 4 and 8 requests at once (\measured). A raw request that offers no tools gets its answer streamed as text, and Flash-Next captures the context classes ahead of a conversation and keeps a verify block's records across a draft step.

On one DGX Spark the 27B ThinkingCap EXL3 4.0 pack decodes 1, 2, 4 and 8 requests of 512 tokens at 42.5, 73.5, 108.9 and 155.9~tok/s in total greedy and 36.5, 62.9, 99.2 and 146.2 sampled at temperature 1.0; one prompt of 2,048, 8,192 and 32,768 tokens reads in 1.90, 6.76 and 29.4~s (\measured, 2026-10-26, 8 lanes, suffix store off).

\section{0.14.2, an Ornith test report}\label{sec:changelog-0142}

0.14.2 fixes what an outside test of Ornith found. With thinking off, the text a model writes before the first \code{</think>} that stands alone on its line, outside a code block, goes to \code{reasoning_content}, and usage counts it as reasoning; a tag in code or in a line of text stays. Of the report's four long contest answers, two now carry their working-out in the reasoning field and only the answer in \code{content} (\measured). Ornith's profiles answer up to the 131,072 tokens its card evaluates with when a request sets no limit; every other profile keeps 32,768.

Ornith's price file prices the MTP head's backlog, the prompt rows the head reads before a turn's first draft (Chapter~\ref{ch:ornith}). A fresh server's first long turn drafted and paid it: at 112,068 prompt tokens greedy decode ran at 18.04~tok/s against 25.16 plain, and now that turn takes the plain step at 25.11; at 220,038 it runs at 18.85, where drafting ran at 10.95 (\measured, one run each). The Setup Guide stops Ornith at one parallel chat, as its engine builds no side slots yet, and \code{/health} says so under \code{--parallel N}.

\code{cache_prompt: false} on a request is a cold read: the whole prompt is read, no prompt cache answers it or keeps it, and the conversation in place is where it was for the next request (Chapter~\ref{ch:caches}). The lookup corpus is built for the served model's tokenizer after its download, one store a tokenizer under \code{CORPUS/tokenizers/}; its effect on speed is not measured yet. A model whose prompt has no place for the conductor's notes, Flint, gets none, and its turns record memory off with the reason instead of an error (Chapter~\ref{ch:memory}). \code{tandem-dash serve} starts without the engine installed (TandemApp's manual). The Setup Guide sets the thinking effort for every model, and the cards name each number's effort. Several requests that arrive together as a fresh server's first traffic no longer answer 500 for a tool-call format not yet registered, and the page check closes its Chrome in about 0.1~s instead of up to 5~s. No fingerprint moves: every model writes the same tokens as in 0.14.1 and keeps its saved conversations.

\section{0.14.1, adapters}\label{sec:changelog-0141}

0.14.1 matches a tool call's argument names to the offered tool's schema before the call leaves the server, on both paths and in every tool-call form, so Flint's \code{read} with \code{path} reaches opencode as \code{filePath}. A call to a tool the request did not offer goes to the offered tool of its kind; a search or a listing with no such tool becomes a quoted \code{grep} or \code{ls} for the client's bash tool, and a write never becomes a shell command. Through opencode's checks at Flint 4.0's card sampling, 8 runs an arm, the checks passed went from 28 and 29 of 32 to 32 of 32 (\measured).

With thinking off, Flint's prompt gets one line before the conversation that asks it not to show its reasoning and to answer or call a tool: greedy, it opened its answer with reasoning on 0 of 10 hard asks instead of 7, and the median answer took 153 tokens instead of 582. \code{FLINT_KV_FP8=1} keeps Flint's full layers' keys and values in e4m3; it changes the output, so it is off by default. Turn it on in the Setup Guide (the Tools step, \emph{Long-context cache (FP8)} under Features) or with the line \code{QWEN38_FLINT_KV_FP8=1} at the end of \code{~/TandemLLM/run/ops/serve.env} and \code{~/TandemLLM/bin/tandem restart}; Chapter~\ref{ch:flint} gives its cost and how to check it took.

Each model family now registers itself in its own package and declares how it is served: its tool-call form, its reasoning, where system messages may stand, the files others build on and its kernel launch options. Guards hold v0.14.0's fingerprints, and shared code names no model (Chapter~\ref{ch:adding-model}). The other models write the same text as in 0.14.0 and keep their saved conversations; Flint's fingerprint moves once, so its saved prompt-cache entries from 0.14.0 are rebuilt once.

\section{0.14.0, Flint}\label{sec:changelog-0140}

0.14.0 adds Recursive's Flint from our EXL3 packs at 4.0 and 6.0 bits a weight and from the BF16 checkpoint (Chapter~\ref{ch:flint}), on a new family, \code{engine/models/flint}: sliding-window layers on a ring, a learned sink per head, full layers without positional encoding, and a sigmoid router with a shared expert. Every profile drafts with the lookup, which StairCut cuts from a price file with a plain price for each context class. On a DGX Spark, one request, 512 greedy tokens, thinking off, one run each, the 4.0 pack decodes at 128.2 to 136.6~tok/s plain and up to 480.0 with the lookup, every drafted text the plain one. Flint writes tool calls in its own arg-key form: the conductor's checks read them in that form, a string argument streams as it is written, and the note after a call that does not parse shows the model's own form with an offered tool, never a placeholder.

The other models write the same text as in 0.13.0 and keep their saved conversations: Flint's edits to shared files take the old paths unless Flint asks for the new ones, and \code{SAME_ROWS} names the new bytes as the old. A turn resumed from a session snapshot keeps its end anchor, end alignment serves a tool call the template wrote back differently as the model wrote it, and every prompt attention reads the KV heads as they are, where LFM2.5's prompt chunks, Ornith's and NemotronH's MTP head had repeated them per query head. After a lost CUDA context \code{/health} fails and the watchdog restarts the engine. In the Chat tab, screenshots fold in blocks, plan mode keeps the chat's effort word, and the run monitor steers a run that circles the same files, then stops it. \code{tools/release.sh} makes a release with one command.

\section{0.13.0, Ornith-1.5-122B}\label{sec:changelog-0130}

0.13.0 adds Ornith AI's Ornith-1.5-122B from our EXL3 pack at 4.0 bits a weight (Chapter~\ref{ch:ornith}). It is a Qwen3.5 mixture-of-experts model that reads images, on a new trunk, \code{engine/models/qwen3_5_moe}. The profile drafts with the model's MTP layer and the lookup: on a DGX Spark, one request, 512 greedy tokens, thinking off, one run each, 37.2 to 38.4~tok/s plain and 85.3 to 124.8 drafted, every drafted text the plain one. The engine's top token agrees with exllamav3's own forward on the pack at 498 of 512 positions.

The other models write the same text as in 0.12.2. The router kernel's new bytes moved the fingerprints of Flash-Next, LFM2.5 and NemotronH; at 512 experts it launches as before, \code{SAME_ROWS} names the new bytes as the old, and their first start adopts 0.12.2's saved conversations once. The release also carries the response cache's key on the resolved effort level, the tool-eval-bench batch tools, the deliverable keep-out in shadow, \code{grep} on one file's path, \code{read_file} past 4~MiB and a grammar test that passes every run.

Ornith's caches have room: the profile's limits are 86 and 90~GiB with a 2~GiB pin, where the first ceiling of 74 refused every end anchor and store snapshot. A refused end alignment serves the part it aligned instead of reading the turn again. In the Chat tab, writes and long commands stay whole in later requests until a fold point, plan mode keeps the request's head, and a run stopped at its call cap leaves the chat idle. Every test file now runs its tests as a script, and a file that runs none fails.

\section{0.12.2, large reads and a streaming preview}\label{sec:changelog-0122}

0.12.2 fixes two Chat-tab bugs that users reported, issues \#18 and \#86, and the adoption of LFM2.5's older saved conversations. A run that reads a large file now ends with an answer instead of a request over the window: one step's tool results share the room the model has left, \code{read_file} stops at whole lines and names the offset to read on from, and near the window chatd counts the request on our engine before it sends it (TandemApp's manual). The port preview passes an app's Server-Sent Events and long-polls on as they arrive, and a browser that leaves ends the relay in the sandbox.

No model, card or profile changed, and every profile writes the same text as in 0.12.1. \code{read_file}'s schema changes, so each chat's first request after the update reads its prompt again. LFM2.5 adopts the saved conversations of 0.9.x to 0.11.1 again; a box that ran 0.11.2 to 0.12.1 already lost them.

\section{0.12.1, a community pack on Super VL}\label{sec:changelog-0121}

0.12.1 serves the first community pack on Nemotron 3.5 Super VL, styles01's EXL3 pack at 3.5 bits a weight (pull request \#63, Section~\ref{sec:community-supervl}). It is served with our 4.0 profile's limits, LiLiCorr and the lookup, text only, since the pack holds no image tower. NemotronH now reads a stored head when the config says the head is tied, as this pack's remapped config does. On a DGX Spark, one request, 512 greedy tokens, thinking off, one run each: 25.9~tok/s plain, and 33.2 prose, 53.9 code, 37.4 chat with LiLiCorr and the lookup; the top token agrees with exllamav3's own forward on the pack at 510 of 512 positions.

A request with images now takes a side slot (Section~\ref{sec:parallel-requests}); \code{QWEN38_SIDE_IMAGES=off} keeps it on slot 0. The Chat tab's port preview keeps an app's redirects inside the preview. Pull requests \#78 to \#85 by joeynyc and \#34 and \#43 by AgentSparko fix chatd's attachments and port proxy, the Setup Guide's remove and its timers, the GGUF reader's size checks and a long thin image's patch grid. Every profile writes the same text as in 0.12.0. The NemotronH profiles' fingerprint changes, so the saved conversations of Bolt, Lightning and Super VL from 0.12.0 read their prompt again; the other models' carry over.

\section{0.12.0, tool clients}\label{sec:changelog-0120}

0.12.0 serves API clients that offer tools through a profile of their own. An OpenAI SDK, curl or an unnamed client that sends tools and names no profile gets \code{tools} (Section~\ref{sec:profiles}): the prompt as the client sent it, token for token the raw path's, with the loop guards and the tool-call checks acting. A request with tools can ask for a JSON answer, and the model's first answer token picks a tool call or the document. An answer that makes two or more changes outside the conversation in one turn is written again one call at a time (C5), and a key the model read in a tool result is masked where the answer repeats it. On the raw path a plain API client gets the thinking in \code{reasoning_content}.

On tool-eval-bench's 69 scenarios, on one DGX Spark with Qwen3.8-27B in the EXL3 6.0 pack, the raw path scored 92 with thinking off and 95 to 97 with it on, and the \code{tools} profile 95. Two runs of the same prompt scored 95 and 97, so rows within 2 points are the same. Profiles that add packages took twice the time and scored no higher. In the bench's hard mode, 28 scenarios, the conductor scored 91 against 84 for raw, one run a side. \code{docs/tool-eval-bench.md} has every row and its command. No change touches the decoding of a request without tools and without a JSON format; the served text against 0.11.4 was not measured.

\section{0.11.4, Agents mode}\label{sec:changelog-0114}

0.11.4 adds Agents mode to the Chat tab (TandemApp's manual). With the chat setting \code{agents} on, the main agent plans and reads, and gives every file write and every multi-step build or fix to a worker through the \code{task} tool. A worker is a full run on a hidden conversation of its own that sees the task, its brief and the user's request, and nothing else of the chat; its report comes back with chatd's check of what it changed and ran. A Task card shows each worker's run, report and transcript, and the Context panel shows each running worker apart, with its live prefill and decode figures. Stop, a queued message and Continue reach the worker. Agents mode is off by default.

The engine keeps a worker's rows and its parent's while they take turns. A request with \code{X-Tandem-Worker} or \code{X-Tandem-Parent-Session} is never a guest; for 600~s after a worker's last request its parent's rows go to the NVMe rather than away, a returning parent tries slot~0 first, and the park and the NVMe spill keep a conversation they used to drop. A request with images resumes at the end anchor (\code{QWEN38_MM_END}), and the live row has a prefill rate from its first mark. These restores copy rows the engine computed before and compute none; the tests hold every resume's logits to a cold prefill's within $10^{-4}$. The identity gate was not run against 0.11.3.

In a chat's own requests, a written file's text goes to the model as a stub from the next request on, a long read stays whole until a compaction, and a compaction cuts to under half the limit that set it off. A memory refusal compacts the chat and the run goes on, and a run on a cloud model takes no run slot. To upgrade, run the installer again, then \code{tandem restart} and \code{tandem chat restart}; chatd's store moves to v8 in place.

\section{0.11.3, the Setup Guide hardened}\label{sec:changelog-0113}

0.11.3 makes a Setup Guide install safe to retry. \code{tandem doctor} checks an install and repairs what is safe (file modes against \code{ops/executables.txt}, the \code{run/ops} copies, stale pid files, the secrets file's mode); it runs before every \code{tandem start} and in the install's check step, and as \code{tandem doctor [--check] [--json]}. A step before the restart checks the scripts, \code{serve.env}, the weights, the memory and the port, and a failure there stops nothing. The scripts call each other through \code{bash}, never through the execute bit. A failed step says what broke in plain words, with Retry from that step and Copy diagnostics. A restart is one card with the job's step and the load's phases, served by the tree's own status responder while the engine is away; a finished install has its own screen with Start chatting. The weights step lists the weights on the box and removes a model's weights, never the running profile's or a shared folder. \code{tools/release\_gate\_setup.sh} installs a candidate twice in scratch folders beside the live engine and drives the whole flow in a phone-sized browser.

Super VL reads images through its C-RADIO tower. A cloud model's turns in the Chat tab have their speed, the first token's wait, the provider's cache and the cost, and Performance and Usage keep the engine's requests apart from relayed ones. A conversation stays on the side slot that holds its rows, the narrator is off by default, the decider's residency claim is 3.0~GiB, and a chat sandbox has 4~GiB of memory. Requests on the raw path write the same text as in 0.11.2; under the conductor a turn can get packages it did not get before, and its text changes with them. To upgrade, run the installer again, then \code{tandem restart} and \code{tandem chat restart}; the conductor's migrations run at the engine's first start.

\section{0.11.2, chatd with the engine}\label{sec:changelog-0112}

0.11.2 starts chatd, the Chat tab's service, with the engine (\code{tandem start}, the systemd unit, the Setup Guide) and stops it with \code{tandem stop}; the Chat tab names chatd when it is down instead of switching to the in-browser mock. The installer copies \code{roles.sh}, \code{watchdog.sh} and \code{hold.sh}, \code{bin/tandem} loads its config on an install without \code{ops/active-profile}, and the Setup Guide's reset keeps the secrets file the engine reads. A second chatd on the same Docker no longer adopts another's sandbox containers, and a sandbox container removed outside chatd starts again at the next tool call. The preview pane opens the running app, the Context panel shows a live run's figures, and the Chat tab has no horizontal overflow on a phone. The model picker lists models, one row each, and DeepSeek V4 Flash and V4 Pro through the proxy can be picked. The Setup Guide sets the number of chats and both windows with steppers, and Nemotron 3.5 Super VL takes a window up to 1,048,576 tokens, checked to 249,983 prompt tokens. Update every box that runs the Chat tab: run the installer again, then \code{tandem chat restart}. Every model writes the same text as in 0.11.1.

\section{0.11.1, the Chat tab's sandbox folders}\label{sec:changelog-0111}

0.11.1 is a security fix: chatd reads and writes the chat's workspace and the sandbox's home without following a symlink, for uploads, the file tools, snapshots and the Open WebUI runner's files. Update every box that runs the Chat tab, and restart chatd (\code{ops/chatd.sh stop}, then \code{start}). Bolt's EXL3 packs 2.5 to 6.0 get their StairCut price files, so every pack drafts at least as fast as it decodes plain (Table~\ref{tab:nh-profiles}). Two contributions: a pinned asset downloads at its revision, and NemotronH's fingerprint covers its GGUF reader and file, so Bolt's and Super VL's saved conversations from 0.10.0 and 0.11.0 read their prompt again once. The setup guide names each weight pack's maker, and the install goes through the setup guide. Every model writes the same text as in 0.11.0.

\section{0.11.0, Nemotron 3.5 Super VL}\label{sec:changelog-0110}

0.11.0 adds a fifth model, NVIDIA's Nemotron 3.5 Super VL 120B-A12B (Section~\ref{sec:nh-supervl}). The other models' text and speed do not change, and their saved conversations carry over.

\paragraph{Super VL.} Our EXL3 packs at 4.0 and 2.0 bits a weight and NVIDIA's NVFP4 release, read as stored (NVFP4 experts, FP8 Mamba2 and attention projections and head), each with NVIDIA's LiLiCorr drafter for Super VL beside the lookup. The 4.0 pack is the default: KL 0.049 against BF16 (NVIDIA's NVFP4 0.041), 31.1~tok/s for one request (NVFP4 19.6), and with LiLiCorr and the lookup 53.6~tok/s on code and 101.0 on edit asks (the mean of 3 runs). Text only: the vision tower is not served yet.

\section{0.10.0, Nemotron 3.5 Bolt}\label{sec:changelog-0100}

0.10.0 adds a fourth model, NVIDIA's Nemotron 3.5 Bolt 74B-A9B (Chapter~\ref{ch:nemotron}). The other models' text and speed do not change, and their saved conversations carry over.

\paragraph{Bolt.} The NemotronH engine (Mamba2, attention without positional encoding, relu$^2$ experts behind a latent projection), our EXL3 packs at 2.0 to 6.0 bits a weight, NVIDIA's IQ2 GGUF read as stored, and NVIDIA's LiLiCorr drafter beside the lookup. The 2.0 pack is the default: KL 0.170 against BF16 (the GGUF 0.259), 43.3~tok/s for one request, and with LiLiCorr and the lookup 65.7~tok/s on prose and 102.8 on code (the mean of 3 runs). A verify block's EXL3 projections decode each trellis tile once for 16 rows, every output byte the one-row call's. Text only: the image tower is not served yet.

\paragraph{Saved conversations.} Fingerprint v2 leaves out modules under \code{engine/} that only another model imports, so the 27B's code fingerprints stay those of 0.9.x; the first Flash-Next or LFM2.5 start adopts the files 0.9.x wrote, once.

\section{0.9.2, install fixes}\label{sec:changelog-092}

0.9.2 fixes two install bugs on a new DGX Spark. The engine, the conductor and the dashboard do not change.

\paragraph{The checksum step.} \code{fetch_weights.py} compared a download's file globs to file names literally. A model whose list is all globs, Qwen3.5-0.8B for the shadow reader or Qwen3.5-2B for the narrator, checked no file and failed with \enquote{0 files checked: FAILED}, and Laya's files under \code{multilingual/} and \code{typed-decisions/} went unchecked. The globs now select the files.

\paragraph{The EXL3 base after a setup install.} The installer wrote \code{QWEN38_MODEL}, the Hugging Face cache's base folder, after the weight profile, so the 27B's EXL3 profiles kept that empty folder instead of \code{base-plain} and \code{use-profile.sh} refused to switch. The installer's own keys now come before the profile, an EXL3 profile takes \code{base-plain} when the folder holds no snapshot, and a standalone community pack reads its own folder. An install with \code{HF_HOME} writes the cache to \code{serve.env} as \code{HF_HUB_CACHE}. To upgrade, run the 0.9.2 installer again: it rewrites \code{serve.env} and the scripts beside it.

\section{0.9.1, web search, media and the loop guard}\label{sec:changelog-091}

0.9.1 changes the Chat tab, the conductor's shipped profiles and the Setup Guide. The decode path does not change: under \code{engine/} only the small models' slot changed, so that a failed load is tried again.

\paragraph{The loop guard.} The same-call guard (C3) is on in the shipped \code{coding}, \code{chat-app}, \code{hermes}, \code{overnight} and \code{chat} profiles. It counts only the current user turn, treats a tool result that took the client 10~s or more as a wait, and holds calls only on a request whose turn already ends with the repeat. After a repeated or near-repeated call the session samples with the model card's settings (the checkpoint's \code{generation_config.json}, a fresh seed) until it makes progress, then decodes greedy again; a client's own sampling wins. Output changes on those turns only. \code{--reasoning-effort none} and \code{off} turn thinking off for every request that says nothing about thinking.

\paragraph{Memory notes.} The shipped \code{chat}, \code{coding}, \code{hermes} and \code{chat-app} profiles have \code{memory: on}. A one-time migration adds the switch to each stored profile without it; an explicit value stays. The memory card tells the model the notes are its memory between chats, a profile with memory off tells it so, and the Context panel shows a turn's notes.

\paragraph{Fixes.} The Setup Guide's install page keeps a finished job final and its whole log after a reload. The change list names the keys a model switch drops from \code{run/setup.env}. The fetch step downloads the models of the tools the selection turns on, and a failed shadow reader load is tried again after 300~s in an idle window. The 27B's EXL3 profiles read their base from \code{base-plain}, fetched by HTTP range requests: 3.54~GB instead of 30.41, so a fresh \code{exl3-6.0} downloads 36.48~GB, not 63.34. The dashboard's service worker shows the cached shell when the reverse proxy answers 502, 503 or 504, and reloads once \code{/health} answers.

\section{0.9.0, LFM2.5-24B-A2B}\label{sec:changelog-090}

0.9.0 adds a third model, LiquidAI's LFM2.5-24B-A2B, the first here that is not a Qwen (Chapter~\ref{ch:lfm25}). The Qwen models' text and speed do not change: the shared EXL3 kernels learned half-integer rates, and \code{SAME_ROWS} names their new bytes as the old ones, so the first Qwen start adopts its saved conversations once.

\paragraph{Profiles.} Liquid's BF16 checkpoint and six EXL3 sizes of our own, 2.5 to 6.0 bits a weight, each at the model's whole 131,072-token context with three slots. The decode step runs at 188.2~tok/s on 2.5, 153.5 on 3.5, 106.9 on 6.0 and 44.8 in BF16 (\measured, 2026-10-14), with each pack's own dense splits. The default is \code{lfm25-exl3-3.5}: the fastest profile in chat as it serves whose chat top-1 agreement with BF16 is at least 85~\% (143.7~tok/s at 86.7~\%; 3.0 is faster at 83.6~\%). An exllamav3 \code{-hq} 4.0 was slower than 4.0 and worse than 4.5, and is not offered.

\paragraph{The lookup without a draft head.} A round with no lookup match is the plain step, priced as one; the lookup's confidence carries its match length; every profile has its own StairCut price file. With the store off, edit and copy asks on 4.0 ran at 291.1 and 326.4~tok/s against 141.3 and 141.9 plain, every text equal.

\paragraph{The conductor.} On for LFM2.5 by default, with Laya and the narrator, through the model's fit to the router: Laya alone answers the class question, a turn nobody classified gets nothing added, and a class sets the level and budget with only its core skill in the note. On 2.5, four short chat asks thought 2,234 reasoning tokens against 2,155 raw and 9,098 under the router as written; seven Chat-tab tasks run twice succeeded 14 of 14, as raw did, against 12 of 14 (Section~\ref{sec:lfm25-conductor}).

\paragraph{Fixes.} A pack stored as one \code{model.safetensors} counts as complete in \code{ops/use-profile.sh} and the Setup Guide. Reading a slot's \code{kv.length} commits nothing. Attention's split buffers are kept per attention class. A suffix store of another tokenizer turns the store off instead of stopping the start. The dashboard's disk tier default reads 16~GB, as the server's. Wren loads its own Qwen3.5-2B beside LFM2.5 (20~s, 3.62~GiB).

\section{0.8.3, repository hygiene}\label{sec:changelog-083}

Repository hygiene for the public release: no private hostnames, home paths or personal notes in the files; no code change. The engine's modules are the bytes of 0.8.2, and the text does not change. Comments, docs and tests no longer name a person who gave an instruction, and tests use \code{192.168.1.x} for a LAN address and \code{tandem.example.lan} for a host behind a proxy. \code{install.sh} installs v0.8.3 by default.

\section{0.8.2, community weights and permission modes}\label{sec:changelog-082}

0.8.2 adds community weight packs, permission modes in the Chat tab and standalone EXL3 packs of the 27B family, and fixes two bugs found on v0.8.0 in Open WebUI on a phone. The engine's modules are the bytes of 0.8.1.

\paragraph{Community weights.} One JSON file per pack under \code{models/<model>/community/}, checked by the registry against \code{models/schema/community.schema.json} and gathered in \code{models/community-index.json}. The setup guide lists the packs after our profiles as \code{community/<id>}, with search and five sort orders. It fetches the index from the public repository when its weights step opens or on Refresh, never on a timer, and falls back to its last copy, then to the one in the release; a 404 shows as the note \enquote{no published list yet}. \code{ops/use-profile.sh community/<id>} downloads a pack at its pinned commit, checks the hub's sha256 and runs it with our drafters. An issue form and a check on pull requests come with it. The first entry is ThinkingCap Qwen3.8-27B Kearuga by 0xWhiteMage: EXL3 at 4.0 bits a weight, standalone, 39.6~tok/s prose, 95.9 code and 51.3 chat on 0.8.2 (\measured). Chapter~\ref{ch:community}.

\paragraph{Standalone EXL3 packs.} A complete BF16-origin \code{mul1} EXL3 pack of the 27B family loads as \code{--model}, with no separate base checkpoint (pull request \#1 by 0xWhiteMage). The path is text only, and speculative parallel serving still needs our drafters (Section~\ref{sec:exl3-standalone}). The 27B's own profiles take the path they took before. \code{engine/loader.py}'s new bytes count as the old ones in the disk tier's fingerprint (\code{SAME_ROWS} in \code{engine/methods/manifest.py}), so the first 27B start adopts v0.8.0's saved files once, and the 27B's code fingerprint stays 0.7.2's.

\paragraph{The slot-0 wait.} A new Open WebUI chat waited 437~s for slot 0 behind another chat's 496~s generation while both side slots were free. It shared only Open WebUI's system prompt and tool list with an entry on the NVMe, and the wait for slot 0 had no limit. Now the spill wait needs the request's own conversation, and every wait for slot 0 ends after 3~s (\code{QWEN38_SIDE_SLOT0_WAIT_S}) when a side slot the request fits in is free. \code{/health} \code{parallel.max_running} is the limit with its reason, not the high-water mark, and the Lanes view says why a request is queued. On the reference box a turn that slot 0's chat had held moved to slot 2 after its 3.0~s wait, while another chat ran 80~s on slot 0 (\measured).

\paragraph{Identity.} Asked who it was, Flash-Next named another assistant and its maker: the reference box served it with the conductor off, so the identity note never went in. The note now says the model is TandemLLM running the served model (Qwen3.8-Flash-Next or Qwen3.8-27B, from the pack's model card) on the user's own machine, for every client the conductor serves, and the reference box runs Flash-Next with the conductor on. With Open WebUI's own system prompt and 49 tools, 20 of 20 answers said so, against 0 of 20 on v0.8.0 (\measured).

\paragraph{Smaller fixes.} The Flash-Next profiles load Laya, so the memory features work there: the remember, forget and recall intent took 20.6~ms on 3.05. A deleted chat's sandbox stops at once; a restore brings the chat back, and the purge after 7 days or on \code{?forever=1} removes the container and the chat's folders. A 404 from the community sync is a note, not an error.

\section{0.8.1, a new box fetches every file}

A new DGX Spark runs every 27B EXL3 profile and Flash-Next from public files, with no Hugging Face token (OPS-95). \code{ops/use-profile.sh} downloads what a 27B profile reads besides its weights before it switches: the K3 drafters and the K4 draft head, the BF16 drafters into \code{train/}, and the FP8 base checkpoint into the Hugging Face cache. It pins each download to the hub's commit, checks each weight file's sha256, checks the disk first, and continues an interrupted run; \code{--dry-run} prints every URL and size. The EXL3 profiles read the BF16 drafters from \code{train/} on every box. A switch from setup mode removes the marker first and puts it back if the profile does not start. \code{install.sh} installs v0.8.1 by default, and \code{--repo} names the repository it clones. The engine's modules are the bytes of 0.8.0. Section~\ref{sec:new-box} has the steps for a new box.

\paragraph{Public weights.} The six EXL3 sizes of the 27B and their drafters are public on Hugging Face: \code{0xBakeer/TandemLLM-Qwen3.8-27B-EXL3-<bpw>} for 2.5, 3.0, 3.5, 4.0, 4.5 and 6.0 bits a weight, and \code{0xBakeer/TandemLLM-Qwen3.8-27B-EXL3-drafters}, in one collection. Each card has the size's KL against BF16 and a speed table; KL runs from 0.1112 at 2.5 bits to 0.0018 at 6.0 (Table~\ref{tab:exl3-27b}). The calibration repository stays private, and nothing that runs reads it.


\section{0.8.0, Qwen3.8-Flash-Next and the Chat tab}\label{sec:changelog-080}

A second model, Qwen3.8-Flash-Next, served from turboderp's EXL3 packs through the same server, caches, ops scripts and setup guide as the 27B (Chapter~\ref{ch:flashnext}). The 27B's text does not change: its modules, flags and disk-tier fingerprint are as in 0.7.2, and every switch added for Flash-Next is off for it. On Flash-Next, greedy and seeded answers with speculation equal plain decoding token for token.

\paragraph{Profiles and windows.} Five weight profiles, \code{flashnext-2.05} to \code{flashnext-6.05}, and a second tile in the setup guide. Each pack has its largest window in the registry: 262,144 tokens for 2.05 and 3.05, 208,896 for 4.05 with the vision tower on; 5.05 does not load yet (its experts use 5 and 7 bits a weight), and 6.05 fits no window on one DGX Spark. 3.05 reads its n-gram table cold. A start with a longer \code{--max-len} is refused before the load.

\paragraph{Speculation.} The pack's own MTP layer drafts, cut by the shared StairCut, with the lookup in the same tree. StairCut keeps its prices, the plain step's cost and a drafted round's cost per context class, takes the plain step where drafting does not pay, and captures missing graphs while the server is idle. One greedy request a length and its repeat, second round on a fresh server: 56.6 and 57.3~tok/s at 1k against 41.0 and 45.4 plain, 42.1 and 52.4 at 140k against 33.9 and 41.4; at 260,800 tokens the class takes plain turns (\measured).

\paragraph{Caches and slots.} Every cache tier holds Flash-Next's state on a 4,096-token grid. The KV pool, side slots and a joint verify pass work on it without changing any answer, and its profiles turn them on with three slots. Three requests arriving together take as long as one after the other (27.18 against 27.03~s), and the later first tokens come at 2.2 to 6.1~s instead of up to about 20~s (\measured). The mixers are int8 with a scale a group of 128 weights: KL 0.0031 on code and 0.0188 on prose against the fp16 mixers.

\paragraph{Images and saved chats.} Flash-Next reads images: the pack's vision tower runs on the EXL3 kernels, and its profiles turn it on (VIS-97). An image chat answers the same with the MTP drafter and without it, and the text after 140 image rows stays at KL 0.019 against transformers (\measured). The disk tier now knows its files by their content (the \code{fp4} fingerprint), not by their folder. Saved chats of 0.7.2 carry over in place, or from the old folder that \code{QSE_DISK_ADOPT_FROM} names; on the reference box Flash-Next kept its 34 saved chats through the switch.

\section{What 0.7 added}

\paragraph{0.7.0: the setup guide.} The dashboard's \texttt{\#/setup} walks through an install in eight steps, with a memory bar that compares the selection's estimate with the box's free memory. One formula in \code{server/setup/estimate.py} gives the numbers, and the page's copy must pass the same cases. The models come from \code{models/<id>/model.yaml}. Wren, Qwen3.5-2B on the engine's own forward, answers from the docs, the manual and the registry and can fill fields, but cannot apply, reset, load or unload anything. Finish runs an install job of fixed steps that brings back the old profile when its check fails; reset has five levels; \code{install.sh --setup} starts the dashboard and Wren alone on a box with no engine yet. Every POST under \code{/v1/setup} must be same-origin JSON. The text does not change: with Wren loaded on the narrator's copy, three greedy prose requests gave the same 411 tokens as without it, at 148.6 against 147.5~tok/s (\measured). The KV figure is corrected to 104.25~KiB a token.

\paragraph{0.7.1 and 0.7.2.} Wren is drawn as the DGX Spark box, with eight states and a still face for reduced motion. It answers a question in words again: a tool call may name only what the page can highlight, an invalid change gets one repair, and the last round sends no tools. On ten change requests and five questions it was right 15 times of 15 (\measured). Neither release touches the engine's modules.

\section{What 0.6 added}

\paragraph{0.6.0: one shared KV pool.} Each slot reserves addresses for up to 262,144 tokens and maps 2~MiB pages as rows are written (\code{QSE_KV_POOL=on}); pages past what a slot keeps go back when its request ends. Four slots of 262,144 tokens fit beside Laya and the narrator, with 58 to 63~GiB of MemAvailable after warm-up against 26 with the old 32k side windows, and the 4-request load test went from 28.4 to 31.5~tok/s on 2 slots to 34.2 to 35.6 on 4 (\measured). A request that does not fit waits, says so in its stream and gets a 503 \code{memory_wait_timeout} only after the queue timeout. The logits, every KV row and the recurrent state match the static buffers bit for bit.

\paragraph{0.6.1 and 0.6.2.} When the pool runs short, the caches give memory back first, then the shadow reader, the narrator, idle side conversations and slot~0's idle conversation; the models come back through the residency planner once memory stays free. An EXL3 prefill offers the GPU between its layers when another slot waits: in a test of three streamed chats beside a 37k-token prefill, the seconds without a decoded token went from 42 of 147 to 0 of 145 (\measured). The dashboard's memory split counts committed pages, and its Perf tab reads the 0.6 live feed again.

\section{What 0.5 added}

\paragraph{0.5.0: weight profiles.} The engine serves the NVFP4 set or an EXL3 quant of Qwen3.8-27B from 2.5 to 6.0 bits a weight, one at a time, through \code{ops/use-profile.sh}, which fetches missing weights, starts the profile, runs a smoke request and brings the old profile back when it fails. EXL3 linear layers have one kernel per bit width from 2 to 8 for 1 to 32 verify rows, and the trellis decode matches exllamav3 bit for bit; the drafters run in EXL3 at 3 bits. Against BF16 on 16,376 held-out positions, NVFP4 has a KL of 0.0582 and agrees on the top token 93.35~\% of the time, EXL3 4.5 0.0217 and 95.78~\% (\measured). An EXL3 profile serves other weights, so its text differs from NVFP4's; speculation stays exact against whichever weights run.

\paragraph{0.5.1.} The dashboard's cache hit counts only the client's own prompt tokens, and the live prefill rate only the rows the engine computed. The reference box runs four parallel requests.

\section{What 0.4.8 adds}

0.4.8 is 0.4.7 plus the 0.4.8-test builds of 2026-10-08; \code{CHANGELOG.md} has them with their live checks.

\paragraph{Parallel requests.} A side slot takes a request whose window leaves it at least 16,384 rows of output, with \code{max_tokens} cut to the room (0.4.7 compared the window with the whole served limit, so a chat past about 24,000 tokens never took the side slot and two chats took turns on slot~0). A side slot has a resident prefix of its own: the next turn of the conversation it holds resumes at the end anchor there, beside slot~0's conversation; the router's question in between leaves it in place. A new conversation beside one a side slot could never hold takes the idle side slot first, and a short one waits up to 20~s for a busy side slot before it sends the long conversation to the NVMe.

\paragraph{Caches.} The end alignment's third pass serves a block the client moved from an earlier user message to its last one (Open WebUI's retrieval block) from the rows where the model saw it, token-aligned on the bare bytes and only when the same bytes are in the last message; a message the user edited is served as sent. The router's notes stay live by their message when the prompt changed before them. The RAM park keeps the end anchor, so a parked Open WebUI chat comes back at its end. A long prompt the end anchor does not serve says where the copy parted and what was refused, in the log and the journal.

\paragraph{Speed.} The speed guard's short gate on the 0.4.8 engine code: 50.64 against release 0.3.0's 49.75~tok/s (\measured, one run a side, store off).

\section{What 0.4.7 added}

0.4.7 is 0.4.0 plus the test builds 0.4.1-test to 0.4.7-test; \code{CHANGELOG.md} has each one with its live checks.

\paragraph{Packages.} The system set (\code{tandem-core}, \code{render-tools}, \code{eval-judge}), four library packages (\code{creative-coding}, \code{charts}, \code{voxel-worlds}, \code{explainer-video}) and 34 everyday packages with \code{everyday-core} and \code{reply-format}. Every prose class offers the whole everyday set; the package question runs once per class (\code{package_sticky: class}) on an idle side slot. Measured against no package, 91 pairs judged by the served model: a mean of +1.19 of 5, indicative only; five everyday packages ship off. Everyday answers draft first and never stop on a form.

\paragraph{Requests.} Regenerate samples (temperature 0.7, top-p 0.95, top-k 20, a fresh seed; \code{X-Tandem-Regenerate}); every other answer stays greedy, and the ledger's \code{sampling} column (schema 9) says which. Continue Response continues the partial answer. A memory note the store refuses is never shown, and a project note without a working directory lands at user scope.

\paragraph{Caches.} A long Open WebUI chat comes back from the NVMe after other chats ran (a 48k-token chat in 1.0~s, TTFT 2.3~s). The session spill is compressed without loss (1.71$\times$ on a real entry's rows), the spill budget is 64~GB and the disk tier 16~GB, and the NVMe writes at most 200~GB a day. A conversation that ran on a side slot is carried into slot~0. Open WebUI chats own their NVMe entries, and a deploy that only registers a knob keeps the disk's files.

\paragraph{Speed.} The speed guard's short gate on the 0.4.7 engine code: 49.93 against release 0.3.0's 50.38~tok/s (\measured, one run a side, store off).

\section{What 0.4.0 brought}

\paragraph{The conductor.} The engine serves one model, \code{TandemLLM}, with a conductor in front (Part IV): profiles; a router that decides each new user message's task class with a decision model and then Qwen itself; effort levels that are real caps on the reasoning; hidden notes with packages and skills, replayed at the same position; thinking budgets in every shape clients send and an answer reserve; side requests; the side lane; waiting streams with an estimate; status lines and the voice; the session monitor with its loop guards, tool checks and claim check; the tool-argument loop guard; the shadow reader; memory notes; session transcripts, turn snapshots and outcomes; session addresses. \code{X-Tandem: raw} turns all of it off for a request. Every client is the one user whose memory notes the engine keeps; \code{X-Tandem-Memory: off} keeps a request out of them. The testing packages name the engine as \code{{tandem_url}} and keep scratch files in \code{{tandem_scratch}}, out of the deliverable, put the deliverable before any test tooling and stop a test that fails the same way twice, and check a page with the kit's own headless Chrome (\code{laya_drive.py page}) instead of a driver of the model's own. The turn budget (S5) ends an agent turn that loops on its own test, sets up before it writes, or keeps a chat waiting. Asked who it is, the model names its model and TandemLLM. A prose request is not routed to a code class on a weak guess. A class's core skill goes in on every turn of the class, and an unsure deliverable is built as picked and stated in one line instead of asked. Each Open WebUI chat is its own session (its chat id as a connection header, or its first turn without one), and a new folder never replaces the user's files (the \code{files-care} package).

\paragraph{Requests.} A malformed request is a 400 that names the field; \code{developer} is the system prompt; \code{usage.prompt_tokens} counts what the client sent and the conductor's tokens are \code{prompt_tokens_details.hidden_tokens}; a JSON answer has its thinking in \code{reasoning_content}, and so has every client but the Hermes agent under a conductor plan (no literal \code{</think>} in Open WebUI's live view, none in opencode or an SDK's answer); stop strings apply to the answer only; a request without tools never gets \code{tool_calls}; \code{cached_tokens} counts the client's tokens below the point where the request resumed. A class's packages fit the request's tools, a reply-only instruction is chat, and the follow-up of an answer that thought resumes at its end. While the engine is away, a status responder answers the reverse proxy's fallback with a page and a 503 that say what is going on.

\paragraph{Parallel requests.} \code{--parallel N}, off by default: side slots beside the engine, a packed draft and a joint verify of at most 32 rows for the decoding requests, admission by bytes; two requests at once 1.38$\times$ the total of one, four 1.42$\times$ (Section~\ref{sec:parallel-requests}); a turn that continues slot~0's conversation, judged on the client's own tokens, waits for slot~0 and resumes in place, so an agent session keeps its cache with parallel requests on; the side slots' bytes stay in the memory limits, and the router's short questions and a client's title tasks take an idle side slot, so a new user message resumes the conversation in place; Open WebUI's live working directory moves from the tool block into a note. A request that shares only its client's head with slot~0's conversation takes the side slot, and a side slot reads the disk tier.

\paragraph{Caches.} The end anchor; the parked session; the engine's memory limits and the memory gate; the disk tier on the NVMe; live conversations on the NVMe through a restart or an eviction; compaction as a continuation with a summary written ahead; images left out of long conversations in steps; a canonical skills list for opencode; a tokenised-prompt cache.

The raw path runs 0.3.0's arithmetic: every build of the series passed the identity gate against the one before it (\code{tools/flagoff_identity.py}, a fixed scenario of fixed shapes, logits and state bit for bit). With the same tree shapes it writes 0.3.0's text, with one exception: the raw path no longer closes a reasoning block that repeats itself (0.3.0's stall close, on by default); under a conductor plan the session monitor's reasoning-loop task (D1) does that job. The length router sizes blocks from timings, so two runs of either version can part at a near tie (Section~\ref{sec:two-runs}). The last benchmark row recorded in the series, with \code{--no-learn-cost}\index{no-learn-cost@\texttt{-\kern0pt-no-learn-cost}} and 50 requests, read 50.09 tok/s (\measured); the paper reports 49.89. The release's speed check, 0.3.0 and 0.4.0 on the same box in the same hour with one run each, read 50.50 and 51.38 tok/s (\measured, \code{release/speed-guard.json}). The upgrade notes are in Chapter~\ref{ch:install-engine}.

\paragraph{Not verified in this release.} OpenAI's adapter with a paid call (it is tested on frames written from OpenAI's documented shapes); the conductor's second requests on a real vLLM or SGLang (tested on a fake backend); the plugins' check against running servers; the per-process GPU memory on a discrete GPU; the docker image (the wheel it installs is built and tested).

\section{Measured in the 0.4.0 series}

Table~\ref{tab:series-measurements} lists the measurements this manual cites from the series, all on the published weights on one DGX Spark.

\begin{table}[htbp]
\centering\small
\caption{Measurements from the 0.4.0 series that this manual cites. All \measured, on the published weights on one DGX Spark.}\label{tab:series-measurements}
\begin{tabularx}{\linewidth}{@{}YL{.36\linewidth}@{}}
\toprule
\hd{what} & \hd{value} \\
\midrule
side lane: 10 short requests fired into a 5,000-token turn & 9 served at 1.2~s median; turn +7.4~s; the turn's state comes back bit for bit, its text only while the tree shapes match \\
compaction of a 57,008-token request & 115.5~s as sent, 40.7~s as a continuation, 0.07~s prepared \\
session spill: a 41,587-token conversation & written in 1.03~s (5.20~GB), read back in 0.94~s \\
disk tier: a 15,360-token head & restored in 334~ms; TTFT 23.1~s down to 2.7~s \\
the voice, Qwen3.5-2B & decode $-1.90$~\% (63.47 to 62.26 tok/s) \\
the shadow reader, Qwen3.5-0.8B, 16,350-token prompt & decode $-2.05$~\% (36.02 to 35.28 tok/s) \\
memory notes on and off & 99.7 against 99.6~ms a block \\
the proxy against a fake backend & +0.33~ms to TTFT at the median \\
the host collector, one sample with a process scan & 7.46~ms of CPU \\
\bottomrule
\end{tabularx}
\end{table}


\backmatter
\chapter{Glossary}\label{ch:glossary}

\begin{description}[style=unboxed, leftmargin=0pt, itemsep=.5ex, font=\normalfont\bfseries]

\item[adapter]\index{adapter} The proxy's code for one backend kind: how it relays a route, where it reads cached and reasoning tokens, how it switches thinking. \code{proxy/adapters/}.

\item[anchor]\index{anchor} A copy of the 48 recurrent states and their convolution tails at one position, 147~MiB. The resident prefix keeps one at every 1,024-token line of a prefill; the disk tier keeps them as files. A request resumes at the largest anchor its prompt shares.

\item[backend]\index{backend} An engine or a cloud API behind the proxy: another TandemLLM, vLLM, SGLang, any OpenAI-compatible server, OpenAI or DeepSeek. Each has a kind, a URL, models, an optional key, and switches.

\item[benchmark row]\index{benchmark row} \code{serve-single-i256-o256-v1}: one request at a time, 256 prompt tokens, 256 generated, thinking off, greedy, 50 requests after 3 warm-ups. Every speed in this book that says \enquote{the benchmark row} is this.

\item[block drafter]\index{block drafter} The drafter that proposes a whole block of tokens in one pass: DFlash2, five decoder layers reading five of the target's hidden states, fine-tuned at block 8 and block 16.

\item[board]\index{board} The DGX Spark the engine runs on: one GB10 GPU and 128~GB of memory shared by the CPU and the GPU. \enquote{One engine per board} is the first rule of operations.

\item[cache source]\index{cache source} Where a request's prompt resumed: \code{resident}, \code{end}, \code{prefix}, \code{session}, \code{disk}, \code{spill}, \code{continuation}, \code{prepared} or \code{none}. In the ledger and on the dashboard.

\item[claim]\index{claim} A role's declared memory peak, which the residency planner checks before a load.

\item[class]\index{task class} The kind of task the router decides a turn is (\code{code_edit}, \code{chat} and eight more). It sets the effort, the packages and the rules.

\item[commit]\index{commit} The step that keeps the accepted path's KV rows and recurrent state after a verify. For the recurrent layers it is one low-rank update per layer, folded into the next verify.

\item[compaction]\index{compaction} A client's request to summarise a long conversation. The engine serves it as a continuation of the conversation it holds, or from a summary written ahead.

\item[conductor]\index{conductor} The harness in front of the model: profiles, the router, packages, the session monitor, status lines, memory notes. It makes the model \code{TandemLLM}.

\item[decider]\index{decider} A decision model that answers a typed question in one forward pass, such as Laya; the router's decider stage, off in the shipped router, asks one.

\item[draft tree]\index{draft tree} The guesses of one round, merged so that a shared prefix is a shared node, stored in depth-first pre-order, at most 32 nodes.

\item[drafter]\index{drafter} A cheap model or procedure that guesses the next tokens for the target to verify: the block drafter and the lookup drafter.

\item[effort]\index{effort} A turn's thinking level: off, low, medium, high or xhigh, with a reasoning cap of 0, 1,024, 8,192, 16,384 or none below the backstop.

\item[end anchor]\index{end anchor} An anchor at the end of the last generation, so the next turn that sends the answer back resumes there instead of re-reading it.

\item[feature group]\index{feature group} One switchable part of the harness in front of a backend: class, effort, packages, side, loops, e1, tools, compaction, states, memory.

\item[grid]\index{grid} The 1,024-token lines of a prefill. Anchors, chunk files and prefix snapshots are taken on the grid, and a resumed request forwards the rest of its prompt on the same grid, so its state equals a cold prefill's bit for bit.

\item[guest]\index{guest} A short request between two turns of a longer resident conversation; the rows it can reach are stashed and put back. Also the short request the side lane serves inside a running turn, which gets its engine state back bit for bit.

\item[harness]\index{harness} The layer around the model that plans each turn and watches it run; in TandemLLM it is the conductor.

\item[hold]\index{hold} Taking the board for a measurement or a second model: \code{ops/hold.sh} stops the service, runs the command, and restores the service.

\item[host collector]\index{host collector} The code that reads a machine's memory, CPU, GPUs, disks and the process behind each backend's port; \code{tandem-app hostcollect} on a remote machine.

\item[journal]\index{journal} The conductor's record of every decision, acted on or not, in the usage ledger's file, kept 30 days.

\item[L1, L2]\index{L1, L2} The two layers of engine figures behind the proxy: L1 is the engine's \code{/metrics} page, L2 the plugin's per-request numbers.

\item[lookup drafter]\index{lookup drafter} The drafter that copies text: the longest match of the last tokens in the request, a corpus and the suffix store.

\item[memory gate]\index{memory gate} The check before every prefill that gives caches back, cheapest first, until the prompt's workspace fits under the declared peak, else a 503.

\item[memory note]\index{memory note} One fact the model keeps as a file, written from a block in its answer or by the background writer, read back in later turns.

\item[memory priority]\index{memory priority} A chat's place in the line for memory, \code{tandem.priority} in the request: \code{"high"}, \code{"normal"} or \code{"low"}. Set per chat in the Chat tab (Section~\ref{sec:chat-priority}).

\item[memory wait]\index{memory wait} A request that does not fit above the floor waits before its prefill and starts by itself once memory frees; \code{waiting_memory} on the dashboard, \code{X-Tandem-Wait: memory} on the response (Section~\ref{sec:memory-wait}).

\item[narrator]\index{narrator} The small model that writes status lines: rewritten templates, thought summaries, or the voice.

\item[near tie]\index{near tie} A position where the two best tokens score within one unit in the last place, where a batched verify and a one-token step may choose differently.

\item[NVMe]\index{NVMe} The box's SSD. The disk tier and the session spill keep their files on it.

\item[note]\index{note} The router's hidden text at the end of a turn's last message: the effort sentence, rules, packages. Replayed at the same position on later turns.

\item[outcome]\index{outcome} The structural reading of a tool-result turn: progress, no visible change, regression, error, stuck and others.

\item[package]\index{package} A text the router may put into a note: a task skill, a test skill, a card, a coding guide, or a library skill.

\item[park]\index{park} Copying another session's rows aside when a request of a new session replaces the resident prompt, so that session's next turn swaps them back.

\item[pin]\index{pin} Copying the head of the resident prompt aside before a long unrelated request cuts it.

\item[plugin]\index{plugin} A small package inside vLLM or SGLang that posts each request's own times to the proxy.

\item[profile]\index{profile} Two things. A conductor profile is a document that sets what the harness does for a client (\code{coding}, \code{chat} and others). A serving profile is a \code{serve.env} file that sets the weights.

\item[proxy]\index{proxy} The router of requests by model id to backends, with one ledger and one dashboard; in the engine's server with \code{--backends}, or standalone as \code{tandem-app}.

\item[raw]\index{raw} The engine without the conductor: the release's raw path. \code{X-Tandem: raw} or the \code{raw} and \code{bench} profiles.

\item[residency planner]\index{residency planner} The code that admits a role's load only when the board keeps its floor after the engine's growth to its peak.

\item[resident prefix]\index{resident prefix} The last prompt's rows, left in the engine's own buffers, with anchors, an end anchor, a stash, a pin and a park.

\item[role]\index{role} Anything beside the big model the conductor asks, with a contract: its kind, claim, window, latency budget, fallback and what it may write.

\item[round]\index{round} One draft, verify, accept and commit: about 93~ms and 4.58 tokens on the benchmark row.

\item[router]\index{router} Two things. The conductor's router decides a turn's class and plan. The speculation router (StairCut, the length router) decides each round's draft tree.

\item[safe point]\index{safe point} One of the four moments the session monitor may act: before the prefill, after a block, at the end token, at the held tool calls.

\item[session resume]\index{session resume} Continuing a conversation from state that verify rounds wrote (the end anchor, the session spill) rather than prefill chunks. It is held to the greedy gate (the same tokens), not to bit equality; the \emph{taint} row marks where it begins.

\item[session monitor]\index{session monitor} The conductor's watch over a running turn: loop guards, the empty-turn continue, tool-call checks, the claim check, session tasks.

\item[session spill]\index{session spill} Whole live conversations written to the NVMe at a planned stop, while idle or before an eviction, and read back: the swap of long chats to the SSD (Section~\ref{sec:spill}).

\item[shadow]\index{shadow mode} A mode in which a part decides and journals what it would do and changes nothing.

\item[shadow reader]\index{shadow reader} A small model that keeps its own exact state of each session, read in idle time, and answers typed questions over the whole session.

\item[side lane]\index{side lane} Serving a short request at a safe point inside the running generation, which then gets its engine state back bit for bit. Its text equals an uninterrupted run's only while the tree shapes are the same (Section~\ref{sec:two-runs}).

\item[side request]\index{side request} A client's title or summary call, known by its shape and served small.

\item[spill]\index{spill} See session spill.

\item[SSD]\index{SSD} Two things. The box's NVMe drive, where the session spill swaps long chats out (Section~\ref{sec:spill}). In NemotronH, Mamba2's chunked state-space form of the prefill (Chapter~\ref{ch:nemotron}).

\item[StairCut]\index{StairCut} The router that cuts each round's draft tree to the node count with the most expected tokens per millisecond on the measured verify-cost staircase.

\item[state store]\index{state store} Whole state snapshots, for the session cache and the prefix cache.

\item[status line]\index{status line} One short line in the client for a step the conductor takes, stripped from whatever the client sends back.

\item[suffix store]\index{suffix store} An append-only log of the token ids the engine read and wrote, with a suffix array, for the lookup drafter.

\item[swap]\index{swap} Two things. The swap of long chats to the SSD is the session spill (Section~\ref{sec:spill}). The park swaps two sessions' rows in the engine's memory (Section~\ref{sec:resident}).

\item[taint]\index{taint} The lowest row that a session resume wrote. Below it the caches are a cold prefill's bit for bit.

\item[TandemApp]\index{TandemApp} The dashboard and the proxy without the engine, as one installable command, \code{tandem-app} (Chapter~\ref{ch:install-app}).

\item[time to first token]\index{time to first token} The time from a request's arrival to its first streamed token: the prefill, plus any wait in the queue. TTFT in the tables.

\item[ulp]\index{ulp} A unit in the last place: one rounding step of BF16 at a value.

\item[verify]\index{verify} One forward pass of the target model over the whole draft tree, reading the model's own token at every node.

\item[voice]\index{voice} The narrator's live lines: facts from the engine and a short checked phrase from the small model.

\item[writer]\index{memory writer} The background turn that writes a session's facts as memory notes at its end or after a compaction.

\end{description}

\printindex

\end{document}
